% Generated mechanically from frozen input p01/ko/sites.tex.
% Do not edit this generated file; edit the bound locale source instead.

% OK, start here.
%

\setcounter{chapter}{6}
\chapter{사이트와 층}



\phantomsection
\label{sites-section-phantom}


\section{서론}
\label{sites-section-introduction}

\noindent
사이트라는 개념은 Grothendieck이 스킴의 \'etale 위상에서 층을 연구할 수
있도록 도입하였다. 이 개념의 기본 참고문헌은 아마도 \cite{SGA4}일 것이다.
여기서 사용하는 사이트의 개념은 \cite{SGA4}의 것과 다르다. 우리가
사이트라고 부르는 것은 \cite[Expos\'e II, D\'efinition 1.3]{SGA4}에서
전위상을 갖춘 범주라고 부르는 것이다. 이렇게 하는 까닭은 대수기하학에서
주어진 덮개들의 모임을 사용하는 편이 흔히 편리하기 때문이다. 예를 들어
스킴의 어떤 성질이 주어진 위상에 관하여 언제 국소적인지를 정의할 때가
그렇다. Descent, 절 \ref{descent-section-descending-properties}을 보라.
이 장의 서술은 \cite{ArtinTopologies}를 충실히 따를 것이다.
우주는 사용하지 않는다.










\section{준층}
\label{sites-section-presheaves}

\noindent
$\mathcal{C}$를 범주라 하자.
{\it 집합의 준층}은 $\mathcal{C}$에서 $\textit{Sets}$로 가는 반변 함자
$\mathcal{F}$이다(Categories, Remark
\ref{categories-remark-functor-into-sets} 참조).
따라서 $\mathcal{C}$의 각 대상 $U$에 대하여 집합
$\mathcal{F}(U)$가 있다. 이 집합의 원소를 $U$ 위에서의
$\mathcal{F}$의 {\it 단면}이라 한다. 각 사상 $f : V \to U$에 대하여
사상 $\mathcal{F}(f) : \mathcal{F}(U) \to \mathcal{F}(V)$를
{\it 제한 사상}이라 하고, 흔히
$f^\ast : \mathcal{F}(U) \to \mathcal{F}(V)$로 나타낸다.
달리 말하면 $f^*(s)$는 $f$에 의한 $s$의 {\it 당김}이다.
함자성은 $g^* f^* (s) = (f \circ g)^*(s)$를 뜻한다.
때로는 $s|_V := f^\ast(s)$라는 표기를 사용한다.
이 표기는 위상수학에서 함수의 제한이라는 개념과 양립한다. 실제로
$W \to V \to U$가 $\mathcal{C}$의 사상들이고 $s$가 $U$ 위에서의
$\mathcal{F}$의 단면이면, $\mathcal{F}$의 함자성에 따라
$s|_W = (s|_V)|_W$이다. 물론 사상 $V \to U$가 둘 이상 존재하는 일은
충분히 가능하고, 그에 대응하는 당김 사상들이 같을 리는 없으므로
주의해야 한다.

\begin{definition}
\label{sites-definition-presheaves-sets}
$\mathcal{C}$ 위의 {\it 집합의 준층}은 $\mathcal{C}$에서
$\textit{Sets}$로 가는 반변 함자이다. {\it 준층의 사상}은 함자들의
변환이다. 집합의 준층들이 이루는 범주를
$\textit{PSh}(\mathcal{C})$로 나타낸다.
\end{definition}

\noindent
$\mathcal{C}$의 임의의 대상 $U$에 대하여 점 함자 $h_U$는
준층임에 유의하라(Categories, Example
\ref{categories-example-hom-functor} 참조). 이들을
{\it 표현 가능한 준층}이라 한다. 이 준층들은 임의의 준층
$\mathcal{F}$에 대하여
\begin{equation}
\label{sites-equation-map-representable-into-presheaf}
\Mor_{\textit{PSh}(\mathcal{C})}(h_U, \mathcal{F})
=
\mathcal{F}(U).
\end{equation}
이라는 좋은 성질을 갖는다. 이것이 Yoneda 보조정리이다
(Categories, Lemma \ref{categories-lemma-yoneda}).

\medskip\noindent
마찬가지로 아벨군의 준층, 환의 준층 등의 개념을 정의할 수 있다.
더 일반적으로는 어떤 범주에 값을 갖는 준층을 정의할 수 있다.

\begin{definition}
\label{sites-definition-presheaf}
$\mathcal{C}$와 $\mathcal{A}$를 범주라 하자.
$\mathcal{C}$ 위에서 $\mathcal{A}$에 값을 갖는 {\it 준층}
$\mathcal{F}$는 $\mathcal{C}$에서 $\mathcal{A}$로 가는 반변 함자,
즉 $\mathcal{F} : \mathcal{C}^{opp} \to \mathcal{A}$이다.
$\mathcal{C}$ 위에서 $\mathcal{A}$에 값을 갖는 준층들의
{\it 사상} $\mathcal{F} \to \mathcal{G}$는 $\mathcal{F}$에서
$\mathcal{G}$로 가는 함자들의 변환이다.
\end{definition}

\noindent
이들은 $\mathcal{C}$ 위에서 $\mathcal{A}$에 값을 갖는 준층들의
범주에서 각각 대상과 사상을 이룬다.

\begin{remark}
\label{sites-remark-big-presheaves}
앞서 지적했듯이 Categories, Remark
\ref{categories-remark-big-categories}에 열거된 임의의 “큰” 범주에
값을 갖는 준층들의 범주를 생각할 수 있다. 이들 또한 “큰” 범주이며,
논의를 진행하면서 위의 Remark에 열거할 것이다.
\end{remark}















\section{준층의 단사 사상과 전사 사상}
\label{sites-section-injective-surjective}

\begin{definition}
\label{sites-definition-presheaves-injective-surjective}
$\mathcal{C}$를 범주라 하고, $\varphi : \mathcal{F}
\to \mathcal{G}$를 집합의 준층들의 사상이라 하자.
\begin{enumerate}
\item $\mathcal{C}$의 모든 대상 $U$에 대하여 사상
$\varphi_U : \mathcal{F}(U) \to \mathcal{G}(U)$가 단사이면
$\varphi$가 {\it 단사}라고 한다.
\item $\mathcal{C}$의 모든 대상 $U$에 대하여 사상
$\varphi_U : \mathcal{F}(U) \to \mathcal{G}(U)$가 전사이면
$\varphi$가 {\it 전사}라고 한다.
\end{enumerate}
\end{definition}

\begin{lemma}
\label{sites-lemma-mono-epi}
위에서 정의한 단사인 사상(각각 전사인 사상)은 정확히
$\textit{PSh}(\mathcal{C})$의 단사사상(각각 전사사상)이다.
사상이 동형사상일 필요충분조건은 단사이면서 전사인 것이다.
\end{lemma}








\begin{proof}
$\varphi : \mathcal{F} \to \mathcal{G}$가 단사일 필요충분조건이
$\textit{PSh}(\mathcal{C})$의 단사사상이라는 것을 보이자.
단사이면 단사사상인 방향은 명백하므로, 그 역을 보이면 된다.
$\varphi$가 단사사상이라고 가정하자.
$U \in \Ob(\mathcal{C})$라 하자. $\varphi_U$가 단사임을 보여야 한다.
$a, b \in \mathcal{F}(U)$가
$\varphi_U (a) = \varphi_U (b)$를 만족한다고 하자. $a = b$임을 확인해야 한다.
동형사상 (\ref{sites-equation-map-representable-into-presheaf}) 아래에서
$\mathcal{F}(U)$의 원소 $a$와 $b$는 두 자연변환
$a', b' \in \Mor_{\textit{PSh}(\mathcal{C})}(h_U, \mathcal{F})$에
대응한다. 마찬가지로 비슷한 동형사상
$\Mor_{\textit{PSh}(\mathcal{C})}(h_U, \mathcal{G})
= \mathcal{G}(U)$ 아래에서 $\mathcal{G}(U)$의 두 같은 원소
$\varphi_U (a)$와 $\varphi_U (b)$는 두 자연변환
$\varphi \circ a', \varphi \circ b'
\in \Mor_{\textit{PSh}(\mathcal{C})}(h_U, \mathcal{G})$에 대응하므로,
이 두 자연변환도 같아야 한다. 따라서
$\varphi \circ a' = \varphi \circ b'$이고, $\varphi$가 단사사상이므로
$a' = b'$이며, 이에 따라 $a = b$이다. 이것으로 (1)이 끝난다.

\medskip\noindent
$\varphi : \mathcal{F} \to \mathcal{G}$가 전사일 필요충분조건이
$\textit{PSh}(\mathcal{C})$의 전사사상이라는 것을 보이자.
전사이면 전사사상인 방향은 명백하므로, 그 역을 보이면 된다.
$\varphi$가 전사사상이라고 가정하자.

\medskip\noindent
범주 $\textit{Sets}$의 임의의 두 사상 $f : A \to B$와 $g : A \to C$에
대하여, $B$와 $C$에서 $B \coprod_A C$로 가는 두 표준 사상을 각각
$\text{inl}_{f,g}$와 $\text{inr}_{f,g}$로 나타낸다.
(여기서 밂은 $\textit{Sets}$에서 계산한다.)

\medskip\noindent
이제 $\mathcal{C}$ 위의 집합의 준층 $\mathcal{H}$를 다음과 같이 정의하자.
각 $U \in \Ob(\mathcal{C})$에 대하여
$\mathcal{H}(U)
= \mathcal{G}(U) \coprod_{\mathcal{F}(U)} \mathcal{G}(U)$로 놓는다
(여기서 밂은 $\textit{Sets}$에서 사상
$\varphi_U : \mathcal{F}(U) \to \mathcal{G}(U)$에 의해 계산한다).
사상들에 대한 작용은 밂의 함자성에 의해 명백한 방식으로 정의한다.
그러면 두 자연변환 $i_1 : \mathcal{G} \to \mathcal{H}$와
$i_2 : \mathcal{G} \to \mathcal{H}$가 있으며, 대상
$U \in \Ob(\mathcal{C})$에서의 그 성분은 각각 사상
$\text{inl}_{\varphi_U, \varphi_U}$와
$\text{inr}_{\varphi_U, \varphi_U}$로 주어진다.
밂의 정의에 따라 $i_1 \circ \varphi
= i_2 \circ \varphi$이다. $\varphi$가 전사사상이므로 $i_1 = i_2$이다.
따라서 모든 $U \in \Ob(\mathcal{C})$에 대하여
$\text{inl}_{\varphi_U, \varphi_U}
= \text{inr}_{\varphi_U, \varphi_U}$이다. 그러므로 $\varphi_U$는
전사여야 한다. 실제로 간단한 조합론적 논증에 따라
$f : A \to B$가 $\textit{Sets}$의 사상이면
$\text{inl}_{f,f} = \text{inr}_{f,f}$일 필요충분조건이
$f$가 전사인 것이다. 다시 말해 $\varphi$는 전사이며, (2)가 증명되었다.

\medskip\noindent
$\varphi : \mathcal{F} \to \mathcal{G}$가 단사이면서 전사일 필요충분조건이
$\textit{PSh}(\mathcal{C})$의 동형사상이라는 것을 보이자.
이번에는 동형사상이면 단사이면서 전사인 방향이 명백하다.
그 역을 보이기 위해서는 다음에 유의하면 충분하다. $\varphi$가 단사이면서
전사이면 모든 $U \in \Ob(\mathcal{C})$에 대하여 $\varphi_U$는 가역 사상이고,
모든 $U$에 대한 이 사상들의 역들을 합치면 $\varphi$의 역인 자연변환
$\mathcal{G} \to \mathcal{F}$를 얻는다.
\end{proof}

\begin{definition}
\label{sites-definition-sub-presheaf}
모든 대상 $U \in \Ob(\mathcal{C})$에 대하여 집합
$\mathcal{F}(U)$가 $\mathcal{G}(U)$의 부분집합이고, 이것이 제한 사상들과
양립하면 $\mathcal{F}$를 $\mathcal{G}$의 {\it 부분준층}이라고 한다.
\end{definition}

\noindent
달리 말하면, 포함 사상들 $\mathcal{F}(U) \to \mathcal{G}(U)$를
합치면 준층들의 (단사인) 사상
$\mathcal{F} \to \mathcal{G}$가 된다.

\begin{lemma}
\label{sites-lemma-image}
$\mathcal{C}$를 범주라 하자.
$\varphi : \mathcal{F} \to \mathcal{G}$가 $\mathcal{C}$ 위의
집합의 준층들의 사상이라고 하자. $\varphi$가
$\mathcal{F} \to \mathcal{G}' \to \mathcal{G}$로 분해되고,
첫 번째 사상이 전사가 되게 하는 부분준층
$\mathcal{G}' \subset \mathcal{G}$가 유일하게 존재한다.
\end{lemma}

\begin{proof}
존재성을 보이려면 모든 $U \in \Ob(C)$에 대하여
% P01-ERRATA: P01-E0057 -- frozen source drops \mathcal from C.
$\mathcal{G}'(U) = \varphi_U \left(\mathcal{F}(U)\right)$로 놓고
(사상들에 대한 작용은 $\mathcal{G}$에서 물려받는다), 이것이
$\mathcal{G}$의 부분준층을 정의하며 $\varphi$가
$\mathcal{F} \to \mathcal{G}' \to \mathcal{G}$로 분해되고
첫 번째 사상이 전사임을 보이면 된다. 유일성은 명백하다.
\end{proof}

\begin{definition}
\label{sites-definition-image}
표기는 Lemma \ref{sites-lemma-image}와 같이 하자.
$\mathcal{G}'$를 {\it $\varphi$의 상}이라고 한다.
\end{definition}

\section{준층의 극한과 여극한}
\label{sites-section-limits-colimits-PSh}

\noindent
$\mathcal{C}$를 범주라 하자.
범주 $\textit{PSh}(\mathcal{C})$에는 극한과 여극한이 존재한다.
더 나아가 모든 $U \in \Ob(\mathcal{C})$에 대하여 함자
$$
\textit{PSh}(\mathcal{C})
\longrightarrow
\textit{Sets}, \quad
\mathcal{F}
\longmapsto
\mathcal{F}(U)
$$
는 극한 및 여극한과 가환한다. 이 명제들을 증명하는 가장 쉬운 방법은
아마 다음과 같다. 도식
$
\mathcal{F} :
\mathcal{I}
\to
\textit{PSh}(\mathcal{C})
$
가 주어지면 준층
$$
\mathcal{F}_{\lim} :
U
\longmapsto
\lim_{i \in \mathcal{I}} \mathcal{F}_i(U)
\text{  and  }
\mathcal{F}_{\colim} :
U
\longmapsto
\colim_{i \in \mathcal{I}} \mathcal{F}_i(U)
$$
을 정의한다. 사영 사상 $\mathcal{F}_{\lim} \to \mathcal{F}_i$와
표준 사상 $\mathcal{F}_i \to \mathcal{F}_{\colim}$가 분명히 존재한다.
이 사상들은 Categories, Definition
\ref{categories-definition-limit}의 극한 사상들(각각 Categories,
Definition \ref{categories-definition-colimit}의 여극한 사상들)에
요구되는 조건을 만족한다. 실제로 이들은 $\mathcal{F}$ 위의
원뿔, 각각 쌍대원뿔을 분명히 이룬다. 더 나아가
$(\mathcal{G}, q_i : \mathcal{G} \to \mathcal{F}_i)$가
다른 한 계라면(극한의 정의에서와 같이), 모든 $U$에 대하여 적절한
함자성 조건을 만족하는 사상들의 계
$\mathcal{G}(U) \to \mathcal{F}_i(U)$를 얻는다. 따라서 유일한 사상
$\mathcal{G}(U) \to \mathcal{F}_{\lim}(U)$를 얻는다.
$U$를 변화시킬 때 이 사상들이 양립하며 원하는 사상
$\mathcal{G} \to \mathcal{F}_{\lim}$에서 나옴은 쉽게 확인된다.
여극한의 경우에도 같은 논증이 성립한다.

\section{준층 범주의 함자성}
\label{sites-section-functoriality-PSh}

\noindent
$u : \mathcal{C} \to \mathcal{D}$를 범주들 사이의 함자라 하자.
이 경우
$$
u^p :
\textit{PSh}(\mathcal{D})
\longrightarrow
\textit{PSh}(\mathcal{C})
$$
는 $\mathcal{D}$ 위의 $\mathcal{G}$에 준층
$u^p\mathcal{G} = \mathcal{G} \circ u$를 대응시키는 함자를 나타낸다.
앞 절에 따라 이 함자는 모든 극한 및 여극한과 가환함에 유의하라.

\medskip\noindent
$V \in \Ob(\mathcal{D})$에 대하여 $\mathcal{I}^u_V$로 다음 범주를
나타내자.
\begin{equation}
\label{sites-equation-colim-category}
\begin{matrix}
\Ob(\mathcal{I}^u_V)
&
=
&
\{
(U, \phi)
\mid
U \in \Ob(\mathcal{C}),
\phi : V \to u(U)
\}
\\
\Mor_{\mathcal{I}^u_V}((U, \phi), (U', \phi'))
&
=
&
\{
f : U \to U' \text{ in }\mathcal{C}
\mid
u(f) \circ \phi = \phi'
\}
\end{matrix}
\end{equation}
때로는 표기에서 위첨자 $ {}^u$를 생략하고 단순히
$\mathcal{I}_V$라고 쓴다. 이 범주들을 사용하여 함자 $u^p$의
왼쪽 수반을 정의할 것이다. 그 전에 몇 가지 기술적 보조정리를 증명한다.

\begin{lemma}
\label{sites-lemma-almost-directed}
$u : \mathcal{C} \to \mathcal{D}$를 범주들 사이의 함자라 하자.
$\mathcal{C}$가 섬유곱과 등화자를 가지며 $u$가 이들과 가환한다고 하자.
그러면 범주 $(\mathcal{I}_V)^{opp}$는 Categories, Lemma
\ref{categories-lemma-split-into-directed}의 가정을 만족한다.
\end{lemma}

\begin{proof}
확인할 조건은 두 가지이다.

\medskip\noindent
먼저 세 대상
$\phi : V \to u(U)$, $\phi' : V \to u(U')$, $\phi'' : V \to u(U'')$와,
$u(a) \circ \phi' = \phi$ 및 $u(b) \circ \phi'' = \phi$를 만족하는 사상
$a : U' \to U$, $b : U'' \to U$가 주어졌다고 하자.
다른 대상 $\phi''' : V \to u(U''')$와 사상
$c : U''' \to U'$, $d : U''' \to U''$로서
$u(c) \circ \phi''' = \phi'$, $u(d) \circ \phi''' = \phi''$ 및
$a \circ c = b \circ d$를 만족하는 것들이 존재함을 보여야 한다.
$U''' = U' \times_U U''$로 놓고 $c$와 $d$를 사영 사상으로 취한다.
$u$가 섬유곱과 가환하므로 이것으로 충분하다. 확인은 생략한다.

\medskip\noindent
둘째, 두 대상 $\phi : V \to u(U)$와 $\phi' : V \to u(U')$, 그리고 사상
$a, b : (U, \phi) \to (U', \phi')$가 주어졌다고 하자.
$a$와 $b$를 등화하는 사상
$c : (U'', \phi'') \to (U, \phi)$를 찾아야 한다.
$c : U'' \to U$를 범주 $\mathcal{C}$에서 $a$와 $b$의 등화자라 하자.
$u$가 등화자와 가환하고
$u(a) \circ \phi = u(b) \circ \phi = \phi'$이므로
사상 $\phi'' : V \to u(U'')$를 얻는다.
\end{proof}

\begin{lemma}
\label{sites-lemma-directed}
$u : \mathcal{C} \to \mathcal{D}$를 범주들 사이의 함자라 하자.
다음을 가정하자.
\begin{enumerate}
\item 범주 $\mathcal{C}$가 끝 대상 $X$를 갖고,
$u(X)$가 $\mathcal{D}$의 끝 대상이다.
\item 범주 $\mathcal{C}$가 섬유곱을 가지며 $u$가 이들과 가환한다.
\end{enumerate}
그러면 지표 범주 $(\mathcal{I}^u_V)^{opp}$는 여과되어 있다
(Categories, Definition \ref{categories-definition-directed} 참조).
\end{lemma}

\begin{proof}
가정들로부터 Lemma \ref{sites-lemma-almost-directed}의 가정들이 성립한다
(Categories, Section \ref{categories-section-finite-limits}의 논의를 보라).
Categories, Lemma \ref{categories-lemma-split-into-directed}에 따라
$\mathcal{I}_V$는 유향 범주들의 (어쩌면 빈) 서로소 합집합이다.
따라서 $\mathcal{I}_V$가 연결임을 보이면 충분하다.

\medskip\noindent
먼저 $\mathcal{I}_V$가 공집합이 아님을 보이자.
가정에 따라 존재하는 $\mathcal{C}$의 끝 대상을 $X$라 하자.
가정에 따라 $u(X)$가 $\mathcal{D}$에서 끝 대상이라는 사실로부터
사상 $V \to u(X)$를 얻는다. 이것은 $\mathcal{I}_V$의 대상을 준다.

\medskip\noindent
둘째, $\mathcal{I}_V$가 연결임을 보이자.
$\phi_1 : V \to u(U_1)$과 $\phi_2 : V \to u(U_2)$가
$\Ob(\mathcal{I}_V)$에 속한다고 하자. 가정에 따라 $U_1\times U_2$가
존재하고 $u(U_1\times U_2) = u(U_1)\times u(U_2)$이다.
곱의 보편 성질에 의해 $(\phi_1, \phi_2)$에 대응하는 사상
$\phi : V \to u(U_1\times U_2)$를 생각하자.
대상 $\phi : V \to u(U_1\times U_2)$는 분명히
$\phi_1 : V \to u(U_1)$과 $\phi_2 : V \to u(U_2)$ 모두로 사상된다.
\end{proof}

\noindent
$\mathcal{D}$의 사상 $g : V' \to V$가 주어지면, 대상 위에서
$\overline{g}(U, \phi) = (U, \phi \circ g)$로 놓아 함자
$\overline{g} : \mathcal{I}_V \to \mathcal{I}_{V'}$를 얻는다.
$\mathcal{C}$ 위의 준층 $\mathcal{F}$가 주어지면 함자
$$
\mathcal{F}_V :
\mathcal{I}_V^{opp}
\longrightarrow
\textit{Sets}, \quad
(U, \phi)
\longmapsto
\mathcal{F}(U).
$$
를 얻는다. 달리 말하면 $\mathcal{F}_V$는 $\mathcal{I}_V$ 위의
집합의 준층이다.
$\mathcal{F}_{V'} \circ \overline{g} = \mathcal{F}_V$임에 유의하라.
다음과 같이 정의한다.
$$
u_p\mathcal{F}(V) =
\colim_{\mathcal{I}_V^{opp}} \mathcal{F}_V
$$
여극한이므로 각 $(U, \phi) \in \Ob(\mathcal{I}_V)$에 대하여 표준 사상
$\mathcal{F}(U)\xrightarrow{c(\phi)}u_p\mathcal{F}(V)$를 얻는다.
위와 같은 $g : V' \to V$에 대하여
$\mathcal{F}_{V'} \circ \overline{g} = \mathcal{F}_V$와 양립하는
표준 제한 사상
$g^* : u_p\mathcal{F}(V) \to u_p\mathcal{F}(V')$가
Categories, Lemma \ref{categories-lemma-functorial-colimit}에 의해 존재한다.
이는 모든 $(U, \phi) \in \Ob(\mathcal{I}_V)$에 대하여 도식
$$
\xymatrix{
\mathcal{F}(U) \ar[r]^{c(\phi)} \ar[d]_{\text{id}}
&
u_p\mathcal{F}(V) \ar[d]^{g^*}
\\
\mathcal{F}(U) \ar[r]^{c(\phi \circ g)}
&
u_p\mathcal{F}(V')
}
$$
이 가환하게 하는 유일한 사상이다. 이 사상들의 유일성에 따라 준층을 얻는다.
이 준층을 $u_p\mathcal{F}$로 나타낼 것이다.

\begin{lemma}
\label{sites-lemma-recover}
제한 사상들($\mathcal{F}$의 것과 $u_p\mathcal{F}$의 것)과 양립하는
표준 사상
$\mathcal{F}(U) \to u_p\mathcal{F}(u(U))$
이 존재한다.
\end{lemma}

\begin{proof}
이는 위에서 도입한 사상 $c(\text{id}_{u(U)})$일 뿐이다.
\end{proof}

\noindent
준층들의 임의의 사상 $\mathcal{F} \to \mathcal{F}'$은 함자들 사이의
양립하는 사상들의 계 $\mathcal{F}_V \to \mathcal{F}'_V$를 주고,
따라서 준층들의 사상
$u_p\mathcal{F} \to u_p\mathcal{F}'$을 준다는 점에 유의하라.
달리 말하면 함자
$$
u_p :
\textit{PSh}(\mathcal{C})
\longrightarrow
\textit{PSh}(\mathcal{D})
$$
를 정의하였다.

\begin{lemma}
\label{sites-lemma-adjoints-u}
함자 $u_p$는 함자 $u^p$의 왼쪽 수반이다. 달리 말하면 공식
$$
\Mor_{\textit{PSh}(\mathcal{C})}(\mathcal{F}, u^p\mathcal{G})
=
\Mor_{\textit{PSh}(\mathcal{D})}(u_p\mathcal{F}, \mathcal{G})
$$
은 $\mathcal{F}$와 $\mathcal{G}$에 대하여 쌍함자적으로 성립한다.
\end{lemma}

\begin{proof}
$\mathcal{G}$를 $\mathcal{D}$ 위의 준층이라 하고,
$\mathcal{F}$를 $\mathcal{C}$ 위의 준층이라 하자.
양방향의 사상을 구성하여 표시한 공식이 성립함을 보이겠다.
% P01-ERRATA: P01-E1329 -- frozen English needs “each other's inverses”.
그 사상들이 서로의 역임을 확인하는 일은 독자에게 맡긴다.

\medskip\noindent
사상 $\alpha : u_p \mathcal{F} \to \mathcal{G}$가 주어지면
$u^p\alpha : u^p u_p \mathcal{F} \to u^p \mathcal{G}$를 얻는다.
Lemma \ref{sites-lemma-recover}에 따라 사상
$\mathcal{F} \to u^p u_p \mathcal{F}$가 존재한다.
이 둘의 합성은 원하는 사상이다.
(이 구성의 좋은 점은 보이는 모든 것에 대하여 명백히 함자적이라는 것이다.)

\medskip\noindent
반대로 사상 $\beta : \mathcal{F} \to u^p\mathcal{G}$가 주어지면 사상
$u_p\beta : u_p\mathcal{F} \to u_p u^p\mathcal{G}$를 얻는다.
$V \in \Ob(\mathcal{D})$라 하자.
$\mathcal{I}_V$ 위의 함자 $u^p\mathcal{G}_V$에서 값이
$\mathcal{G}(V)$인 상수 함자로 가는 표준 사상이 있다고 주장한다.
실제로 $\mathcal{I}_V$의 모든 대상 $(U, \phi)$에 대하여
이 대상에서 $u^p\mathcal{G}_V$의 값은 $\mathcal{G}(u(U))$이며,
이는 $\mathcal{G}(\phi) = \phi^*$에 의해 $\mathcal{G}(V)$로 사상된다.
$\mathcal{G}$ 자체가 함자이므로 이것은 함자들의 변환이다.
따라서 사상 $u_p u^p \mathcal{G}(V) \to \mathcal{G}(V)$를 얻는다.
다시 자명하게 확인하면 이것은 $V$에 대하여 함자적이므로,
준층들의 사상 $u_p u^p \mathcal{G} \to \mathcal{G}$를 얻는다.
합성
$u_p\mathcal{F} \to u_p u^p\mathcal{G} \to
\mathcal{G}$가 원하는 사상이다.
\end{proof}

\begin{remark}
\label{sites-remark-functoriality-presheaves-values}
$\mathcal{A}$가, 임의의 도식 $\mathcal{I}_Y \to \mathcal{A}$가
$\mathcal{A}$ 안에서 여극한을 갖게 하는 범주라고 하자.
이 경우 $\mathcal{A}$에 값을 갖는 준층들의 범주 위에서도
위와 완전히 같은 방식으로 정의한 함자 $u^p$와 $u_p$가 있음은 명백하다.
또한 쌍 $u^p$와 $u_p$의 수반성은 이 상황에서도 계속 성립한다.
\end{remark}

\begin{lemma}
\label{sites-lemma-pullback-representable-presheaf}
$u : \mathcal{C} \to \mathcal{D}$를 범주들 사이의 함자라 하자.
$\mathcal{C}$의 임의의 대상 $U$에 대하여 $u_ph_U = h_{u(U)}$이다.
\end{lemma}

\begin{proof}
$u_p$와 $u^p$의 수반성에 따라
$$
\Mor_{\textit{PSh}(\mathcal{D})}(u_ph_U, \mathcal{G})
=
\Mor_{\textit{PSh}(\mathcal{C})}(h_U, u^p\mathcal{G})
=
u^p\mathcal{G}(U) =
\mathcal{G}(u(U))
$$
이다. 따라서 Yoneda 보조정리에 의해 준층으로서
$u_ph_U = h_{u(U)}$임을 알 수 있다.
\end{proof}























\section{사이트}
\label{sites-section-sites-definitions}

\noindent
여기서 사용하는 사이트의 개념에는 다음과 같은 구조가 쓰인다.

\begin{definition}
\label{sites-definition-family-morphisms-fixed-target}
$\mathcal{C}$를 범주라 하자. Conventions, Section
\ref{conventions-section-categories}을 보라.
$\mathcal{C}$에서 {\it 공역이 고정된 사상족}은 대상
$U \in \Ob(\mathcal{C})$, 집합 $I$, 그리고 각 $i\in I$에 대하여 공역이
$U$인 $\mathcal{C}$의 사상 $U_i \to U$로 주어진다.
이를 $\{U_i \to U\}_{i\in I}$로 나타낸다.
\end{definition}

\noindent
집합 $I$가 공집합일 수도 있다! 이 표기는 위상수학에서의 열린 덮개를
연상시키려는 것이다.

\begin{definition}
\label{sites-definition-site}
{\it 사이트}\footnote{서론에서 설명했듯이 이 표기는 \cite{SGA4}의 표기와
다르다.}는 범주 $\mathcal{C}$와, {\it $\mathcal{C}$의 덮개}라고 부르는
공역이 고정된 사상족 $\{U_i \to U\}_{i \in I}$들의 집합
$\text{Cov}(\mathcal{C})$로 주어지며, 다음 공리들을 만족한다.
\begin{enumerate}
\item $V \to U$가 동형사상이면 $\{V \to U\} \in
\text{Cov}(\mathcal{C})$이다.
\item $\{U_i \to U\}_{i\in I} \in \text{Cov}(\mathcal{C})$이고 각
$i$에 대하여 $\{V_{ij} \to U_i\}_{j\in J_i} \in
\text{Cov}(\mathcal{C})$이면
$\{V_{ij} \to U\}_{i \in I, j\in J_i} \in \text{Cov}(\mathcal{C})$이다.
\item $\{U_i \to U\}_{i\in I}\in \text{Cov}(\mathcal{C})$이고
$V \to U$가 $\mathcal{C}$의 사상이면, 모든 $i$에 대하여
$U_i \times_U V$가 존재하고
$\{U_i \times_U V \to V \}_{i\in I} \in \text{Cov}(\mathcal{C})$이다.
\end{enumerate}
\end{definition}

\noindent
몇 가지를 분명히 하자. 공리 (1)에서는 $I = \{i\}$가 한원소집합이고
사상 $U_i \to U$가 (1)에서 주어진 사상 $V \to U$와 같은
$\mathcal{C}$의 덮개 $\{U_i \to U\}_{i \in I}$가 존재할 것을 요구한다.
이하에서는 지표집합이 한원소집합인 공역이 고정된 사상족을 흔히
$\{V \to U\}$로 나타낸다. 공리 (3)에서는 $i \in I$ 각각에 대한 섬유곱
$U_i \times_U V$의 어떤 선택에 관하여 그 덮개가 존재할 것을 요구한다.

\begin{remark}
\label{sites-remark-no-big-sites}
(집합론적 문제에 관하여---처음 읽을 때는 건너뛰어도 된다.)
사이트를 도입하는 주된 이유는 사이트 위의 층 범주를 연구하기 위해서이다.
이 범주는 대수기하학에서 매우 중요한 역할을 해 온 위상공간 위 층 범주의
일반화이기 때문이다. ``클래스들의 클래스'' 같은 것을 생각하지 않도록,
범주와 $2$-범주에 대해서와 달리 사이트가 ``큰'' 범주인 것은 허용하지 않는다.

\medskip\noindent
$\mathcal{C}$가 범주이고 $\text{Cov}(\mathcal{C})$가 위의 (1), (2), (3)을
만족하는 덮개들의 고유 모임이라고 하자. 이것도 사이트로 허용하지 않는데,
주된 이유는 앞으로 덮개들을 지표로 하는 극한을 취할 것이기 때문이다.
그러나 $\text{Cov}(\mathcal{C})$를 같은 층 범주를 주는 덮개들의 집합이나
조금 다른 구조로 바꾸는 자연스러운 방법이 여러 가지 있다. 예를 들면 다음과 같다.
\begin{enumerate}
\item Sets, Section \ref{sets-section-coverings-site}에서 같은 층 범주를 주는
적당한 덮개들의 집합을 고르는 방법을 보인다.
\item 연관된 위상을 취할 수도 있다(Definition
\ref{sites-definition-topology-associated-site}을 보라). 이렇게 얻은
$\mathcal{C}$ 위의 위상은 같은 층 범주를 갖는다. 두 위상이 같은 층 범주를
가질 필요충분조건은 두 위상이 같은 것이며, Theorem
\ref{sites-theorem-topology-and-topos}을 보라. 범주 위의 위상은 각 대상 위의
시브들을 선택하여 주어진다. $\mathcal{C}$ 위의 가능한 모든 시브들의 모임,
나아가 가능한 모든 위상들의 모임도 집합이다.
\item 사이트의 개념을 조금 바꿀 수도 있다. 아래 Remark
\ref{sites-remark-shrink-coverings}을 보라. 그러면 표준적인 덮개들의 집합을 얻는다.
\end{enumerate}
이 해법들은 각각 작은 단점이 있다. 첫 번째 방법에서는 나중의 구성들이
덮개집합의 선택에 의존하지 않는지 확인해야 한다. 두 번째 방법에서는 위상을
배우고 사이트에 관한 많은 논증을 다시 해야 한다. 세 번째 방법에 관해서는
Remark \ref{sites-remark-shrink-coverings}의 마지막 문장을 보라.

\medskip\noindent
우리는 위의 Definition \ref{sites-definition-site}에서 말한 사이트를 사용할 것이다.
위와 같은 덮개들의 고유 모임을 갖는 범주 $\mathcal{C}$가 주어지면,
Sets, Lemma \ref{sets-lemma-coverings-site}을 사용하여 이를 사이트를 이루는
덮개들의 집합으로 바꾼다. 아래 Lemma
\ref{sites-lemma-choice-set-coverings-immaterial}에서 이렇게 얻은 층 범주(토포스)가
이 선택에 무관함을 보인다. 원한다면 독자는 이 문제를 다루기 위해 나머지 두
전략 가운데 하나를 사용해도 된다.
\end{remark}

\begin{example}
\label{sites-example-site-topological}
$X$를 위상공간이라 하자. $X_{Zar}$를 다음 범주라 하자. 그 대상은 $X$의 모든
열린집합 $U$이고 사상은 포함사상뿐이다. 즉 $X_{Zar}$의 임의의 두 대상 사이에는
사상이 많아야 하나뿐이다. 이제
$\{U_i \to U\}_{i \in I}\in \text{Cov}(X_{Zar})$인 것을
$\bigcup U_i = U$인 것과 동치가 되도록 정의한다. 위의 조건 (1)과 (2)는
명백하며, $X_{Zar}$에서 $U \times V = U \cap V$임을 알면 (3)도 명백하다.
특히 공집합은 $X_{Zar}$의 원소여야 한다. 그렇지 않으면 일반적으로 이 구성이
성립하지 않는다. 더욱이 뒤에서 보듯이, {\it 공집합의 빈 덮개를 덮개로 허용하는
것}도 마찬가지로 중요하다! Sets, Lemma
\ref{sets-lemma-coverings-site}에서처럼 적당한 덮개들의 집합
$\text{Cov}(X_{Zar})_{\kappa, \alpha}$를 골라 $X_{Zar}$를 사이트로 만든다.
% P01-ERRATA: P01-E0058 -- the frozen source has the ungrammatical plural phrase "a presheaves".
사이트 $X_{Zar}$ 위의 준층과 층(아래에서 정의한다)은 Sheaves, Section
\ref{sheaves-section-introduction}에서 정의한 위상공간 위의 통상적인 준층과
층의 개념과 정확히 일치한다.
\end{example}

\begin{example}
\label{sites-example-site-on-group}
$G$를 군이라 하자. 대상이 왼쪽 $G$-작용을 갖는 집합 $X$이고 사상이
$G$-등변사상인 범주 $G\textit{-Sets}$를 생각하자. 중요한 예는 ${}_GG$이다.
이는 밑집합이 $G$이고 작용이 왼쪽 곱셈으로 주어진 $G$-집합이다. 이 범주는
섬유곱을 갖는다. Categories, Section
\ref{categories-section-example-fibre-products}을 보라.
$\bigcup_{i\in I} \varphi_i(U_i) = U$일 때
$\{\varphi_i : U_i \to U\}_{i\in I}$를 덮개라고 선언한다.
이는 $G\textit{-Sets}$ 위의 덮개들의 클래스를 주며 Definition
\ref{sites-definition-site}의 조건 (1), (2), (3)을 만족함을 쉽게 알 수 있다.
밑범주의 대상들의 모임과 덮개들의 모임이 모두 고유 모임을 이루므로
그 결과는 사이트가 아니다.
% P01-ERRATA: P01-E0059 -- the frozen source says "replace by ... by".
먼저 Sets, Lemma \ref{sets-lemma-sets-with-group-action}에서처럼
$G\textit{-Sets}$를 충만부분범주 $G\textit{-Sets}_\alpha$로 바꾼다.
그 뒤 Sets, Lemma \ref{sets-lemma-coverings-site}에서처럼 덮개들의 모임을
적절히 제한하여 얻은 사이트
$(G\textit{-Sets}_\alpha,
\text{Cov}_{\kappa, \alpha'}(G\textit{-Sets}_\alpha))$
를 $\mathcal{T}_G$로 나타낸다.

\medskip\noindent
특히 군 $G$가 가산이면 $\mathcal{T}_G$를 가산 $G$-집합들의 범주로 삼고,
$I$가 가산인 공동 전사 사상족
$\{\varphi_i : U_i \to U\}_{i \in I}$를 덮개로 삼을 수 있다.
\end{example}

\begin{example}
\label{sites-example-indiscrete}
$\mathcal{C}$를 범주라 하자. $U$의 덮개가
$\{f : V \to U \mid f\text{ is an isomorphism}\}$인 사이트로 이 범주를
만드는 표준적인 방법이 있다. 이 사이트 위의 층은 $\mathcal{C}$ 위의 준층이다.
% P01-ERRATA: P01-E0060 -- the frozen source says "This corresponding topology".
대응하는 위상을 {\it 카오스 위상} 또는 {\it 비이산 위상}이라고 한다.
\end{example}















\section{층}
\label{sites-section-sheaves}

\noindent
$\mathcal{C}$를 사이트라 하자. 한 범주에 값을 갖는 층의 개념을 도입하기 전에,
집합의 준층이 층이라는 것이 무엇인지 설명하자. $\mathcal{F}$를
$\mathcal{C}$ 위의 집합의 준층이라 하고
$\{U_i \to U\}_{i\in I}$를 $\text{Cov}(\mathcal{C})$의 원소라 하자.
가정에 따라 모든 섬유곱 $U_i \times_U U_j$가 $\mathcal{C}$ 안에 존재한다.
다음 두 자연스러운 사상이 있다.
$$
\xymatrix{
\prod\nolimits_{i\in I}
\mathcal{F}(U_i)
\ar@<1ex>[r]^-{\text{pr}_0^*} \ar@<-1ex>[r]_-{\text{pr}_1^*}
&
\prod\nolimits_{(i_0, i_1) \in I \times I}
\mathcal{F}(U_{i_0} \times_U U_{i_1})
}
$$
표시된 식과 같이 이들을 $\text{pr}^*_i$, $i = 0, 1$로 나타낼 것이다.
실제로 왼쪽의 한 원소는 사상족 $(s_i)_{i\in I}$에 대응하며, 각 $s_i$는
$U_i$ 위의 $\mathcal{F}$의 단면이다. 각 쌍
$(i_0, i_1) \in I \times I$에 대하여 사영사상
$$
\text{pr}^{(i_0, i_1)}_{i_0} :
U_{i_0} \times_U U_{i_1}
\longrightarrow
U_{i_0}
\text{ and }
\text{pr}^{(i_0, i_1)}_{i_1} :
U_{i_0} \times_U U_{i_1}
\longrightarrow
U_{i_1}.
$$
이 있다. 따라서 단면 $s_{i_0}$를 이들 중 첫 번째 사상을 따라 당길 수도 있고,
단면 $s_{i_1}$을 두 번째 사상을 따라 당길 수도 있다. 구체적으로 위에서 말한
사상들은
$$
\text{pr}_0^* :
(s_i)_{i\in I}
\longmapsto
\Big(
\text{pr}^{(i_0, i_1), *}_{i_0}(s_{i_0})
\Big)_{(i_0, i_1) \in I \times I}
$$
및
$$
\text{pr}_1^* :
(s_i)_{i\in I}
\longmapsto
\Big(
\text{pr}^{(i_0, i_1), *}_{i_1}(s_{i_1})
\Big)_{(i_0, i_1) \in I \times I}.
$$
이다. 마지막으로 자연스러운 사상
$$
\mathcal{F}(U)
\longrightarrow
\prod\nolimits_{i\in I}
\mathcal{F}(U_i), \quad
s
\longmapsto
(s|_{U_i})_{i \in I}
$$
을 생각하자. 여기서 $s|_{U_i}$는 $U_i \to U$를 따라 $s$를 당긴 것을
나타낸다. $\mathcal{F}$의 함자성과 섬유곱 도식들의 가환성으로부터
$\text{pr}_0^*( (s|_{U_i})_{i \in I} ) =
\text{pr}_1^*( (s|_{U_i})_{i \in I} )$임은 명백하다.

\begin{definition}
\label{sites-definition-sheaf-sets}
$\mathcal{C}$를 사이트라 하고 $\mathcal{F}$를 $\mathcal{C}$ 위의 집합의
준층이라 하자. 모든 덮개
$\{U_i \to U\}_{i \in I} \in \text{Cov}(\mathcal{C})$에 대하여 도식
\begin{equation}
\label{sites-equation-sheaf-condition}
\xymatrix{
\mathcal{F}(U) \ar[r]
&
\prod\nolimits_{i\in I}
\mathcal{F}(U_i)
\ar@<1ex>[r]^-{\text{pr}_0^*} \ar@<-1ex>[r]_-{\text{pr}_1^*}
&
\prod\nolimits_{(i_0, i_1) \in I \times I}
\mathcal{F}(U_{i_0} \times_U U_{i_1})
}
\end{equation}
에서 첫 번째 화살표가 $\text{pr}_0^*$와 $\text{pr}_1^*$의 등화자로
나타나면 $\mathcal{F}$를 {\it 층}이라고 한다.
\end{definition}

\noindent
느슨하게 말하면 다음을 뜻한다. 단면들 $s_i \in \mathcal{F}(U_i)$가 주어져
모든 쌍 $(i, j) \in I \times I$에 대하여
$\mathcal{F}(U_i \times_U U_j)$ 안에서
$$
s_i|_{U_i \times_U U_j} = s_j|_{U_i \times_U U_j}
$$
이면 $s_i = s|_{U_i}$를 만족하는 $s \in \mathcal{F}(U)$가 유일하게 존재한다.

\begin{remark}
\label{sites-remark-sheaf-condition-empty-covering}
덮개 $\{U_i \to U\}_{i \in I}$가 빈 사상족(즉 $I = \emptyset$)이면,
층 조건은 $\mathcal{F}(U) = \{*\}$가 한원소집합임을 뜻한다. 이는
(\ref{sites-equation-sheaf-condition})의 두 번째와 세 번째 집합이 집합의 범주에서
빈 곱이고, 이들이 집합의 범주의 끝 대상, 따라서 한원소집합이기 때문이다.
\end{remark}

\begin{example}
\label{sites-example-sheaves-topological}
$X$를 위상공간이라 하자. $X_{Zar}$를 Example
\ref{sites-example-site-topological}에서 구성한 사이트라 하자.
$X_{Zar}$ 위의 층 개념은 Sheaves, Definition
\ref{sheaves-definition-sheaf}에서 도입한 $X$ 위의 층 개념과 일치한다.
\end{example}

\begin{example}
\label{sites-example-topological-wrong}
$X$를 위상공간이라 하자. Example \ref{sites-example-site-topological}의 사이트
$X_{Zar}$와 같되 공집합의 빈 덮개를 허용하지 않는 사이트 $X'_{Zar}$를
생각하자. 달리 말하면 덮개 $\{\emptyset \to \emptyset\}$는 허용하지만
지표집합이 빈 덮개는 허용하지 않는다. 이것도 사이트를 정의함은 쉽게 보인다.
그러나 $X'_{Zar}$ 위의 층들은 $X_{Zar}$ 위의 층들과 다르다고 주장한다.
극단적인 예로 $X = \{p\}$가 한원소집합인 경우를 생각하자.
그러면 $X'_{Zar}$의 대상들은 $\emptyset, X$이고, $\emptyset$의 모든 덮개는
$\{\emptyset \to \emptyset\}$로 세분되며 $X$의 모든 덮개는
$\{X \to X\}$로 세분된다. 명백히 이 위의 한 층은 집합
$\mathcal{F}(\emptyset)$의 임의의 선택, 집합 $\mathcal{F}(X)$의 임의의 선택,
그리고 제한 사상 $\mathcal{F}(X) \to \mathcal{F}(\emptyset)$의 임의의 선택으로
주어진다. 따라서 $X'_{Zar}$ 위의 층들은 열린집합들이
$\{\emptyset, \{\eta\}, \{p, \eta\}\}$인 두 점 공간 $\{\eta, p\}$ 위의
통상적인 층들과 같다. 일반적으로 $X'_{Zar}$ 위의 층들은
$X \amalg \{\eta\}$ 위의 층들과 같다. 여기서 열린집합은 공집합과,
열린 $U \subset X$에 대하여 $U \cup \{\eta\}$ 꼴인 모든 집합이다.
\end{example}


\begin{definition}
\label{sites-definition-category-sheaves-sets}
집합의 층들이 이루는 범주 {\it $\Sh(\mathcal{C})$}는 대상들이 집합의 층인
범주 $\textit{PSh}(\mathcal{C})$의 충만부분범주이다.
\end{definition}

\noindent
$\mathcal{A}$를 범주라 하자. 덮개 사상족들의 지표집합으로 나타나는 임의의
$I$에 대하여 $I$ 및 $I \times I$로 지표화된 곱들이 $\mathcal{A}$ 안에
존재하면 위의 Definition \ref{sites-definition-sheaf-sets}은 의미를 가지며,
$\mathcal{A}$에 값을 갖는 $\mathcal{C}$ 위의 층이라는 개념을 정의한다.
$\mathcal{A}$ 안의 도식
$$
\xymatrix{
\mathcal{F}(U) \ar[r]
&
\prod\nolimits_{i\in I}
\mathcal{F}(U_i)
\ar@<1ex>[r]^-{\text{pr}_0^*} \ar@<-1ex>[r]_-{\text{pr}_1^*}
&
\prod\nolimits_{(i_0, i_1) \in I \times I}
\mathcal{F}(U_{i_0} \times_U U_{i_1})
}
$$
이 등화자 도식일 필요충분조건은 $\mathcal{A}$의 모든 대상 $X$에 대하여
집합들의 도식
$$
\xymatrix{
\Mor_\mathcal{A}(X, \mathcal{F}(U)) \ar[r]
&
\prod
\Mor_\mathcal{A}(X, \mathcal{F}(U_i))
\ar@<1ex>[r]^-{\text{pr}_0^*} \ar@<-1ex>[r]_-{\text{pr}_1^*}
&
\prod
\Mor_\mathcal{A}(X, \mathcal{F}(U_{i_0} \times_U U_{i_1}))
}
$$
이 등화자 도식인 것이다.

\medskip\noindent
$\mathcal{A}$가 임의의 범주라고 하자. $\mathcal{F}$를 $\mathcal{A}$에 값을
갖는 준층이라 하자. 임의의 대상 $X\in \Ob(\mathcal{A})$를 택하면 규칙
$$
\mathcal{F}_X(U) = \Mor_\mathcal{A}(X, \mathcal{F}(U)).
$$
으로 정의되는 집합의 준층 $\mathcal{F}_X$를 얻는다. 위의 논의에 따라 모든
준층 $\mathcal{F}_X$가 집합의 층일 것을 요구하면 좋은 정의를 얻게 된다.

\begin{definition}
\label{sites-definition-sheaf}
$\mathcal{C}$를 사이트, $\mathcal{A}$를 범주라 하고, $\mathcal{F}$를
$\mathcal{A}$에 값을 갖는 $\mathcal{C}$ 위의 준층이라 하자.
$\mathcal{A}$의 모든 대상 $X$에 대하여 위에서 정의한 집합의 준층
$\mathcal{F}_X$가 층이면 $\mathcal{F}$를 {\it 층}이라고 한다.
\end{definition}







\section{공역이 고정된 사상족}
\label{sites-section-refinements}

\noindent
이 절에서는 공역이 같은 사상족과 관련된 몇 가지 개념을 도입한다.

\begin{definition}
\label{sites-definition-morphism-coverings}
$\mathcal{C}$를 범주라 하자.
$\mathcal{U} = \{U_i \to U\}_{i\in I}$를 공역이 고정된
$\mathcal{C}$의 사상족이라 하자.
$\mathcal{V} = \{V_j \to V\}_{j\in J}$를 또 다른 그러한 사상족이라 하자.
\begin{enumerate}
\item
{\it $\mathcal{C}$의 공역이 고정된 사상족의 $\mathcal{U}$에서
$\mathcal{V}$로 가는 사상}, 또는 간단히
{\it $\mathcal{U}$에서 $\mathcal{V}$로 가는 사상}은 사상 $U \to V$,
집합의 사상 $\alpha : I \to J$, 그리고 각 $i\in I$에 대한 사상
$U_i \to V_{\alpha(i)}$로 주어지며, 이때 도식
$$
\xymatrix{
U_i \ar[r] \ar[d]
&
V_{\alpha(i)} \ar[d]
\\
U \ar[r]
&
V
}
$$
은 가환이다.
\item 특히 $U = V$이고 $U \to V$가 항등사상인 경우에는
$\mathcal{U}$를 사상족 $\mathcal{V}$의 {\it 세분}이라고 한다.
\end{enumerate}
\end{definition}

\noindent
자명하지만 중요한 사실은
$\mathcal{V} = \{V_j \to V\}_{j \in J}$가 {\it 공집합인 사상족},
즉 $J = \emptyset$이면 $I \not = \emptyset$인 어떤 사상족
$\mathcal{U} = \{U_i \to V\}_{i \in I}$도 $\mathcal{V}$를
세분할 수 없다는 것이다!

\begin{definition}
\label{sites-definition-combinatorial-tautological}
$\mathcal{C}$를 범주라 하자.
$\mathcal{U} = \{\varphi_i : U_i \to U\}_{i\in I}$와
$\mathcal{V} = \{\psi_j : V_j \to U\}_{j\in J}$를 공역이 고정된
두 사상족이라 하자.
\begin{enumerate}
\item $\varphi_i = \psi_{\alpha(i)}$이고
$\psi_j = \varphi_{\beta(j)}$이 되게 하는 사상
$\alpha : I \to J$와 $\beta : J\to I$가 존재하면
$\mathcal{U}$와 $\mathcal{V}$가 {\it 조합적으로 동치}라고 한다.
\item 사상 $\alpha : I \to J$와 $\beta : J\to I$가 존재하고,
모든 $i\in I$와 $j \in J$에 대하여 가환도식
$$
\xymatrix{
U_i \ar[rd] \ar[rr] & &
V_{\alpha(i)} \ar[ld] & &
V_j \ar[rd] \ar[rr] & &
U_{\beta(j)} \ar[ld] \\
&
U & & & &
U &
}
$$
에서 수평 화살표들이 동형사상이면 $\mathcal{U}$와 $\mathcal{V}$가
{\it 자명하게 동치}라고 한다.
\end{enumerate}
\end{definition}

\begin{lemma}
\label{sites-lemma-tautological-combinatorial}
$\mathcal{C}$를 범주라 하자.
$\mathcal{U} = \{\varphi_i : U_i \to U\}_{i\in I}$와
$\mathcal{V} = \{\psi_j : V_j \to U\}_{j\in J}$를 동일하게 고정된
공역을 갖는 두 사상족이라 하자.
\begin{enumerate}
\item $\mathcal{U}$와 $\mathcal{V}$가 조합적으로 동치이면
자명하게 동치이다.
\item $\mathcal{U}$와 $\mathcal{V}$가 자명하게 동치이면
$\mathcal{U}$는 $\mathcal{V}$의 세분이고 $\mathcal{V}$는
$\mathcal{U}$의 세분이다.
\item ``조합적으로 동치임''이라는 관계는 공역이 고정된 모든
사상족 위의 동치관계이다.
\item ``자명하게 동치임''이라는 관계는 공역이 고정된 모든
사상족 위의 동치관계이다.
\item ``$\mathcal{U}$가 $\mathcal{V}$를 세분하고 $\mathcal{V}$가
$\mathcal{U}$를 세분함''이라는 관계는 공역이 고정된 모든
사상족 위의 동치관계이다.
\end{enumerate}
\end{lemma}

\begin{proof}
생략한다.
\end{proof}


\noindent
다음 보조정리에서 범주 $\mathcal{C}$, $\mathcal{C}$ 위의 준층
$\mathcal{F}$, 그리고 모든 올곱 $U_i \times_U U_{i'}$가 존재하는 사상족
$\mathcal{U} = \{U_i \to U\}_{i\in I}$가 주어졌을 때 도식
(\ref{sites-equation-sheaf-condition})이 등화자 도식이면
{\it $\mathcal{U}$에 대한 $\mathcal{F}$의 층 조건}이 성립한다고 한다.

\begin{lemma}
\label{sites-lemma-tautological-same-sheaf}
$\mathcal{C}$를 범주라 하자. $\mathcal{U} =
\{\varphi_i : U_i \to U\}_{i\in I}$와
$\mathcal{V} = \{\psi_j : V_j \to U\}_{j\in J}$를 동일하게 고정된
공역을 갖는 두 사상족이라 하자. 올곱
$U_i \times_U U_{i'}$와 $V_j \times_U V_{j'}$가 존재한다고 가정하자.
$\mathcal{U}$와 $\mathcal{V}$가 자명하게 동치이면 $\mathcal{C}$ 위의
임의의 준층 $\mathcal{F}$에 대하여 $\mathcal{U}$에 대한
$\mathcal{F}$의 층 조건과 $\mathcal{V}$에 대한 $\mathcal{F}$의 층 조건은
동치이다.
\end{lemma}

\begin{proof}
먼저 $\varphi : A \to B$가 범주 $\mathcal{C}$의 동형사상이면
$\varphi^* : \mathcal{F}(B) \to \mathcal{F}(A)$도 동형사상임에 유의하자.
$\beta : J \to I$를 사상이라 하고, 가정에 의해 존재하는 $U$ 위의
동형사상을 $\chi_j : V_j \to U_{\beta(j)}$라 하자.
$\mathcal{V}$에 대한 층 조건이 $\mathcal{U}$에 대한 층 조건을
함의함을 보이자. 모든 $(i, i') \in I \times I$에 대하여
$\mathcal{F}(U_i \times_U U_{i'})$ 안에서
$$
s_i|_{U_i \times_U U_{i'}} = s_{i'}|_{U_i \times_U U_{i'}}
$$
이 되도록 절단 $s_i \in \mathcal{F}(U_i)$들이 주어졌다고 하자.
$s_j = \chi_j^*s_{\beta(j)}$로 정의하자. 임의의
$(j, j') \in J \times J$에 대하여 사상
$\chi_j \times_{\text{id}_U} \chi_{j'} : V_j \times_U V_{j'} \to
U_{\beta(j)} \times_U U_{\beta(j')}$도 동형사상이다.
그러므로 구조의 운반에 의해
$$
s_j|_{V_j \times_U V_{j'}} = s_{j'}|_{V_j \times_U V_{j'}}
$$
임을 알 수 있다. $\mathcal{V}$에 대한 층 조건으로부터 모든
$j \in J$에 대하여 $s|_{V_j} = s_j$가 되는 유일한 $s$가 존재한다.
증명의 첫 번째 관찰에 의해 이는 모든 $i \in \Im(\beta)$에 대하여
$s|_{U_i} = s_i$임도 함의한다. 이제 $i \in I$이고
$i \not \in \Im(\beta)$라고 하자. 그러한 $i$에 대하여 $U$ 위의 동형사상
$U_i \to V_{\alpha(i)} \to U_{\beta(\alpha(i))}$들이 있다. 이들은 사상
$U_i \to U_i \times_U U_{\beta(\alpha(i))}$를 주며, 이는 사영의 절단이다.
$s_i$와 $s_{\beta(\alpha(i))}$가 올곱 위에서 같은 원소로 제한되므로
$s_{\beta(\alpha(i))}$를 $U_i \to U_{\beta(\alpha(i))}$에 따라
당기면 $s_i$가 됨을 알 수 있다. 따라서 원하는 대로
$s_i = s|_{U_i}$이다.
\end{proof}

\begin{lemma}
\label{sites-lemma-compare-separated-presheaf-condition}
$\mathcal{C}$를 범주라 하자. $\mathcal{V} = \{V_j \to U\}_{j \in J} \to
\mathcal{U} = \{U_i \to U\}_{i \in I}$를 $\text{id} : U \to U$,
$\alpha : J \to I$, 그리고 $f_j : V_j \to U_{\alpha(j)}$로 주어지는
$\mathcal{C}$의 공역이 고정된 사상족의 사상이라 하자.
$\mathcal{F}$를 $\mathcal{C}$ 위의 준층이라 하자. 사상
$\mathcal{F}(U) \to \prod_{j \in J} \mathcal{F}(V_j)$가 단사이면
$\mathcal{F}(U) \to \prod_{i \in I} \mathcal{F}(U_i)$도 단사이다.
\end{lemma}

\begin{proof}
생략한다.
\end{proof}

\begin{lemma}
\label{sites-lemma-compare-sheaf-condition}
$\mathcal{C}$를 범주라 하자. $\mathcal{V} = \{V_j \to U\}_{j \in J} \to
\mathcal{U} = \{U_i \to U\}_{i \in I}$를 $\text{id} : U \to U$,
$\alpha : J \to I$, 그리고 $f_j : V_j \to U_{\alpha(j)}$로 주어지는
$\mathcal{C}$의 공역이 고정된 사상족의 사상이라 하자.
$\mathcal{F}$를 $\mathcal{C}$ 위의 준층이라 하자. 다음을 가정하자.
\begin{enumerate}
\item 올곱 $U_i \times_U U_{i'}$, $U_i \times_U V_j$,
$V_j \times_U V_{j'}$가 존재한다.
\item $\mathcal{F}$가 $\mathcal{V}$에 대한 층 조건을 만족한다.
\item 모든 $i \in I$에 대하여 사상
$\mathcal{F}(U_i) \to \prod_{j \in J} \mathcal{F}(V_j \times_U U_i)$가
단사이다.
\end{enumerate}
그러면 $\mathcal{F}$는 $\mathcal{U}$에 대한 층 조건을 만족한다.
\end{lemma}

\begin{proof}
Lemma \ref{sites-lemma-compare-separated-presheaf-condition}에 의해 사상
$\mathcal{F}(U) \to \prod \mathcal{F}(U_i)$는 단사이다. 모든
$i, i' \in I$에 대하여
$s_i|_{U_i \times_U U_{i'}} = s_{i'}|_{U_i \times_U U_{i'}}$가 되도록
$s_i \in \mathcal{F}(U_i)$들이 주어졌다고 하자.
$s_j = f_j^*(s_{\alpha(j)}) \in \mathcal{F}(V_j)$로 놓자.
사상 $f_j$들이 $U$ 위의 사상이므로 사영사상을 통해 $f_j, f_{j'}$와
호환되는 유도사상을 얻는다.
% P01-ERRATA: P01-E0061 -- frozen codomain uses undefined i,i' instead of j,j'.
$f_{jj'} : V_j \times_U V_{j'} \to
U_{\alpha(i)} \times_U U_{\alpha(i')}$
따라서
$$
s_j|_{V_j \times_U V_{j'}}
= f_{jj'}^*(s_{\alpha(j)}|_{U_{\alpha(j)} \times_U U_{\alpha(j')}})
= f_{jj'}^*(s_{\alpha(j')}|_{U_{\alpha(j)} \times_U U_{\alpha(j')}})
= s_{j'}|_{V_j \times_U V_{j'}}
$$
가 모든 $j, j' \in J$에 대하여 성립한다. 그러므로 $\mathcal{V}$에 대한
$\mathcal{F}$의 층 조건에 의해 각 $V_j$ 위에서 $s_j$로 제한되는 절단
$s \in \mathcal{F}(U)$를 얻는다. 임의의 $i \in I$에 대하여 $s$가
$U_i$ 위에서 $s_i$로 제한됨을 보이면 충분하다. $\mathcal{F}$가 (3)을
만족하므로 모든 $j \in J$에 대하여 $s$와 $s_i$가
$U_i \times_U V_j$ 위에서 같은 원소로 제한됨을 보이면 충분하다.
이를 위해
$$
s|_{U_i \times_U V_j} = s_j|_{U_i \times_U V_j} =
(\text{id} \times f_j)^*s_{\alpha(j)}|_{U_i \times_U U_{\alpha(j)}} =
(\text{id} \times f_j)^*s_i|_{U_i \times_U U_{\alpha(j)}} =
s_i|_{U_i \times_U V_j}
$$
를 사용하면 된다.
\end{proof}

\begin{lemma}
\label{sites-lemma-refine-same-topology}
$\mathcal{C}$를 범주라 하자. $\text{Cov}_i$, $i = 1, 2$를 각각
$\mathcal{C}$ 위에 사이트 구조를 정의하는, 공역이 고정된 사상족들의
두 집합이라 하자.
\begin{enumerate}
\item 모든 $\mathcal{U} \in \text{Cov}_1$가 어떤
$\mathcal{V} \in \text{Cov}_2$와 자명하게 동치이면
$\Sh(\mathcal{C}, \text{Cov}_2) \subset
\Sh(\mathcal{C}, \text{Cov}_1)$이다.
또한 모든 $\mathcal{U} \in \text{Cov}_2$가 어떤
$\mathcal{V} \in \text{Cov}_1$와 자명하게 동치이면
% P01-ERRATA: P01-E0062 -- frozen English has singular/plural disagreement.
두 층의 범주는 같다.
\item 각 $\mathcal{U} \in \text{Cov}_1$에 대하여
$\mathcal{V} \in \text{Cov}_2$가 존재하고 $\mathcal{V}$가
$\mathcal{U}$를 세분한다고 가정하자. 이 경우
$\Sh(\mathcal{C}, \text{Cov}_2) \subset
\Sh(\mathcal{C}, \text{Cov}_1)$이다.
또한 모든 $\mathcal{U} \in \text{Cov}_2$에 대하여
$\mathcal{V} \in \text{Cov}_1$가 존재하고 $\mathcal{V}$가
$\mathcal{U}$를 세분하면 두 층의 범주는 같다.
\end{enumerate}
\end{lemma}

\begin{proof}
(1)은 Lemma \ref{sites-lemma-tautological-same-sheaf}와 정의에서 바로 따른다.

\medskip\noindent
(2)를 증명하자. $\mathcal{F}$를 사이트
$(\mathcal{C}, \text{Cov}_2)$의 집합의 층이라 하자.
$\mathcal{U} \in \text{Cov}_1$라 하고
$\mathcal{U} = \{U_i \to U\}_{i \in I}$라고 쓰자. 가정에 의해
$\mathcal{U}$의 세분 $\mathcal{V} \in \text{Cov}_2$를 택할 수 있다.
% P01-ERRATA: P01-E1330 -- frozen English omits an article and finite verb in “and refinement given”.
$\mathcal{V} = \{V_j \to U\}_{j \in J}$라고 쓰고, 그 세분은
$\alpha : J \to I$와 $f_j : V_j \to U_{\alpha(j)}$로 주어진다고 하자.
$\mathcal{F}$는 $\mathcal{V}$에 대한 층 조건을 만족하고, 덮개
$\{V_j \times_U U_i \to U_i\}_{j \in J}$들도 $\text{Cov}_2$에 속하므로
이 덮개들에 대한 층 조건도 만족한다. 따라서 Lemma
\ref{sites-lemma-compare-sheaf-condition}에 의해 $\mathcal{F}$는
$\mathcal{U}$에 대한 층 조건을 만족한다.
\end{proof}

\begin{lemma}
\label{sites-lemma-choice-set-coverings-immaterial}
$\mathcal{C}$를 범주라 하자. $\text{Cov}(\mathcal{C})$를 Definition
\ref{sites-definition-site}의 조건 (1), (2), (3)을 만족하는 덮개들의
진클래스라 하자. $\text{Cov}_1, \text{Cov}_2 \subset
\text{Cov}(\mathcal{C})$를 $\text{Cov}(\mathcal{C})$의 두 부분집합으로서
$\mathcal{C}$에 사이트 구조를 부여하는 것들이라 하자. 모든 덮개
$\mathcal{U} \in \text{Cov}(\mathcal{C})$가 $\text{Cov}_1$의 어떤 덮개와
조합적으로 동치이고 $\text{Cov}_2$의 어떤 덮개와도 조합적으로 동치이면
$\Sh(\mathcal{C}, \text{Cov}_1) =
\Sh(\mathcal{C}, \text{Cov}_2)$이다.
\end{lemma}

\begin{proof}
가정에 의해 모든 덮개
$\mathcal{U} \in \text{Cov}_1 \subset \text{Cov}(\mathcal{C})$가
$\text{Cov}_2$의 한 원소와 조합적으로 동치이고, $\text{Cov}_1$과
$\text{Cov}_2$의 역할을 바꾸어도 마찬가지이므로, 위의 Lemmas
\ref{sites-lemma-refine-same-topology}와
\ref{sites-lemma-tautological-combinatorial}에서 바로 따른다.
\end{proof}

\section{G-집합의 예}
\label{sites-section-example-sheaf-G-sets}

\noindent
예로서 Example \ref{sites-example-site-on-group}의 사이트 $\mathcal{T}_G$를
생각하자. $\mathcal{T}_G$ 위의 층들의 범주를 기술할 것이다. 그 답은
$\mathcal{T}_G$를 정의할 때 한 선택들과 무관함이 드러날 것이다. 실제로
증명에서는 사이트 $\mathcal{T}_G$의 다음 성질들만 필요하다.
\begin{enumerate}
\item[(a)] $\mathcal{T}_G$는 $G\textit{-Sets}$의 충만부분범주이다.
\item[(b)] $\mathcal{T}_G$는 $G$-집합 ${}_GG$를 포함한다.
\item[(c)] $\mathcal{T}_G$에는 올곱들이 존재하며, 이들은
$G\textit{-Sets}$에서의 올곱들과 같다.
\item[(d)] $U \in \Ob(\mathcal{T}_G)$와 $G$-안정 부분집합
$O \subset U$가 주어지면 $O$와 동형인 $\mathcal{T}_G$의 대상이 존재한다.
% P01-ERRATA: P01-E1264 -- candidate concerns the frozen reduced-relative wording.
\item[(e)] $U, U_i \in \Ob(\mathcal{T}_G)$인 임의의 전사 사상족
$\{U_i \to U\}_{i \in I}$는 $\mathcal{T}_G$의 한 덮개와
조합적으로 동치이다.
\end{enumerate}
이 성질들은 Sets, Lemmas \ref{sets-lemma-what-is-in-it-G-sets}와
\ref{sets-lemma-coverings-site}에 의해 성립한다.

\medskip\noindent
사상
$$
\Hom_G({}_GG, {}_GG)
\longrightarrow
G^{opp},
\varphi
\longmapsto
\varphi(1)
$$
이 군의 동형사상임에 유의하자. 그 역은 $g \in G$를 사상
$R_g : s \mapsto sg$(즉 오른쪽 곱셈)로 보낸다.
$R_{g_1g_2} = R_{g_2} \circ R_{g_1}$이므로 반대군이 필요하다.

\medskip\noindent
이로부터 $\mathcal{T}_G$ 위의 모든 준층 $\mathcal{F}$에 대하여 값
$\mathcal{F}({}_GG)$는 다음과 같이 $G$-집합 구조를 물려받는다.
$g \in G$와 $s \in \mathcal{F}({}_GG)$에 대하여 $g \cdot s$를
$\mathcal{F}(R_g)(s)$로 정의한다. 이것은 왼쪽 작용인데,
$$
(g_1g_2) \cdot s  = \mathcal{F}(R_{g_1g_2})(s) =
\mathcal{F}(R_{g_2}\circ R_{g_1})(s) =
\mathcal{F}(R_{g_1})( \mathcal{F}(R_{g_2})(s)) =
g_1 \cdot (g_2 \cdot s).
$$
이기 때문이다. 여기서는 $\mathcal{F}$가 반변임을 사용했다.
$\mathcal{F} \to \mathcal{G}$가 $\mathcal{T}_G$ 위의 집합의 준층들의
사상이면 방금 정의한 $G$-작용들과 호환되는 사상
$\mathcal{F}({}_GG) \to \mathcal{G}({}_GG)$를 얻음에 유의하자.
종합하면 함자
$$
\textit{PSh}(\mathcal{T}_G)
\longrightarrow
G\textit{-Sets}, \quad
\mathcal{F}
\longmapsto
\mathcal{F}({}_GG).
$$
를 구성했다. 이 구성이 표현가능 준층 $h_U$를 $U$와 정준적으로 동형인
대상으로 보낸다는 좋은 성질의 확인은 독자에게 맡긴다. 공식으로 쓰면
$h_U({}_GG) = \Hom_G({}_GG, U) \cong U$이다.

\medskip\noindent
$S$가 $G$-집합이라고 하자. 공식\footnote{이것은 $S$로 정의되는 표현가능
준층처럼 보일 수 있다. 그러나 충분히 큰 $G$-집합들의 집합이 되도록 선택한
$\mathcal{T}_G$의 대상이 $S$가 아닐 수도 있으므로 반드시 그렇지는 않다.}
$$
\mathcal{F}_S(U)
=
\Mor_{G\textit{-Sets}}(U, S).
$$
으로 준층 $\mathcal{F}_S$를 정의한다. 이는 명백히 준층이다. 한편
$\{U_i \to U\}_{i\in I}$가 $\mathcal{T}_G$의 덮개라고 하자. 그러면
$\coprod_i U_i \to U$는 전사이다. 따라서 사상
$$
\mathcal{F}_S(U)
=
\Mor_{G\textit{-Sets}}(U, S)
\longrightarrow
\prod \mathcal{F}_S(U_i)
=
\prod \Mor_{G\textit{-Sets}}(U_i, S)
$$
이 단사임은 명백하다. 또한 모든 도식
$$
\xymatrix{
U_i \times_U U_j \ar[d] \ar[r]
&
U_j \ar[d]^{s_j}
\\
U_i \ar[r]^{s_i}
&
S
}
$$
이 가환이 되도록 $G$-등변사상들의 족 $s_i : U_i \to S$가 주어지면,
$s_i$가 합성 $U_i \to U \to S$가 되게 하는 유일한 $G$-등변사상
$s : U \to S$가 존재한다. 실제로 $U_i \to U$에 의해 $u$로 가는 어떤
$u_i \in U_i$가 존재하는 임의의 지표 $i\in I$를 택하여
$s(u) = s_i(u_i)$로 정의한다. 위 도식들의 가환성은 정확히 이 구성이
잘 정의됨을 뜻한다. 종합하면 함자
$$
G\textit{-Sets}
\longrightarrow
\Sh(\mathcal{T}_G), \quad
S
\longmapsto
\mathcal{F}_S
.
$$
를 구성했다.

\medskip\noindent
이제 범주들과 함자들의 다음 도식을 얻는다.
$$
\xymatrix{
\textit{PSh}(\mathcal{T}_G) \ar[rr]^{\mathcal{F} \mapsto \mathcal{F}({}_GG)}
&
&
G\textit{-Sets} \ar[ld]_{S \mapsto \mathcal{F}_S}
\\
&
\Sh(\mathcal{T}_G) \ar[lu]
&
}
$$
정의에서 바로
$\mathcal{F}_S({}_GG) = \Mor_G({}_GG, S) = S$임을 알 수 있으며,
마지막 등식은 $1$에서 값을 취하여 얻는다. 이것으로 다음 명제는 거의
증명된다.

\begin{proposition}
\label{sites-proposition-sheaves-on-group}
함자 $\mathcal{F} \mapsto \mathcal{F}({}_GG)$와
$S \mapsto \mathcal{F}_S$는 $\Sh(\mathcal{T}_G)$와
$G\textit{-Sets}$ 사이의 서로 준역인 동치를 정의한다.
\end{proposition}

\begin{proof}
두 함자를 한 순서로 합성하면 항등함자와 동형임은 이미 보였다.
다른 순서에서는 임의의 층 $\mathcal{H}$에 대하여 층들의 자연스러운 사상
$$
can :
\mathcal{H}
\longrightarrow
\mathcal{F}_{\mathcal{H}({}_GG)}.
$$
이 존재한다. 실제로 $\mathcal{T}_G$의 임의의 대상 $U$에 대하여
$can_U$를 사상
$$
\begin{matrix}
\mathcal{H}(U)
&
\longrightarrow
&
\mathcal{F}_{\mathcal{H}({}_GG)}(U)
=
\Mor_G(U, \mathcal{H}({}_GG))
\\
s
&
\longmapsto
&
(u \mapsto \alpha_u^*s).
\end{matrix}
$$
으로 둔다. 여기서 $\alpha_u : {}_GG \to U$는 사상
$\alpha_u(g) = gu$이고 $\alpha_u^* : \mathcal{H}(U)
\to \mathcal{H}({}_GG)$는 당김사상이다. 자명하지만 혼동하기 쉬운 확인을
통해 이것이 실제로 준층들의 사상임을 알 수 있다. $can$이 동형사상임을
보여야 한다. 모든 $U \in \Ob(\mathcal{T}_G)$에 대하여 $can_U$가
동형사상임을 보이겠다. $U = {}_GG$인 중요하지만 쉬운 경우는 독자에게
맡긴다. $\mathcal{T}_G$의 일반적인 대상 $U$는 $G$-궤도들의 서로소 합집합
$U = \coprod_{i\in I} O_i$이다. 사상족
$\{O_i \to U\}_{i \in I}$는 $\mathcal{T}_G$의 한 덮개와 자명하게 동치이다
(이 절의 시작에서 열거한 $\mathcal{T}_G$의 성질들에 의한다). 따라서 Lemma
\ref{sites-lemma-tautological-same-sheaf}에 의해 층 $\mathcal{H}$는
$\{O_i \to U\}_{i \in I}$에 대한 층 성질을 만족한다. 이 덮개에 대한 층
성질로부터 $\mathcal{H}(U) = \prod_i \mathcal{H}(O_i)$를 얻는다.
따라서 $U$가 하나의 $G$-궤도로 이루어진 경우에 $can_U$가 동형사상임을
보이면 충분하다. $u \in U$라 하고 $H \subset G$를 그 안정자군이라 하자.
명백히 $\Mor_G(U, \mathcal{H}({}_GG)) = \mathcal{H}({}_GG)^H$는
$H$-불변 원소들의 부분집합과 같다. 한편 $g \mapsto gu$로 주어지는 덮개
$\{{}_GG \to U\}$를 생각하자(역시 $\mathcal{T}_G$의 어떤 덮개와
조합적으로 동치일 뿐이며, 역시 이 차이는 중요하지 않다). 올곱
$({}_GG)\times_U ({}_GG)$는
$\{(g, gh), g\in G, h\in H\} \cong \coprod_{h\in H} {}_GG$와 같다.
따라서 이 덮개에 대한 층 성질은
$$
\xymatrix{
\mathcal{H}(U) \ar[r]
&
\mathcal{H}({}_GG)
\ar@<1ex>[r]^-{\text{pr}_0^*} \ar@<-1ex>[r]_-{\text{pr}_1^*}
&
\prod_{h \in H}
\mathcal{H}({}_GG).
}
$$
이다. 인자 $\mathcal{H}({}_GG)$로 가는 두 사상 $\text{pr}_i^*$는
$h$를 곱하는 것만큼 다르다. 이제 이 사실과 $U = {}_GG$일 때 $can$이
동형사상이라는 사실로부터 결론이 따른다.
\end{proof}

\section{층화}
\label{sites-section-sheafification}

\noindent
층화를 정의하기 위해 덮개의 0차 {\v C}ech 코호몰로지 군과 그 함자적
성질들을 연구한다.
% P01-ERRATA: P01-E0063 -- set-valued H^0 need not be a group.

\medskip\noindent
$\mathcal{F}$를 $\mathcal{C}$ 위의 집합의 준층이라 하고
$\mathcal{U} = \{U_i \to U\}_{i \in I}$를 $\mathcal{C}$의 덮개라 하자.
다음 등화자를 나타내기 위해 $\mathcal{F}(\mathcal{U})$라는 표기를 쓰자.
$$
H^0(\mathcal{U}, \mathcal{F})
=
\{
(s_i)_{i\in I} \in \prod\nolimits_i \mathcal{F}(U_i)
\mid
s_i|_{U_i \times_U U_j} = s_j|_{U_i \times_U U_j}
\ \forall i, j \in I
\}.
$$
나중에 보겠지만 이는 덮개 $\mathcal{U}$에 대한 $U$ 위의
$\mathcal{F}$의 0차 {\v C}ech 코호몰로지이다. 작은 관찰 하나는 모든
사상 $U_i \to U$가 표현가능이기만 하면
$H^0(\mathcal{U}, \mathcal{F})$를 정의할 수 있다는 것이다. 즉
$\mathcal{U}$가 사이트의 덮개일 필요는 없다. 표준적인 사상
$\mathcal{F}(U) \to H^0(\mathcal{U}, \mathcal{F})$가 존재한다.
덮개들의 사상 $\mathcal{U} \to \mathcal{V}$가 가환도식
$$
\xymatrix{
& U_i \ar[rr] & & V_{\alpha(i)} \\
U_i \times_U U_j \ar[rr] \ar[ur] \ar[dr] & &
V_{\alpha(i)} \times_V V_{\alpha(j)} \ar[ur] \ar[dr] & \\
& U_j \ar[rr] & & V_{\alpha(j)}
}.
$$
을 유도함은 명백하다. 다시 이것은 사상
$H^0(\mathcal{V}, \mathcal{F}) \to
H^0(\mathcal{U}, \mathcal{F})$를 만들며, 이는 사상
$\mathcal{F}(V) \to \mathcal{F}(U)$와 호환된다.

\medskip\noindent
구성에 의해 준층 $\mathcal{F}$가 층일 필요충분조건은
$\mathcal{C}$의 모든 덮개 $\mathcal{U}$에 대하여 자연사상
$\mathcal{F}(U) \to H^0(\mathcal{U}, \mathcal{F})$가 전단사인 것이다.
이 개념을 사용하여 층의 극한에 관한 다음 간단한 보조정리를 증명하겠다.

\begin{lemma}
\label{sites-lemma-limit-sheaf}
\begin{slogan}
층들의 도표의 극한은 존재하며 준층으로서의 극한과 일치한다
\end{slogan}
$\mathcal{F} : \mathcal{I} \to \Sh(\mathcal{C})$를 도표라 하자.
그러면 $\lim_\mathcal{I} \mathcal{F}$는 존재하고 준층의 범주에서의
극한과 같다.
\end{lemma}

\begin{proof}
$\lim_i \mathcal{F}_i$를 준층으로서의 극한이라 하자. 이것이 층임을
보이면 층의 범주에서 극한이라는 사실은 자명하게 따른다. 층 성질을
증명하기 위해 $\mathcal{V} = \{V_j \to V\}_{j\in J}$를 덮개라 하고
$(s_j)_{j\in J}$를 $H^0(\mathcal{V}, \lim_i \mathcal{F}_i)$의 원소라 하자.
사영사상들을 사용하면 $H^0(\mathcal{V}, \mathcal{F}_i)$의 원소
$(s_{j, i})_{j\in J}$를 얻는다. $\mathcal{F}_i$의 층 성질에 의해
$s_{j, i} = s_i|_{V_j}$가 되는 유일한 $s_i \in \mathcal{F}_i(V)$가
존재함을 알 수 있다. $\phi : i \to i'$를 지표범주의 사상이라 하자.
$\mathcal{F}(\phi) : \mathcal{F}_i \to \mathcal{F}_{i'}$가 $s_i$를
$s_{i'}$로 보냄을 보이고자 한다.
% P01-ERRATA: P01-E1265 -- frozen proof reverses the indices s_{j,i} to s_{i,j}.
모든 $j$에 대하여 절단 $s_{i, j}$와 $s_{i', j}$에 대해서는 이것이
성립함을 알고 있으므로 $\mathcal{F}_{i'}$의 층 성질에 의해 원하는
결론이 성립한다. 이제 $(\lim_i \mathcal{F}_i)(V)$의 원소
$s = (s_i)_{i \in \Ob(\mathcal{I})}$를 얻었다. 이 원소가
$s_j = s|_{V_j}$라는 필요한 성질을 가짐을 확인하는 일은 독자에게 맡긴다.
\end{proof}

\begin{example}
\label{sites-example-singleton-sheaf}
특별한 예는 공도표 위의 극한이다. 이는 (준)층들의 범주에서 끝대상을 준다.
각 $\mathcal{C}$의 대상 $U$에 한원소집합을 대응시키고 유일한 제한사상을
갖는 준층이며, 이 준층은 또한 층이다. 이 층을 흔히 $*$로 나타낸다.
\end{example}

\noindent
$\mathcal{J}_U$를 $U$의 모든 덮개들의 범주라 하자. 다시 말해
$\mathcal{J}_U$의 대상들은 $\mathcal{C}$에서의 $U$의 덮개들이고 사상들은
세분들이다. 사이트에 관한 우리의 규약에 의해 이는 실제로 범주이다. 즉
대상들과 사상들의 모음은 집합을 이룬다. $\{\text{id}_U\}$가 대상이므로
$\Ob(\mathcal{J}_U)$는 공집합이 아님에 유의하자. 위 관찰들에 의해 구성
$\mathcal{U} \mapsto H^0(\mathcal{U}, \mathcal{F})$는
$\mathcal{J}_U$ 위의 반변함자이다. 다음과 같이 정의한다.
$$
\mathcal{F}^{+}(U)
=
\colim_{\mathcal{J}_U^{opp}}
H^0(\mathcal{U}, \mathcal{F})
$$
극한과 여극한에 관한 논의는 Categories, Section
\ref{categories-section-limits}를 보라. 나중에 $\mathcal{F}^{+}(U)$가
$U$ 위의 $\mathcal{F}$의 0차 {\v C}ech 코호몰로지임을 보게 된다는
점을 지적해 둔다.

\medskip\noindent
여극한의 구조에 관해 더 말하기 전에 집합들의 모음
$\mathcal{F}^{+}(U)$, $U \in \Ob(\mathcal{C})$를 준층으로 만들자.
이를 위해 $V \to U$를 $\mathcal{C}$의 사상이라 하자. 사이트의 공리들에
의해 함자\footnote{이 구성은 실제로 올곱 $U_i \times_U V$의 선택을
수반하므로 선택공리를 사용한다. 아래에서 보듯이 얻어진 사상은 한 선택들에
의존하지 않는다.}
$$
\mathcal{J}_U
\longrightarrow
\mathcal{J}_V, \quad
\{U_i \to U\}
\longmapsto
\{U_i \times_U V \to V\}.
$$
가 존재한다. 사영사상들은 덮개들의 함자적인 사상
$\{U_i \times_U V \to V\} \to \{U_i \to U\}$를 주며, 따라서 위의
구성에 의해 집합들의 함자적인 사상
$H^0(\{U_i \to U\}, \mathcal{F}) \to
H^0(\{U_i \times_U V \to V\}, \mathcal{F})$를 준다.
달리 말하면 함자
$H^0(-, \mathcal{F}) : \mathcal{J}_U^{opp} \to \textit{Sets}$에서 합성
$\mathcal{J}_U^{opp} \to \mathcal{J}_V^{opp}
\xrightarrow{H^0(-, \mathcal{F})} \textit{Sets}$로 가는 함자의 변환이
존재한다. 따라서 여극한의 일반론에 의해 표준적인 사상
$\mathcal{F}^+(U) \to \mathcal{F}^+(V)$를 얻는다. 위에서 기술한 집합
$\mathcal{F}^+(U)$의 관점에서 이는
$s = (s_i) \in H^0(\{U_i \to U\}, \mathcal{F})$에 대응하는 원소를
$(s_i|_{V \times_U U_i})
\in H^0(\{U_i \times_U V \to V\}, \mathcal{F})$에 대응하는 원소로
보낼 뿐이다.

\begin{lemma}
\label{sites-lemma-plus-presheaf}
위 구성들은 준층 $\mathcal{F}^+$와 준층들의 표준적인 사상
$\mathcal{F} \to \mathcal{F}^+$를 정의한다.
\end{lemma}

\begin{proof}
사상 $W \to V \to U$가 주어졌을 때 합성
$\mathcal{F}^+(U) \to \mathcal{F}^+(V) \to \mathcal{F}^+(W)$가 사상
$\mathcal{F}^+(U) \to \mathcal{F}^+(W)$와 같음을 보이기만 하면 된다.
덮개 $\{U_i \to U\}$와
$s = (s_i) \in H^0(\{U_i \to U\}, \mathcal{F})$가 주어졌을 때
표준적인 동형
$W \times_U U_i \cong W \times_V (V \times_U U_i)$가 존재하고,
$s_i|_{W \times_U U_i}$가 이 동형을 통해
$(s_i|_{V \times_U U_i})|_{W \times_V (V \times_U U_i)}$와
대응함을 직접 확인하여 이를 보일 수 있다.
\end{proof}

\noindent
좀 더 간접적으로는 Lemma \ref{sites-lemma-independent-refinement}의 결과에 의해
위와 같은 원소 $s$를 $W$의 임의의 덮개에서
$\{U_i \to U\}$로 가는 임의의 사상으로 당길 수 있고, 언제나
$\mathcal{F}^+(W)$의 같은 원소를 얻는다.

\begin{lemma}
\label{sites-lemma-plus-functorial}
대응 $\mathcal{F} \mapsto (\mathcal{F} \to \mathcal{F}^+)$는 함자이다.
\end{lemma}

\begin{proof}
이를 증명하는 대신 무엇을 증명해야 하는지 정확히 진술하겠다. 준층들의 사상
$\mathcal{F} \to \mathcal{G}$가 주어졌을 때 다음 도식의 가환성을 증명하라.
$$
\xymatrix{
\mathcal{F} \ar[r] \ar[d]
&
\mathcal{F}^{+} \ar[d]
\\
\mathcal{G} \ar[r]
&
\mathcal{G}^{+}
}
$$
\end{proof}

\noindent
다음 두 보조정리는 위의 여극한들이 유향집합 위의 여극한임을 함의한다.

\begin{lemma}
\label{sites-lemma-common-refinement}
사이트 $\mathcal{C}$의 주어진 대상 $U$의 두 덮개
$\{U_i \to U\}$와 $\{V_j \to U\}$가 주어지면 이 둘의 공통 세분인
덮개가 존재한다.
\end{lemma}

\begin{proof}
$\mathcal{C}$가 사이트이므로 모든 $i$에 대하여 사상족
$\{V_j \times_U U_i \to U_i\}_j$는 덮개이다. 이어서 다른 공리에 의해
$\{V_j \times_U U_i \to U\}_{i, j}$는 $U$의 덮개이다. 명백히 이 덮개는
주어진 두 덮개를 모두 세분한다.
\end{proof}

\begin{lemma}
\label{sites-lemma-independent-refinement}
덮개들의 두 사상 $f, g: \mathcal{U} \to \mathcal{V}$가 같은 사상
$U \to V$를 유도하면 이들은 같은 사상
$H^0(\mathcal{V}, \mathcal{F}) \to  H^0(\mathcal{U}, \mathcal{F})$를
유도한다.
\end{lemma}

\begin{proof}
$\mathcal{U} = \{U_i \to U\}_{i\in I}$와
$\mathcal{V} = \{V_j \to V\}_{j\in J}$라 하자. 사상 $f$는 사상
$U\to V$, 사상 $\alpha : I \to J$, 그리고 사상들
$f_i : U_i \to V_{\alpha(i)}$로 이루어진다. 마찬가지로 $g$는 사상
$\beta : I \to J$와 사상들 $g_i : U_i \to V_{\beta(i)}$를 정한다.
$f$와 $g$가 같은 사상 $U\to V$를 유도하므로 모든 $i\in I$에 대하여 도식
$$
\xymatrix{
&
V_{\alpha(i)} \ar[dr]
\\
U_i \ar[ur]^{f_i} \ar[dr]_{g_i}
&
&
V
\\
&
V_{\beta(i)} \ar[ur]
}
$$
은 가환이다. 따라서 $f$와 $g$는 올곱을 통하여 분해된다.
$$
\xymatrix{
&
V_{\alpha(i)}
\\
U_i \ar[r]^-\varphi \ar[ur]^{f_i} \ar[dr]_{g_i}
&
V_{\alpha(i)} \times_V V_{\beta(i)} \ar[u]_{\text{pr}_1} \ar[d]^{\text{pr}_2}
\\
&
V_{\beta(i)}.
}
$$
이제 $s = (s_j)_j \in H^0(\mathcal{V}, \mathcal{F})$라 하자. 그러면 모든
$i\in I$에 대하여
$$
(f^*s)_i
=
f_i^*(s_{\alpha(i)})
=
\varphi^*\text{pr}_1^*(s_{\alpha(i)})
=
\varphi^*\text{pr}_2^*(s_{\beta(i)})
=
g_i^*(s_{\beta(i)})
=
(g^*s)_i,
$$
이며, 가운데 등식은 $H^0(\mathcal{V}, \mathcal{F})$의 정의에서 나온다.
이는 $f$와 $g$가 유도하는 사상
$H^0(\mathcal{V}, \mathcal{F}) \to H^0(\mathcal{U}, \mathcal{F})$가
같음을 보인다.
\end{proof}

\begin{remark}
\label{sites-remark-both-refine-same-H0}
특히 이 보조정리는 $\mathcal{U}$가 $\mathcal{V}$의 세분이고
$\mathcal{V}$가 $\mathcal{U}$의 세분이면 표준적인 동일시
$H^0(\mathcal{U}, \mathcal{F}) =
H^0(\mathcal{V}, \mathcal{F})$가 존재함을 보인다.
\end{remark}

\noindent
이 두 보조정리와 $\mathcal{J}_U$가 공집합이 아니라는 사실로부터 도표
$H^0(-, \mathcal{F}) : \mathcal{J}_U^{opp}
\to \textit{Sets}$가 여과됨이 따른다. Categories, Definition
\ref{categories-definition-directed}를 보라.
따라서 Categories, Section \ref{categories-section-directed-colimits}에 의해
여극한 $\mathcal{F}^{+}(U)$를 다음과 같이 간단히 기술할 수 있다. 즉 집합
$\mathcal{F}^{+}(U)$의 모든 원소는 $U$의 어떤 덮개 $\mathcal{U}$에 대한
$s \in H^0(\mathcal{U}, \mathcal{F})$에서 나온다. 두 번째 원소
$s' \in H^0(\mathcal{U}', \mathcal{F})$가 주어졌을 때 $s$와 $s'$가
여극한의 같은 원소를 정할 필요충분조건은 $U$의 덮개 $\mathcal{V}$와 세분
$f : \mathcal{V} \to \mathcal{U}$ 및
$f' : \mathcal{V} \to \mathcal{U}'$가 존재하여
$H^0(\mathcal{V}, \mathcal{F})$에서 $f^*s = (f')^*s'$가 되는 것이다.
자명한 덮개 $\{\text{id}_U\}$가 $\mathcal{J}_U$의 대상이므로 표준적인
사상 $\mathcal{F}(U) \to \mathcal{F}^+(U)$를 얻는다.

\begin{lemma}
\label{sites-lemma-plus-surjective}
사상 $\theta : \mathcal{F} \to \mathcal{F}^+$는 다음 성질을 갖는다.
$\mathcal{C}$의 모든 대상 $U$와 모든 절단
$s \in \mathcal{F}^+(U)$에 대하여 덮개 $\{U_i \to U\}$가 존재하여
$s|_{U_i}$는 $\theta : \mathcal{F}(U_i)
\to \mathcal{F}^{+}(U_i)$의 상에 속한다.
\end{lemma}

\begin{proof}
$s$가 $(s_i) \in H^0(\{U_i \to U\}, \mathcal{F})$에서 나오게 하는 덮개
$\{U_i \to U\}$를 택하자. Lemma \ref{sites-lemma-independent-refinement}에 의해
$s$를 $\mathcal{F}^+(U_i)$의 원소로 당기는 것을 계산할 때 덮개
$\{U_i \to U_i\}$와 덮개들의 (명백한) 사상
$\{U_i \to U_i\} \to \{U_i \to U\}$를 사용할 수 있다. 실제로 이 덮개를
사용하면 $s$를 $U_i$에 제한한 것으로 정확히 $\theta(s_i)$를 얻는다.
\end{proof}

\begin{definition}
\label{sites-definition-separated}
사이트 $\mathcal{C}$ 위의 집합의 준층 $\mathcal{F}$에 대하여 모든 덮개
$\{U_i \rightarrow U\}$에서 사상
$\mathcal{F}(U) \to \prod \mathcal{F}(U_i)$가 단사이면
그 준층이 {\it 분리되어 있다}고 한다.
% P01-ERRATA: P01-E0064 -- frozen English quantifies malformedly over coverings of a family.
\end{definition}

\begin{theorem}
\label{sites-theorem-plus}
위와 같은 $\mathcal{F}$에 대하여 다음이 성립한다.
\begin{enumerate}
\item
\label{sites-item-sep}
준층 $\mathcal{F}^+$는 분리되어 있다.
\item
\label{sites-item-sheaf}
$\mathcal{F}$가 분리되어 있으면 $\mathcal{F}^+$는 층이고 준층들의 사상
$\mathcal{F} \to \mathcal{F}^+$는 단사이다.
\item
\label{sites-item-plus-iso}
$\mathcal{F}$가 층이면 $\mathcal{F} \to \mathcal{F}^+$는 동형사상이다.
\item
\label{sites-item-plusplus}
준층 $\mathcal{F}^{++}$는 언제나 층이다.
\end{enumerate}
\end{theorem}

\begin{proof}
(\ref{sites-item-sep})의 증명. $s, s' \in \mathcal{F}^+(U)$라 하고 모든 $i$에
대하여 $s|_{U_i} = s'|_{U_i}$가 되게 하는 어떤 덮개
$\{U_i \to U\}$가 존재한다고 하자. 이제 $U$의 덮개 세 개가 있다. 하나는
위의 덮개 $\{U_i \to U\}$이고, 하나는 Lemma
\ref{sites-lemma-plus-surjective}에서처럼 $s$에 대한 덮개 $\mathcal{U}$이며,
나머지는 $s'$에 대한 비슷한 덮개 $\mathcal{U}'$이다. Lemma
\ref{sites-lemma-common-refinement}에 의해 이들의 공통 세분
$\{W_j \to U\}$를 찾을 수 있다. 이는
$s|_{W_j} = \theta(s_j)$이고 $s'|_{W_j}$에 대해서도 마찬가지이며
$\theta(s_j) = \theta(s'_j)$가 되게 하는
$s_j, s'_j \in \mathcal{F}(W_j)$가 존재한다는 뜻이다. 마지막 등식은 어떤
덮개 $\{W_{jk} \to W_j\}$가 존재하여
$s_j|_{W_{jk}} = s'_j|_{W_{jk}}$가 된다는 뜻이다. 그러면
$\{W_{jk} \to U\}$가 덮개이므로 원하는 대로 $s, s'$는
$H^0(\{W_{jk} \to U\}, \mathcal{F})$의 같은 원소로 간다.

\medskip\noindent
(\ref{sites-item-sheaf})의 증명. 모든 사상
$\mathcal{F}(U) \to H^0(\mathcal{U}, \mathcal{F})$가 단사이므로
$\mathcal{F} \to \mathcal{F}^+$가 단사임은 명백하다. 또한
$\mathcal{U} \to \mathcal{U}'$가 세분이면
$H^0(\mathcal{U}', \mathcal{F})
\to H^0(\mathcal{U}, \mathcal{F})$가 단사임도 명백하다. 이제
$\{U_i \to U\}$가 덮개이고 $(s_i)$가 모든 $i, i' \in I$에 대하여 층 조건
$s_i|_{U_i \times_U U_{i'}} = s_{i'}|_{U_i \times_U U_{i'}}$를 만족하는
$\mathcal{F}^+(U_i)$의 원소족이라 하자. Lemma
\ref{sites-lemma-plus-surjective}에서처럼 각 $s_i|_{U_{ij}}$가 유일한 원소
$s_{ij} \in \mathcal{F}(U_{ij})$의 상이 되도록 덮개
$\{U_{ij} \to U_i\}$를 택하자. 이 올곱은
$U_i \times_U U_{i'}$로 가고 그 위에서 등식이 성립하므로, 층 조건은
$s_{ij}$와 $s_{i'j'}$가 $U_{ij} \times_U U_{i'j'}$ 위에서 일치함을
함의한다. 따라서 $(s_{ij}) \in H^0(\{U_{ij} \to U\}, \mathcal{F})$는
원소 $s \in \mathcal{F}^+(U)$를 준다. $s|_{U_i} = s_i$임을 확인하는
일은 독자에게 맡긴다.

\medskip\noindent
(\ref{sites-item-plus-iso})의 증명. 층 성질은 $U$의 모든 덮개 $\mathcal{U}$에
대하여 모든 사상
% P01-ERRATA: P01-E0065 -- frozen domain omits (U).
$\mathcal{F} \to H^0(\mathcal{U}, \mathcal{F})$가 전단사임을 정확히
뜻하므로 정의에서 바로 따른다.

\medskip\noindent
명제 (\ref{sites-item-plusplus})는 이제 명백하다.
\end{proof}

\begin{definition}
\label{sites-definition-associated-sheaf}
$\mathcal{C}$를 사이트라 하고 $\mathcal{F}$를 $\mathcal{C}$ 위의 집합의
준층이라 하자. 표준적인 사상 $\mathcal{F} \to \mathcal{F}^\#$와 함께
층 $\mathcal{F}^\# := \mathcal{F}^{++}$를
{\it $\mathcal{F}$에 연관된 층}이라고 한다.
\end{definition}

\begin{proposition}
\label{sites-proposition-sheafification-adjoint}
표준적인 사상 $\mathcal{F} \to \mathcal{F}^\#$는 다음 보편 성질을 갖는다.
$\mathcal{G}$가 집합의 층인 임의의 사상 $\mathcal{F} \to \mathcal{G}$에
대하여 합성 $\mathcal{F} \to \mathcal{F}^\# \to \mathcal{G}$가 주어진
사상과 같게 하는 유일한 사상 $\mathcal{F}^\# \to \mathcal{G}$가 존재한다.
\end{proposition}

\begin{proof}
Lemma \ref{sites-lemma-plus-functorial}에 의해 가환도식
$$
\xymatrix{
\mathcal{F} \ar[r] \ar[d]
&
\mathcal{F}^{+} \ar[r] \ar[d]
&
\mathcal{F}^{++} \ar[d]
\\
\mathcal{G} \ar[r]
&
\mathcal{G}^{+} \ar[r]
&
\mathcal{G}^{++}
}
$$
을 얻고, Theorem \ref{sites-theorem-plus}에 의해 아래쪽 수평 사상들은
동형사상이다. 유일성은 $\mathcal{F}^\#$의 모든 절단이 국소적으로
$\mathcal{F}$의 절단들에서 나온다는 Lemma
\ref{sites-lemma-plus-surjective}로부터 따른다.
\end{proof}

\noindent
이 결과에서 함자 $\mathcal{F} \mapsto (\mathcal{F} \to \mathcal{F}^\#)$가
함자들의 유일한 동형을 제외하면 유일함은 명백하다. 실제로 포함 함자를
잠시 $i : \Sh(\mathcal{C}) \to \textit{PSh}(\mathcal{C})$라고 쓰자.
위 결과는 실제로
$$
\Mor_{\textit{PSh}(\mathcal{C})}(\mathcal{F}, i(\mathcal{G}))
=
\Mor_{\Sh(\mathcal{C})}(\mathcal{F}^\#, \mathcal{G}).
$$
라고 말한다. 다시 말해 층화 함자는 포함 함자 $i$의 왼쪽 수반이다.

\begin{example}
\label{sites-example-constant-sheaf}
$\mathcal{C}$를 사이트라 하고 $E$를 집합이라 하자.
{\it 값이 $E$인 상수층 $\underline{E}$}은 값이 $E$인 상수함자로 주어지는
준층의 층화이다. 모든 $\mathcal{C}$ 위의 층 $\mathcal{F}$에 대하여
$\Mor_{\Sh(\mathcal{C})}(\underline{E}, \mathcal{F}) =
\Mor_{\textit{Sets}}(E, \Gamma(\mathcal{C}, \mathcal{F}))$라는 규칙으로
특징지어지며, 이 표기는 Definition \ref{sites-definition-global-sections}에서
가져왔다.
\end{example}

\noindent
두 보조정리로 이 절을 마친다.

\begin{lemma}
\label{sites-lemma-colimit-sheaves}
\begin{slogan}
층의 범주에서의 여극한은 준층의 범주에서의 여극한을 층화한 것과 같다.
\end{slogan}
$\mathcal{F} : \mathcal{I} \to \Sh(\mathcal{C})$를 도표라 하자. 그러면
$\colim_\mathcal{I} \mathcal{F}$는 존재하고 준층의 범주에서의 여극한을
층화한 것이다.
\end{lemma}

\begin{proof}
층화 함자는 왼쪽 수반이므로 모든 여극한과 가환한다. Categories, Lemma
\ref{categories-lemma-adjoint-exact}를 보라. 따라서
$\textit{PSh}(\mathcal{C})$가 여극한들을 가지므로
$\Sh(\mathcal{C})$도 여극한들을 가진다(이들은 준층에서의 여극한들을
층화한 것이다).
\end{proof}

\begin{lemma}
\label{sites-lemma-sheafification-exact}
함자 $\textit{PSh}(\mathcal{C}) \to \Sh(\mathcal{C})$,
$\mathcal{F} \mapsto \mathcal{F}^\#$는 완전하다.
\end{lemma}

\begin{proof}
이는 왼쪽 수반이므로 오른쪽 완전하다. Categories, Lemma
\ref{categories-lemma-exact-adjoint}를 보라. 한편 Lemmas
\ref{sites-lemma-common-refinement}와 Lemma
\ref{sites-lemma-independent-refinement}에 의해 $\mathcal{F}^+$의 구성에서
여극한들은 실제로 유향집합 $\Ob(\mathcal{J}_U)$ 위의 여극한들이다. 여기서
$\mathcal{U}$가 $\mathcal{U}'$의 세분일 필요충분조건으로
$\mathcal{U} \geq \mathcal{U}'$라 둔다. 그러므로 Categories, Lemma
\ref{categories-lemma-directed-commutes}에 의해 함자
$\mathcal{F} \to \mathcal{F}^+$가 유한 극한들과 가환함을 알 수 있다
(준층에서 준층으로 가는 함자로서). 이제 Lemma
\ref{sites-lemma-limit-sheaf}를 사용하여 결론을 얻는다.
\end{proof}

\begin{lemma}
\label{sites-lemma-sections-sheafification}
$\mathcal{C}$를 사이트라 하고 $\mathcal{F}$를 $\mathcal{C}$ 위의 집합의
준층이라 하자. $\theta^2 : \mathcal{F} \to \mathcal{F}^\#$로
$\mathcal{F}$에서 그 층화로 가는 표준적인 사상을 나타내자.
$U$를 $\mathcal{C}$의 대상이라 하고 $s \in \mathcal{F}^\#(U)$라 하자.
다음을 만족하는 덮개 $\{U_i \to U\}$와 절단들
$s_i \in \mathcal{F}(U_i)$가 존재한다.
\begin{enumerate}
\item $s|_{U_i} = \theta^2(s_i)$이다.
\item 모든 $i, j$에 대하여 $\mathcal{C}$의 덮개
$\{U_{ijk} \to U_i \times_U U_j\}$가 존재하여 $s_i$와 $s_j$를 각
$U_{ijk}$로 당긴 것들이 일치한다.
\end{enumerate}
역으로 임의의 덮개 $\{U_i \to U\}$와 (2)를 만족하는 원소들
$s_i \in \mathcal{F}(U_i)$가 주어지면, (1)을 만족하는 유일한 절단
$s \in \mathcal{F}^\#(U)$가 존재한다.
\end{lemma}

\begin{proof}
생략한다.
\end{proof}

\begin{lemma}
\label{sites-lemma-isomorphism-sheafifications}
$\mathcal{C}$를 사이트라 하고 $\mathcal{F} \to \mathcal{G}$를
$\mathcal{C}$ 위의 집합의 준층들의 사상이라 하자.
$\mathcal{B}$로 $\mathcal{F}(U) \to \mathcal{G}(U)$가 전단사인
$U \in \Ob(\mathcal{C})$들의 집합을 나타내자. $\mathcal{C}$의 모든 대상이
$\mathcal{B}$의 원소들로 이루어진 덮개를 가지면
$\mathcal{F}^\# \to \mathcal{G}^\#$는 동형사상이다.
\end{lemma}

\begin{proof}
$U \in \Ob(\mathcal{C})$라 하자.
$\mathcal{F}^\#(U) \to \mathcal{G}^\#(U)$가 전사임을 증명하자.
이를 위해 Lemma \ref{sites-lemma-sections-sheafification}를 더 언급하지 않고
사용하겠다. 임의의 $s \in \mathcal{G}^\#(U)$에 대하여 덮개
$\{U_i \to U\}$와 절단들 $s_i \in \mathcal{G}(U_i)$가 존재하여
\begin{enumerate}
\item $s|_{U_i}$는 $\mathcal{G} \to \mathcal{G}^\#$에 의한 $s_i$의 상이다.
\item 모든 $i, j$에 대하여
% P01-ERRATA: P01-E0066 -- frozen source names undefined site D instead of C.
$\mathcal{D}$의 덮개 $\{U_{ijk} \to U_i \times_U U_j\}$가 존재하여
$s_i$와 $s_j$를 각 $U_{ijk}$로 당긴 것들이 일치한다.
\end{enumerate}
가정에 의해 각 $i$에 대하여 $U_{i, a} \in \mathcal{B}$인 덮개
$\{U_{i, a} \to U_i\}$를 택할 수 있다. 그러면
$\{U_{i, a} \to U\}$는 덮개이다. $s_i$의 $\mathcal{G}(U_{i, a})$에서의
상을 $s_{i, a}$로 나타내면 덮개
$$
\{(U_{i, a} \times_U U_{j, b}) \times_{U_i \times_U U_j} U_{ijk}
\to U_{i, a} \times_U U_{j, b}\}
$$
의 각 원소로 $s_{i, a}$와 $s_{j, b}$를 당긴 것들이 일치함을 알 수 있다.
따라서 $U_i \in \mathcal{B}$라고 가정해도 된다. 논증을 반복하면 모든
$i, j, k$에 대하여 $U_{ijk} \in \mathcal{B}$라고 가정해도 된다
(세부사항은 생략한다). 보조정리의 진술에서 $\mathcal{B}$를 정의한 방식에
따라 $\mathcal{F}(U_i) \to \mathcal{G}(U_i)$가 전단사이므로 $s_i$로 가는
유일한 $t_i \in \mathcal{F}(U_i)$를 얻는다. $U_{ijk} \in \mathcal{B}$이고
% P01-ERRATA: P01-E0067 -- frozen rationale repeats the conclusion instead of citing s_i,s_j.
$t_i$와 $t_j$에 대한 일치가 있으므로 $t_i$와 $t_j$를 각 $U_{ijk}$로
당긴 것들은 $\mathcal{F}(U_{ijk})$에서 일치한다. 그러면 보조정리에 의해
$t|_{U_i}$가 $t_i$의 상이 되게 하는 유일한
$t \in \mathcal{F}^\#(U)$가 존재한다. 따라서 $t$는 $s$로 간다.
$\mathcal{F}^\#(U) \to \mathcal{G}^\#(U)$가 단사라는 증명은 생략한다.
\end{proof}

\section{층들의 단사 사상과 전사 사상}
\label{sites-section-sheaves-injective}

\begin{definition}
\label{sites-definition-sheaves-injective-surjective}
$\mathcal{C}$를 사이트라 하고 $\varphi : \mathcal{F}
\to \mathcal{G}$를 집합의 층들의 사상이라 하자.
\begin{enumerate}
\item $\mathcal{C}$의 모든 대상 $U$에 대하여 사상
$\varphi : \mathcal{F}(U) \to \mathcal{G}(U)$가 단사이면
$\varphi$가 {\it 단사}라고 한다.
\item $\mathcal{C}$의 모든 대상 $U$와 모든 절단
$s\in \mathcal{G}(U)$에 대하여 덮개 $\{U_i \to U\}$가 존재하여 모든
$i$에 대해 제한 $s|_{U_i}$가 사상
$\varphi : \mathcal{F}(U_i) \to \mathcal{G}(U_i)$의 상에 속하면
$\varphi$가 {\it 전사}라고 한다.
\end{enumerate}
\end{definition}

\begin{lemma}
\label{sites-lemma-mono-epi-sheaves}
위에서 정의한 단사 사상(각각 전사 사상)은 정확히 범주
$\Sh(\mathcal{C})$의 단사사상(각각 전사사상)이다. 층들의 사상이
동형사상일 필요충분조건은 그 사상이 단사이면서 전사인 것이다.
\end{lemma}

\begin{proof}
생략한다.
\end{proof}

\begin{lemma}
\label{sites-lemma-coequalizer-surjection}
$\mathcal{C}$를 사이트라 하자. $\mathcal{F} \to \mathcal{G}$를
집합의 층들의 전사 사상이라 하자. 그러면 도식
$$
\xymatrix{
\mathcal{F} \times_\mathcal{G} \mathcal{F}
\ar@<1ex>[r] \ar@<-1ex>[r]
&
\mathcal{F} \ar[r]
&
\mathcal{G}}
$$
은 $\mathcal{G}$를 여등화자로 나타낸다.
\end{lemma}

\begin{proof}
$\mathcal{H}$를 집합의 층이라 하고
$\varphi : \mathcal{F} \to \mathcal{H}$를 두 사상
$\mathcal{F} \times_\mathcal{G} \mathcal{F} \to \mathcal{F}$를
등화하는 층들의 사상이라 하자. $\mathcal{G}' \subset \mathcal{G}$를
사상 $\mathcal{F} \to \mathcal{G}$의 준층으로서의 상이라 하자. 곱
$\mathcal{F} \times_\mathcal{G} \mathcal{F}$는 준층의 범주에서 계산할 수
있으므로 준층의 곱 $\mathcal{F} \times_{\mathcal{G}'} \mathcal{F}$와
같음을 알 수 있다. 따라서 $\varphi$는 준층들의 유일한 사상
$\psi' : \mathcal{G}' \to \mathcal{H}$를 유도한다. Lemma
\ref{sites-lemma-mono-epi-sheaves}에 의해 $\mathcal{G}$는 $\mathcal{G}'$의
층화이므로 $\psi'$가 층들의 사상
$\psi : \mathcal{G} \to \mathcal{H}$로 유일하게 확장된다고 결론 내린다.
$\varphi$가 $\psi$와 주어진 사상의 합성과 같다는 확인은 생략한다.
\end{proof}

\section{표현가능 층}
\label{sites-section-representable-sheaves}

\noindent
$\mathcal{C}$를 범주라 하자. 표준 위상은 모든 표현가능 준층이 층이 되게 하는
가장 세밀한 위상이다(Definition \ref{sites-definition-canonical-topology}에서
형식적으로 정의하지만 여기서는 필요하지 않다). 이 위상은
$\mathcal{C}$ 위의 사이트 구조에 연관된 위상과 언제나 같지는 않다.
$\mathcal{C}$가 올곱을 가지는 경우 이 위상을 생성하는 덮개들의 모음을
제시할 것이다. 먼저 다음 일반적인 정의를 둔다.

\begin{definition}
\label{sites-definition-universal-effective-epimorphisms}
$\mathcal{C}$를 범주라 하자. 사상족 $\{U_i \to U\}_{i \in I}$의 모든 사상
$U_i \to U$가 표현가능이고(Categories, Definition
\ref{categories-definition-representable-morphism}를 보라), 모든
$X\in \Ob(\mathcal{C})$에 대하여 열
$$
\xymatrix{
\Mor_\mathcal{C}(U, X) \ar[r]
&
\prod\nolimits_{i \in I} \Mor_\mathcal{C}(U_i, X)
\ar@<1ex>[r] \ar@<-1ex>[r]
&
\prod\nolimits_{(i, j) \in I^2} \Mor_\mathcal{C}(U_i \times_U U_j, X)
}
$$
이 등화자 도식이면 이 사상족을 {\it 유효 전사}라고 한다. 사상족
$\{U_i \to U\}$에 대하여 임의의 사상 $V \to U$에 의한 밑변환
$\{U_i \times_U V \to V\}$가 유효 전사이면 원래 사상족을
{\it 보편 유효 전사}라고 한다.
\end{definition}

\noindent
보편 유효 전사인 사상족들의 클래스는 Definition
\ref{sites-definition-site}의 공리들을 만족한다. $\mathcal{C}$가 올곱들을 가지면
연관된 위상은 표준 위상이다(이 경우 사이트를 얻으려면 Sets, Lemma
\ref{sets-lemma-coverings-site}에서와 같이 논증한다).

\medskip\noindent
역으로 $\mathcal{C}$가 모든 표현가능 준층이 층인 사이트라고 하자.
그러면 명백히 모든 덮개는 보편 유효 전사이다. 따라서 사이트의 맥락에서는
다음 정의가 ``올바른'' 정의이다.

\begin{definition}
\label{sites-definition-weaker-than-canonical}
사이트 $\mathcal{C}$의 위상이 {\it 표준 위상보다 약하다}고 하거나
{\it 준표준}이라고 하는 것은 $\mathcal{C}$의 모든 덮개가 보편 유효
전사인 경우이다.
\end{definition}

\noindent
표현가능 층은 층이기도 한 표현가능 준층이다. 위상이 준표준이 아닐 때는
이 용어를 피하는 편이 나을 수 있으므로, 그 경우에만 이를 형식적으로 정의한다.

\begin{definition}
\label{sites-definition-representable-sheaf}
$\mathcal{C}$를 그 위상이 준표준인 사이트라 하자. 요네다 매입 $h$는
(Categories, Section \ref{categories-section-opposite}를 보라)
$\mathcal{C}$를 $\mathcal{C}$ 위의 층들의 범주의 충만부분범주로 나타낸다.
이 경우 $U \in \Ob(\mathcal{C})$인 꼴 $h_U$의 층들을
$\mathcal{C}$ 위의 {\it 표현가능 층}이라고 한다. 표기: 때로는 $U$에
연관된 표현가능 층 $h_U$를 {\it $\underline{U}$}로 나타낸다.
\end{definition}

\noindent
정의의 상황에서는 준층에 대해 성립하는 사실로부터 모든 층
$\mathcal{F}$에 대하여
$$
\Mor_{\Sh(\mathcal{C})}(h_U, \mathcal{F}) = \mathcal{F}(U)
$$
가 성립함에 유의하자. (\ref{sites-equation-map-representable-into-presheaf})를
보라. 일반적으로 준층 $h_U$들은 층이 아니며 층을 얻으려면 이들을 층화해야
한다. 이 경우에도 모든 층 $\mathcal{F}$에 대하여
\begin{equation}
\label{sites-equation-map-representable-into-sheaf}
\Mor_{\Sh(\mathcal{C})}(h_U^\#, \mathcal{F}) =
\Mor_{\textit{PSh}(\mathcal{C})}(h_U, \mathcal{F}) =
\mathcal{F}(U)
\end{equation}
가 성립한다. 실제로 첫 번째 등식은 $\#$의 수반성에서, 두 번째 등식은
(\ref{sites-equation-map-representable-into-presheaf})에서 나온다.

\begin{lemma}
\label{sites-lemma-covering-surjective-after-sheafification}
\begin{slogan}
덮개는 층화한 뒤 전사가 된다.
\end{slogan}
$\mathcal{C}$를 사이트라 하자. $\{U_i \to U\}_{i \in I}$가 사이트
$\mathcal{C}$의 덮개이면 집합의 준층들의 사상
$$
\coprod\nolimits_{i \in I} h_{U_i} \to h_U
$$
은 층화한 뒤 전사가 된다.
\end{lemma}

\begin{proof}
위의 Lemma \ref{sites-lemma-mono-epi-sheaves}에 의해
$\coprod\nolimits_{i \in I} h_{U_i}^\# \to h_U^\#$가 전사사상임을
보여야 한다. $\mathcal{F}$를 집합의 층이라 하자. 사상
$h_U^\# \to \mathcal{F}$는 절단 $s \in \mathcal{F}(U)$에 대응한다.
따라서 $\Mor(h_U^\#, \mathcal{F})
\to \prod_i \Mor(h_{U_i}^\#, \mathcal{F})$의 단사성은
$\mathcal{F}$의 층 성질에서 바로 따른다.
\end{proof}

\noindent
다음 보조정리는 위상이 표준 위상보다 약한 경우에 모든 층이 어떤 의미에서
표현가능 층들로 이루어진다고 말한다.

\begin{lemma}
\label{sites-lemma-sheaf-coequalizer-representable}
$\mathcal{C}$를 사이트라 하자. $E \subset \Ob(\mathcal{C})$를
$\mathcal{C}$의 모든 대상이 $E$의 원소들로 이루어진 덮개를 가지게 하는
부분집합이라 하자. $\mathcal{F}$를 집합의 층이라 하자. $\mathcal{F}$를
여등화자로 나타내는 집합의 층들의 도식
$$
\xymatrix{
\mathcal{F}_1 \ar@<1ex>[r] \ar@<-1ex>[r] &
\mathcal{F}_0 \ar[r] &
\mathcal{F}
}
$$
이 존재하며, $\mathcal{F}_i$, $i = 0, 1$은 $U \in E$인 꼴
$h_U^\#$의 층들의 쌍대곱들이다.
\end{lemma}

\begin{proof}
먼저 원하는 꼴의 전사사상 $\mathcal{F}_0 \to \mathcal{F}$가 존재함을
보이자. 다음 사상을 취하면 된다.
$$
\mathcal{F}_0 =
\coprod\nolimits_{U \in E, s \in \mathcal{F}(U)}
(h_U)^\# \longrightarrow \mathcal{F}
$$
여기서 $(U, s)$에 대응하는 성분에 제한된 화살표는 원소
$\text{id}_U \in h_U^\#(U)$를 절단 $s \in \mathcal{F}(U)$로 보낸다.
Lemma \ref{sites-lemma-mono-epi-sheaves}와 $E$에 관한 조건에 의해 이는
전사사상이다. $\mathcal{F}_1$을 구성하기 위해 먼저
$\mathcal{G} = \mathcal{F}_0 \times_\mathcal{F} \mathcal{F}_0$로 놓고,
위에서와 같이 전사사상 $\mathcal{F}_1 \to \mathcal{G}$를 구성한다.
Lemma \ref{sites-lemma-coequalizer-surjection}를 보라.
\end{proof}

\begin{lemma}
\label{sites-lemma-colimit-representable}
$\mathcal{C}$를 사이트라 하고 $\mathcal{F}$를 $\mathcal{C}$ 위의 집합의
층이라 하자. 다음을 만족하는 도표 $\mathcal{I} \to \mathcal{C}$,
$i \mapsto U_i$가 존재한다.
$$
\mathcal{F} = \colim_{i \in \mathcal{I}} h_{U_i}^\#
$$
또한 $E \subset \Ob(\mathcal{C})$가 $\mathcal{C}$의 모든 대상이
$E$의 원소들로 이루어진 덮개를 가지게 하는 부분집합이면 모든
$i \in \Ob(\mathcal{I})$에 대하여 $U_i$가 $E$의 원소라고 가정할 수 있다.
\end{lemma}

\begin{proof}
$\mathcal{I}$를 대상이 $U \in \Ob(\mathcal{C})$와
$s \in \mathcal{F}(U)$의 순서쌍 $(U, s)$이고, 사상
$(U, s) \to (U', s')$이 $f^*s' = s$를 만족하는 $\mathcal{C}$의 사상
$f : U \to U'$인 범주라 하자. $\mathcal{I}$의 각 대상 $(U, s)$에 대하여
원소 $s$는 요네다 보조정리에 의해 준층들의 사상
$s : h_U \to \mathcal{F}$를 정의한다. 따라서 층화의 보편 성질에 의해 사상
$h_U^\# \to \mathcal{F}$를 정의한다. 이 사상들이 범주 $\mathcal{I}$의
사상들과 호환됨은 바로 알 수 있다. 따라서 준층들의 사상
$\colim_{(U, s)} h_U \to \mathcal{F}$(여기서 여극한은 준층의 범주에서
취한다)와 층들의 사상
$\colim_{(U, s)} (h_U)^\# \to \mathcal{F}$(여기서 여극한은 층의 범주에서
취한다)를 얻는다. 층화는 포함 함자
$\Sh(\mathcal{C}) \to \textit{PSh}(\mathcal{C})$의 왼쪽 수반이므로
(Proposition \ref{sites-proposition-sheafification-adjoint}) Categories, Lemma
\ref{categories-lemma-adjoint-exact}에 의해
$\colim (h_U)^\# = (\colim h_U)^\#$이다. 따라서 준층들의 사상
$\colim_{(U, s)} h_U \to \mathcal{F}$가 동형사상임을 보이면 충분하다.
이를 위해 $\mathcal{C}$의 모든 대상 $X$에 대하여 사상
$\colim_{(U, s)} h_U(X) \to \mathcal{F}(X)$가 전단사임을 보인다.
실제로 이 사상의 역은 원소 $t \in \mathcal{F}(X)$를 $(X, t)$와
$\text{id}_X \in h_X(X)$에 대응하는 집합
$\colim_{(U, s)} h_U(X)$의 원소로 보낸다.

\medskip\noindent
마지막 명제의 증명은 생략한다.
\end{proof}


\section{연속 함자}
\label{sites-section-continuous-functors}

\begin{definition}
\label{sites-definition-continuous}
$\mathcal{C}$와 $\mathcal{D}$를 사이트라 하자.
함자 $u : \mathcal{C} \to \mathcal{D}$를 {\it 연속}이라고 하는 것은 모든
$\{V_i \to V\}_{i\in I} \in \text{Cov}(\mathcal{C})$에 대하여
다음이 성립하는 경우이다.
\begin{enumerate}
\item $\{u(V_i) \to u(V)\}_{i\in I}$가 $\text{Cov}(\mathcal{D})$의
원소이고,
\item $\mathcal{C}$의 임의의 사상 $T \to V$에 대하여 사상
$u(T \times_V V_i) \to u(T) \times_{u(V)} u(V_i)$가 동형사상이다.
\end{enumerate}
\end{definition}

\noindent
위와 같은 함자 $u$와 $\mathcal{D}$ 위의 집합의 준층 $\mathcal{F}$가
주어졌을 때, $u^p\mathcal{F}$를 단순히 준층 $\mathcal{F} \circ u$로
정의했음을 상기하자. 다시 말해 $\mathcal{C}$의 모든 대상 $V$에 대하여
$$
u^p\mathcal{F} (V) = \mathcal{F}(u(V))
$$
이다.

\begin{lemma}
\label{sites-lemma-pushforward-sheaf}
$\mathcal{C}$와 $\mathcal{D}$를 사이트라 하자.
$u : \mathcal{C} \to \mathcal{D}$를 연속 함자라 하자.
$\mathcal{F}$가 $\mathcal{D}$ 위의 층이면 $u^p\mathcal{F}$도 층이다.
\end{lemma}

\begin{proof}
$\{V_i \to V\}$를 덮개라 하자. 가정에 의해
$\{u(V_i) \to u(V)\}$는 $\mathcal{D}$의 덮개이고
$u(V_i \times_V V_j) = u(V_i)\times_{u(V)}u(V_j)$이다. 따라서
$u^p\mathcal{F}$와 덮개 $\{V_i \to V\}$에 대한 층 조건은
$\mathcal{F}$와 덮개 $\{u(V_i) \to u(V)\}$에 대한 층 조건과 정확히 같다.
\end{proof}

\noindent
혼동을 피하기 위해 집합의 층들의 부분범주에 제한한 함자 $u^p$를 때로
$$
u^s :
\Sh(\mathcal{D})
\longrightarrow
\Sh(\mathcal{C})
$$
로 나타낸다. $u^p$가 왼쪽 수반
$u_p : \textit{PSh}(\mathcal{C}) \to \textit{PSh}(\mathcal{D})$를
가짐을 상기하자. Section \ref{sites-section-functoriality-PSh}를 보라.

\begin{lemma}
\label{sites-lemma-adjoint-sheaves}
Lemma \ref{sites-lemma-pushforward-sheaf}의 상황에서 함자
$u_s : \mathcal{G} \mapsto (u_p \mathcal{G})^\#$는 $u^s$의 왼쪽 수반이다.
\end{lemma}

\begin{proof}
Lemma \ref{sites-lemma-adjoints-u}와 Proposition
\ref{sites-proposition-sheafification-adjoint}에서 바로 따른다.
\end{proof}

\noindent
다음은 기술적인 보조정리이다.

\begin{lemma}
\label{sites-lemma-technical-up}
Lemma \ref{sites-lemma-pushforward-sheaf}의 상황에서 $\mathcal{C}$ 위의 임의의
준층 $\mathcal{G}$에 대하여
$(u_p\mathcal{G})^\# = (u_p(\mathcal{G}^\#))^\#$이다.
\end{lemma}

\begin{proof}
$\mathcal{D}$ 위의 임의의 층 $\mathcal{F}$에 대하여
\begin{eqnarray*}
\Mor_{\Sh(\mathcal{D})}(u_s(\mathcal{G}^\#), \mathcal{F})
& = &
\Mor_{\Sh(\mathcal{C})}(\mathcal{G}^\#, u^s\mathcal{F}) \\
& = &
\Mor_{\textit{PSh}(\mathcal{C})}(\mathcal{G}^\#, u^p\mathcal{F}) \\
& = &
\Mor_{\textit{PSh}(\mathcal{C})}(\mathcal{G}, u^p\mathcal{F}) \\
& = &
\Mor_{\textit{PSh}(\mathcal{D})}(u_p\mathcal{G}, \mathcal{F}) \\
& = &
\Mor_{\Sh(\mathcal{D})}((u_p\mathcal{G})^\#, \mathcal{F})
\end{eqnarray*}
이고, 결과는 요네다 보조정리에서 따른다.
\end{proof}

\begin{lemma}
\label{sites-lemma-pullback-representable-sheaf}
$u : \mathcal{C} \to \mathcal{D}$를 사이트 사이의 연속 함자라 하자.
$\mathcal{C}$의 임의의 대상 $U$에 대하여
$u_sh_U^\# = h_{u(U)}^\#$이다.
\end{lemma}

\begin{proof}
Lemmas \ref{sites-lemma-pullback-representable-presheaf}와
\ref{sites-lemma-technical-up}에서 따른다.
\end{proof}


\begin{remark}
\label{sites-remark-quasi-continuous}
(처음 읽을 때는 건너뛰어도 된다.)
$\mathcal{C}$와 $\mathcal{D}$를 사이트라 하자. Definition
\ref{sites-definition-combinatorial-tautological}의 사상족의 자명한 동치 개념을
사용하여 연속성을 정의하는 조건을 (약간) 약화하자.
$u : \mathcal{C} \to \mathcal{D}$를 함자라 하자. 모든
$\mathcal{V} = \{V_i \to V\}_{i\in I} \in \text{Cov}(\mathcal{C})$에
대하여 다음이 성립하면 $u$를 {\it 준연속}이라고 하자.
\begin{enumerate}
\item[(1')] 사상족 $\{u(V_i) \to u(V)\}_{i\in I}$가
$\text{Cov}(\mathcal{D})$의 한 원소와 자명하게 동치이고,
\item[(2)] $\mathcal{C}$의 임의의 사상 $T \to V$에 대하여 사상
$u(T \times_V V_i) \to u(T) \times_{u(V)} u(V_i)$가 동형사상이다.
\end{enumerate}
$u$가 준연속인 경우에도 Lemmas \ref{sites-lemma-pushforward-sheaf}와
\ref{sites-lemma-adjoint-sheaves}가 성립함을 보이겠다.

\medskip\noindent
먼저 (첫 번째 조건에 의해) 사상 $u(V_i) \to u(V)$는 표현가능 사상과
동형이므로 표현가능이다. 특히 사상족
$u(\mathcal{V}) = \{u(V_i) \to u(V)\}_{i\in I}$는 $\mathcal{D}$ 위의
임의의 준층 $\mathcal{F}$에 대하여 0차 {\v C}ech 코호몰로지 군
% P01-ERRATA: P01-E0063 -- set-valued H^0 need not be a group.
$H^0(u(\mathcal{V}), \mathcal{F})$을 정의한다.
$\mathcal{U} = \{U_j \to u(V)\}_{j \in J}$를
$\{u(V_i) \to u(V)\}_{i \in I}$와 자명하게 동치인
$\text{Cov}(\mathcal{D})$의 원소라 하자. $u(\mathcal{V})$는
$\mathcal{U}$의 세분이고 그 역도 마찬가지임에 유의하자. 따라서 Remark
\ref{sites-remark-both-refine-same-H0}에 의해
$H^0(u(\mathcal{V}), \mathcal{F}) = H^0(\mathcal{U}, \mathcal{F})$이다.
특히 $\mathcal{F}$가 층이면, 0차 {\v C}ech 코호몰로지 군으로 표현된 층
성질에 의해
$\mathcal{F}(u(V)) = H^0(u(\mathcal{V}), \mathcal{F})$이다.
$H^0(\mathcal{V}, u^p\mathcal{F}) =
H^0(u(\mathcal{V}), \mathcal{F})$이고, 방금 본 대로 이는
$\mathcal{F}(u(V)) = u^p\mathcal{F}(V)$와 같으므로, $\mathcal{F}$가 층이면
$u^p\mathcal{F}$도 층이라고 결론 내린다. 따라서 Lemma
\ref{sites-lemma-pushforward-sheaf}가 성립한다. Lemma
\ref{sites-lemma-adjoint-sheaves}는 바로 따른다.
\end{remark}


\section{사이트의 사상}
\label{sites-section-morphism-sites}

\begin{definition}
\label{sites-definition-morphism-sites}
$\mathcal{C}$와 $\mathcal{D}$를 사이트라 하자.
{\it 사이트의 사상} $f : \mathcal{D} \to \mathcal{C}$는 함자 $u_s$가
완전하도록 하는 연속 함자 $u : \mathcal{C} \to \mathcal{D}$로 주어진다.
\end{definition}

\noindent
함자 $u$의 방향은 사상 $f$의 방향과 {\it 반대}임에 유의하자.
$f \leftrightarrow u$가 사이트의 사상이면
$f^{-1} = u_s$와 $f_* = u^s$라는 표기를 사용한다.
함자 $f^{-1}$을 {\it 역상 함자}라 하고, 함자 $f_*$를
{\it 직접상 함자}라 한다. 위상수학에서와 마찬가지로 다음 수반성이 있다.
$$
\Mor_{\Sh(\mathcal{D})}(f^{-1}\mathcal{G}, \mathcal{F})
=
\Mor_{\Sh(\mathcal{C})}(\mathcal{G}, f_*\mathcal{F})
$$
이 정의의 동기는 다음 예에서 나온다.

\begin{example}
\label{sites-example-continuous-map}
$f : X  \to Y$를 위상공간들의 연속사상이라 하자. Example
\ref{sites-example-site-topological}의 사이트 $X_{Zar}$와 $Y_{Zar}$를
상기하자. 함자 $u : Y_{Zar} \to X_{Zar}$,
$V \mapsto f^{-1}(V)$를 생각하자. 열린 덮개의 역상은 열린 덮개이므로 이
함자는 분명히 연속이다. (실제로 이는 $X_{Zar}$와 $Y_{Zar}$의 덮개들의
집합을 어떻게 택했는지에 달려 있다. 그러나 어느 경우든 이 함자는
준연속이다. Remark \ref{sites-remark-quasi-continuous}를 보라.) 함자 $u^s$가
위상수학에서 통상적인 직접상 함자 $f_*$와 같음은 쉽게 확인된다. 따라서
$u_s$가 수반이고 통상적인 위상수학적 역상 함자 $f^{-1}$도 수반이므로,
표준 동형사상 $f^{-1} \cong u_s$를 얻는다. $f^{-1}$은 완전하므로
$u_s$도 완전하다고 결론 내린다. 따라서 $u$는 사이트의 사상
$f : X_{Zar} \to Y_{Zar}$를 정의한다. 이미 함자 $u_s, u^s$가 각각의
통상적인 개념과 일치함을 보았으므로 이 사상도 $f$로 나타낼 수 있다.
\end{example}

\begin{example}
\label{sites-example-finer-topology}
$\mathcal{C}$를 범주라 하자.
$$
\text{Cov}(\mathcal{C}) \supset \text{Cov}'(\mathcal{C})
$$
를 $\mathcal{C}$에 각각 사이트 구조를 주는, 공역이 고정된 사상족들의 두
집합이라 하자. $\text{Cov}(\mathcal{C})$에 대응하는 사이트를
$\mathcal{C}_\tau$로, $\text{Cov}'(\mathcal{C})$에 대응하는 사이트를
$\mathcal{C}_{\tau'}$로 나타내자. $\mathcal{C}$의 항등 함자가 사이트의
사상
$$
\epsilon : \mathcal{C}_\tau \longrightarrow \mathcal{C}_{\tau'}
$$
을 정의한다고 주장한다. 실제로 모든 $\tau'$-덮개는 $\tau$-덮개이므로
$\text{id} : \mathcal{C}_{\tau'} \to \mathcal{C}_\tau$는 연속이다.
따라서 함자 $\epsilon_* = \text{id}^s$는 바탕 준층들 위의 항등 함자이다.
그러므로 $\epsilon_*$의 왼쪽 수반 $\epsilon^{-1}$은 $\tau'$-층
$\mathcal{F}$를 $\mathcal{F}$의 $\tau$-층화로 보낸다(층화의 보편 성질에
의한다). $\tau'$-층들의 유한 극한은 준층들의 유한 극한과 일치하고(Lemma
\ref{sites-lemma-limit-sheaf}), $\tau$-층화는 유한 극한과 가환한다(Lemma
\ref{sites-lemma-sheafification-exact}). 따라서 $\epsilon^{-1}$은 좌완전이다.
$\epsilon^{-1}$은 왼쪽 수반이므로 우완전이기도 하다(Categories, Lemma
\ref{categories-lemma-exact-adjoint}). 따라서 $\epsilon^{-1}$은 완전하며,
Definition \ref{sites-definition-morphism-sites}의 모든 조건을 확인하였다.
\end{example}

\begin{lemma}
\label{sites-lemma-composition-morphisms-sites}
\begin{slogan}
사이트 함자들의 합성은 사이트 함자들의 연속성을 보존한다.
\end{slogan}
$\mathcal{C}_i$, $i = 1, 2, 3$을 사이트라 하자.
$u : \mathcal{C}_2 \to \mathcal{C}_1$과
$v : \mathcal{C}_3 \to \mathcal{C}_2$를 사이트의 사상들을 유도하는 연속
함자라 하자. 그러면 함자 $u \circ v : \mathcal{C}_3 \to \mathcal{C}_1$은
연속이며 사이트의 사상 $\mathcal{C}_1 \to \mathcal{C}_3$을 정의한다.
\end{lemma}

\begin{proof}
정의에서 $u \circ v$가 연속 함자임은 바로 따른다. 또한 분명히
$(u \circ v)^p = v^p \circ u^p$이고, 따라서
$(u \circ v)^s = v^s \circ u^s$이다. 그러므로 함자
$(u \circ v)_s$와 $u_s \circ v_s$는 모두 $(u \circ v)^s$의 왼쪽
수반이다. 따라서 $(u \circ v)_s \cong u_s \circ v_s$이고,
$(u \circ v)_s$는 완전 함자들의 합성이므로 완전하다고 결론 내린다.
\end{proof}

\begin{definition}
\label{sites-definition-composition-morphisms-sites}
$\mathcal{C}_i$, $i = 1, 2, 3$을 사이트라 하자.
$f : \mathcal{C}_1 \to \mathcal{C}_2$와
$g : \mathcal{C}_2 \to \mathcal{C}_3$을 각각 연속 함자
$u : \mathcal{C}_2 \to \mathcal{C}_1$과
$v : \mathcal{C}_3 \to \mathcal{C}_2$로 주어지는 사이트의 사상이라 하자.
{\it 합성} $g \circ f$는 함자 $u \circ v$에 대응하는 사이트의 사상이다.
\end{definition}

\noindent
이 상황에서 $(g \circ f)_* = g_* \circ f_*$이고
$(g \circ f)^{-1} = f^{-1} \circ g^{-1}$이다(Lemma
\ref{sites-lemma-composition-morphisms-sites}의 증명을 보라).

\begin{lemma}
\label{sites-lemma-directed-morphism}
$\mathcal{C}$와 $\mathcal{D}$를 사이트라 하자.
$u : \mathcal{C} \to \mathcal{D}$를 연속이라 하자. Section
\ref{sites-section-functoriality-PSh}의 모든 범주
$(\mathcal{I}_V^u)^{opp}$가 여과되었다고 가정하자. 그러면 $u$는 사이트의
사상 $\mathcal{D} \to \mathcal{C}$를 정의한다. 다시 말해 $u_s$는
완전하다.
\end{lemma}

\begin{proof}
$u_s$는 $u^s$의 왼쪽 수반이므로 $u_s$는 우완전이다. Categories, Lemma
\ref{categories-lemma-exact-adjoint}를 보라. 따라서 $u_s$가 좌완전임을
보이면 충분하다. 다시 말해 $u_s$가 유한 극한과 가환함을 보여야 한다.
범주들 $\mathcal{I}_Y^{opp}$가 여과되었으므로 $u_p$는 유한 극한과
가환한다. Categories, Lemma \ref{categories-lemma-directed-commutes}를
보라(이는 $\textit{PSh}$에서의 극한에 대한 기술도 사용한다. Section
\ref{sites-section-limits-colimits-PSh}를 보라). 층화도 유한 극한과 가환하므로
(Lemma \ref{sites-lemma-sheafification-exact}), $u_s = \# \circ u_p$에 의해
결론을 얻는다.
\end{proof}

\begin{proposition}
\label{sites-proposition-get-morphism}
$\mathcal{C}$와 $\mathcal{D}$를 사이트라 하자.
$u : \mathcal{C} \to \mathcal{D}$를 연속이라 하자. 또한 다음을 가정하자.
\begin{enumerate}
\item 범주 $\mathcal{C}$가 종대상 $X$를 가지며 $u(X)$가
$\mathcal{D}$의 종대상이다.
\item 범주 $\mathcal{C}$가 섬유곱들을 가지며 $u$가 이들과 가환한다.
\end{enumerate}
그러면 $u$는 사이트의 사상 $\mathcal{D} \to \mathcal{C}$를 정의한다.
다시 말해 $u_s$는 완전하다.
\end{proposition}

\begin{proof}
Lemmas \ref{sites-lemma-directed}와 \ref{sites-lemma-directed-morphism}에서 따른다.
\end{proof}

\begin{remark}
\label{sites-remark-explain-left-exact}
위 Proposition \ref{sites-proposition-get-morphism}의 조건들은 $u$가 좌완전,
즉 유한 극한과 가환한다는 것과 동치이다. Categories, Lemmas
\ref{categories-lemma-finite-limits-exist}와
\ref{categories-lemma-characterize-left-exact}를 보라. 예들에서는 종대상과
섬유곱으로 표현하는 편이 더 기하학적인 의미를 지니는 듯하므로 이 표현이 더
자연스러워 보인다.
\end{remark}

\noindent
Lemma \ref{sites-lemma-continuous-with-continuous-left-adjoint}는 연속 함자가
사이트의 사상을 유도함을 증명하는 또 다른 방법을 제공할 것이다.

\begin{remark}
\label{sites-remark-quasi-continuous-morphism-sites}
(처음 읽을 때는 건너뛰어도 된다.)
$\mathcal{C}$와 $\mathcal{D}$를 사이트라 하자. Definition
\ref{sites-definition-morphism-sites}와 유사하게, {\it 사이트의 준사상
$f : \mathcal{D} \to \mathcal{C}$}은 $u_s$가 완전하도록 하는 준연속 함자
$u : \mathcal{C} \to \mathcal{D}$로 주어진다고 한다(Remark
\ref{sites-remark-quasi-continuous}를 보라). 이 상황에서 Proposition
\ref{sites-proposition-get-morphism}의 유사 명제는 ``연속''을 ``준연속''으로,
``사상''을 ``준사상''으로 바꾸면 얻어진다. 증명은 글자 그대로 같다.
\end{remark}

\noindent
Definition \ref{sites-definition-morphism-sites}에서 $u_s$가 완전하다는 조건은
생략할 수 없다. 예를 들어 $u$가 연속이라고만 가정하면 다음 보조정리의
결론이 성립하지 않을 수 있다.

\begin{lemma}
\label{sites-lemma-morphism-of-sites-covering}
$f : \mathcal{D} \to \mathcal{C}$를 함자
$u : \mathcal{C} \to \mathcal{D}$로 주어지는 사이트의 사상이라 하자.
$\mathcal{D}$의 임의의 대상 $V$가 주어졌을 때, 모든 $j$에 대하여
$\mathcal{C}$의 어떤 대상 $U_j$로 가는 사상 $V_j \to u(U_j)$가 존재하는
덮개 $\{V_j \to V\}$가 존재한다.
\end{lemma}

\begin{proof}
$f^{-1} = u_s$가 완전하므로 $f^{-1}* = *$이다. 여기서 $*$는 층들의
범주의 종대상을 나타낸다(Example \ref{sites-example-singleton-sheaf}).
$f^{-1}* = u_s*$는 $u_p*$의 층화이므로, $(u_p*)(V_j)$가 공집합이 아닌
덮개 $\{V_j \to V\}$가 존재한다. $(u_p*)(V_j)$는 대상이 사상
$V_j \to u(U)$인 범주 $\mathcal{I}^u_{V_j}$ 위의 여극한이므로 보조정리가
따른다.
\end{proof}


\section{토포스}
\label{sites-section-topoi}

\noindent
여기서는 우리의 목적에 알맞은 토포스의 정의를 제시한다. 즉 토포스는 한
사이트 위의 층들의 범주이다. 토포스를 지정하려면 사이트만 지정하면 된다.
토포스와 사이트의 진정한 차이는 사상의 정의에 있다. 실제로 바탕 사이트의
사상에서 오지 않는 토포스의 사상이 많이 존재한다.

\begin{definition}[토포스]
\label{sites-definition-topos}
{\it 토포스}는 사이트 $\mathcal{C}$ 위의 층들의 범주
$\Sh(\mathcal{C})$이다.
\begin{enumerate}
\item $\mathcal{C}$, $\mathcal{D}$를 사이트라 하자.
$\Sh(\mathcal{D})$에서 $\Sh(\mathcal{C})$로 가는
{\it 토포스의 사상} $f$는 다음을 만족하는 두 함자
$f_* : \Sh(\mathcal{D}) \to \Sh(\mathcal{C})$와
$f^{-1} : \Sh(\mathcal{C}) \to \Sh(\mathcal{D})$로 주어진다.
\begin{enumerate}
\item 이중함자적으로
$$
\Mor_{\Sh(\mathcal{D})}(f^{-1}\mathcal{G}, \mathcal{F})
=
\Mor_{\Sh(\mathcal{C})}(\mathcal{G}, f_*\mathcal{F})
$$
가 성립하고,
\item 함자 $f^{-1}$은 유한 극한과 가환한다. 즉 좌완전이다.
\end{enumerate}
\item $\mathcal{C}$, $\mathcal{D}$, $\mathcal{E}$를 사이트라 하자.
토포스의 사상 $f :\Sh(\mathcal{D}) \to \Sh(\mathcal{C})$와
$g :\Sh(\mathcal{E}) \to \Sh(\mathcal{D})$가 주어졌을 때,
{\it 합성 $f\circ g$}는 함자
$(f \circ g)_* = f_* \circ g_*$와
$(f \circ g)^{-1} = g^{-1} \circ f^{-1}$로 정의되는 토포스의 사상이다.
\end{enumerate}
\end{definition}

\noindent
$\alpha : \mathcal{S}_1 \to \mathcal{S}_2$가 (아마도 ``큰'') 범주들의
동치라고 하자. $\mathcal{S}_1$, $\mathcal{S}_2$가 토포스라면
$f_* = \alpha$로 놓고 $f^{-1}$을 $\alpha$의 준역으로 놓아서 토포스의
사상 $f : \mathcal{S}_1 \to \mathcal{S}_2$를 얻는다. 더욱이 이 사상은
토포스들의 $2$-범주에서의 동치이다(Section
\ref{sites-section-2-category}를 보라). 따라서 $\mathcal{S}$가 어떤 사이트
위의 층들의 범주와 동치이면(반드시 그 범주와 같을 필요는 없다)
``$\mathcal{S}$는 토포스이다''라고 말해도 뜻이 통한다. 때때로 이와 같이
표기를 남용할 것이다.

\medskip\noindent
{\it 공 토포스}는 $\mathcal{C}$가 공범주인 사이트 $\mathcal{C}$ 위의
층들의 토포스이다. 때로 이 사이트를 $\emptyset$으로 쓴다. 이 사이트에는
유일한 층이 존재한다($\emptyset$은 대상을 갖지 않기 때문이다). 따라서
$\Sh(\emptyset)$은 대상 하나와 사상 하나를 갖는 범주와 동치이다.

\medskip\noindent
{\it 점 토포스}는 대상 하나 $pt$와 사상 하나 $\text{id}_{pt}$를 가지며
유일한 덮개가 $\{\text{id}_{pt}\}$인 사이트 $\mathcal{C}$ 위의 층들의
토포스이다. 이 사이트는 단순히 $pt$로 쓸 것이다. 층의 범주$ = $준층의
범주$ = $집합의 범주임은 분명하다. 식으로 쓰면
$\Sh(pt) = \textit{Sets}$이다.

\medskip\noindent
$\mathcal{C}$와 $\mathcal{D}$를 사이트라 하자.
$f : \Sh(\mathcal{D}) \to \Sh(\mathcal{C})$를 토포스의 사상이라 하자.
$f_*$는 모든 극한과 가환하고 $f^{-1}$은 모든 여극한과 가환함에 유의하자.
Categories, Lemma \ref{categories-lemma-adjoint-exact}를 보라. 특히 위 정의의
$f^{-1}$에 대한 조건은 $f^{-1}$이 완전함을 보장한다. 토포스의 사상은 흔히
Lemma \ref{sites-lemma-cocontinuous-morphism-topoi}나 다음 보조정리를 이용하여
구성한다.

\begin{lemma}
\label{sites-lemma-morphism-sites-topoi}
함자 $u : \mathcal{C} \to \mathcal{D}$에 대응하는 사이트의 사상
$f : \mathcal{D} \to \mathcal{C}$가 주어지면, 함자의 쌍
$(f^{-1} = u_s, f_* = u^s)$는 토포스의 사상
$\Sh(\mathcal{D}) \to \Sh(\mathcal{C})$이다.
\end{lemma}

\begin{proof}
Definition \ref{sites-definition-morphism-sites}에서 명백하다.
\end{proof}

\begin{remark}
\label{sites-remark-pt-topos}
토포스 $\Sh(pt)$를 주는 사이트는 많이 있다. 다음은 유용한 예이다.
$S$가 적어도 하나의 공집합이 아닌 원소를 포함하는 집합(들의 집합)이라고
하자. $\mathcal{S}$를 대상이 $S$의 원소들이고 사상이 임의의 집합 사상인
범주라 하자. $\mathcal{S}$가 섬유곱을 가진다고 가정하자. 예를 들어
$S = \mathcal{P}(\text{infinite set})$가 임의의 무한집합의 멱집합이면
이 조건이 성립한다(집합론 연습문제). 전사인 사상족을 덮개로 선언하여
$\mathcal{S}$에 사이트 구조를 주자(그리고 Sets, Section
\ref{sets-section-coverings-site}에서와 같이 충분히 큰 적절한 덮개족들의
집합을 택한다). $\Sh(\mathcal{S})$가 집합들의 범주와 동치라고 주장한다.

\medskip\noindent
먼저 $S$가 한원소집합인 $e \in S$를 포함하는 경우를 증명하자. 이 경우
토포스의 동치 $i : \Sh(pt) \to \Sh(\mathcal{S})$가 함자들
\begin{equation}
\label{sites-equation-sheaves-pt-sets}
i^{-1}\mathcal{F} = \mathcal{F}(e), \quad
i_*E = (U \mapsto \Mor_{\textit{Sets}}(U, E))
\end{equation}
로 주어진다. 실제로 $\mathcal{F}$를 $\mathcal{S}$ 위의 층이라 하자.
임의의 $U \in \Ob(\mathcal{S}) = S$에 대하여 덮개
$\{\varphi_u : e \to U\}_{u \in U}$를 찾을 수 있다. 여기서
$\varphi_u$는 $e$의 유일한 원소를 $u \in U$로 보낸다. 이 경우 층 조건은
$\mathcal{F}(U) = \prod_{u \in U} \mathcal{F}(e)$를 함의한다. 다시 말해
$\mathcal{F}(U) = \Mor_{\textit{Sets}}(U, \mathcal{F}(e))$이다. 더욱이 이
규칙은 제한 사상들과 호환된다. 따라서 함자
$$
i_* :
\textit{Sets} = \Sh(pt)
\longrightarrow
\Sh(\mathcal{S}), \quad
E \longmapsto (U \mapsto \Mor_{\textit{Sets}}(U, E))
$$
는 범주들의 동치이고, 그 역은 위에서 주어진 함자 $i^{-1}$이다.

\medskip\noindent
$\mathcal{S}$가 한원소집합을 포함하지 않는 경우에도 위에서 정의한 함자
$i_*$는 여전히 의미를 갖는다. 이 경우에도 동치임을 보이기 위해 공집합이 아닌
임의의 $\tilde e \in S$와 상이 한원소집합인 사상
$\varphi : \tilde e \to \tilde e$를 택하자. 임의의 층 $\mathcal{F}$에
대하여
$$
\mathcal{F}(e) :=
\Im(
\mathcal{F}(\varphi) :
\mathcal{F}(\tilde e)
\longrightarrow
\mathcal{F}(\tilde e)
)
$$
로 놓고 이것이 $i_*$의 준역임을 보이면 된다. 세부사항은 생략한다.
\end{remark}

\begin{remark}
\label{sites-remark-morphism-topoi-big}
(토포스의 사상과 관련된 집합론적 문제이다. 처음 읽을 때는 건너뛰어도 된다.)
위에서 정의한 토포스의 사상은 집합이 아니라 클래스이다. 다시 말해 이는
수학적 대상이 아니라 수학적 공식으로 주어진다. 주어진 두 토포스 사이의 모든
사상들의 모음을 생각할 수는 있지만, 이를 수학적 대상으로 도입하는 것은 좋은
생각이 아니다. 한편 $\mathcal{C}$와 $\mathcal{D}$가 주어진 사이트라고
하자. 함자 $\Phi : \mathcal{C} \to \Sh(\mathcal{D})$를 생각하자. 이러한
것은 집합, 다시 말해 수학적 대상이다. $\Phi$에 관하여 다음 질문들을
차례로 물을 수 있다.
\begin{enumerate}
\item $\mathcal{D}$ 위의 층 $\mathcal{F}$가 주어졌을 때 규칙
$U \mapsto \Mor_{\Sh(\mathcal{D})}(\Phi(U), \mathcal{F})$가
$\mathcal{C}$ 위의 층을 정의하는가? 그렇다면 이는 함자
$\Phi_* : \Sh(\mathcal{D}) \to \Sh(\mathcal{C})$를 정의한다.
\item $\Phi_*$가 왼쪽 수반을 가지는가? 그렇다면 이 왼쪽 수반을
$\Phi^{-1}$로 쓰자.
\item $\Phi^{-1}$은 완전한가?
\end{enumerate}
마지막 질문의 답도 여전히 ``그렇다''이면 토포스의 사상
$(\Phi_*, \Phi^{-1})$을 얻는다. 더욱이 토포스의 임의의 사상
$(f_*, f^{-1})$이 주어지면 $\Phi(U) = f^{-1}(h_U^\#)$로 놓아서 위와 같은
함자 $\Phi$를 얻으며, $f_* \cong \Phi_*$와
$f^{-1} \cong \Phi^{-1}$가 성립한다(수반성과 호환된다). 결론적으로
토포스의 사상 대신 $\Phi$들의 모음을 사용함으로써 (a) 토포스의 사상이라는
개념을 수학적 대상으로 바꾸었고, (b) $\Phi$들의 모음이 클래스들의 모음이
아닌 하나의 클래스를 이룬다. 물론 더 많은 것을 말할 수 있다. 예를 들어 위
조건 (2)와 (3)의 의미를 더 정확히 규명할 수 있다. Section
\ref{sites-section-points}에서 토포스의 점에 대하여 이를 수행한다.
\end{remark}

\begin{remark}
\label{sites-remark-quasi-continuous-morphism-topoi}
(처음 읽을 때는 건너뛰어도 된다.)
$\mathcal{C}$와 $\mathcal{D}$를 사이트라 하자. 사이트의 준사상
$f : \mathcal{D} \to \mathcal{C}$는(Remark
\ref{sites-remark-quasi-continuous-morphism-sites}를 보라) Lemma
\ref{sites-lemma-morphism-sites-topoi}에서와 정확히 같은 방식으로
$\Sh(\mathcal{D})$에서 $\Sh(\mathcal{C})$로 가는 토포스의 사상 $f$를
유도한다.
\end{remark}


\section{G-집합과 사상}
\label{sites-section-G-sets-morphisms}

\noindent
$\varphi : G \to H$를 군의 준동형이라 하자. Example
\ref{sites-example-site-on-group}과 Section
\ref{sites-section-example-sheaf-G-sets}에서와 같이 (적절한) 사이트
$\mathcal{T}_G$와 $\mathcal{T}_H$를 택하자.
$u : \mathcal{T}_H \to \mathcal{T}_G$를 $H$-집합 $U$에 바탕 집합은 같고,
$g \in G$와 $\xi \in U$에 대하여 $G$-작용이
$g \cdot \xi = \varphi(g)\xi$로 정의된 $G$-집합 $U_\varphi$를 대응시키는
함자라 하자. $u$가 유한 극한과 가환하며 연속임은
분명하다\footnote{집합론적 주석: 먼저 $\mathcal{T}_H$를 택한다. 다음으로
$\mathcal{T}_G$가 $u(\mathcal{T}_H)$를 포함하고 $\mathcal{T}_H$의 모든
덮개가 $\mathcal{T}_G$의 덮개에 대응하도록 택한다.
이는 Sets, Lemmas
\ref{sets-lemma-sets-with-group-action},
\ref{sets-lemma-what-is-in-it-G-sets},
\ref{sets-lemma-coverings-site}에 의해 가능하다.}.
Proposition \ref{sites-proposition-get-morphism}과 Lemma
\ref{sites-lemma-morphism-sites-topoi}를 적용하면 $\varphi$에 연관된 토포스의 사상
$$
f : \Sh(\mathcal{T}_G) \longrightarrow \Sh(\mathcal{T}_H)
$$
을 얻는다. Proposition \ref{sites-proposition-sheaves-on-group}을 사용하면 수반
함자의 쌍
$$
f_* : G\textit{-Sets} \longrightarrow H\textit{-Sets}, \quad
f^{-1} : H\textit{-Sets} \longrightarrow G\textit{-Sets}.
$$
을 얻음을 알 수 있다. 이 경우 이 함자들이 무엇인지 계산해 보자.

\medskip\noindent
먼저 $f_*$의 공식을 계산한다. $G$-집합 $S$가 주어졌을 때 이에 대응하는
$\mathcal{T}_G$ 위의 층 $\mathcal{F}_S$는 규칙
$\mathcal{F}_S(U) = \Mor_G(U, S)$로 주어짐을 상기하자. 한편
$\mathcal{T}_H$ 위의 층 $\mathcal{G}$가 주어졌을 때 이에 대응하는
$H$-집합은 규칙 $\mathcal{G}({}_HH)$로 주어진다. 따라서
$$
f_*S = \Mor_{G\textit{-Sets}}(({}_HH)_\varphi, S).
$$
이다. 이를 조금 더 구체적으로 계산하면
% P01-ERRATA: P01-E0068 -- the frozen equivariance formula has ill-typed a(gh).
$$
f_*S = \{ a : H \to S \mid a(gh) = ga(h) \}
$$
를 얻으며, 왼쪽 $H$-작용은 임의의 원소 $a \in f_*S$에 대하여
$(h \cdot a)(h') = a(h'h)$로 주어진다.

\medskip\noindent
다음으로 $f^{-1}$을 명시적으로 계산한다. $\mathcal{T}_G$와
$\mathcal{T}_H$의 위상은 준표준이므로 모든 표현가능 준층이 층임에
유의하자. 더욱이 $\mathcal{T}_H$의 대상 $V$가 주어지면
$f^{-1}h_V$는 $h_{u(V)}$와 같다(Lemma
\ref{sites-lemma-pullback-representable-sheaf}를 보라). 따라서 표현가능 층에
대해서는 $f^{-1}S = S_\varphi$이다. $\mathcal{T}_H$ 위의 모든 층은
표현가능 층들의 쌍대곱이므로 일반적으로도 참이라고 결론 내린다. 따라서
임의의 $H$-집합 $T$에 대하여
$$
f^{-1}T = T_\varphi.
$$
이다. $f^{-1}$과 $f_*$ 사이의 수반성은 위와 같은 $f_*S$에 대한 공식
$$
\Mor_{G\textit{-Sets}}(T_\varphi, S) =
\Mor_{H\textit{-Sets}}(T, f_*S)
$$
으로 나타난다. 이는 직접 증명할 수 있다. 더욱이 이로부터
$(f^{-1}, f_*)$가 수반쌍을 이루고 $f^{-1}$이 완전함은 분명하다. 따라서
위의 방법 대신 토포스의 사상 $f : G\textit{-Sets} \to H\textit{-Sets}$를
이 방식으로 직접 정의할 수도 있다.


\section{준콤팩트 대상과 여극한}
\label{sites-section-quasi-compact}

\noindent
위상공간의 경우와 같은 언어를 사용할 수 있도록 다음 용어를 도입한다.

\begin{definition}
\label{sites-definition-quasi-compact}
$\mathcal{C}$를 사이트라 하자. $\mathcal{C}$의 대상 $U$에 대하여,
$\mathcal{C}$의 덮개 $\mathcal{U} = \{U_i \to U\}_{i \in I}$가 주어질
때마다 다른 덮개 $\mathcal{V} = \{V_j \to U\}_{j \in J}$와
$\text{id} : U \to U$, $\alpha : J \to I$, $V_j \to U_{\alpha(j)}$로
주어지는, 공역이 고정된 사상족들의 사상
$\mathcal{V} \to \mathcal{U}$가 존재하고(Definition
\ref{sites-definition-morphism-coverings}을 보라), $\alpha$의 상이 $I$의 유한
부분집합이면 이 대상을 {\it 준콤팩트}라고 한다.
\end{definition}

\noindent
물론 통상적인 개념만으로도 $U$가 준콤팩트임을 결론 내리기에 충분하다.

\begin{lemma}
\label{sites-lemma-conclude-quasi-compact}
$\mathcal{C}$를 사이트라 하고 $U$를 $\mathcal{C}$의 대상이라 하자.
다음 조건들을 생각하자.
\begin{enumerate}
\item $U$는 준콤팩트이다.
\item $\mathcal{C}$의 모든 덮개 $\{U_i \to U\}_{i \in I}$에 대하여
$\mathcal{U}$를 세분하는 $\mathcal{C}$의 유한 덮개
$\{V_j \to U\}_{j = 1, \ldots, m}$가 존재한다.
\item $\mathcal{C}$의 모든 덮개 $\{U_i \to U\}_{i \in I}$에 대하여
$\{U_i \to U\}_{i \in I'}$가 $\mathcal{C}$의 덮개가 되는 유한 부분집합
$I' \subset I$가 존재한다.
\end{enumerate}
항상 (3) $\Rightarrow$ (2) $\Rightarrow$ (1)이지만, 역함의들은 일반적으로
성립하지 않는다.
\end{lemma}

\begin{proof}
함의들은 정의에서 바로 따른다. $X = [0, 1] \subset \mathbf{R}$을
(통상적인 $\epsilon$-$\delta$ 위상을 갖춘) 위상공간으로 보자.
$\mathcal{C}$를 $X$의 열린 부분공간들을 대상으로 하고 포함사상을 사상으로
하며 통상적인 열린 덮개를 갖는 범주라 하자(Example
\ref{sites-example-site-topological}과 비교하라). 다만 $\mathcal{C}$의 덮개
개념을 바꾸어 $\{U \to U\}$ 꼴의 덮개를 제외한 모든 유한 덮개를
배제한다. 이것이 Definition \ref{sites-definition-site}의 의미에서 사이트가
됨은 쉽게 알 수 있다. 이 사이트에서 대상 $X = U = [0, 1]$은 Definition
\ref{sites-definition-quasi-compact}의 의미에서 준콤팩트이지만 $U$는 (2)를
만족하지 않는다. (2)는 만족하지만 (3)은 만족하지 않는 예를 만드는 일은
독자에게 맡긴다.
\end{proof}

\noindent
다음은 준콤팩트 대상의 토포스 이론적 의미이다.

\begin{lemma}
\label{sites-lemma-quasi-compact}
$\mathcal{C}$를 사이트라 하고 $U$를 $\mathcal{C}$의 대상이라 하자.
다음은 서로 동치이다.
\begin{enumerate}
\item $U$는 준콤팩트이다.
\item 층들의 모든 전사사상
$\coprod_{i \in I} \mathcal{F}_i \to h_U^\#$에 대하여
$\coprod_{i \in J} \mathcal{F}_i \to h_U^\#$가 전사인 유한 부분집합
$J \subset I$가 존재한다.
\end{enumerate}
\end{lemma}

\begin{proof}
(1)을 가정하고 $\coprod_{i \in I} \mathcal{F}_i \to h_U^\#$를
전사사상이라 하자. 그러면 $\text{id}_U$는 $U$ 위에서 $h_U^\#$의
절단이다. 따라서 덮개 $\{U_a \to U\}_{a \in A}$와 각 $a \in A$에
대하여 $U_a$ 위의 $\coprod_{i \in I} \mathcal{F}_i$의 절단 $s_a$가
존재하여 $\text{id}_U$로 사상된다. 준층의 쌍대곱을 층화하여 쌍대곱을
구성하는 법에 의해(Lemma \ref{sites-lemma-colimit-sheaves}), 각 $a$에 대하여
덮개 $\{U_{ab} \to U_a\}_{b \in B_a}$가 존재하고, 모든 $b \in B_a$에
대하여 $\iota(b) \in I$와 $U_{ab}$ 위의 $\mathcal{F}_{\iota(b)}$의 절단
$s_{b}$가 존재하여 $\text{id}_U|_{U_{ab}}$로 사상된다. 따라서 덮개
$\{U_a \to U\}_{a \in A}$를
$\{U_{ab} \to U\}_{a \in A, b \in B_a}$로 바꾼 뒤에는 사상
$\iota : A \to I$와 각 $a \in A$에 대하여 $U_a$ 위의
$\mathcal{F}_{\iota(a)}$의 절단 $s_a$가 존재하여 $\text{id}_U$로
사상된다고 가정할 수 있다. $U$가 준콤팩트이므로 덮개
$\{V_c \to U\}_{c \in C}$, 유한 상을 갖는 사상 $\alpha : C \to A$,
그리고 $U$ 위의 사상 $V_c \to U_{\alpha(c)}$가 존재한다. 그러면
$J = \Im(\iota \circ \alpha) \subset I$가 원하는 부분집합이다. 실제로
$\coprod_{c \in C} h_{V_c}^\# \to h_U^\#$는 전사이고(Lemma
\ref{sites-lemma-covering-surjective-after-sheafification}),
$\coprod_{i \in J} \mathcal{F}_i \to h_U^\#$를 통하여 분해된다. (여기서는
% P01-ERRATA: P01-E0069 -- the frozen intermediate representable lacks sheafification.
합성 $h_{V_c}^\# \to h_{U_{\alpha(c)}}
\xrightarrow{s_{\alpha(c)}} \mathcal{F}_{\iota(\alpha(c))} \to h_U^\#$가
$s_{\alpha(c)}$가 $\text{id}_U|_{U_{\alpha(c)}}$로 사상된다는 사실 때문에
사상 $V_c \to U$에서 오는 사상 $h_{V_c}^\# \to h_U^\#$임을 사용한다.)

\medskip\noindent
(2)를 가정하자. $\{U_i \to U\}_{i \in I}$를 덮개라 하자. Lemma
\ref{sites-lemma-covering-surjective-after-sheafification}에 의해
$\coprod_{i \in I} h_{U_i}^\# \to h_U^\#$는 전사이다. 따라서
$\coprod_{j \in J} h_{U_j}^\# \to h_U^\#$가 전사인 유한 부분집합
$J \subset I$를 찾을 수 있다. 위와 같이 논증하면 $\mathcal{C}$에서
$U$의 덮개 $\{V_c \to U\}_{c \in C}$와 사상 $\iota : C \to J$를 찾아
% P01-ERRATA: P01-E0070 -- the frozen phrase "a section of s_c of" is malformed.
$\text{id}_U$가 $V_c$ 위에서 $h_{U_{\iota(c)}}^\#$의 절단 $s_c$로
올라가게 할 수 있다. 덮개를 더 세분하면 $\text{id}_U$로 사상되는
$s_c \in h_{U_{\iota(c)}}(V_c)$라고 가정할 수 있다. 그러면
$s_c : V_c \to U_{\iota(c)}$는 $U$ 위의 사상이므로 결론을 얻는다.
\end{proof}

\noindent
위 보조정리는 다음 정의의 동기를 준다.

\begin{definition}
\label{sites-definition-quasi-compact-topos}
토포스 $\Sh(\mathcal{C})$의 대상 $\mathcal{F}$에 대하여,
$\Sh(\mathcal{C})$의 임의의 전사사상
$\coprod_{i \in I} \mathcal{F}_i \to \mathcal{F}$가 주어질 때마다
$\coprod_{i \in J} \mathcal{F}_i \to \mathcal{F}$가 전사인 유한 부분집합
$J \subset I$가 존재하면 이 대상을 {\it 준콤팩트}라고 한다.
토포스 $\Sh(\mathcal{C})$의 종대상 $*$가 준콤팩트 대상이면 이 토포스를
{\it 준콤팩트}라고 한다.
\end{definition}

\noindent
Lemma \ref{sites-lemma-quasi-compact}에 의해, 사이트 $\mathcal{C}$가 종대상
$X$를 가지면 $\Sh(\mathcal{C})$가 준콤팩트일 필요충분조건은 $X$가
준콤팩트인 것이다.

\begin{lemma}
\label{sites-lemma-image-of-quasi-compact}
\begin{slogan}
층의 전사사상은 준콤팩트성을 전달한다.
\end{slogan}
$\mathcal{C}$를 사이트라 하자.
\begin{enumerate}
\item $U \to V$가 $\mathcal{C}$의 사상이고
$h_U^\# \to h_V^\#$가 전사이며 $U$가 준콤팩트이면 $V$도 준콤팩트이다.
\item $\mathcal{F} \to \mathcal{G}$가 집합의 층들의 전사사상이고
$\mathcal{F}$가 준콤팩트이면 $\mathcal{G}$도 준콤팩트이다.
\end{enumerate}
\end{lemma}

\begin{proof}
생략한다.
\end{proof}

\begin{lemma}
\label{sites-lemma-coproduct-quasi-compact}
\begin{slogan}
% P01-ERRATA: P01-E0071 -- the frozen source misspells "coproducts".
층들의 유한 쌍대곱은 준콤팩트성을 보존한다.
\end{slogan}
$\mathcal{C}$를 사이트라 하자. $n \geq 1$이고
$\mathcal{F}_1, \ldots, \mathcal{F}_n$이 $\mathcal{C}$ 위의 준콤팩트
층들이면 $\coprod_{i = 1, \ldots, n} \mathcal{F}_i$도 준콤팩트이다.
\end{lemma}

\begin{proof}
생략한다.
\end{proof}

\noindent
다음 두 보조정리는 사이트에 대한 Sheaves, Lemma
\ref{sheaves-lemma-directed-colimits-sections}의 유사 명제이다.

\begin{lemma}
\label{sites-lemma-directed-colimits-sections}
$\mathcal{C}$를 사이트라 하자.
$\mathcal{I} \to \Sh(\mathcal{C})$, $i \mapsto \mathcal{F}_i$를 집합의
층들의 여과된 도표라 하자. $U \in \Ob(\mathcal{C})$라 하고 표준 사상
$$
\Psi :
\colim_i \mathcal{F}_i(U)
\longrightarrow
\left(\colim_i \mathcal{F}_i\right)(U)
$$
을 생각하자. 위에서 도입한 용어를 사용하면 다음이 성립한다.
\begin{enumerate}
\item 모든 전이사상이 단사이면 임의의 $U$에 대하여 $\Psi$는 단사이다.
\item $U$가 준콤팩트이면 $\Psi$는 단사이다.
\item $U$가 준콤팩트이고 모든 전이사상이 단사이면 $\Psi$는 동형사상이다.
\item $U$가 $J$가 유한이고 모든 $j, j' \in J$에 대하여
$U_j \times_U U_{j'}$가 준콤팩트인 덮개
$\{U_j \to U\}_{j \in J}$들의 공종계를 가지면 $\Psi$는 전단사이다.
\end{enumerate}
\end{lemma}

\begin{proof}
모든 전이사상이 단사라고 가정하자. 이 경우 준층
$\mathcal{F}' : V \mapsto \colim_i \mathcal{F}_i(V)$는 분리되어 있다
(Definition \ref{sites-definition-separated}를 보라). Lemma
\ref{sites-lemma-colimit-sheaves}에 의해
$(\mathcal{F}')^\# = \colim_i \mathcal{F}_i$이다. Theorem
\ref{sites-theorem-plus}에 의해 $\mathcal{F}' \to (\mathcal{F}')^\#$가 단사임을
알 수 있다. 이로써 (1)이 증명된다.

\medskip\noindent
$U$가 준콤팩트라고 가정하자. $s \in \mathcal{F}_i(U)$와
$s' \in \mathcal{F}_{i'}(U)$가 왼쪽의 원소들을 주며, 이들이 $\Psi$에
의해 같은 상으로 간다고 하자. 이는 덮개
$\{U_a \to U\}_{a \in A}$를 택할 수 있고, 각 $a \in A$에 대하여
$i_a \in I$, $i_a \geq i$, $i_a \geq i'$인 지표가 존재하여
$\varphi_{ii_a}(s) = \varphi_{i'i_a}(s')$가 된다는 뜻이다. $U$가
준콤팩트이므로 덮개 $\{V_b \to U\}_{b \in B}$, 유한 상을 갖는 사상
$\alpha : B \to A$, 그리고 $U$ 위의 사상 $V_b \to U_{\alpha(b)}$를
택할 수 있다. $i''\in I$를 모든 $i_{\alpha(b)}$에 대하여 $\geq$인 것으로 택하자.
$\alpha$의 상이 유한이므로 이는 가능하다. 그러면 모든 $b \in B$에 대하여
% P01-ERRATA: P01-E0072 -- the frozen second terms use s instead of s-prime.
$\varphi_{ii''}(s)$와 $\varphi_{i'i''}(s)$가 $V_b$ 위에서 일치하고, 따라서
$\varphi_{ii''}(s) = \varphi_{i'i''}(s)$라고 결론 내린다. 이로써 (2)가
증명된다.

\medskip\noindent
$U$가 준콤팩트이고 모든 전이사상이 단사라고 가정하자. $s$를 $\Psi$의
공역의 원소라 하자. 덮개 $\{U_a \to U\}_{a \in A}$가 존재하고 각
$a \in A$에 대하여 지표 $i_a \in I$와 절단
$s_a \in \mathcal{F}_{i_a}(U_a)$가 존재하여 모든 $a \in A$에 대하여
$s|_{U_a}$가 $s_a$에서 온다. $U$가 준콤팩트이므로 덮개
$\{V_b \to U\}_{b \in B}$, 유한 상을 갖는 사상 $\alpha : B \to A$,
그리고 $U$ 위의 사상 $V_b \to U_{\alpha(b)}$를 택할 수 있다.
$i \in I$를 모든 $i_{\alpha(b)}$에 대하여 $\geq$인 것으로 택하자. $\alpha$의 상이
유한이므로 이는 가능하다. (1)에 의해 절단
$s_b = \varphi_{i_{\alpha(b)} i}(s_{\alpha(b)})|_{V_b}$들은
$V_b \times_U V_{b'}$ 위에서 일치한다. 따라서 이들은
$s' \in \mathcal{F}_i(U)$인 절단으로 붙고, 이 절단은 $\Psi$에 의해
$s$로 간다. 이로써 (3)이 증명된다.

\medskip\noindent
(4)의 가정을 하자. Lemma \ref{sites-lemma-conclude-quasi-compact}에 의해 대상
$U$는 준콤팩트이므로 (2)에 따라 $\Psi$는 단사이다. 전사성을 증명하기 위해
$s$를 $\Psi$의 공역의 원소라 하자. 가정에 따라 유한 덮개
$\{U_j \to U\}_{j = 1, \ldots, m}$가 존재하며, 모든
$1 \leq j, j' \leq m$에 대하여 $U_j \times_U U_{j'}$가 준콤팩트이고, 각
$j$에 대하여 지표 $i_j \in I$와 $s_j \in \mathcal{F}_{i_j}(U_j)$가
존재하여 모든 $j$에 대해 $s|_{U_j}$가 $s_j$의 상이다.
$U_j \times_U U_{j'}$가 준콤팩트이므로 (2)를 적용하면
$i_{jj'} \in I$, $i_{jj'} \geq i_j$, $i_{jj'} \geq i_{j'}$인 지표가
존재하여 $\varphi_{i_ji_{jj'}}(s_j)$와
$\varphi_{i_{j'}i_{jj'}}(s_{j'})$가 $U_j \times_U U_{j'}$ 위에서
일치한다. 모든 $i_{jj'}$ 이상인 지표 $i \in I$를 택하자. 그러면
$\mathcal{F}_i$의 절단들 $\varphi_{i_ji}(s_j)$가 $U$ 위의
$\mathcal{F}_i$의 절단으로 붙는다. 이 절단은 요구한 대로 원소 $s$로 간다.
\end{proof}

\begin{lemma}
\label{sites-lemma-directed-colimits-global-sections}
$\mathcal{C}$를 사이트라 하자.
$\mathcal{I} \to \Sh(\mathcal{C})$, $i \mapsto \mathcal{F}_i$를 집합의
층들의 여과된 도표라 하자. 표준 사상
$$
\Psi :
\colim_i \Gamma(\mathcal{C}, \mathcal{F}_i)
\longrightarrow
\Gamma(\mathcal{C}, \colim_i \mathcal{F}_i)
$$
을 생각하자. 다음이 성립한다.
\begin{enumerate}
\item 모든 전이사상이 단사이면 $\Psi$는 단사이다.
\item $\Sh(\mathcal{C})$가 준콤팩트이면 $\Psi$는 단사이다.
\item $\Sh(\mathcal{C})$가 준콤팩트이고 모든 전이사상이 단사이면
$\Psi$는 동형사상이다.
\item 다음 성질들을 갖는 집합 $S \subset \Ob(\Sh(\mathcal{C}))$가
존재한다고 가정하자.
\begin{enumerate}
\item 모든 전사사상 $\mathcal{F} \to *$에 대하여
$\mathcal{K} \in S$와 사상 $\mathcal{K} \to \mathcal{F}$가 존재하여
$\mathcal{K} \to *$가 전사이다.
\item $\mathcal{K} \in S$에 대하여 곱
$\mathcal{K} \times \mathcal{K}$는 준콤팩트이다.
\end{enumerate}
그러면 $\Psi$는 전단사이다.
\end{enumerate}
\end{lemma}

\begin{proof}
(1)의 증명. 모든 전이사상이 단사라고 가정하자. 이 경우 준층
$\mathcal{F}' : V \mapsto \colim_i \mathcal{F}_i(V)$는 분리되어 있다
(Definition \ref{sites-definition-separated}를 보라). Lemma
\ref{sites-lemma-colimit-sheaves}에 의해
$(\mathcal{F}')^\# = \colim_i \mathcal{F}_i$이다. Theorem
\ref{sites-theorem-plus}에 의해 $\mathcal{F}' \to (\mathcal{F}')^\#$가
단사임을 알 수 있다. 이로써 (1)이 증명된다.

\medskip\noindent
(2)의 증명. $\Sh(\mathcal{C})$가 준콤팩트라고 가정하자.
$\Sh(\mathcal{C})$의 모든 $\mathcal{F}$에 대하여
$\Gamma(\mathcal{C}, \mathcal{F}) = \Mor(*, \mathcal{F})$임을 상기하자.
$a_i, b_i : * \to \mathcal{F}_i$라 하고, $i' \geq i$에 대하여 계의
전이사상과 합성한 사상들을 $a_{i'}, b_{i'} : * \to \mathcal{F}_{i'}$로
나타내자. $a = \colim_{i' \geq i} a_{i'}$로 놓고 $b$도 마찬가지로 놓자.
$i' \geq i$에 대하여
$$
E_{i'} = \text{Equalizer}(a_{i'}, b_{i'}) \subset *
\quad\text{and}\quad
E = \text{Equalizer}(a, b) \subset *
$$
로 나타내자. Categories, Lemma
\ref{categories-lemma-directed-commutes}에 의해
$E = \colim_{i' \geq i} E_{i'}$이다. 따라서
$\coprod_{i' \geq i} E_{i'} \to E$는 층들의 전사사상이다. 그러므로
$E = *$, 즉 $a = b$이면 $*$가 준콤팩트이므로 어떤 $i' \geq i$에 대하여
$E_{i'} = *$이고, 어떤 $i' \geq i$에 대하여
$a_{i'} = b_{i'}$라고 결론 내린다. 이로써 (2)가 증명된다.

\medskip\noindent
(3)의 증명. $\Sh(\mathcal{C})$가 준콤팩트이고 모든 전이사상이 단사라고
가정하자. $a : * \to \colim \mathcal{F}_i$를 사상이라 하자. 그러면
$E_i = a^{-1}(\mathcal{F}_i) \subset *$는 부분층이고
$\colim E_i = *$이다(위의 인용에 의한다). 따라서 어떤 $i$에 대하여
$E_i = *$이고, 요구한 대로 $a$의 상은 $\mathcal{F}_i$에 포함된다.

\medskip\noindent
(4)의 증명. $S \subset \Ob(\Sh(\mathcal{C}))$가 (4)(a), (b)를
만족한다고 하자. (4)(a)를 $\text{id} : * \to *$에 적용하면
$\mathcal{K} \to *$가 전사인 $\mathcal{K} \in S$가 존재함을 알 수 있다.
사상 $\mathcal{K} \times \mathcal{K} \to \mathcal{K} \to *$는 전사이다.
(4)(b)와 Lemma \ref{sites-lemma-image-of-quasi-compact}에 의해
$\mathcal{K}$와 $\Sh(\mathcal{C})$가 준콤팩트라고 결론 내린다. 따라서
(2)에 의해 $\Psi$는 단사이다. $\mathcal{F} = \colim \mathcal{F}_i$로
놓자. $s : * \to \mathcal{F}$를 여극한의 대역 절단이라 하자.
$\coprod \mathcal{F}_i \to \mathcal{F}$가 전사이므로 사영사상
$$
\coprod\nolimits_{i \in I} * \times_{s, \mathcal{F}} \mathcal{F}_i \to *
$$
은 전사이다. (4)(a)에 의해 $\mathcal{K} \in S$와 사상
$\mathcal{K} \to \coprod_{i \in I} * \times_{s, \mathcal{F}} \mathcal{F}_i$를
얻으며 $\mathcal{K} \to *$는 전사이다. $\mathcal{K}$가 준콤팩트이므로
어떤 유한 부분집합 $I' \subset I$에 대한 분해
$\mathcal{K} \to \coprod_{i' \in I'} * \times_{s, \mathcal{F}} \mathcal{F}_{i'}$를
얻는다. $i \in I$를 유한 부분집합 $I'$의 상계라 하자. 전이사상들은 사상
$\coprod_{i' \in I'} \mathcal{F}_{i'} \to \mathcal{F}_i$를 정의한다.
이는 다시 사상 $\mathcal{K} \to * \times_{s, \mathcal{F}} \mathcal{F}_i$를
준다. 다시 말해 $\mathcal{K} \to *$가 전사인 $\mathcal{K} \in S$와
가환도식
$$
\xymatrix{
\mathcal{K} \times \mathcal{K} \ar@<1ex>[r] \ar@<-1ex>[r] &
\mathcal{K} \ar[d] \ar[r] & {*} \ar[d]^s \\
& \mathcal{F}_i \ar[r] & \mathcal{F} \ar@{=}[r] & \colim \mathcal{F}_i
}
$$
을 얻는다. 이 도식의 윗줄이 여등화자임에 유의하자. 따라서 $i$를 늘린 뒤
유도되는 두 사상 $a_i, b_i : \mathcal{K} \times \mathcal{K} \to \mathcal{F}_i$가
같음을 보이면 충분하다.
% P01-ERRATA: P01-E0073 -- the frozen phrase "This is done shown" is malformed.
이는 (2)의 증명에서와 정확히 같은 논증으로 다음 문단에서 보인다. 독자는
나머지 증명을 건너뛰기를 권한다.

\medskip\noindent
$i' \geq i$에 대하여 $a_i, b_i$와 계의 전이사상의 합성을
$a_{i'}, b_{i'} : \mathcal{K} \times \mathcal{K} \to \mathcal{F}_{i'}$로
나타내자. $a = \colim_{i' \geq i} a_{i'} :
\mathcal{K} \times \mathcal{K} \to \mathcal{F}$로 놓고 $b$도 마찬가지로
놓자. 위 도식의 가환성에 의해 $a = b$이다. $i' \geq i$에 대하여
$$
E_{i'} = \text{Equalizer}(a_{i'}, b_{i'}) \subset
\mathcal{K} \times \mathcal{K}
\quad\text{and}\quad
E = \text{Equalizer}(a, b) \subset
\mathcal{K} \times \mathcal{K}
$$
로 나타내자. Categories, Lemma
\ref{categories-lemma-directed-commutes}에 의해
$E = \colim_{i' \geq i} E_{i'}$이다. 따라서
$\coprod_{i' \geq i} E_{i'} \to E$는 층들의 전사사상이다. $a = b$이므로
$E = \mathcal{K} \times \mathcal{K}$이다.
$\mathcal{K} \times \mathcal{K}$는 (4)(b)에 의해 준콤팩트이므로 어떤
$i' \geq i$에 대하여 $E_{i'} = \mathcal{K} \times \mathcal{K}$이고,
어떤 $i' \geq i$에 대하여 $a_{i'} = b_{i'}$라고 결론 내린다.
\end{proof}

\begin{remark}
\label{sites-remark-stronger-conditions}
$\mathcal{C}$를 사이트라 하자. Lemma
\ref{sites-lemma-directed-colimits-global-sections}의 (4)의 가정들을 보장하는
방법은 여럿 있다. 몇 가지를 제시한다.
\begin{enumerate}
\item 다음 성질들을 갖는 집합 $\mathcal{B} \subset \Ob(\mathcal{C})$가
존재한다고 가정하자.
\begin{enumerate}
\item 모든 전사사상 $\mathcal{F} \to *$에 대하여 $m \geq 0$과
$U_1, \ldots, U_m \in \mathcal{B}$가 존재하여 $\mathcal{F}(U_j)$가
공집합이 아니고 $\coprod h_{U_j}^\# \to *$가 전사이다.
\item $U, U' \in \mathcal{B}$에 대하여 층
$h_U^\# \times h_{U'}^\#$는 준콤팩트이다.
\end{enumerate}
\item 다음 성질들을 갖는 집합 $\mathcal{B} \subset \Ob(\mathcal{C})$가
존재한다고 가정하자.
\begin{enumerate}
\item $m \geq 0$과 $U_1, \ldots, U_m \in \mathcal{B}$가 존재하여
$\coprod h_{U_j}^\# \to *$가 전사이다.
\item $U \in \mathcal{B}$에 대하여 $U$의 임의의 덮개는
$U_j \in \mathcal{B}$인 유한 덮개
$\{U_j \to U\}_{j = 1, \ldots, m}$로 세분된다.
\item $U, U' \in \mathcal{B}$에 대하여 $m \geq 0$,
$U_1, \ldots, U_m \in \mathcal{B}$와 사상 $U_j \to U$, $U_j \to U'$가
존재하여 $\coprod h_{U_j}^\# \to h_U^\# \times h_{U'}^\#$가 전사이다.
\end{enumerate}
\item 다음을 가정하자.
\begin{enumerate}
\item $\Sh(\mathcal{C})$는 준콤팩트이다.
\item $\mathcal{C}$의 모든 대상은 준콤팩트 대상들로 이루어진 덮개를 가진다.
\item $U$와 $U'$가 준콤팩트이면 층
$h_U^\# \times h_{U'}^\#$는 준콤팩트이다.
\end{enumerate}
\end{enumerate}
(1)과 (2)의 경우에는 $S \subset \Ob(\Sh(\mathcal{C}))$를
$U \in \mathcal{B}$인 층들 $h_U^\#$의 유한 쌍대곱들의 집합과 같게 놓는다.
(3)의 경우에는 $S \subset \Ob(\Sh(\mathcal{C}))$를 준콤팩트인
$U \in \Ob(\mathcal{C})$에 대한 층들 $h_U^\#$의 유한 쌍대곱들의 집합과
같게 놓는다.
\end{remark}

\noindent
나중에 다음과 같이 여극한에서 일어날 수 있는 일에 대한 상계가 필요하다.

\begin{lemma}
\label{sites-lemma-colimit-over-ordinal-sections}
$\mathcal{C}$를 사이트라 하고 $\beta$를 서수라 하자.
$\beta \to \Sh(\mathcal{C})$, $\alpha \mapsto \mathcal{F}_\alpha$를
$\beta$ 위의 층들의 계라 하자. $U \in \Ob(\mathcal{C})$에 대하여 표준 사상
$$
\colim_{\alpha < \beta} \mathcal{F}_\alpha(U)
\longrightarrow
\left(\colim_{\alpha < \beta} \mathcal{F}_\alpha \right)(U)
$$
을 생각하자. $\beta$의 공종성이 충분히 크면 이 사상은 모든 $U$에 대하여
전단사이다.
\end{lemma}

\begin{proof}
왼쪽은 준층들의 범주에서 취한 여극한 $\mathcal{F}_{\colim}$의 $U$에서의
값이다. Section \ref{sites-section-limits-colimits-PSh}를 보라.
$\colim_{\alpha < \beta} \mathcal{F}_\alpha$가
$\mathcal{F}_{\colim}$의 층화 $\mathcal{F}_{\colim}^\#$임을 상기하자.
Lemma \ref{sites-lemma-colimit-sheaves}를 보라.
$\mathcal{U} = \{U_i \to U\}_{i \in I}$를 $\mathcal{C}$의 덮개들의 집합
$\text{Cov}(\mathcal{C})$의 원소라 하자. $\beta$의 공종성이 $I$의
기수보다 크면
$$
H^0(\mathcal{U}, \mathcal{F}_{\colim}) =
\colim H^0(\mathcal{U}, \mathcal{F}_\alpha) =
\colim \mathcal{F}_\alpha(U) = \mathcal{F}_{\colim}(U)
$$
라고 주장한다. 두 번째와 세 번째 등호는 명백하다. 첫 번째를 위해
$s = (s_i) \in H^0(\mathcal{U}, \mathcal{F}_{\colim})$라 하자. 그러면
각 $i$에 대하여 원소 $s_i$는 어떤 $\alpha_i < \beta$에 대한 원소
$s_{i, \alpha_i} \in \mathcal{F}_{\alpha_i}(U_i)$에서 온다. 공종성에 대한
가정에 의해
% P01-ERRATA: P01-E0074 -- the frozen proof claims the original alpha_i are equal instead of transporting to a common upper stage.
$i$와 무관하게 $\alpha_i = \alpha$로 택할 수 있다. 그러면 어떤
$\alpha_{i, j} < \beta$에 대하여 $s_i$와 $s_j$는
$\mathcal{F}_{\alpha_{i, j}}(U_i \times_U U_j)$의 같은 원소로 사상된다.
% P01-ERRATA: P01-E0075 -- the frozen source mistypes "cardinality of" as "cardinality if".
$I \times I$의 기수도 $\beta$의 공종성보다 작으므로, $\alpha$를 늘린 뒤
모든 $i, j$에 대하여 $\alpha_{i, j} = \alpha$라고 가정할 수 있다. 이로써
자연스러운 사상
$\colim H^0(\mathcal{U}, \mathcal{F}_\alpha)
\to H^0(\mathcal{U}, \mathcal{F}_{\colim})$가 전사임을 증명하였다. 매우
유사한 논증으로 단사임을 보일 수 있다. 특히 $\mathcal{F}_{\colim}$은
$\mathcal{U}$에 대한 층 조건을 만족한다. 따라서 $\beta$의 공종성이 덮개들의
지표집합들 $I$의 기수들의 상한보다 크면 결론을 얻는다.
\end{proof}

\section{사이트들의 여극한}
\label{sites-section-colimit-sites}

\noindent
사이트가 사이트들의 역계의 극한인 경우에는
Lemma \ref{sites-lemma-directed-colimits-sections}의 유사물이 필요하다. 간단히 하기
위하여 덮개들의 지표집합이 유한한 경우에만 이 구성을 설명한다.

\begin{situation}
\label{sites-situation-inverse-limit-sites}
다음이 주어졌다고 하자.
\begin{enumerate}
\item 공여과 지표범주 $\mathcal{I}$,
\item 각 $i \in \Ob(\mathcal{I})$에 대하여, $\mathcal{C}_i$의 모든 덮개가
유한한 지표집합을 갖는 사이트 $\mathcal{C}_i$,
\item $\mathcal{I}$의 사상 $a : i \to j$에 대하여, 연속 함자
$u_a : \mathcal{C}_j \to \mathcal{C}_i$로 주어지는 사이트의 사상
$f_a : \mathcal{C}_i \to \mathcal{C}_j$.
\end{enumerate}
이들은 $\mathcal{I}$에서 $c = a \circ b$일 때마다
$f_a \circ f_b = f_c$를 만족한다고 하자.
\end{situation}

\begin{lemma}
\label{sites-lemma-colimit-sites}
Situation \ref{sites-situation-inverse-limit-sites}에서 다음과 같이 사이트
$(\mathcal{C}, \text{Cov}(\mathcal{C}))$를 구성할 수 있다.
\begin{enumerate}
\item 범주로서 $\mathcal{C} = \colim \mathcal{C}_i$이고,
\item $\text{Cov}(\mathcal{C})$는
$u_i : \mathcal{C}_i \to \mathcal{C}$에 의한
$\text{Cov}(\mathcal{C}_i)$의 상들의 합집합이다.
\end{enumerate}
\end{lemma}

\begin{proof}
사이트의 사상들의 합성에 대한 우리의 정의에 따르면 $\mathcal{I}$에서
$c = a \circ b$일 때마다 $u_b \circ u_a = u_c$이다.
공식 $\mathcal{C} = \colim \mathcal{C}_i$는
$\Ob(\mathcal{C}) = \colim \Ob(\mathcal{C}_i)$이고
$\text{Arrows}(\mathcal{C}) = \colim \text{Arrows}(\mathcal{C}_i)$임을 뜻한다.
그러면 원천, 목표 및 합성은 $\text{Arrows}(\mathcal{C}_i)$ 위의 원천,
목표 및 합성에서 유도된다. 이 방법으로 범주를 얻는다.
명백한 함자를 $u_i : \mathcal{C}_i \to \mathcal{C}$로 나타내자.
$\mathcal{C}$ 안의 임의의 유한 도표가 주어지면, 이 도표가
$\mathcal{C}_i$ 안의 어떤 도표의 상이 되는 $i$가 존재함을 주목하자.

\medskip\noindent
$\{U^t \to U\}$를 $\mathcal{C}$의 덮개라 하자. 먼저
$V \to U$가 $\mathcal{C}$의 사상이면 $U^t \times_U V$가 존재함을 보인다.
위의 관찰과 덮개의 정의에 의해, $\mathcal{C}_i$의 덮개
$\{U_i^t \to U_i\}$와 사상 $V_i \to U_i$가 어떤 $i$에 대하여 존재하고,
$u_i$에 의한 그 상은 주어진 자료가 된다. $U^t \times_U V$는
$u_i$에 의한 $U^t_i \times_{U_i} V_i$의 상이라고 주장한다.
실제로 $\mathcal{I}$의 모든 $a : j \to i$에 대하여 함자 $u_a$는
연속이므로
$u_a(U^t_i \times_{U_i} V_i) = u_a(U^t_i) \times_{u_a(U_i)} u_a(V_i)$이다.
특히 원한다면 $i$를 $j$로 바꿀 수 있다. 따라서 $W$가
$\mathcal{C}$의 또 다른 대상이면 $W = u_i(W_i)$라고 가정할 수 있고,
다음을 얻는다.
\begin{align*}
& \Mor_\mathcal{C}(W, u_i(U^t_i \times_{U_i} V_i)) \\
& =
\colim_{a : j \to i}
\Mor_{\mathcal{C}_j}(u_a(W_i), u_a(U^t_i \times_{U_i} V_i)) \\
& =
\colim_{a : j \to i}
\Mor_{\mathcal{C}_j}(u_a(W_i), u_a(U^t_i))
\times_{\Mor_{\mathcal{C}_j}(u_a(W_i), u_a(U_i))}
\Mor_{\mathcal{C}_j}(u_a(W_i), u_a(V_i)) \\
& =
\Mor_\mathcal{C}(W, U^t)
\times_{\Mor_\mathcal{C}(W, U)}
\Mor_\mathcal{C}(W, V)
\end{align*}
마지막 등호는 여과된 여극한이 유한 극한과 가환하기 때문이다
(Categories, Lemma \ref{categories-lemma-directed-commutes}). 또한
$\{U^t \times_U V \to V\}$가 $\mathcal{C}$의 덮개임이 따른다.
이로써 Definition \ref{sites-definition-site}의 공리 (3)이 성립함을 알 수 있다.

\medskip\noindent
Definition \ref{sites-definition-site}의 공리 (2)를 확인하기 위하여
$\{U^t \to U\}_{t \in T}$를 $\mathcal{C}$의 덮개라 하고, 각 $t$에 대하여
$\{U^{ts} \to U^t\}$를 $\mathcal{C}$의 덮개라 하자. 그러면 어떤 $i$와
$\mathcal{C}_i$의 덮개 $\{U^t_i \to U_i\}_{t \in T}$를 찾아서,
$u_i$에 의한 그 상이 $\{U^t \to U\}$가 되게 할 수 있다. $T$가
{\bf 유한}하므로 $\mathcal{I}$의 어떤 $a : j \to i$와
$\mathcal{C}_j$의 덮개들 $\{U^{ts}_j \to u_a(U^t_i)\}$를 택하여,
$u_j$에 의한 그 상들이 $\{U^{ts} \to U^t\}$를 주게 할 수 있다.
이제 사이트 $\mathcal{C}_j$에 공리 (2)를 적용하면
$\{U^{ts} \to U\}$가 $\mathcal{C}$의 덮개임을 얻는다.

\medskip\noindent
Definition \ref{sites-definition-site}의 공리 (1)의 증명은 생략한다.
\end{proof}

\begin{lemma}
\label{sites-lemma-compute-pullback-to-limit}
Situation \ref{sites-situation-inverse-limit-sites}에서
$u_i : \mathcal{C}_i \to \mathcal{C}$를
Lemma \ref{sites-lemma-colimit-sites}에서 구성한 것이라 하자. 그러면 $u_i$는
사이트의 사상 $f_i : \mathcal{C} \to \mathcal{C}_i$를 정의한다.
$U_i \in \Ob(\mathcal{C}_i)$이고 $\mathcal{F}$가 $\mathcal{C}_i$ 위의
층이면 다음이 성립한다.
\begin{equation}
\label{sites-equation-compute-pullback-to-limit}
f_i^{-1}\mathcal{F}(u_i(U_i)) =
\colim_{a : j \to i} f_a^{-1}\mathcal{F}(u_a(U_i))
\end{equation}
\end{lemma}

\begin{proof}
Lemma \ref{sites-lemma-colimit-sites}의 증명에 있는 논증으로부터 함자 $u_i$들이
연속임은 즉시 알 수 있다. 증명을 끝내려면
$f_i^{-1} := u_{i, s}$가 완전 함자
$\Sh(\mathcal{C}_i) \to \Sh(\mathcal{C})$임을 보여야 한다. 실제로
$f_i^{-1}$가 좌완전임을 보이면 충분하다. 왼쪽 수반이므로 우완전이기 때문이다
(Categories, Lemma \ref{categories-lemma-exact-adjoint}). 먼저
(\ref{sites-equation-compute-pullback-to-limit})을 증명한 다음 완전성을 도출한다.

\medskip\noindent
% P01-ERRATA: P01-E0076 -- the frozen source duplicates the indefinite article before a.
% P01-ERRATA: P01-E0077 -- the frozen source places V_j in C although u_j has domain C_j.
$\mathcal{C}$의 임의의 대상 $V$에 대하여 $a : j \to i$와
$V = u_j(V_j)$를 만족하는 대상 $V_j \in \Ob(\mathcal{C})$를 택할 수 있다.
그러면 다음과 같이 놓을 수 있다.
$$
\mathcal{G}(V) = \colim_{b : k \to j} f_{a \circ b}^{-1}\mathcal{F}(u_b(V_j))
$$
% P01-ERRATA: P01-E0078 -- the frozen prose changes the fixed outer arrow a to b and collides with the colimit variable.
여극한의 값 $\mathcal{G}(V)$는 $b : j \to i$와
$u_j(V_j) = V$인 대상 $V_j$의 선택과 무관하다. 확인은 생략한다.
더 나아가 $\alpha : V \to V'$가 $\mathcal{C}$의 사상이면,
$u_j(\alpha_j) = \alpha$를 만족하는 $b : j \to i$와 사상
$\alpha_j : V_j \to V'_j$를 택할 수 있다. 사상
$u_b(\alpha_j) : u_b(V_j) \to u_b(V'_j)$를 따른 제한을 사용하면 이로부터
사상 $\mathcal{G}(V') \to \mathcal{G}(V)$가 유도된다. 확인해 보면
$\mathcal{G}$는 준층이다(생략한다). 사실 $\mathcal{G}$는 층 조건을 만족한다.
실제로 $\mathcal{C}$의 임의의 덮개
$\mathcal{U} = \{U^t \to U\}$는 유한한 단계에서 온다. 어떤
$\mathcal{I}$의 $a : j \to i$에 대하여
$\mathcal{U}_j = \{U^t_j \to U_j\}$가 $u_j$에 의해
$\mathcal{U}$로 사상된다고 하자. 그러면 다음을 얻는다.
% P01-ERRATA: P01-E0079 -- the frozen display reverses the composable order a o b to the ill-typed b o a twice.
$$
H^0(\mathcal{U}, \mathcal{G}) =
\colim_{b : k \to j} H^0(u_b(\mathcal{U}_j), f_{b \circ a}^{-1}\mathcal{F}) =
\colim_{b : k \to j} f_{b \circ a}^{-1}\mathcal{F}(u_b(U_j)) =
\mathcal{G}(U)
$$
이는 원하는 바이다. 첫 번째 등호는 여과된 여극한이 유한 극한과 가환하기
때문에 성립한다
(Categories, Lemma \ref{categories-lemma-directed-commutes}). 구성에 의해
$\mathcal{G}(U)$는 (\ref{sites-equation-compute-pullback-to-limit})의 우변으로
주어진다. 따라서 $\mathcal{G}$가 $f_i^{-1}\mathcal{F}$와 같음을 보이면
(\ref{sites-equation-compute-pullback-to-limit})이 성립한다.

\medskip\noindent
이 문단에서는 $\mathcal{G}$가 $f_i^{-1}\mathcal{F}$와 표준적으로
동형임을 확인한다. 독자에게 이 문단을 건너뛸 것을 강력히 권한다. 이를
확인하려면 $\mathcal{C}$ 위의 층 $\mathcal{H}$에 함자적인 전단사
$\Mor_{\Sh(\mathcal{C})}(\mathcal{G}, \mathcal{H}) =
\Mor_{\Sh(\mathcal{C}_i)}(\mathcal{F}, f_{i, *}\mathcal{H})$가 존재함을
보여야 한다. 여기서 $f_{i, *} = u_i^p$이다. 사상
$\mathcal{G} \to \mathcal{H}$를 주는 것은 모든 $a : j \to i$,
$b : k \to j$ 및 $V_j \in \Ob(\mathcal{C}_j)$에 대하여 양립하는 사상계
$$
\varphi_{a, b, V_j} :
f_{a \circ b}^{-1}\mathcal{F}(u_b(V_j))
\longrightarrow
\mathcal{H}(u_j(V_j))
$$
를 주는 것과 같다. 양립성에 의해 사상 $\varphi_{a, b, V_j}$는
$\varphi_{a \circ b, \text{id}, u_b(V_j)}$와 같아야 한다.
$a : j \to i$가 주어지면 사상족 $\varphi_{a, \text{id}, V_j}$는 층의 사상
$\varphi_a : f_a^{-1}\mathcal{F} \to f_{j, *}\mathcal{H}$에 대응한다.
% P01-ERRATA: P01-E0080 -- the frozen English gives the plural subject compatibilities a singular verb.
$\varphi_{a, \text{id}, u_a(V_i)}$와
$\varphi_{\text{id}, \text{id}, V_i}$ 사이의 양립성에 따르면
$\varphi_a$는 다음을 통한 사상 $\varphi_{id}$의 수반이다.
$$
\Mor_{\Sh(\mathcal{C}_j)}(f_a^{-1}\mathcal{F}, f_{j, *}\mathcal{H}) =
\Mor_{\Sh(\mathcal{C}_i)}(\mathcal{F}, f_{a, *}f_{j, *}\mathcal{H}) =
\Mor_{\Sh(\mathcal{C}_i)}(\mathcal{F}, f_{i, *}\mathcal{H})
$$
따라서 마침내 전체 사상계 $\varphi_{a, b, V_j}$가 사상
$\varphi_{id} : \mathcal{F} \to f_{i, *}\mathcal{H}$에 의해 결정됨을 안다.
역으로 그러한 사상 $\psi : \mathcal{F} \to f_{i, *}\mathcal{H}$가 주어지면,
방금 한 논증을 거꾸로 읽어 사상족 $\varphi_{a, b, V_j}$를 구성할 수 있다.
이로써 $\mathcal{G} = f_i^{-1}\mathcal{F}$임을 증명하였다.

\medskip\noindent
(\ref{sites-equation-compute-pullback-to-limit})이 성립한다고 하자. 그러면
층의 유한 극한은 준층의 범주에서 계산되고
(Lemma \ref{sites-lemma-limit-sheaf}), 함자 $f_a^{-1}$들은 유한 극한과 가환하며,
여과된 여극한은 유한 극한과 가환하므로, 함자
$\mathcal{F} \mapsto f_i^{-1}\mathcal{F}(U)$는 유한 극한과 가환한다.
$\mathcal{C}$의 일반적인 대상 $V$에 대하여 함자
$\mathcal{F} \mapsto f_i^{-1}\mathcal{F}(V)$가 유한 극한과 가환함을 보이려면
위에서 준 $f_i^{-1}\mathcal{F}(V) = \mathcal{G}(V)$의 공식을 사용하여 같은
논증을 적용하면 된다. 따라서 $f_i^{-1}$는 좌완전이고 보조정리의 증명이
끝난다.
\end{proof}

\begin{lemma}
\label{sites-lemma-colimit}
Situation \ref{sites-situation-inverse-limit-sites}에서 다음이 주어졌다고 하자.
\begin{enumerate}
\item 모든 $i \in \Ob(\mathcal{I})$에 대하여 $\mathcal{C}_i$ 위의 층
$\mathcal{F}_i$,
\item $a : j \to i$에 대하여 $\mathcal{C}_j$ 위의 층의 사상
$\varphi_a : f_a^{-1}\mathcal{F}_i \to \mathcal{F}_j$.
\end{enumerate}
$c = a \circ b$일 때마다
$\varphi_c = \varphi_b \circ f_b^{-1}\varphi_a$를 만족한다고 하자.
Lemma \ref{sites-lemma-colimit-sites}의 사이트 $\mathcal{C}$ 위에서
$\mathcal{F} = \colim f_i^{-1}\mathcal{F}_i$로 놓는다.
$i \in \Ob(\mathcal{I})$이고 $X_i \in \text{Ob}(\mathcal{C}_i)$이면
$$
\colim_{a : j \to i} \mathcal{F}_j(u_a(X_i)) = \mathcal{F}(u_i(X_i))
$$
이다.
\end{lemma}

\begin{proof}
형식적인 논증에 의해 다음이 성립한다.
% P01-ERRATA: P01-E0081 -- the frozen proof evaluates F_i at an object of C_j instead of using F_j as in the statement.
$$
\colim_{a : j \to i} \mathcal{F}_i(u_a(X_i)) =
\colim_{a : j \to i} \colim_{b : k \to j}
f_b^{-1}\mathcal{F}_j(u_{a \circ b}(X_i))
$$
(\ref{sites-equation-compute-pullback-to-limit})에 의해 안쪽 여극한은
$f_j^{-1}\mathcal{F}_j(u_i(X_i))$와 같으므로
Lemma \ref{sites-lemma-directed-colimits-sections}에 의해 결론을 얻는다.
\end{proof}

\begin{lemma}
\label{sites-lemma-colimit-push-pull}
Situation \ref{sites-situation-inverse-limit-sites}에서 $\mathcal{C}$ 위의 층
$\mathcal{F}$가 주어졌다고 하자. 그러면
$$
\mathcal{F} = \colim f_i^{-1}f_{i, *}\mathcal{F}
$$
이다. 여기서 전이사상은 $a : j \to i$에 대한
$f_j^{-1}\varphi_a$이고,
$\varphi_a : f_a^{-1}f_{i, *}\mathcal{F} \to f_{j, *}\mathcal{F}$는
Lemma \ref{sites-lemma-colimit}에서와 같은 여순환 조건을 만족하는 표준적인 사상이다.
\end{lemma}

\begin{proof}
사상
$$
\varphi_a : f_a^{-1}f_{i, *}\mathcal{F} \to f_{j, *}\mathcal{F}
$$
으로 항등사상의 수반을 택한다.
$$
f_{i, *}\mathcal{F} \to f_{a, *}f_{j, *}\mathcal{F}
$$
따라서 $\varphi_a$는 $(f_a^{-1}, f_{a, *})$가 주는 수반의 여단위원이다.
% P01-ERRATA: P01-E0082 -- the frozen source gives b target i although a o b requires target j.
$a : j \to i$와 $b : k \to i$가 $c = a \circ b$를 합성으로 가질 때마다
$\varphi_c = \varphi_b \circ f_b^{-1}\varphi_a$임을 증명해야 한다.
이는 Categories, Lemma \ref{categories-lemma-compose-counits}에서 따른다.
마지막으로 수반에 의해 주어지는 사상
$$
\colim_{i \in I} f_i^{-1}f_{i, *}\mathcal{F}
\longrightarrow
\mathcal{F}
$$
이 동형임을 증명해야 한다. $\mathcal{C}$의 대상 $U$에 대하여 사상
% P01-ERRATA: P01-E0083 -- the frozen proof uses undefined F_i instead of f_{i,*}F in both evaluated colimits.
$$
(\colim_{i \in I} f_i^{-1}\mathcal{F}_i)(U) \to \mathcal{F}(U)
$$
이 전단사임을 보이면 된다. $u_i(U_i) = U$를 만족하는 어떤 $i$와
$\mathcal{C}_i$의 대상 $U_i$를 택하자. 그러면 좌변은
$$
(\colim_{i \in I} f_i^{-1}\mathcal{F}_i)(U) =
\colim_{a : j \to i} f_{j, *}\mathcal{F}(u_a(U_i))
$$
와 같다. 이는 Lemma \ref{sites-lemma-colimit}에서 따른다.
$u_j(u_a(U_i)) = U$이므로 정의에 의해 모든 $a : j \to i$에 대하여
$f_{j, *}\mathcal{F}(u_a(U_i)) = \mathcal{F}(U)$이다. 따라서 여극한의 값은
$\mathcal{F}(U)$이고 증명이 끝난다.
\end{proof}

\section{준층의 함자성에 대한 추가 사항}
\label{sites-section-more-functoriality-PSh}

\noindent
이 절에서는 Section \ref{sites-section-functoriality-PSh}의 내용을 다시 살펴본다.
$u : \mathcal{C} \to \mathcal{D}$를 범주 사이의 함자라 하자. 함자
$$
u^p :
\textit{PSh}(\mathcal{D})
\longrightarrow
\textit{PSh}(\mathcal{C})
$$
는 $\mathcal{D}$ 위의 $\mathcal{G}$에 준층
$u^p\mathcal{G} = \mathcal{G} \circ u$를 대응시킴을 상기하자. 이 함자는
왼쪽 수반(즉 $u_p$)뿐 아니라 오른쪽 수반도 갖는다.

\medskip\noindent
구체적으로 각 $V \in \Ob(\mathcal{D})$에 대하여 범주
${}_V\mathcal{I} = {}_V^u\mathcal{I}$를 정의한다. 그 대상은 쌍
$(U, \psi : u(U) \to V)$이다. 화살표의 방향은
Section \ref{sites-section-functoriality-PSh}에서 범주 $\mathcal{I}_V^u$를 정의할
때 사용한 화살표의 방향과 반대임을 주목하자. 사상
$(U, \psi) \to (U', \psi')$는
$\psi = \psi' \circ u(\alpha)$를 만족하는 사상
$\alpha : U \to U'$로 주어진다. 또한 $\mathcal{C}$ 위의 임의의 집합값 준층
$\mathcal{F}$에 대하여 규칙
${}_V\mathcal{F}(U, \psi) = \mathcal{F}(U)$로 정의되는 함자
${}_V\mathcal{F} : {}_V\mathcal{I}^{opp} \to \textit{Sets}$를 도입한다.
다음과 같이 정의하자.
$$
{}_pu(\mathcal{F})(V) := \lim_{{}_V\mathcal{I}^{opp}} {}_V\mathcal{F}
$$
이는 극한이므로 ${}_V\mathcal{I}$의 각 대상 $(U, \psi)$에 대하여 사영사상
$c(\psi) : {}_pu(\mathcal{F})(V) \to \mathcal{F}(U)$가 있다. 실제로
$$
{}_pu(\mathcal{F})(V)
=
\left\{
\begin{matrix}
\text{collections }
s_{(U, \psi)} \in \mathcal{F}(U) \\
\forall \beta : (U_1, \psi_1) \to (U_2, \psi_2)
\text{ in }{}_V\mathcal{I} \\
\text{ we have } \beta^*s_{(U_2, \psi_2)} = s_{(U_1, \psi_1)}
\end{matrix}
\right\}
$$
이고, 대응은 $s \mapsto s_{(U, \psi)} = c(\psi)(s)$로 주어진다.
$\mathcal{D}$의 임의의 사상 $V' \to V$에 결부된 제한사상
${}_pu(\mathcal{F})(V) \to {}_pu(\mathcal{F})(V')$의 정의는 독자에게
맡긴다. 이렇게 얻은 준층을 ${}_pu\mathcal{F}$로 나타낸다.

\begin{lemma}
\label{sites-lemma-recover-pu}
제한사상들과 양립하는 표준적인 사상
${}_pu\mathcal{F}(u(U)) \to \mathcal{F}(U)$가 있다.
\end{lemma}

\begin{proof}
이는 위의 사영사상 $c(\text{id}_{u(U)})$일 뿐이다.
\end{proof}

\noindent
준층의 임의의 사상 $\mathcal{F} \to \mathcal{F}'$는 함자 사이의 양립하는
사상계 ${}_V\mathcal{F} \to {}_V\mathcal{F}'$를 낳고, 따라서 준층의 사상
${}_pu\mathcal{F} \to {}_pu\mathcal{F}'$를 낳는다. 다시 말해 함자
$$
{}_pu :
\textit{PSh}(\mathcal{C})
\longrightarrow
\textit{PSh}(\mathcal{D})
$$
를 정의하였다.

\begin{lemma}
\label{sites-lemma-adjoints-pu}
함자 ${}_pu$는 함자 $u^p$의 오른쪽 수반이다. 다시 말해 공식
$$
\Mor_{\textit{PSh}(\mathcal{C})}(u^p\mathcal{G}, \mathcal{F})
=
\Mor_{\textit{PSh}(\mathcal{D})}(\mathcal{G}, {}_pu\mathcal{F})
$$
는 $\mathcal{F}$와 $\mathcal{G}$에 대하여 두 변수 함자적으로 성립한다.
\end{lemma}

\begin{proof}
이는 Lemma \ref{sites-lemma-adjoints-u}의 증명과 정확히 같은 방법으로 증명된다.
Lemma \ref{sites-lemma-recover-pu}의 사상
$u^p{}_pu \mathcal{F} \to \mathcal{F}$가 우변에서 좌변으로 가는 데 사용되는
사상임을 주목한다.

\medskip\noindent
다른 방법으로, $\mathcal{C}$ 위의 집합값 준층 $\mathcal{F}$를
$\textit{Sets}^{opp}$에 값을 갖는 $\mathcal{C}^{opp}$ 위의 준층
$\mathcal{F}'$로 생각하고, $\mathcal{D}$ 위에서도 마찬가지로 생각하자.
$({}_pu \mathcal{F})' = u_p(\mathcal{F}')$이고
$(u^p\mathcal{G})' = u^p(\mathcal{G}')$임을 확인하라.
Remark \ref{sites-remark-functoriality-presheaves-values}에 의해
$\textit{Sets}^{opp}$에 값을 갖는 준층에 대하여 $u_p$와 $u^p$의 수반성이
성립한다. 이제 결과는 형식적으로 따른다.
\end{proof}

\noindent
따라서 범주의 함자 $u : \mathcal{C} \to \mathcal{D}$가 주어지면 준층들의
범주 사이의 함자열
$$
u_p, u^p, {}_pu
$$
를 얻으며, 연속한 각 쌍에서 첫째 함자는 둘째 함자의 왼쪽 수반이다.

\begin{lemma}
\label{sites-lemma-adjoint-functors}
$u : \mathcal{C} \to \mathcal{D}$와 $v : \mathcal{D} \to \mathcal{C}$를
범주의 함자라 하자. $v$가 $u$의 오른쪽 수반이라고 가정하자. 그러면
다음이 성립한다.
\begin{enumerate}
\item $\mathcal{D}$의 임의의 $V$에 대하여 $u^ph_V = h_{v(V)}$이다.
\item 범주 $\mathcal{I}^v_U$는 시초대상을 갖는다.
\item 범주 ${}_V^u\mathcal{I}$는 종대상을 갖는다.
\item ${}_pu = v^p$이다.
\item $u^p = v_p$이다.
\end{enumerate}
\end{lemma}

\begin{proof}
(1)의 증명. $V$를 $\mathcal{D}$의 대상이라 하자. 가정에 의해
$u^ph_V(U) = \Mor_\mathcal{D}(u(U), V) = \Mor_\mathcal{C}(U, v(V))$이므로
$u^ph_V = h_{v(V)}$이다.

\medskip\noindent
(2)의 증명. $U$를 $\mathcal{C}$의 대상이라 하자. 사상
$\text{id} : u(U) \to u(U)$의 수반인 사상을
$\eta : U \to v(u(U))$라 하자. 그러면 $(u(U), \eta)$가
$\mathcal{I}_U^v$의 시초대상이라고 주장한다. 실제로
$\mathcal{I}_U^v$의 대상 $(V, \phi : U \to v(V))$가 주어지면 사상
$\phi$는 사상 $\psi : u(U) \to V$의 수반이고, 이 사상은 다시 사상
$(u(U), \eta) \to (V, \phi)$를 정의한다.

\medskip\noindent
(3)의 증명. $V$를 $\mathcal{D}$의 대상이라 하자. 사상
$\text{id} : v(V) \to v(V)$의 수반인 사상을
$\xi : u(v(V)) \to V$라 하자. 그러면 $(v(V), \xi)$가
${}_V^u\mathcal{I}$의 종대상이라고 주장한다. 실제로
${}_V^u\mathcal{I}$의 대상 $(U, \psi : u(U) \to V)$가 주어지면 사상
$\psi$는 사상 $\phi : U \to v(V)$의 수반이고, 이 사상은 다시 사상
$(U, \psi) \to (v(V), \xi)$를 정의한다.

\medskip\noindent
따라서 $\mathcal{C}$ 위의 임의의 준층 $\mathcal{F}$에 대하여 다음을 얻는다.
\begin{eqnarray*}
v^p\mathcal{F}(V)
& = &
\mathcal{F}(v(V)) \\
& = &
\Mor_{\textit{PSh}(\mathcal{C})}(h_{v(V)}, \mathcal{F}) \\
& = &
\Mor_{\textit{PSh}(\mathcal{C})}(u^ph_V, \mathcal{F}) \\
& = &
\Mor_{\textit{PSh}(\mathcal{D})}(h_V, {}_pu\mathcal{F}) \\
& = &
{}_pu\mathcal{F}(V)
\end{eqnarray*}
이로써 (4)가 증명된다. (5)는 수반함자의 유일성에서 따른다.
\end{proof}

\section{쌍대연속 함자}
\label{sites-section-cocontinuous-functors}

\noindent
토포스의 사상을 구성하는 또 다른 방법이 있다. 이 방법은 다음과 같이
정의되는 사이트 사이의 쌍대연속 함자를 사용한다.

\begin{definition}
\label{sites-definition-cocontinuous}
$\mathcal{C}$와 $\mathcal{D}$를 사이트라 하자.
$u : \mathcal{C} \to \mathcal{D}$를 함자라 하자. 모든
$U \in \Ob(\mathcal{C})$와 $\mathcal{D}$의 모든 덮개
$\{V_j \to u(U)\}_{j \in J}$에 대하여, $\mathcal{C}$의 덮개
$\{U_i \to U\}_{i\in I}$가 존재하여 사상족
$\{u(U_i) \to u(U)\}_{i \in I}$가 덮개
$\{V_j \to u(U)\}_{j \in J}$를 세분하면 함자 $u$를
{\it 쌍대연속}이라 한다.
\end{definition}

\noindent
일반적으로 $\{u(U_i) \to u(U)\}_{i \in I}$는 사이트 $\mathcal{D}$의
덮개가 {\it 아님}에 유의하자.

\begin{lemma}
\label{sites-lemma-pu-sheaf}
$\mathcal{C}$와 $\mathcal{D}$를 사이트라 하자.
$u : \mathcal{C} \to \mathcal{D}$가 쌍대연속이라고 하자.
$\mathcal{F}$를 $\mathcal{C}$ 위의 층이라 하자. 그러면
${}_pu\mathcal{F}$는 $\mathcal{D}$ 위의 층이며, 이를
${}_su\mathcal{F}$로 나타낸다.
\end{lemma}

\begin{proof}
$\{V_j \to V\}_{j \in J}$를 사이트 $\mathcal{D}$의 덮개라 하자. 도표
$$
\xymatrix{
&\ \phantom{}_pu\mathcal{F}(V) \ar[r] &
\prod {}_pu\mathcal{F}(V_j) \ar@<1ex>[r] \ar@<-1ex>[r] &
\prod {}_pu\mathcal{F}(V_j \times_V V_{j'}) &
}
$$
가 등화자 도표임을 보여야 한다. ${}_pu$는 $u^p$의 오른쪽 수반이므로
$$
{}_pu\mathcal{F}(V) =
\Mor_{\textit{PSh}(\mathcal{D})}(h_V, {}_pu\mathcal{F}) =
\Mor_{\textit{PSh}(\mathcal{C})}(u^ph_V, \mathcal{F}) =
\Mor_{\Sh(\mathcal{C})}((u^ph_V)^\#, \mathcal{F})
$$
이다. 따라서
\begin{equation}
\label{sites-equation-coequalizer}
\xymatrix{
\coprod u^p h_{V_j \times_V V_{j'}} \ar@<1ex>[r] \ar@<-1ex>[r] &
\coprod u^p h_{V_j} \ar[r] &
u^p h_V
}
\end{equation}
가 층화한 뒤 쌍대등화자 도표가 됨을 보이면 충분하다. (층의 범주에서의
쌍대곱은 준층의 범주에서의 쌍대곱을 층화한 것임을 상기하라.
Lemma \ref{sites-lemma-colimit-sheaves}를 보라.)

\medskip\noindent
먼저 (\ref{sites-equation-coequalizer})의 두 번째 화살표가 층화한 뒤 전사가 됨을
보인다. 이를 위해 Lemma \ref{sites-lemma-mono-epi-sheaves}를 사용한다. 따라서
$U$ 위의 $u^ph_V$의 절단 $s$가 $U$의 어떤 덮개를 이루는 대상들 위에서
$\coprod u^p h_{V_j}$의 절단으로 들어 올려짐을 보이면 충분하다. $s$는 사상
$s : u(U) \to V$임을 주목하자. 그러면
$\{V_j \times_{V, s} u(U) \to u(U)\}$는 $\mathcal{D}$의 덮개이다.
$u$가 쌍대연속이므로, $\{u(U_i) \to u(U)\}$가
$\{V_j \times_{V, s} u(U) \to u(U)\}$를 세분하도록 하는 덮개
$\{U_i \to U\}$가 존재한다. 이는 각 제한
$s|_{U_i} : u(U_i) \to V$가 어떤 $j$에 대한 사상
$s_i : u(U_i) \to V_j$를 통해 분해됨을 뜻한다. 즉, 원하는 대로
$s|_{U_i}$는 $u^ph_{V_j}(U_i) \to u^ph_V(U_i)$의 상에 속한다.

\medskip\noindent
$s, s' \in (\coprod u^ph_{V_j})^\#(U)$가
$(u^ph_V)^\#(U)$의 같은 원소로 사상된다고 하자. 보조정리의 증명을
끝내기 위하여, $U$를 어떤 덮개의 대상들로 바꾼 뒤 $s, s'$가
(\ref{sites-equation-coequalizer})의 두 사상에 의한
$\coprod u^p h_{V_j \times_V V_{j'}}$의 같은 절단의 상임을 보인다.
먼저 $U$를 덮개의 대상들로 바꾸어
$s \in u^ph_{V_j}(U)$이고 $s' \in u^ph_{V_{j'}}(U)$라고 가정할 수 있다.
다시 한 번 그렇게 바꾸면 $s$와 $s'$가 층화에서가 아니라
$u^ph_V(U)$에서 같은 상을 갖도록 할 수 있다. 따라서
$s : u(U) \to V_j$와 $s' : u(U) \to V_{j'}$는 다음 도표가 가환하도록
하는 $\mathcal{D}$의 사상이다.
$$
\xymatrix{
u(U) \ar[r]_{s'} \ar[d]_s & V_{j'} \ar[d] \\
V_j \ar[r] & V
}
$$
따라서 $t = (s, s') : u(U) \to V_j \times_V V_{j'}$를 얻는다. 즉, 원하는
대로 $s, s'$로 사상되는 절단
$t \in u^ph_{V_j \times_V V_{j'}}(U)$를 얻는다.
\end{proof}

\begin{lemma}
\label{sites-lemma-exact-cocontinuous}
$\mathcal{C}$와 $\mathcal{D}$를 사이트라 하고
$u : \mathcal{C} \to \mathcal{D}$를 쌍대연속이라 하자. 함자
$\Sh(\mathcal{D}) \to \Sh(\mathcal{C})$,
$\mathcal{G} \mapsto (u^p\mathcal{G})^\#$는 위의
Lemma \ref{sites-lemma-pu-sheaf}에서 도입한 함자 ${}_su$의 왼쪽 수반이다.
더 나아가 이 함자는 완전하다.
\end{lemma}

\begin{proof}
수반성을 다음과 같이 증명하자.
\begin{eqnarray*}
\Mor_{\Sh(\mathcal{C})}
((u^p\mathcal{G})^\#, \mathcal{F})
& = &
\Mor_{\textit{PSh}(\mathcal{C})}
(u^p\mathcal{G}, \mathcal{F}) \\
& = &
\Mor_{\textit{PSh}(\mathcal{D})}
(\mathcal{G}, {}_pu\mathcal{F}) \\
& = &
\Mor_{\Sh(\mathcal{D})}
(\mathcal{G}, {}_su\mathcal{F}).
\end{eqnarray*}
따라서 이는 왼쪽 수반이고, 그러므로 우완전이다. Categories,
Lemma \ref{categories-lemma-exact-adjoint}를 보라. 층화가 좌완전임은 이미
보았다. Lemma \ref{sites-lemma-sheafification-exact}를 보라. 또한 포함함자
$i : \Sh(\mathcal{D}) \to \textit{PSh}(\mathcal{D})$는
Lemma \ref{sites-lemma-limit-sheaf}에 의해 좌완전이다. 마지막으로 $u^p$는
오른쪽 수반(즉 $u_p$의 오른쪽 수반)이므로 좌완전이다. 따라서 이 함자는
좌완전 함자들의 합성 ${}^\# \circ u^p \circ i$이고, 그러므로 좌완전이다.
\end{proof}

\noindent
이 절을 기술적인 보조정리로 마친다.

\begin{lemma}
\label{sites-lemma-technical-pu}
% P01-ERRATA: P01-E0084 -- the frozen lemma opens with a sentence fragment.
Lemma \ref{sites-lemma-exact-cocontinuous}의 상황에서이다. $\mathcal{D}$ 위의
임의의 준층 $\mathcal{G}$에 대하여
$(u^p\mathcal{G})^\# = (u^p(\mathcal{G}^\#))^\#$이다.
\end{lemma}

\begin{proof}
$\mathcal{C}$ 위의 임의의 층 $\mathcal{F}$에 대하여 다음이 성립한다.
\begin{eqnarray*}
\Mor_{\Sh(\mathcal{C})}((u^p(\mathcal{G}^\#))^\#, \mathcal{F})
& = &
\Mor_{\Sh(\mathcal{D})}(\mathcal{G}^\#, {}_su\mathcal{F}) \\
& = &
\Mor_{\Sh(\mathcal{D})}(\mathcal{G}^\#, {}_pu\mathcal{F}) \\
& = &
\Mor_{\textit{PSh}(\mathcal{D})}(\mathcal{G}, {}_pu\mathcal{F}) \\
& = &
\Mor_{\textit{PSh}(\mathcal{C})}(u^p\mathcal{G}, \mathcal{F}) \\
& = &
\Mor_{\Sh(\mathcal{C})}((u^p\mathcal{G})^\#, \mathcal{F})
\end{eqnarray*}
이고 결과는 요네다 보조정리에서 따른다.
\end{proof}

\begin{remark}
\label{sites-remark-cartesian-cocontinuous}
$u : \mathcal{C} \to \mathcal{D}$를 범주 사이의 함자라 하자.
$\mathcal{D}$의 사상 $g : u(U) \to V$와 $f : W \to V$가 주어지면 함자
$$
\mathcal{C}^{opp} \longrightarrow \textit{Sets},\quad
T \longmapsto
\Mor_\mathcal{C}(T, U)
\times_{\Mor_\mathcal{D}(u(T), V)}
\Mor_\mathcal{D}(u(T), W)
$$
를 생각할 수 있다. 이 함자가 표현가능하면 이에 대응하는
$\mathcal{C}$의 대상을 $U \times_{g, V, f} W$로 나타낸다.
$\mathcal{C}$와 $\mathcal{D}$가 사이트라고 가정하자. 다음 성질 $P$를
생각하자. $\mathcal{D}$의 모든 덮개 $\{f_j : V_j \to V\}$와 임의의 사상
$g : u(U) \to V$에 대하여 다음이 성립한다.
\begin{enumerate}
% P01-ERRATA: P01-E0085 -- the frozen property switches from covering index j to undefined i in both clauses.
\item 모든 $i$에 대하여 $U \times_{g, V, f_i} V_i$가 존재하고,
\item $\{U \times_{g, V, f_i} V_i \to U\}$는 $\mathcal{C}$의 덮개이다.
\end{enumerate}
연속 함자의 정의와의 유사성에 유의하라. $u$가 $P$를 가지면 $u$는
쌍대연속이다(세부사항은 생략한다). 앞으로 만날 쌍대연속 함자들 중 많은
것이 $P$를 만족한다.
\end{remark}

\section{쌍대연속 함자와 토포스의 사상}
\label{sites-section-cocontinuous-morphism-topoi}

\noindent
위의 결과로부터 쌍대연속 함자 $u$가 $u$와 같은 방향의 토포스의 사상을
준다는 것은 명백하다. 따라서 이는 Definition \ref{sites-definition-morphism-sites},
Proposition \ref{sites-proposition-get-morphism}, 그리고
Lemma \ref{sites-lemma-morphism-sites-topoi}에서와 같이 (어떤 조건 아래)
연속인 $u$에 결부되는 토포스의 사상과 반대 방향이다.

\begin{lemma}
\label{sites-lemma-cocontinuous-morphism-topoi}
$\mathcal{C}$와 $\mathcal{D}$를 사이트라 하자.
$u : \mathcal{C} \to \mathcal{D}$가 쌍대연속이라고 하자. 함자
$g_* = {}_su$와 $g^{-1} = (u^p\ )^\#$는
$\Sh(\mathcal{C})$에서 $\Sh(\mathcal{D})$로 가는 토포스의 사상
$g$를 정의한다.
\end{lemma}

\begin{proof}
이는 정확히 Lemma \ref{sites-lemma-exact-cocontinuous}의 내용이다.
\end{proof}

\begin{lemma}
\label{sites-lemma-composition-cocontinuous}
\begin{slogan}
사이트 함자의 합성은 사이트 함자의 쌍대연속성을 보존한다.
\end{slogan}
$u : \mathcal{C} \to \mathcal{D}$와 $v : \mathcal{D} \to \mathcal{E}$를
쌍대연속 함자라 하자. 그러면 $v \circ u$는 쌍대연속이고 $h = g \circ f$이다.
여기서 $f : \Sh(\mathcal{C}) \to \Sh(\mathcal{D})$, 각각
$g : \Sh(\mathcal{D}) \to \Sh(\mathcal{E})$, 각각
$h : \Sh(\mathcal{C}) \to \Sh(\mathcal{E})$는 $u$, 각각 $v$, 각각
$v \circ u$에 결부된 토포스의 사상이다.
\end{lemma}

\begin{proof}
$U \in \Ob(\mathcal{C})$라 하자.
% P01-ERRATA: P01-E0086 -- the frozen source calls E_i a covering of U in E although its target is v(u(U)).
$\{E_i \to v(u(U))\}$를 $\mathcal{E}$에서 $U$의 덮개라 하자. 가정에 의해
$\{v(D_j) \to v(u(U))\}$가 $\{E_i \to v(u(U))\}$를 세분하도록 하는
$\mathcal{D}$의 덮개 $\{D_j \to u(U)\}$가 존재한다. 다시 가정에 의해
$\{u(C_l) \to u(U)\}$가 $\{D_j \to u(U)\}$를 세분하도록 하는
$\mathcal{C}$의 덮개 $\{C_l \to U\}$가 존재한다. 그러면
$\{v(u(C_l)) \to v(u(U))\}$가 덮개 $\{E_i \to v(u(U))\}$를 세분한다.
이로써 $v \circ u$가 쌍대연속임이 증명된다. 마지막 주장을 증명하려면
${}_sv \circ {}_su = {}_s(v \circ u)$임을 보이면 충분하다.
Lemma \ref{sites-lemma-pu-sheaf}에 의해
${}_pv \circ {}_pu = {}_p(v \circ u)$임을 보이면 충분하다.
${}_pu$, 각각 ${}_pv$, 각각 ${}_p(v \circ u)$는 $u^p$, 각각 $v^p$, 각각
$(v \circ u)^p$의 오른쪽 수반이므로, $u^p \circ v^p = (v \circ u)^p$임을
보이면 충분하다. 이는 정의에서 곧바로 따른다.
\end{proof}

\begin{example}
\label{sites-example-open-immersion-cocontinuous}
$X$를 위상공간이라 하자. $j : U  \to X$를 열린 부분공간의 포함이라 하자.
사이트 $X_{Zar}$와 $U_{Zar}$가 있음을 상기하자.
Example \ref{sites-example-site-topological}을 보라. $j$에 결부된 함자
$u : X_{Zar} \to U_{Zar}$가 연속이고 사이트의 사상
$U_{Zar} \to X_{Zar}$를 낳음을 상기하자.
Example \ref{sites-example-continuous-map}을 보라. 이는 토포스의 사상
$(j_*, j^{-1})$도 준다. 다음으로 함자
$v : U_{Zar} \to X_{Zar}$, $V \mapsto v(V) = V$를 생각하자(같은 열린집합을
이제 $X_{Zar}$의 대상으로 생각할 뿐이다). 이 함자는 쌍대연속이다.
실제로 $v(V) = \bigcup_{j \in J} W_j$가 $X$의 열린 덮개이면, 각 $W_j$는
$U$의 부분집합이어야 하므로 $v(V_j)$ 꼴이고, 자명하게
$V = \bigcup_{j \in J} V_j$는 $U$의 열린 덮개이다. 위의
Lemma \ref{sites-lemma-cocontinuous-morphism-topoi}에 의해 $v$에 결부된 토포스의
사상
$$
\Sh(U) \longrightarrow \Sh(X)
$$
가 존재하며, 이는 ${}_sv$와 $(v^p\ )^\#$로 주어진다. 실제로
$(v^p\ )^\# = j^{-1}$이고 ${}_sv = j_*$라고 주장한다. 다시 말해 이는 위에서
주어진 토포스의 사상과 같다. 이를 보는 가장 쉬운 방법은 아마 $X$ 위의
임의의 층 $\mathcal{G}$에 대하여
$v^p\mathcal{G}(V) = \mathcal{G}(V)$임을 관찰하는 것이다. 이는 Sheaves,
Lemma \ref{sheaves-lemma-j-pullback}에 따르면 $j^{-1}\mathcal{G}$의 기술이다
(따라서 이 경우 층화는 불필요하다). ${}_sv$와 $j_*$의 등식은 수반함자의
유일성에서 따른다(직접 계산할 수도 있다).
\end{example}

\begin{example}
\label{sites-example-open-map-cocontinuous}
이 예는 Example \ref{sites-example-open-immersion-cocontinuous}의 약간의 일반화이다.
$f : X \to Y$를 위상공간의 연속사상이라 하고 $f$가 열린사상이라고
가정하자. 사이트 $X_{Zar}$와 $Y_{Zar}$가 있음을 상기하자.
Example \ref{sites-example-site-topological}을 보라. $f$에 결부된 함자
$u : Y_{Zar} \to X_{Zar}$가 연속이고 사이트의 사상
$X_{Zar} \to Y_{Zar}$를 낳음을 상기하자.
Example \ref{sites-example-continuous-map}을 보라. 이는 토포스의 사상
$(f_*, f^{-1})$도 준다. 다음으로 함자
$v : X_{Zar} \to Y_{Zar}$, $U \mapsto v(U) = f(U)$를 생각하자. 이 함자는
쌍대연속이다. 실제로 $f(U) = \bigcup_{j \in J} V_j$가 $Y$의 열린 덮개이면
$U_j = f^{-1}(V_j) \cap U$로 놓아서 열린 덮개 $U = \bigcup U_j$를 얻으며,
$f(U) = \bigcup f(U_j)$는 $f(U) = \bigcup V_j$의 세분이다. 위의
Lemma \ref{sites-lemma-cocontinuous-morphism-topoi}에 의해 $v$에 결부된 토포스의
사상
$$
\Sh(X) \longrightarrow \Sh(Y)
$$
가 존재하며, 이는 ${}_sv$와 $(v^p\ )^\#$로 주어진다. 실제로
$(v^p\ )^\# = f^{-1}$이고 ${}_sv = f_*$라고 주장한다. 다시 말해 이는 위에서
주어진 토포스의 사상과 같다. $Y$ 위의 임의의 층 $\mathcal{G}$에 대하여
$v^p\mathcal{G}(U) = \mathcal{G}(f(U))$이다. 한편
$u_p\mathcal{G}(U) = \colim_{f(U) \subset V} \mathcal{G}(V)
= \mathcal{G}(f(U))$로 계산할 수 있다. 실제로
$(f(U), U \subset f^{-1}(f(U)))$는 명백히
Section \ref{sites-section-functoriality-PSh}의 범주 $\mathcal{I}_U^u$의
시초대상이다. 따라서 $u_p = v^p$이고
$f^{-1} = u_s = (v^p\ )^\#$를 얻는다. ${}_sv$와 $f_*$의 등식은 수반함자의
유일성에서 따른다(직접 계산할 수도 있다).
\end{example}

\noindent
첫 번째 Example \ref{sites-example-open-immersion-cocontinuous}에서는 함자 $v$도
연속이다. 그러나 두 번째 Example \ref{sites-example-open-map-cocontinuous}에서는
Definition \ref{sites-definition-continuous}의 조건 (2)가 실패할 수 있으므로
일반적으로 연속이 아니다. 따라서 다음 보조정리는 첫 번째 예에는 적용되지만
두 번째 예에는 적용되지 않는다.

\begin{lemma}
\label{sites-lemma-when-shriek}
\begin{slogan}
동시에 연속이고 쌍대연속인 사이트 함자에 결부된 토포스의 사상에서는
층화가 불필요하다.
\end{slogan}
$\mathcal{C}$와 $\mathcal{D}$를 사이트라 하고
$u : \mathcal{C} \to \mathcal{D}$를 함자라 하자. 다음을 가정하자.
\begin{enumerate}
\item[(a)] $u$는 쌍대연속이고,
\item[(b)] $u$는 연속이다.
\end{enumerate}
$g : \Sh(\mathcal{C}) \to \Sh(\mathcal{D})$를 이에 결부된 토포스의
사상이라 하자. 그러면 다음이 성립한다.
\begin{enumerate}
\item 공식 $g^{-1} = (u^p\ )^\#$에서 층화는 불필요하다. 다시 말해
$g^{-1}(\mathcal{G})(U) = \mathcal{G}(u(U))$이다.
\item $g^{-1}$는 왼쪽 수반 $g_{!} = (u_p\ )^\#$를 갖는다.
\item $g^{-1}$는 임의의 극한 및 여극한과 가환한다.
\end{enumerate}
\end{lemma}

\begin{proof}
Lemma \ref{sites-lemma-pushforward-sheaf}에 의해 $\mathcal{D}$ 위의 임의의 층
$\mathcal{G}$에 대하여 준층 $u^p\mathcal{G}$는 $\mathcal{C}$ 위의 층이다
(이를 $u^s\mathcal{G}$로 나타낸다). 함자
$\mathcal{F} \mapsto (u_p\mathcal{F})^\#$는
Lemma \ref{sites-lemma-adjoint-sheaves}에 의해 $u^s$의 왼쪽 수반이다. 극한 및
여극한에 대한 주장은 Categories, Section \ref{categories-section-adjoint}의
논의에서 따른다.
\end{proof}

\noindent
위의 Lemma \ref{sites-lemma-when-shriek}의 상황에서는 수반함자열
$$
g_{!}, \ g^{-1}, \ g_*.
$$
을 얻는다. 일반적으로 함자 $g_!$는 $\Sh(\mathcal{C})$의 종대상을
$\Sh(\mathcal{D})$의 종대상으로 보내지 않으므로 {\it 완전하지 않다}.
Sheaves, Remark \ref{sheaves-remark-j-shriek-not-exact}를 보라. 한편
Example \ref{sites-example-open-immersion-cocontinuous}의 위상적 상황에서는 함자
$j_!$가 아벨 층에 대하여 완전하다. Modules,
Lemma \ref{modules-lemma-j-shriek-exact}를 보라. 다음 보조정리는 이를
사이트의 경우로 일반화한다.

\begin{lemma}
\label{sites-lemma-preserve-equalizers}
$\mathcal{C}$와 $\mathcal{D}$를 사이트라 하고
$u : \mathcal{C} \to \mathcal{D}$를 함자라 하자. 다음을 가정하자.
\begin{enumerate}
\item[(a)] $u$는 쌍대연속이고,
\item[(b)] $u$는 연속이며,
\item[(c)] $\mathcal{C}$에서 섬유곱과 등화자가 존재하고 $u$는 이들과
가환한다.
\end{enumerate}
이 경우 위의 함자 $g_!$는 섬유곱 및 등화자와 가환한다(더 일반적으로는
유한 연결 극한과 가환한다).
\end{lemma}

\begin{proof}
(a), (b), (c)를 가정하자. $g_! = (u_p\ )^\#$이다. 층의 극한은 대응하는
준층의 극한과 같음을 상기하자(Lemma \ref{sites-lemma-limit-sheaf}). 그리고 층화는
유한 극한과 가환한다(Lemma \ref{sites-lemma-sheafification-exact}). 따라서
$u_p$가 섬유곱 및 등화자와 가환함을 보이면 충분하다. 이를 위해서는
Section \ref{sites-section-functoriality-PSh}의 범주
$(\mathcal{I}_V^u)^{opp}$에 걸친 여극한이 섬유곱 및 등화자와 가환하면
충분하다. 이는 Lemma \ref{sites-lemma-almost-directed}와 Categories,
Lemma \ref{categories-lemma-almost-directed-commutes-equalizers}에서 따른다.
\end{proof}

\noindent
다음 보조정리는 위상공간의 열린 몰입에 결부된 사상과 훨씬 더 닮은 경우를
다룬다.

\begin{lemma}
\label{sites-lemma-back-and-forth}
$\mathcal{C}$와 $\mathcal{D}$를 사이트라 하고
$u : \mathcal{C} \to \mathcal{D}$를 함자라 하자. 다음을 가정하자.
\begin{enumerate}
\item[(a)] $u$는 쌍대연속이고,
\item[(b)] $u$는 연속이며,
\item[(c)] $u$는 충실충만하다.
\end{enumerate}
위와 같은 $g_!, g^{-1}, g_*$에 대하여 표준적인 사상
$\mathcal{F} \to g^{-1}g_!\mathcal{F}$와
$g^{-1}g_*\mathcal{F} \to \mathcal{F}$는 $\mathcal{C}$ 위의 모든 층
$\mathcal{F}$에 대하여 동형이다.
\end{lemma}

\begin{proof}
$X$를 $\mathcal{C}$의 대상이라 하자. Lemmas \ref{sites-lemma-pu-sheaf}와
\ref{sites-lemma-when-shriek}에서 함자 $g^{-1} = (u^p\ )^\#$와
$g_{*} = ({}_pu\ )^\#$에는 층화가 필요하지 않음을 보았다. 다음과 같이
계산할 수 있다.
$(g^{-1}g_{*}\mathcal{F})(X) = g_{*}\mathcal{F}(u(X))
= \lim \mathcal{F}(Y)$. 여기서 극한은 쌍 $(Y, u(Y) \to u(X))$들의 범주에
걸쳐 취하며, 사상 $u(Y) \to u(X)$가 $\mathcal{C}$의 사상 $\alpha$에 대한
$u(\alpha)$ 꼴일 필요는 없다. 가정 (c)에 의해 이 사상들은 자동으로
$\mathcal{C}$의 사상들에서 오며, 따라서 극한은
$(X, u(\text{id}_X))$에서의 값, 즉 $\mathcal{F}(X)$이다. 이로써
$g^{-1}g_{*}\mathcal{F} = \mathcal{F}$임이 증명된다.

\medskip\noindent
한편 $(g^{-1}g_{!}\mathcal{F})(X) =
g_{!}\mathcal{F}(u(X)) = (u_p\mathcal{F})^\#(u(X))$이고,
$u_p\mathcal{F}(u(X)) = \colim \mathcal{F}(Y)$이다. 여기서 여극한은 쌍
$(Y, u(X) \to u(Y))$들의 범주에 걸쳐 취하며, 사상
$u(X) \to u(Y)$가 $\mathcal{C}$의 사상 $\alpha$에 대한 $u(\alpha)$ 꼴일
필요는 없다. 가정 (c)에 의해 이 사상들은 자동으로 $\mathcal{C}$의 사상들에서
오며, 따라서 여극한은 $(X, u(\text{id}_X))$에서의 값, 즉
$\mathcal{F}(X)$이다. 그러므로 모든 $X \in \Ob(\mathcal{C})$에 대하여
$u_p\mathcal{F}(u(X)) = \mathcal{F}(X)$이다. $u$가 쌍대연속이고 연속이므로
$\mathcal{D}$에서 $u(X)$의 임의의 덮개는 $\mathcal{D}$의 덮개(!)
$\{u(X_i) \to u(X)\}$로 세분될 수 있고, 여기서
$\{X_i \to X\}$는 $\mathcal{C}$의 덮개이다. $u(X)$에서
$(u_p\mathcal{F})^+$의 값을 정의하는 여극한에서는 위와 같은 덮개들의 공종계
$\{u(X_i) \to u(X)\}$에 국한할 수 있으므로, 이는
$(u_p\mathcal{F})^+(u(X)) = \mathcal{F}(X)$도 함의한다. 따라서
$\mathcal{C}$의 모든 대상 $X$에 대하여
$(u_p\mathcal{F})^+(u(X)) = \mathcal{F}(X)$임도 알 수 있다. 이 논증을 한 번
더 반복하면 $\mathcal{C}$의 모든 대상 $X$에 대하여
$(u_p\mathcal{F})^\#(u(X)) = \mathcal{F}(X)$를 얻는다. 이로써 원하는 등식
$g^{-1}g_!\mathcal{F} = \mathcal{F}$를 얻는다.
\end{proof}

\noindent
마지막으로 대응하는 위상적 예가 없는 경우를 제시한다. 이 보조정리를 사용하여
같은 위상을 유지하면서 스킴들의 ``부분 우주''를 확대할 때 무슨 일이
일어나는지 살펴볼 것이다. 보조정리의 상황에서 토포스의 사상
$g : \Sh(\mathcal{C}) \to \Sh(\mathcal{D})$는 $\Sh(\mathcal{C})$를
$\Sh(\mathcal{D})$의 부분토포스로 식별하고
(Section \ref{sites-section-subtopoi}), 더 나아가 주어진 매장은 수축을 갖는다.

\begin{lemma}
\label{sites-lemma-bigger-site}
$\mathcal{C}$와 $\mathcal{D}$를 사이트라 하고
$u : \mathcal{C} \to \mathcal{D}$를 함자라 하자. 다음을 가정하자.
\begin{enumerate}
\item[(a)] $u$는 쌍대연속이고,
\item[(b)] $u$는 연속이며,
\item[(c)] $u$는 충실충만하고,
\item[(d)] $\mathcal{C}$에서 섬유곱이 존재하며 $u$는 이들과 가환하고,
\item[(e)] $u(e_\mathcal{C}) = e_\mathcal{D}$를 만족하는 종대상
$e_\mathcal{C} \in \Ob(\mathcal{C})$와
$e_\mathcal{D} \in \Ob(\mathcal{D})$가 존재한다.
\end{enumerate}
$g_!, g^{-1}, g_*$를 위와 같이 두자. 그러면 $u$는
$f_* = g^{-1}$, $f^{-1} = g_!$인 사이트의 사상
$f : \mathcal{D} \to \mathcal{C}$를 정의한다. 합성
$$
\xymatrix{
\Sh(\mathcal{C}) \ar[r]^g &
\Sh(\mathcal{D}) \ar[r]^f &
\Sh(\mathcal{C})
}
$$
은 토포스 $\Sh(\mathcal{C})$의 항등사상과 동형이다. 더 나아가 함자
$f^{-1}$는 충실충만하다.
\end{lemma}

\begin{proof}
가정에 의해 함자 $u$는 Proposition \ref{sites-proposition-get-morphism}의 가정을
만족한다. 따라서 $u$는 사이트의 사상을 정의하고, 그러므로
Lemma \ref{sites-lemma-morphism-sites-topoi}에서와 같은 토포스의 사상 $f$를
정의한다. 등식 $f_* = g^{-1}$와 $f^{-1} = g_!$는 인용한 보조정리와
Lemma \ref{sites-lemma-when-shriek}에서 명백하다. Lemma
\ref{sites-lemma-back-and-forth}에 의해
$f_* \circ g_* = g^{-1} \circ g_* \cong \text{id}$이고,
$g^{-1} \circ f^{-1} = g^{-1} \circ g_! \cong \text{id}$이다.

\medskip\noindent
이제 $f^{-1}$가 충실충만함을 보여야 한다.
$\mathcal{F}, \mathcal{G} \in \Ob(\Sh(\mathcal{C}))$라 하자. 사상
$$
\Mor_{\Sh(\mathcal{C})}(\mathcal{F}, \mathcal{G})
\longrightarrow
\Mor_{\Sh(\mathcal{D})}(f^{-1}\mathcal{F}, f^{-1}\mathcal{G})
$$
이 전단사임을 보이면 된다. 그런데 우변은 다음과 같다.
\begin{align*}
\Mor_{\Sh(\mathcal{D})}(f^{-1}\mathcal{F}, f^{-1}\mathcal{G})
& =
\Mor_{\Sh(\mathcal{C})}(\mathcal{F}, f_*f^{-1}\mathcal{G}) \\
& =
\Mor_{\Sh(\mathcal{C})}(\mathcal{F}, g^{-1}f^{-1}\mathcal{G}) \\
& =
\Mor_{\Sh(\mathcal{C})}(\mathcal{F}, \mathcal{G})
\end{align*}
(첫 번째 등호는 수반성에 의한다.) 이것이 원하는 바를 증명한다.
\end{proof}

\begin{example}
\label{sites-example-closed-map-cocontinuous-false}
$X$를 위상공간이라 하자. $i : Z  \to X$를 (유도된 위상을 갖는)
부분집합의 포함이라 하자. 함자 $u : X_{Zar} \to Z_{Zar}$,
$U \mapsto u(U) = Z \cap U$를 생각하자. 언뜻 보면 이 함자도 쌍대연속인
것처럼 보일 수 있다. 결국 $Z$가 유도된 위상을 가지므로
% P01-ERRATA: P01-E0087 -- the frozen interrogative has a spurious pronoun.
$U\cap Z$의 임의의 덮개는 $X$에서 $U$의 덮개로부터 와야 하지 않는가?
그렇지 않다! 실제로 $U \cap Z = \emptyset$이면 어떻게 되는가? 이 경우
공 덮개는 $U \cap Z$의 덮개이고, 공 덮개는 공 덮개로만 세분될 수 있다.
따라서 다음을 얻는다.
% P01-ERRATA: P01-E0088 -- the frozen English predicate omits the copula is.
$u$ 쌍대연속 $\Rightarrow$ $X$의 모든 공집합이 아닌 열린집합 $U$는
$Z$와 공집합이 아닌 교집합을 갖는다. 그러나 이것으로 충분하지 않다.
예를 들어 $X = \mathbf{R}$가 보통 위상을 갖는 실수직선이고
$Z = \mathbf{R} \setminus \{0\}$이면, $Z$의 열린 덮개
$Z = \{x < 0\} \cup \bigcup_n \{1/n < x\}$가 있는데, 이는 $X$의 어떤
열린 덮개의 제한으로도 세분될 수 없다.
\end{example}

\section{오른쪽 수반을 갖는 쌍대연속 함자}
\label{sites-section-cocontinuous-adjoint}

\noindent
$\mathcal{C}$와 $\mathcal{D}$를 사이트라 하자.
$u : \mathcal{C} \to \mathcal{D}$와 $v : \mathcal{D} \to \mathcal{C}$를
기저 범주의 함자라 하고, $v$가 $u$의 오른쪽 수반이라고 하자. 이 경우
$v$가 연속이면 $u$는 쌍대연속이다
(Lemma \ref{sites-lemma-continuous-with-continuous-left-adjoint}). $u$가
쌍대연속이면 흔히 (항상 그렇지는 않다.
Example \ref{sites-example-not-equivalent}를 보라.) $v$가 연속이고, 그 경우
$v$는 사이트의 사상을 정의하며 이에 결부된 토포스의 사상은 $u$가 정의하는
것과 같다.

\begin{lemma}
\label{sites-lemma-have-functor-other-way}
$\mathcal{C}$와 $\mathcal{D}$를 사이트라 하자.
$u : \mathcal{C} \to \mathcal{D}$와 $v : \mathcal{D} \to \mathcal{C}$를
함자라 하자. $u$가 쌍대연속이고 $v$가 $u$의 오른쪽 수반이라고 가정하자.
$g : \Sh(\mathcal{C}) \to \Sh(\mathcal{D})$를 $u$에 결부된 토포스의
사상이라 하자. Lemma \ref{sites-lemma-cocontinuous-morphism-topoi}를 보라.
그러면 다음이 성립한다.
\begin{enumerate}
\item $\mathcal{C}$ 위의 층 $\mathcal{F}$에 대하여 층
$g_*\mathcal{F}$는 준층 $v^p\mathcal{F}$와 같다. 다시 말해
$(g_*\mathcal{F})(V) = \mathcal{F}(v(V))$이다.
\item $\mathcal{D}$ 위의 층 $\mathcal{G}$에 대하여
$g^{-1}\mathcal{G} = (v_p\mathcal{G})^\#$이다.
\end{enumerate}
\end{lemma}

\begin{proof}
(1)의 $\mathcal{F}$에 대하여
$$
g_*\mathcal{F} =
{}_su\mathcal{F} =
{}_pu\mathcal{F} =
v^p\mathcal{F} =
\mathcal{F} \circ v
$$
이다. 첫 번째 등식은 Lemma \ref{sites-lemma-cocontinuous-morphism-topoi}이다.
두 번째 등식은 Lemma \ref{sites-lemma-pu-sheaf}이다. 세 번째 등식은
Lemma \ref{sites-lemma-adjoint-functors}이다. 마지막 등식은
Section \ref{sites-section-functoriality-PSh}에 있는 $v^p$의 정의이다. 이로써
(1)이 증명된다. (2)의 $\mathcal{G}$에 대하여
$$
g^{-1}\mathcal{G} = (u^p\mathcal{G})^\# = (v_p\mathcal{G})^\#
$$
이다. 첫 번째 등식은 Lemma \ref{sites-lemma-cocontinuous-morphism-topoi}이고,
두 번째 등식은 Lemma \ref{sites-lemma-adjoint-functors}이다.
\end{proof}

\begin{lemma}
\label{sites-lemma-have-functor-other-way-morphism}
표기와 가정을 Lemma \ref{sites-lemma-have-functor-other-way}에서와 같이 두자.
추가로 $v$가 연속이면 $v$는 사이트의 사상
$f : \mathcal{C} \to \mathcal{D}$를 정의하고, 이에 결부된 토포스의 사상은
$g$와 같다.
\end{lemma}

\begin{proof}
Lemma \ref{sites-lemma-have-functor-other-way}의 결과를 더 언급하지 않고 사용한다.
$v$가 보조정리의 명제와 같은 사이트의 사상 $f$를 정의함을 증명하려면
$v_s$가 완전 함자임을 보여야 한다
(Definition \ref{sites-definition-morphism-sites}를 보라). 그런데
$v_s\mathcal{G} = (v_p\mathcal{G})^\# = g^{-1}\mathcal{G}$이므로 이는
$g$가 토포스의 사상이라는 사실에서 따른다. 이제
$f^{-1} = v_s = g^{-1}$임을 알 수 있고, 토포스의 사상으로서 $f = g$를
얻는다.
\end{proof}

\begin{example}
\label{sites-example-finer-topology-bis}
이 예는 Example \ref{sites-example-finer-topology}의 논의를 이어가며, 그곳의 표기
$\mathcal{C}, \tau, \tau', \epsilon$을 빌려 쓴다. 항등함자
$v : \mathcal{C}_{\tau'} \to \mathcal{C}_\tau$는 연속 함자이고, 항등함자
$u : \mathcal{C}_\tau \to \mathcal{C}_{\tau'}$는 쌍대연속 함자임을
관찰하자. 또한 $u$는 $v$의 왼쪽 수반이다. 따라서 Lemmas
\ref{sites-lemma-have-functor-other-way}와
\ref{sites-lemma-have-functor-other-way-morphism}의 결과가 적용되고, $v$가
사이트의 사상
$$
\epsilon : \mathcal{C}_\tau \longrightarrow \mathcal{C}_{\tau'}
$$
을 정의하며 이에 대응하는 토포스의 사상은 쌍대연속 함자 $u$에 결부된
토포스의 사상과 같다는 결론을 얻는다.
\end{example}

\begin{lemma}
\label{sites-lemma-continuous-with-continuous-left-adjoint}
$\mathcal{C}$와 $\mathcal{D}$를 사이트라 하고
$v : \mathcal{D} \to \mathcal{C}$를 연속 함자라 하자. $v$가 왼쪽 수반
$u : \mathcal{C} \to \mathcal{D}$를 갖는다고 가정하자. 그러면
\begin{enumerate}
\item $u$는 쌍대연속이다.
\item Lemmas \ref{sites-lemma-have-functor-other-way}와
\ref{sites-lemma-have-functor-other-way-morphism}의 결과가 성립한다.
\end{enumerate}
특히 $v$는 사이트의 사상 $f : \mathcal{C} \to \mathcal{D}$를 정의한다.
\end{lemma}

\begin{proof}
$U$를 $\mathcal{C}$의 대상이라 하고
$\{V_j \to u(U)\}_{j \in J}$를 $\mathcal{D}$의 덮개라 하자. 그러면
% P01-ERRATA: P01-E0089 -- the frozen conclusion changes j in J to undefined i in I.
$\{v(V_j) \to v(u(U))\}_{i \in I}$는 $\mathcal{C}$의 덮개이다. 수반사상
$U \to v(u(U))$를 통해 이를 기저변환하여 덮개
$\{W_j \to U\}_{j \in J}$를 얻으며, 여기서
$W_j = v(V_j) \times_{v(u(U))} U$이다. 첫 번째 사영을
$p_j : W_j \to v(V_j)$로 나타내면 사상
$$
u(W_j) \xrightarrow{u(p_j)} u(v(V_j)) \longrightarrow V_j
$$
를 얻으며, 여기서 두 번째 화살표는 수반사상이다. 이는 공역이 고정된
사상족의 사상
$\{u(W_j) \to u(U)\} \to \{V_j \to u(U)\}$를 정하므로, $u$가 실제로
쌍대연속임을 보인다. 나머지 주장들은 Lemmas
\ref{sites-lemma-have-functor-other-way}와
\ref{sites-lemma-have-functor-other-way-morphism}에서 즉시 따른다.
\end{proof}

\begin{example}
\label{sites-example-not-equivalent}
$\mathcal{C}$와 $\mathcal{D}$를 사이트라 하자.
$u : \mathcal{C} \to \mathcal{D}$와 $v : \mathcal{D} \to \mathcal{C}$를
기저 범주의 함자라 하고, $v$가 $u$의 오른쪽 수반이라고 하자.
Lemma \ref{sites-lemma-continuous-with-continuous-left-adjoint}에 따르면 $v$가
연속이면
% P01-ERRATA: P01-E0090 -- the frozen source misspells cocontinuous.
$u$는 쌍대연속이다. 역으로 $u$가 쌍대연속이어도 $v$가 연속이라고 결론내릴
수는 없다. 스킴의 큰 \'etale 사이트와 큰 매끄러운 사이트를 사용하여 이
현상의 예를 주겠지만, 아마 초보적인 예도 있을 것이다. 구체적으로 스킴
$S$와 사이트 $(\Sch/S)_\etale$ 및 $(\Sch/S)_{smooth}$를 생각하자. 이
사이트들은 같은 기저 범주를 갖는다고 가정해도 된다. Topologies,
Remark \ref{topologies-remark-choice-sites}를 보라. $u = v = \text{id}$로
놓자. $(\Sch/S)_\etale$에서 $(\Sch/S)_{smooth}$로 가는 함자로서 $u$는
쌍대연속이다. 실제로 스킴의 모든 매끄러운 덮개는 \'etale 덮개로 세분될
수 있다. More on Morphisms,
Lemma \ref{more-morphisms-lemma-etale-dominates-smooth}를 보라. 역으로
$(\Sch/S)_{smooth}$에서 $(\Sch/S)_\etale$로 가는 함자 $v$는 연속이 아니다.
일반적으로 매끄러운 덮개는 \'etale 덮개가 아니기 때문이다.
\end{example}

\section{왼쪽 수반을 갖는 쌍대연속 함자}
\label{sites-section-cocontinuous-left-adjoint}

\noindent
쌍대연속 함자 $u$가 왼쪽 수반 $w$를 가질 수도 있다.

\begin{lemma}
\label{sites-lemma-have-left-adjoint}
$\mathcal{C}$와 $\mathcal{D}$를 사이트라 하자.
$g : \Sh(\mathcal{C}) \to \Sh(\mathcal{D})$를 연속이고 쌍대연속인 함자
$u : \mathcal{C} \to \mathcal{D}$에 결부된 토포스의 사상이라 하자.
Lemmas \ref{sites-lemma-cocontinuous-morphism-topoi}와
\ref{sites-lemma-when-shriek}를 보라.
\begin{enumerate}
\item $w : \mathcal{D} \to \mathcal{C}$가 $u$의 왼쪽 수반이면 다음이
성립한다.
\begin{enumerate}
\item $g_!\mathcal{F}$는 준층 $w^p\mathcal{F}$에 결부된 층이다.
\item $g_!$는 완전하다.
\end{enumerate}
\item $w$가 연속인 왼쪽 수반이면 $g_!$는 왼쪽 수반을 갖는다.
\item $w$가 쌍대연속인 왼쪽 수반이면 $g_! = h^{-1}$이고
$g^{-1} = h_*$이다. 여기서 $h : \Sh(\mathcal{D}) \to \Sh(\mathcal{C})$는
$w$에 결부된 토포스의 사상이다.
\end{enumerate}
\end{lemma}

\begin{proof}
$g_!\mathcal{F}$는 $u_p\mathcal{F}$의 층화임을 상기하자. 따라서 (1)(a)는
Lemma \ref{sites-lemma-adjoint-functors}에 의한 $u_p = w^p$에서 따른다.

\medskip\noindent
(1)(b)를 보기 위하여, $g_!$가 모든 여극한과 가환함에 유의하자. 이는
$g_!$가 왼쪽 수반이기 때문이다
(Categories, Lemma \ref{categories-lemma-adjoint-exact}).
$i \mapsto \mathcal{F}_i$를 $\Sh(\mathcal{C})$의 유한 도표라 하자. 그러면
$\lim \mathcal{F}_i$는 준층의 범주에서 계산된다
(Lemma \ref{sites-lemma-limit-sheaf}). $w^p$는 오른쪽 수반이므로
(Lemma \ref{sites-lemma-adjoints-u})
$w^p \lim \mathcal{F}_i = \lim w^p\mathcal{F}_i$임을 안다. 층화가
완전하므로(Lemma \ref{sites-lemma-sheafification-exact}) (1)(a)에 의해 결론을
얻는다.

\medskip\noindent
$w$가 연속이라고 가정하자. 그러면 $g_! = (w^p\ )^\# = w^s$이지만 층화는
필요하지 않고 왼쪽 수반 $w_s$가 있다. Lemmas
\ref{sites-lemma-pushforward-sheaf}와 \ref{sites-lemma-adjoint-sheaves}를 보라.

\medskip\noindent
$w$가 쌍대연속이라고 가정하자. 등식 $g_! = h^{-1}$는 (1)(a)와 정의들에서
따른다. 등식 $g^{-1} = h_*$는 $g_! = h^{-1}$와 수반함자의 유일성에서
따른다. 또는 Lemma \ref{sites-lemma-have-functor-other-way}에서 이를 도출할 수
있다.
\end{proof}









\section{하부 슈리크의 존재}
\label{sites-section-lower-shriek}

\noindent
이 절에서는 $f^{-1}$가 왼쪽 수반 $f_!$를 갖는 토포스의 사상 $f$의 몇 가지
경우를 논의한다.

\begin{lemma}
\label{sites-lemma-existence-lower-shriek}
$\mathcal{C}$, $\mathcal{D}$를 두 사이트라 하자.
$f : \Sh(\mathcal{D}) \to \Sh(\mathcal{C})$를 토포스의 사상이라 하자.
$E \subset \Ob(\mathcal{D})$를 다음 조건을 만족하는 부분집합이라 하자.
\begin{enumerate}
\item $V \in E$에 대하여 $\mathcal{C}$ 위의 층 $\mathcal{G}$가 존재하여,
$\Sh(\mathcal{C})$의 $\mathcal{F}$에 대하여 함자적으로
$f^{-1}\mathcal{F}(V) = 
\Mor_{\Sh(\mathcal{C})}(\mathcal{G}, \mathcal{F})$이다.
\item $\mathcal{D}$의 모든 대상은 $E$의 대상들로 이루어진 덮개를 갖는다.
\end{enumerate}
그러면 $f^{-1}$는 왼쪽 수반 $f_!$를 갖는다.
\end{lemma}

\begin{proof}
요네다 보조정리(Categories, Lemma \ref{categories-lemma-yoneda})에 의해
$V \in E$에 대응하는 층 $\mathcal{G}_V$는 공식
\[
f^{-1}\mathcal{F}(V) = \Mor_{\Sh(\mathcal{C})}(\mathcal{G}_V, \mathcal{F})
\]
에 의해 유일한 동형을 제외하고 정해진다. 다음을 상기하자.
$f^{-1}\mathcal{F}(V) = \Mor_{\Sh(\mathcal{D})}(h_V^\#, f^{-1}\mathcal{F})$.
$\Mor(\mathcal{G}_V, \mathcal{G}_V)$의 $\text{id}$에 대응하는 사상을
$i_V : h_V^\# \to f^{-1}\mathcal{G}_V$로 나타내자. (1)의 함자성에 의해
전단사는 다음과 같이 주어진다.
$$
\Mor_{\Sh(\mathcal{C})}(\mathcal{G}_V, \mathcal{F}) \to
\Mor_{\Sh(\mathcal{D})}(h_V^\#, f^{-1}\mathcal{F}),\quad
\varphi \mapsto f^{-1}\varphi \circ i_V
$$
임의의 $V_1, V_2 \in E$에 대하여 표준적인 사상
$$
\Mor_{\Sh(\mathcal{D})}(h^\#_{V_2}, h^\#_{V_1})
\to
\Hom_{\Sh(\mathcal{C})}(\mathcal{G}_{V_2}, \mathcal{G}_{V_1}),\quad
\varphi \mapsto f_!(\varphi)
$$
이 있으며, 이는
$f^{-1}(f_!(\varphi)) \circ i_{V_2} = i_{V_1} \circ \varphi$로 특징지어진다.
$\varphi \mapsto f_!(\varphi)$가 합성과 양립함에 유의하자. 이는 그 특징에서
직접 알 수 있다. 따라서 $h_V^\# \mapsto \mathcal{G}_V$ 및
$\varphi \mapsto f_!\varphi$는 대상이 $h_V^\#$인 $\Sh(\mathcal{D})$의
충만부분범주에서 정의된 함자이다.

\medskip\noindent
$J$를 집합이라 하고 $J \to E$, $j \mapsto V_j$를 사상이라 하자. 그러면
위의 전단사들의 곱을 사용하여 함자적인 전단사
$$
\Mor_{\Sh(\mathcal{C})}(\coprod \mathcal{G}_{V_j}, \mathcal{F})
\longrightarrow
\Mor_{\Sh(\mathcal{D})}(\coprod h_{V_j}^\#, f^{-1}\mathcal{F})
$$
를 얻는다. 따라서 함자 $f_!$를 대상이 $V \in E$에 대한 $h_V^\#$들의
쌍대곱인 $\Sh(\mathcal{D})$의 충만부분범주로 확장할 수 있다.

\medskip\noindent
$\mathcal{D}$ 위의 임의의 층 $\mathcal{H}$가 주어지면 쌍대등화자 도표
$$
\xymatrix{
\mathcal{H}_1 \ar@<1ex>[r] \ar@<-1ex>[r] &
\mathcal{H}_0 \ar[r] &
\mathcal{H}
}
$$
를 택하자. 여기서 $\mathcal{H}_i = \coprod h_{V_{i, j}}^\#$는
$V_{i, j} \in E$인 쌍대곱이다. 이는 가정 (2)에 의해 가능하다.
Lemma \ref{sites-lemma-sheaf-coequalizer-representable}을 보라(집합론적 문제를
걱정하는 독자는 Lemma \ref{sites-lemma-sheaf-coequalizer-representable}에서 준
구성이 표준적임에 유의하라). $f_!(\mathcal{H})$를 다음 도표가 쌍대등화자
도표가 되게 하는 $\mathcal{C}$ 위의 층으로 정의하자.
$$
\xymatrix{
f_!\mathcal{H}_1 \ar@<1ex>[r] \ar@<-1ex>[r] &
f_!\mathcal{H}_0 \ar[r] &
f_!\mathcal{H}
}
$$
그러면
\begin{align*}
\Mor(f_!\mathcal{H}, \mathcal{F})
& =
\text{Equalizer}(
\xymatrix{
\Mor(f_!\mathcal{H}_0, \mathcal{F}) \ar@<1ex>[r] \ar@<-1ex>[r] &
\Mor(f_!\mathcal{H}_1, \mathcal{F})
}
) \\
& =
\text{Equalizer}(
\xymatrix{
\Mor(\mathcal{H}_0, f^{-1}\mathcal{F}) \ar@<1ex>[r] \ar@<-1ex>[r] &
\Mor(\mathcal{H}_1, f^{-1}\mathcal{F})
}
) \\
& =
\Hom(\mathcal{H}, f^{-1}\mathcal{F})
\end{align*}
이다. 따라서 $f_!$를 $\mathcal{D}$ 위의 층들의 전체 범주로 확장할 수 있음을
알 수 있다.
\end{proof}

\section{국소화}
\label{sites-section-localize}

\noindent
$\mathcal{C}$를 사이트라 하고 $U \in \Ob(\mathcal{C})$라 하자. $U$ 위의
대상들의 범주 $\mathcal{C}/U$의 정의는 Categories,
Example \ref{categories-example-category-over-X}를 보라. $U$ 위의 대상들의
사상족 $\{V_j \to V\}$가 $\mathcal{C}$에서 덮개일 필요충분조건으로
$\mathcal{C}/U$에서 덮개라고 선언하여 $\mathcal{C}/U$를 사이트로 만든다.
망각함자
$$
j_U : \mathcal{C}/U \longrightarrow \mathcal{C}.
$$
를 생각하자. 이는 명백히 쌍대연속이고 연속이다. 따라서
Lemma \ref{sites-lemma-when-shriek}에 의해 $j_U^{-1}$와 $j_{U*}$로 주어지는
토포스의 사상
$$
j_U : \Sh(\mathcal{C}/U) \longrightarrow \Sh(\mathcal{C})
$$
과 $j_U^{-1}$의 왼쪽 수반인 함자 $j_{U!}$를 얻는다.

\begin{definition}
\label{sites-definition-localize}
$\mathcal{C}$를 사이트라 하고 $U \in \Ob(\mathcal{C})$라 하자.
\begin{enumerate}
\item 사이트 $\mathcal{C}/U$를 {\it 대상 $U$에서 사이트 $\mathcal{C}$의
국소화}라 한다.
\item 토포스의 사상 $j_U : \Sh(\mathcal{C}/U) \to \Sh(\mathcal{C})$를
{\it 국소화 사상}이라 한다.
\item 함자 $j_{U*}$를 {\it 직접상 함자}라 한다.
\item $\mathcal{C}$ 위의 층 $\mathcal{F}$에 대하여 층
$j_U^{-1}\mathcal{F}$를 {$\mathcal{C}/U$에 대한 $\mathcal{F}$의
{\it 제한}}이라 한다.
\item $\mathcal{C}/U$ 위의 층 $\mathcal{G}$에 대하여 층
$j_{U!}\mathcal{G}$를 {$\mathcal{G}$의 {\it 공집합에 의한 확장}}이라 한다.
\end{enumerate}
\end{definition}

\noindent
예상대로 제한 $j_U^{-1}\mathcal{F}$는 규칙
$j_U^{-1}\mathcal{F}(X/U) = \mathcal{F}(X)$로 정의되는 층이다. 이 경우
공집합에 의한 확장도 아주 간단히 기술할 수 있다. 다음과 같다.

\begin{lemma}
\label{sites-lemma-describe-j-shriek}
$\mathcal{C}$를 사이트라 하고 $U \in \Ob(\mathcal{C})$라 하자.
$\mathcal{G}$를 $\mathcal{C}/U$ 위의 준층이라 하자. 그러면
$j_{U!}(\mathcal{G}^\#)$는 명백한 제한사상들을 갖는 준층
$$
V
\longmapsto
\coprod\nolimits_{\varphi \in \Mor_\mathcal{C}(V, U)}
\mathcal{G}(V \xrightarrow{\varphi} U)
$$
에 결부된 층이다.
\end{lemma}

\begin{proof}
Lemma \ref{sites-lemma-when-shriek}에 의해
$j_{U!}(\mathcal{G}^\#) = ((j_U)_p\mathcal{G}^\#)^\#$이다.
Lemma \ref{sites-lemma-technical-up}에 의해 이는
$((j_U)_p\mathcal{G})^\#$와 같다. 따라서 $\mathcal{C}/U$ 위의 임의의 준층
$\mathcal{G}$에 대하여 $(j_U)_p$가 위의 공식으로 주어짐을 증명하면 충분하다.
% P01-ERRATA: P01-E0091 -- the frozen proof contains the informal editing remnant "OK,".
Section \ref{sites-section-functoriality-PSh}의 정의에 의해
$$
(j_U)_p\mathcal{G}(V)
=
\colim_{(W/U, V \to W)} \mathcal{G}(W)
$$
이다. 이제 쌍 $(W/U, V \to W)$들의 범주는 각 $\varphi : V \to U$에 대하여
대상 $O_\varphi = (\varphi : V \to U, \text{id} : V \to V)$를 가지며, 더
나아가 모든 대상에 대하여 $O_\varphi$ 중 하나에서 그 대상으로 가는 유일한
사상이 있음은 명백하다. 결과가 따른다.
\end{proof}

\begin{lemma}
\label{sites-lemma-describe-j-shriek-representable}
$\mathcal{C}$를 사이트라 하고 $U \in \Ob(\mathcal{C})$라 하자.
$X/U$를 $\mathcal{C}/U$의 대상이라 하자. 그러면
$j_{U!}(h_{X/U}^\#) = h_X^\#$이다.
\end{lemma}

\begin{proof}
$p : X \to U$를 $X$의 구조사상으로 나타내자.
Lemma \ref{sites-lemma-describe-j-shriek}에 의해
$j_{U!}(h_{X/U}^\#)$는 준층
$$
V
\longmapsto
\coprod\nolimits_{\varphi \in \Mor_\mathcal{C}(V, U)}
\{\psi : V \to X \mid p \circ \psi = \varphi\}
$$
에 결부된 층이다. 이는 명백히 $\Mor_\mathcal{C}(V, X)$와 같다. 따라서
보조정리가 따른다.
\end{proof}

\noindent
위의 두 보조정리 중 어느 것을 사용해도 $j_{U!}(*) = h_U^\#$를 얻는다.
따라서 $\mathcal{C}/U$ 위의 모든 층 $\mathcal{G}$에 대하여 층의 표준적인
사상 $j_{U!}\mathcal{G} \to h_U^\#$가 있다. 이는 $j_{U!}$의 본질적 상에
속하는 층들을 특징짓는다.

\begin{lemma}
\label{sites-lemma-essential-image-j-shriek}
$\mathcal{C}$를 사이트라 하고 $U \in \Ob(\mathcal{C})$라 하자. 함자
$j_{U!}$는 범주의 동치
$$
\Sh(\mathcal{C}/U)
\longrightarrow
\Sh(\mathcal{C})/h_U^\#
$$
를 준다.
\end{lemma}

\begin{proof}
$\mathcal{C}/U$의 대상을 쌍 $(X, a)$로 나타내자. 여기서 $X$는
$\mathcal{C}$의 대상이고 $a : X \to U$는 $\mathcal{C}$의 사상이다.
마찬가지로 $\Sh(\mathcal{C})/h_U^\#$의 대상을 쌍
$(\mathcal{F}, \varphi)$로 나타내자. 함자
$\Sh(\mathcal{C}/U) \to \Sh(\mathcal{C})/h_U^\#$는 $\mathcal{G}$를 쌍
$(j_{U!}\mathcal{G}, \gamma)$로 보내며, 여기서 $\gamma$는
$j_{U!}\mathcal{G} \to j_{U!}*$와 식별 $j_{U!}* = h_U^\#$의 합성이다.

\medskip\noindent
$\Sh(\mathcal{C})/h_U^\#$에서 $\Sh(\mathcal{C}/U)$로 가는 함자를 구성하자.
$(\mathcal{F}, \varphi)$가 주어졌다고 하자. $\mathcal{C}/U$의 대상
$(X, a)$에 대하여 $\mathcal{F}_\varphi(X, a)$를 다음과 같은 원소
$s \in \mathcal{F}(X)$들의 집합으로 정의한다. 즉, 각 원소의 $\varphi$에 의한
상은 $a \in \Mor_\mathcal{C}(X, U) = h_U(X)$의 $h_U^\#(X)$에서의 상이다.
$(X, a) \mapsto \mathcal{F}_\varphi(X, a)$가 $\mathcal{C}/U$ 위의 층임은
쉽게 알 수 있다. 규칙 $(\mathcal{F}, \varphi) \mapsto \mathcal{F}_\varphi$는
명백히 함자 $\Sh(\mathcal{C})/h_U^\# \to \Sh(\mathcal{C}/U)$를 정의한다.

\medskip\noindent
또한 함자
$\textit{PSh}(\mathcal{C})/h_U \to \textit{PSh}(\mathcal{C}/U)$,
$(\mathcal{F}, \varphi) \mapsto \mathcal{F}_\varphi$를 생각하자. 여기서
$\mathcal{F}_\varphi(X, a)$는 $\mathcal{F}(X)$의 원소 중 $a \in h_U(X)$로
사상되는 것들의 집합으로 정의한다. 수직 화살표가 층화로 주어지는 도표
$$
\xymatrix{
\textit{PSh}(\mathcal{C})/h_U \ar[r] \ar[d] &
\textit{PSh}(\mathcal{C}/U) \ar[d] \\
\Sh(\mathcal{C})/h_U^\# \ar[r] &
\Sh(\mathcal{C}/U)
}
$$
가 가환한다고 주장한다. 이를 보려면\footnote{다른 방법으로, 준층에 대해서는
$\mathcal{F}_\varphi$를 데카르트 도표
$$
\vcenter{
\xymatrix{
\mathcal{F}_\varphi \ar[r] \ar[d] & {*} \ar[d] \\
\mathcal{F}|_{\mathcal{C}/U} \ar[r] & h_U|_{\mathcal{C}/U}
}
}
\quad\text{for presheaves and}\quad
\vcenter{
\xymatrix{
\mathcal{F}_\varphi \ar[r] \ar[d] & {*} \ar[d] \\
\mathcal{F}|_{\mathcal{C}/U} \ar[r] & h_U^\#|_{\mathcal{C}/U}
}
}
$$
로 기술하고, 층에 대해서도 마찬가지로 한 뒤 $\mathcal{C}/U$에 대한 제한이
층화와 가환함을 사용할 수 있다.} Lemmas \ref{sites-lemma-plus-presheaf}와
\ref{sites-lemma-plus-functorial} 및 Theorem \ref{sites-theorem-plus}의 함자
$\mathcal{F} \mapsto \mathcal{F}^+$와 이 구성이 가환함을 증명하면 충분하다.
$\mathcal{F} \mapsto \mathcal{F}^+$와의 가환성은, $(X, a)$가 주어졌을 때
$\mathcal{C}/U$에서 $(X, a)$의 덮개들의 범주와 $\mathcal{C}$에서 $X$의
덮개들의 범주가 표준적으로 식별된다는 사실에서 따른다.

\medskip\noindent
다음으로 $\textit{PSh}(\mathcal{C}/U) \to \textit{PSh}(\mathcal{C})/h_U$가
$\mathcal{G}$를 쌍 $(j_{U!}^{PSh}\mathcal{G}, \gamma)$로 보내게 하자.
% P01-ERRATA: P01-E0092 -- the frozen relative clause omits its verb.
여기서 $j_{U!}^{PSh}\mathcal{G}$는 Lemma \ref{sites-lemma-describe-j-shriek}의
공식으로 정의되는 준층이고, $\gamma$는
% P01-ERRATA: P01-E0093 -- the frozen target drops the PSh superscript.
$j_{U!}^{PSh}\mathcal{G} \to j_{U!}*$와 식별
$j_{U!}^{PSh}* = h_U$의 합성이다(공식에서 명백하다). 그러면 수직 화살표가
층화인 도표
$$
\xymatrix{
\textit{PSh}(\mathcal{C}/U) \ar[r] \ar[d] &
\textit{PSh}(\mathcal{C})/h_U \ar[d] \\
\Sh(\mathcal{C}/U) \ar[r] &
\Sh(\mathcal{C})/h_U^\#
}
$$
가 가환함은 즉시 알 수 있다. 모든 것을 종합하면
$\textit{PSh}(\mathcal{C}/U)$의 $\mathcal{G}$에 대하여 함자적인 동형
$(j_{U!}^{PSh}\mathcal{G})_\gamma = \mathcal{G}$가 존재하고,
$\textit{PSh}(\mathcal{C})/h_U$의 $(\mathcal{F}, \varphi)$에 대하여
$j_{U!}^{PSh}\mathcal{F}_\varphi = \mathcal{F}$임을 보이면 충분하다.
준층 $(j_{U!}^{PSh}\mathcal{G})_\gamma$의 $(X, a)$에서의 값은 사상
$$
\coprod\nolimits_{a' : X \to U} \mathcal{G}(X, a') \to \Mor_\mathcal{C}(X, U)
$$
의 $a$ 위의 섬유이며, 이는 $\mathcal{G}(X, a)$이다. 이로써 첫 번째 등식이
증명된다. 준층 $j_{U!}^{PSh}\mathcal{F}_\varphi$의
% P01-ERRATA: P01-E0094 -- the frozen sentence duplicates "is" around "on X".
$X$에서의 값은
$$
\coprod\nolimits_{a : X \to U} \mathcal{F}_\varphi(X, a) =
\mathcal{F}(X)
$$
이다. 실제로 집합의 사상 $S \to S'$가 주어지면 집합 $S$는 그 섬유들의
서로소 합집합이다.
\end{proof}

\noindent
Lemma \ref{sites-lemma-essential-image-j-shriek}에 따르면 함자 $j_{U!}$는 합성
$$
\Sh(\mathcal{C}/U) \rightarrow
\Sh(\mathcal{C})/h_U^\# \rightarrow
\Sh(\mathcal{C})
$$
이며, 여기서 첫 번째 화살표는 동치이다.

\begin{lemma}
\label{sites-lemma-j-shriek-commutes-equalizers-fibre-products}
$\mathcal{C}$를 사이트라 하고 $U \in \Ob(\mathcal{C})$라 하자. 함자
$j_{U!}$는 섬유곱 및 등화자와 가환한다(더 일반적으로 유한 연결 극한과
가환한다). 특히 $\Sh(\mathcal{C}/U)$에서
$\mathcal{F} \subset \mathcal{F}'$이면
$j_{U!}\mathcal{F} \subset j_{U!}\mathcal{F}'$이다.
\end{lemma}

\begin{proof}
Lemma \ref{sites-lemma-essential-image-j-shriek}와 범주의 동치가 모든 극한과
가환한다는 사실을 사용하면, 이는 함자
$\Sh(\mathcal{C})/h_U^\# \rightarrow \Sh(\mathcal{C})$가 섬유곱 및
등화자와 가환한다는 사실로 귀착된다. 또는 층화가 완전하다는 사실과
Lemma \ref{sites-lemma-describe-j-shriek}의 $j_{U!}$에 대한 기술을 사용하여 이를
직접 증명할 수 있다. (또한 $\mathcal{C}$가 섬유곱과 등화자를 갖는 경우에는
Lemma \ref{sites-lemma-preserve-equalizers}에서 결과가 따른다.)
\end{proof}

\begin{lemma}
\label{sites-lemma-j-shriek-reflects-injections-surjections}
$\mathcal{C}$를 사이트라 하고 $U \in \Ob(\mathcal{C})$라 하자. 함자
$j_{U!}$는 단사 및 전사를 반영한다.
\end{lemma}

\begin{proof}
Lemma \ref{sites-lemma-mono-epi-sheaves}에 따라 $j_{U!}$가 모노사상 및
에피사상을 반영함을 보여야 한다. Lemma
\ref{sites-lemma-essential-image-j-shriek}를 사용하면 이는 함자
$\Sh(\mathcal{C})/h_U^\# \to \Sh(\mathcal{C})$가 모노사상 및
에피사상을 반영한다는 사실로 귀착된다.
\end{proof}

\begin{lemma}
\label{sites-lemma-compute-j-shriek-restrict}
$\mathcal{C}$를 사이트라 하고 $U \in \Ob(\mathcal{C})$라 하자.
$\mathcal{C}$ 위의 임의의 층 $\mathcal{F}$에 대하여
$j_{U!}j_U^{-1}\mathcal{F} = \mathcal{F} \times h_U^\#$이다.
\end{lemma}

\begin{proof}
이는 Lemma \ref{sites-lemma-describe-j-shriek}의 $j_{U!}$에 대한 기술에서 명백하다.
\end{proof}

\begin{lemma}
\label{sites-lemma-relocalize}
$\mathcal{C}$를 사이트라 하고 $f : V \to U$를 $\mathcal{C}$의 사상이라 하자.
그러면 연속이고 쌍대연속인 함자들의 가환도표
$$
\xymatrix{
\mathcal{C}/V \ar[rd]_{j_V} \ar[rr]_j & &
\mathcal{C}/U \ar[ld]^{j_U} \\
& \mathcal{C} &
}
$$
가 존재한다. 식별 $(\mathcal{C}/U)/(V/U) = \mathcal{C}/V$ 아래에서 함자
$j : \mathcal{C}/V \to \mathcal{C}/U$,
$(a : W \to V) \mapsto (f \circ a : W \to U)$는 함자
$j_{V/U} : (\mathcal{C}/U)/(V/U) \to \mathcal{C}/U$와 식별된다. 더 나아가
$j_{V!} = j_{U!} \circ j_!$, $j_V^{-1} = j^{-1} \circ j_U^{-1}$ 및
$j_{V*} = j_{U*} \circ j_*$가 성립한다.
\end{lemma}

\begin{proof}
도표의 가환성은 즉시 알 수 있다. $j$와 $j_{V/U}$가 일치한다는 것은
정의에서 따른다. Lemma \ref{sites-lemma-composition-cocontinuous}에 의해 다음
토포스 사상들의 도표가 가환함을 알 수 있다.
\begin{equation}
\label{sites-equation-relocalize}
\vcenter{
\xymatrix{
\Sh(\mathcal{C}/V) \ar[rd]_{j_V} \ar[rr]_j & &
\Sh(\mathcal{C}/U) \ar[ld]^{j_U} \\
& \Sh(\mathcal{C}) &
}
}
\end{equation}
이로써 $j_V^{-1} = j^{-1} \circ j_U^{-1}$ 및
$j_{V*} = j_{U*} \circ j_*$가 증명된다. 등식
$j_{V!} = j_{U!} \circ j_!$는 수반성의 성질에서 형식적으로 따른다.
\end{proof}

\begin{lemma}
\label{sites-lemma-relocalize-explicit}
Lemma \ref{sites-lemma-relocalize}에서와 같이 $\mathcal{C}$, $f : V \to U$,
$j_U$, $j_V$, $j$를 표기하자. Lemma
\ref{sites-lemma-essential-image-j-shriek}의 식별
$\Sh(\mathcal{C}/V) = \Sh(\mathcal{C})/h_V^\#$ 및
$\Sh(\mathcal{C}/U) = \Sh(\mathcal{C})/h_U^\#$ 아래에서 다음이 성립한다.
\begin{enumerate}
\item 함자 $j^{-1}$는 다음과 같이 기술된다.
$$
j^{-1}(\mathcal{H} \xrightarrow{\varphi} h_U^\#)
=
(\mathcal{H} \times_{\varphi, h_U^\#, f} h_V^\# \to h_V^\#).
$$
\item 함자 $j_!$는 다음과 같이 기술된다.
$$
j_!(\mathcal{H} \xrightarrow{\varphi} h_V^\#) =
(\mathcal{H} \xrightarrow{h_f \circ \varphi} h_U^\#)
$$
\end{enumerate}
\end{lemma}

\begin{proof}
(2)를 증명하자. 식별
$\Sh(\mathcal{C}/V) \to \Sh(\mathcal{C})/h_V^\#$는 $\mathcal{G}$를
$j_{V!}\mathcal{G} \to j_{V!}(*) = h_V^\#$로 보내며,
$\Sh(\mathcal{C}/U) \to \Sh(\mathcal{C})/h_U^\#$도 마찬가지임을 상기하자.
따라서 $j_!\mathcal{G}$는
$j_{U!}(j_!\mathcal{G}) \to j_{U!}(*) = h_U^\#$로 사상되고, Lemma
\ref{sites-lemma-relocalize}에 의해 $j_{U!}j_! = j_{V!}$이므로 (2)가 따른다.

\medskip\noindent
이제 독자는 $j^{-1}$가 $j_!$의 오른쪽 수반이라는 사실과 (1)의 규칙이 (2)의
규칙의 오른쪽 수반이라는 사실을 사용하여 (1)을 증명할 수 있다. 직접 증명은
다음과 같다. $\varphi : \mathcal{H} \to h_U^\#$가
$\Sh(\mathcal{C})/h_U^\#$의 대상이라고 하자. Lemma
\ref{sites-lemma-essential-image-j-shriek}의 증명에 의해 이는 $\mathcal{C}/U$ 위에서
규칙
$$
(a : W \to U)
\longmapsto
\{ s \in \mathcal{H}(W) \mid \varphi(s) = a\}
$$
로 정의되는 $\mathcal{C}/U$ 위의 층 $\mathcal{H}_\varphi$에 대응한다.
$\mathcal{C}/V$로의 당김
$j^{-1}\mathcal{H}_\varphi$는 규칙
$$
(a : W \to V)
\longmapsto
\{ s \in \mathcal{H}(W) \mid \varphi(s) = f \circ a\}
$$
로 주어진다.
% P01-ERRATA: P01-E0095 -- the frozen source reverses the localization indices in j_{U/V}^{-1}.
이는 $j^{-1} = j_{U/V}^{-1}$를 $\mathcal{H}_\varphi$의 $\mathcal{C}/V$에 대한
제한으로 기술한 데 따른다. 한편 $\Sh(\mathcal{C})/h_V^\#$의 대상
$$
\xymatrix{
\mathcal{H}' = \mathcal{H} \times_{\varphi, h_U^\#, f} h_V^\#
\ar[rr]^-{\varphi'} & & h_V^\#
}
$$
에 이 규칙을 적용하면
\begin{align*}
(a : W \to V)
\longmapsto
&  \{ s' \in \mathcal{H}'(W) \mid \varphi'(s') = a\} \\
= &
\{ (s, a') \in \mathcal{H}(W) \times h_V^\#(W) \mid
a' = a \text{ and } \varphi(s) = f \circ a'\}
\end{align*}
로 주어지는 $\mathcal{H}'_{\varphi'}$를 얻는다. 이는 위에서
$j^{-1}\mathcal{H}_\varphi$를 기술한 규칙과 정확히 같다.
\end{proof}

\begin{remark}
\label{sites-remark-localize-presheaves}
국소화와 준층. $\mathcal{C}$를 범주라 하고 $U$를 $\mathcal{C}$의 대상이라
하자. 엄밀히 말하면 함자 $j_U^{-1}$, $j_{U*}$ 및 $j_{U!}$는 준층에 대하여
정의되지 않았다. 그러나 물론 준층을 $\mathcal{C}$ 위의 카오스 위상에 대한
층으로 생각할 수 있다(Example \ref{sites-example-indiscrete} 참조). 따라서 함자
$$
j_U^{-1} :
\textit{PSh}(\mathcal{C})
\longrightarrow
\textit{PSh}(\mathcal{C}/U)
$$
와 $j_U^{-1}$의 각각 오른쪽, 왼쪽 수반인 함자
$$
j_{U*}, j_{U!} :
\textit{PSh}(\mathcal{C}/U)
\longrightarrow
\textit{PSh}(\mathcal{C})
$$
를 얻는다. Lemma \ref{sites-lemma-describe-j-shriek}에 의해 $j_{U!}\mathcal{G}$는
준층
$$
V \longmapsto
\coprod\nolimits_{\varphi \in \Mor_\mathcal{C}(V, U)}
\mathcal{G}(V \xrightarrow{\varphi} U)
$$
이다. 또한 함자 $j_{U!}$는 섬유곱 및 등화자와 가환한다.
\end{remark}

\begin{remark}
\label{sites-remark-localization-cartesian-cocontinuous}
$\mathcal{C}$를 사이트라 하고 $U \to V$를 $\mathcal{C}$의 사상이라 하자.
쌍대연속 함자 $\mathcal{C}/U \to \mathcal{C}$ 및
$j : \mathcal{C}/U \to \mathcal{C}/V$ (Lemma \ref{sites-lemma-relocalize})는
Remark \ref{sites-remark-cartesian-cocontinuous}의 성질 $P$를 만족한다. 예를 들어
대상 $(X/U)$, $(W/V)$, 사상 $g : j(X/U) \to (W/V)$ 및 덮개
$\{f_i  : (W_i/V) \to (W/V)\}$가 주어졌다면,
$(X \times_W W_i/U)$는 $(X/U) \times_{g, (W/V), f_i} (W_i/V)$의 한 표현이고
사상족 $\{(X \times_W W_i/U) \to (X/U)\}$는 $\mathcal{C}/U$의 덮개이다.
\end{remark}

\section{층 붙이기}
\label{sites-section-glueing-sheaves}

\noindent
이 절은 Sheaves, Section \ref{sheaves-section-glueing-sheaves}에 대응한다.

\begin{lemma}
\label{sites-lemma-glue-maps}
\begin{slogan}
층의 사상들은 붙일 수 있다.
\end{slogan}
$\mathcal{C}$를 사이트라 하고 $\{U_i \to U\}$를 $\mathcal{C}$의 덮개라 하자.
$\mathcal{F}$, $\mathcal{G}$를 $\mathcal{C}$ 위의 층이라 하자. 층의 사상들의
모음
$$
\varphi_i :
\mathcal{F}|_{\mathcal{C}/U_i}
\longrightarrow
\mathcal{G}|_{\mathcal{C}/U_i}
$$
가 주어지고, 모든 $i, j \in I$에 대하여 사상 $\varphi_i, \varphi_j$가 같은
사상
$\varphi_{ij} : \mathcal{F}|_{\mathcal{C}/U_i \times_U U_j} \to
\mathcal{G}|_{\mathcal{C}/U_i \times_U U_j}$로 제한된다고 하자. 그러면 각
$\mathcal{C}/U_i$에 대한 제한이 $\varphi_i$와 일치하는 유일한 층의 사상
$$
\varphi :
\mathcal{F}|_{\mathcal{C}/U}
\longrightarrow
\mathcal{G}|_{\mathcal{C}/U}
$$
가 존재한다.
\end{lemma}

\begin{proof}
보조정리에서 사용한 제한은 Lemma \ref{sites-lemma-relocalize}의 제한이다.
$V/U$를 $\mathcal{C}/U$의 대상이라 하자. $V_i = U_i \times_U V$라 놓고
$\mathcal{V} = \{V_i \to V\}$로 나타내자. 그러면
$(U_i \times_U U_j) \times_U V = V_i \times_V V_j$임을 주목하라. 또한
$\mathcal{F}|_{\mathcal{C}/U_i}(V_i/U_i) = \mathcal{F}(V_i)$이고
$\mathcal{F}|_{\mathcal{C}/U_i \times_U U_j}(V_i \times_V V_j/U_i \times_U U_j)
= \mathcal{F}(V_i \times_V V_j)$이며, $\mathcal{G}$에 대해서도 마찬가지이다.
따라서 $V$ 위의 절단들에 대한 $\varphi$를 다음 도표의 점선 화살표로 정의할
수 있다.
$$
\xymatrix{
\mathcal{F}(V) \ar@{=}[r] &
H^0(\mathcal{V}, \mathcal{F}) \ar@{..>}[d] \ar[r] &
\prod \mathcal{F}(V_i)
\ar[d]_{\prod \varphi_i}
\ar@<1ex>[r] \ar@<-1ex>[r] &
\prod \mathcal{F}(V_i \times_V V_j) \ar[d]_{\prod \varphi_{ij}} \\
\mathcal{G}(V) \ar@{=}[r] &
H^0(\mathcal{V}, \mathcal{G}) \ar[r] &
\prod \mathcal{G}(V_i)
\ar@<1ex>[r] \ar@<-1ex>[r] &
\prod \mathcal{G}(V_i \times_V V_j)
}
$$
등호들은 층 조건에서 나온다. $H^0(\mathcal{V}, -)$의 표기는 Section
\ref{sites-section-sheafification}을 보라. 이 사상들이 제한사상들과 양립한다는
확인은 생략한다.
\end{proof}

\noindent
앞의 보조정리에 따르면 사이트 $\mathcal{C}$ 위의 두 층 $\mathcal{F}$,
$\mathcal{G}$가 주어졌을 때 규칙
$$
U \longmapsto
\Mor_{\Sh(\mathcal{C}/U)}(
\mathcal{F}|_{\mathcal{C}/U}, \mathcal{G}|_{\mathcal{C}/U})
$$
은 층 $\SheafHom(\mathcal{F},\mathcal{G})$를 정의한다. 이는 일종의
{\it 내부 Hom 층}이다. 집합의 층을 다룰 때에는 거의 사용하지 않으며,
보통 가군의 층을 다룰 때 사용한다. Modules on Sites, Section
\ref{sites-modules-section-internal-hom}을 보라.

\begin{lemma}
\label{sites-lemma-internal-hom-sheaf}
\begin{slogan}
사이트 위의 층들의 범주는 데카르트 닫힌 범주이다.
\end{slogan}
$\mathcal{C}$를 사이트라 하자. $\mathcal{F}$, $\mathcal{G}$ 및
$\mathcal{H}$를 $\mathcal{C}$ 위의 층이라 하자. 세 성분 모두에 대하여
함자적인 표준 전단사
$$
\Mor_{\Sh(\mathcal{C})}(\mathcal{F}\times\mathcal{G},\mathcal{H}) =
\Mor_{\Sh(\mathcal{C})}(\mathcal{F},\SheafHom(\mathcal{G},\mathcal{H}))
$$
가 존재한다.
\end{lemma}

\begin{proof}
보조정리는 함자 $-\times\mathcal{G}$와 $\SheafHom(\mathcal{G},-)$가 서로
수반임을 말한다. 이를 보이기 위해 단위원과 여단위원의 개념을 사용한다.
Categories, Section \ref{categories-section-adjoint}를 보라. 단위원
$$
\eta_\mathcal{F} :
\mathcal{F}
\longrightarrow
\SheafHom(\mathcal{G},\mathcal{F}\times\mathcal{G})
$$
은 $s \in \mathcal{F}(U)$를 다음 사상으로 보낸다.
$\mathcal{G}|_{\mathcal{C}/U} \to
\mathcal{F}|_{\mathcal{C}/U}\times\mathcal{G}|_{\mathcal{C}/U}$
이 사상은 $V/U$ 위에서
$$
\mathcal{G}(V) \longrightarrow \mathcal{F}(V)\times \mathcal{G}(V), \quad
t \longmapsto (s|_{V},t).
$$
로 주어진다. 여단위원
$$
\epsilon_{\mathcal{H}} :
\SheafHom(\mathcal{G}, \mathcal{H}) \times \mathcal{G}
\longrightarrow
\mathcal{H}
$$
은 평가사상이다. 이는 규칙
$$
\Mor_{\Sh(\mathcal{C}/U)}(
\mathcal{G}|_{\mathcal{C}/U}, \mathcal{H}|_{\mathcal{C}/U})
\times \mathcal{G}(U)
\longrightarrow
\mathcal{H}(U),\quad
(\varphi, s) \longmapsto \varphi(s).
$$
로 주어진다. 그러면 각 $\varphi : \mathcal{F} \times \mathcal{G} \to \mathcal{H}$에
대응하는 사상 $\mathcal{F} \to \SheafHom(\mathcal{G},\mathcal{H})$는 각 절단
$s \in \mathcal{F}(U)$를 $\mathcal{C}/U$ 위의 층의 사상으로 보내서 얻는다.
그 사상은 $V/U$ 위의 절단들에 대하여
$$
\mathcal{G}(V) \longrightarrow \mathcal{H}(V),\quad
t \longmapsto \varphi(s|_V, t).
$$
로 주어진다. 역으로 각
$\psi : \mathcal{F} \to \SheafHom(\mathcal{G}, \mathcal{H})$에 대응하는 사상
$\mathcal{F} \times \mathcal{G} \to \mathcal{H}$는 대상 $U$ 위의 절단들에
대하여
$$
\mathcal{F}(U) \times \mathcal{G}(U) \longrightarrow \mathcal{H}(U),\quad
(s, t) \longmapsto \psi(s)(t)
$$
로 주어진다. 이 구성들이 서로 역임을 보이는 증명의 세부사항은 생략한다.
\end{proof}

\begin{lemma}
\label{sites-lemma-hom-sheaf-hU}
$\mathcal{C}$를 사이트라 하고 $U \in \Ob(\mathcal{C})$라 하자.
$\Sh(\mathcal{C})$의 $\mathcal{F}$에 대하여
$\SheafHom(h_U^\#, \mathcal{F}) = j_*(\mathcal{F}|_{\mathcal{C}/U})$이다.
\end{lemma}

\begin{proof}
층의 동형사상을 직접 구성하여 이를 보일 수 있다. 대신 다음과 같이
논증하자. $\mathcal{G}$를 $\mathcal{C}$ 위의 층이라 하자. 그러면
\begin{align*}
\Mor(\mathcal{G}, j_*(\mathcal{F}|_{\mathcal{C}/U}))
& =
\Mor(\mathcal{G}|_{\mathcal{C}/U}, \mathcal{F}|_{\mathcal{C}/U}) \\
& =
\Mor(j_!(\mathcal{G}|_{\mathcal{C}/U}), \mathcal{F}) \\
& =
\Mor(\mathcal{G} \times h_U^\#, \mathcal{F}) \\
& =
\Mor(\mathcal{G}, \SheafHom(h_U^\#, \mathcal{F}))
\end{align*}
이고, 요네다 보조정리에 의해 결론을 얻는다. 여기서는 Lemmas
\ref{sites-lemma-internal-hom-sheaf}와
\ref{sites-lemma-compute-j-shriek-restrict}를 사용했다.
\end{proof}

\noindent
$\mathcal{C}$를 사이트라 하고 $\{U_i \to U\}_{i \in I}$를 $\mathcal{C}$의
덮개라 하자. 각 $i \in I$에 대하여 $\mathcal{F}_i$를 $\mathcal{C}/U_i$ 위의
집합의 층이라 하자. 각 쌍 $i, j \in I$에 대하여
$$
\varphi_{ij} :
\mathcal{F}_i|_{\mathcal{C}/U_i \times_U U_j}
\longrightarrow
\mathcal{F}_j|_{\mathcal{C}/U_i \times_U U_j}
$$
를 집합의 층의 동형사상이라 하자. 또한 모든 세 지표 $i, j, k \in I$에
대하여 다음 도표가 가환한다고 가정하자.
$$
\xymatrix{
\mathcal{F}_i|_{\mathcal{C}/U_i \times_U U_j \times_U U_k}
\ar[rr]_{\varphi_{ik}}
\ar[rd]_{\varphi_{ij}} & &
\mathcal{F}_k|_{\mathcal{C}/U_i \times_U U_j \times_U U_k} \\
&
\mathcal{F}_j|_{\mathcal{C}/U_i \times_U U_j \times_U U_k}
\ar[ru]_{\varphi_{jk}}
}
$$
% P01-ERRATA: P01-E1331 -- the frozen source uses the singular article with plural "glueing data".
이러한 자료의 모음 $(\mathcal{F}_i, \varphi_{ij})$를 덮개
$\{U_i \to U\}_{i \in I}$에 대한 {\it 집합의 층의 붙이기 자료}라 한다.

\begin{lemma}
\label{sites-lemma-glue-sheaves}
$\mathcal{C}$를 사이트라 하고 $\{U_i \to U\}_{i \in I}$를 $\mathcal{C}$의
덮개라 하자. 덮개 $\{U_i \to U\}_{i \in I}$에 대한 집합의 층의 임의의 붙이기 자료
$(\mathcal{F}_i, \varphi_{ij})$가 주어지면, $\mathcal{C}/U$ 위의 집합의 층
$\mathcal{F}$와 동형사상
$$
\varphi_i : \mathcal{F}|_{\mathcal{C}/U_i} \to \mathcal{F}_i
$$
가 존재하여 도표
$$
\xymatrix{
\mathcal{F}|_{\mathcal{C}/U_i \times_U U_j}
\ar[d]_{\text{id}} \ar[r]_{\varphi_i} &
\mathcal{F}_i|_{\mathcal{C}/U_i \times_U U_j} \ar[d]^{\varphi_{ij}} \\
\mathcal{F}|_{\mathcal{C}/U_i \times_U U_j} \ar[r]^{\varphi_j} &
\mathcal{F}_j|_{\mathcal{C}/U_i \times_U U_j}
}
$$
가 가환한다.
\end{lemma}

\begin{proof}
$\mathcal{C}/U$ 위의 층 $\mathcal{F}$를 구성하는 방법을 기술하자.
$a : V \to U$를 $\mathcal{C}/U$의 대상이라 하자. 그러면
$$
\mathcal{F}(V/U) = \{
(s_i)_{i \in I} \in \prod_{i \in I} \mathcal{F}_i(U_i \times_U V/U_i)
\mid
\varphi_{ij}(s_i|_{U_i \times_U U_j \times_U V})
=
s_j|_{U_i \times_U U_j \times_U V}
\}
$$
이다. 제한사상의 구성은 생략한다. 이것이 층이라는 확인도 생략한다.
동형사상 $\varphi_i$의 구성과 보조정리의 도표들이 가환한다는 증명도
생략한다.
\end{proof}

\noindent
$\mathcal{C}$를 사이트라 하고 $\{U_i \to U\}_{i \in I}$를 $\mathcal{C}$의
덮개라 하자. $\mathcal{F}$를 $\mathcal{C}/U$ 위의 층이라 하자.
$\mathcal{F}$에는 제한 $\mathcal{F}|_{\mathcal{C}/U_i}$와 표준 동형사상
$$
\left(\mathcal{F}|_{\mathcal{C}/U_i}\right)|_{\mathcal{C}/U_i \times_U U_j}
=
\left(\mathcal{F}|_{\mathcal{C}/U_j}\right)|_{\mathcal{C}/U_i \times_U U_j}
$$
로 주어지는 {\it 표준 붙이기 자료}가 결부된다. 이 동형사상은 함자들의 합성
$\mathcal{C}/U_i \times_U U_j \to \mathcal{C}/U_i \to \mathcal{C}/U$와
$\mathcal{C}/U_i \times_U U_j \to \mathcal{C}/U_j \to \mathcal{C}/U$가
같다는 사실에서 나온다.

\begin{lemma}
\label{sites-lemma-mapping-property-glue}
$\mathcal{C}$를 사이트라 하고 $\{U_i \to U\}_{i \in I}$를 $\mathcal{C}$의
덮개라 하자. $\mathcal{C}/U$ 위의 $\mathcal{F}$에 그 표준 붙이기 자료를
대응시키는 함자를 통하여 범주 $\Sh(\mathcal{C}/U)$는 붙이기 자료들의
범주와 동치이다.
\end{lemma}

\begin{proof}
Lemma \ref{sites-lemma-glue-maps}에서 이 함자가 충실충만임을 보았고, Lemma
\ref{sites-lemma-glue-sheaves}에서 그것이 본질적으로 전사임을 증명했다(준역함자를
명시적으로 구성하였다).
\end{proof}

\noindent
$\mathcal{C}$를 사이트라 하자. 이제 전체 사이트 $\mathcal{C}$ 위에서 층을
붙이는 한 형태를 논의하겠다. $\mathcal{C}$의 각 대상 $U$에 대하여
$\mathcal{F}_U$를 $\mathcal{C}/U$ 위의 층이라 하자. $\mathcal{C}$의 각 사상
$f : V \rightarrow U$에는 $(a : W \to V) \mapsto (f \circ a : W \to U)$로
주어지는 함자 $j_f : \mathcal{C}/V \rightarrow \mathcal{C}/U$가 결부됨을
상기하자. 이러한 각 $f$에 대하여
$$
c_f : j_f^{-1} \mathcal{F}_U \to \mathcal{F}_V
$$
를 층의 동형사상이라 하자. $\mathcal{C}$의 임의의 두 화살표
$f : V \to U$와 $g : W \to V$가 주어졌을 때 합성
$c_g \circ j_{g}^{-1} c_f$가 $c_{f \circ g}$와 같다고 가정한다. 이러한 자료의
모음 $(\mathcal{F}_{U}, c_f)$를 $\mathcal{C}$ 위의 {\it 집합의 층의 절대
붙이기 자료}라 한다. 절대 붙이기 자료의 사상
$(\mathcal{F}_{U}, c_f) \to (\mathcal{G}_{U}, c'_f)$은 층의 사상들의 모음
$(\varphi_U)$로 주어지며, 여기서
$\varphi_U : \mathcal{F}_U \to \mathcal{G}_U$이고 모든 $\mathcal{C}$의 사상
$f : V \to U$에 대하여 도표
$$
\xymatrix{
j_f^{-1} \mathcal{F}_U \ar[r]_{c_f} \ar[d]_{j_f^{-1} \varphi_U} &
\mathcal{F}_V \ar[d]^{\varphi_V} \\
j_f^{-1} \mathcal{G}_U \ar[r]^{c'_f} &
\mathcal{G}_V
}
$$
가 가환한다.

\medskip\noindent
$\mathcal{C}$ 위의 임의의 층 $\mathcal{F}$에는 그 {\it 표준 절대 붙이기 자료}
$(\mathcal{F}|_{\mathcal{C}/U},c_f)$가 결부된다. 여기서 $f : V \to U$에 대한
표준 동형사상
$c_f : j_f^{-1} \mathcal{F}|_{\mathcal{C}/U} \to
\mathcal{F}|_{\mathcal{C}/V}$는 Lemma \ref{sites-lemma-relocalize}에서처럼 관계
$j_V = j_U \circ j_f$에서 나온다. $\mathcal{C}$의 층들의 임의의 사상
$\varphi : \mathcal{F} \to \mathcal{G}$는 표준 절대 붙이기 자료의 사상
$(\varphi|_{\mathcal{C}/U})$를 유도한다.

\begin{lemma}
\label{sites-lemma-glue-sheaves-absolute}
$\mathcal{C}$를 사이트라 하자. $\mathcal{C}$ 위의 $\mathcal{F}$에 그 표준 절대
붙이기 자료를 대응시키는 함자를 통하여 범주 $\Sh(\mathcal{C})$는 절대
붙이기 자료들의 범주와 동치이다.
\end{lemma}

\begin{proof}
% P01-ERRATA: P01-E1332 -- the frozen source again combines a singular article with plural "data".
절대 붙이기 자료 $(\mathcal{F}_U, c_f)$가 주어졌다고 하자.
$\mathcal{F}(U) = \mathcal{F}_U(U)$로 놓아 $\mathcal{C}$ 위의 층
$\mathcal{F}$를 구성한다.
% P01-ERRATA: P01-E0096 -- the frozen relative clause omits "is" before "given".
$f : V \to U$를 따른 제한은 가환도표
$$
\xymatrix{
\mathcal{F}_U(U) \ar[r] \ar@{=}[d] &
\mathcal{F}_U(V) \ar[r]^{c_f} &
\mathcal{F}_V(V) \ar@{=}[d] \\
\mathcal{F}(U) \ar[rr] &  &
\mathcal{F}(V)
}
$$
로 주어진다. 양립성 조건 $c_g \circ j_{g}^{-1} c_f = c_{f \circ g}$는
$\mathcal{F}$가 준층임을 보장하고, 사상
$c_f : \mathcal{F}_U(V) \to \mathcal{F}(V)$가 각 $U$에 대하여 동형사상
$\mathcal{F}_U \to \mathcal{F}|_{\mathcal{C}/U}$를 정의함도 보장한다. 각
$\mathcal{F}_U$가 층이므로 이는 $\mathcal{F}$도 층임을 뜻한다. 방금 구성한
함자 $(\mathcal{F}_U, c_f) \mapsto \mathcal{F}$는 $\mathcal{C}$ 위의 층을 그
표준 붙이기 자료로 보내는 함자의 준역함자이다. 나머지 세부사항은 생략한다.
\end{proof}

\begin{remark}
\label{sites-remark-variant-glue-sheaves-absolute}
Lemma \ref{sites-lemma-glue-sheaves-absolute}에는 대수기하학에서 나타나는 다음
변형이 있다. 즉, $\mathcal{C}$가 모든 섬유곱을 갖는 사이트이고 각
$U \in \Ob(\mathcal{C})$에 대하여 다음 성질을 갖는 충만부분범주
$U_\tau \subset \mathcal{C}/U$가 주어졌다고 하자.
\begin{enumerate}
\item $U/U$는 $U_\tau$에 속한다.
\item $U_\tau$의 $V/U$와 $\mathcal{C}$의 덮개 $\{V_j \to V\}$에 대하여
$V_j/U$는 $U_\tau$에 속한다.
\item $\mathcal{C}$의 사상 $U' \to U$와 $U_\tau$의 $V/U$에 대하여
기저변환 $V' = V \times_U U'$는 $U'_\tau$에 속한다.
\end{enumerate}
이 상황에서 $\mathcal{C}$의 모든 $U$에 대하여 $U_\tau$는 사이트이고,
$\mathcal{C}$의 모든 사상 $f : U' \to U$에 대하여 기저변환 함자
$U_\tau \to U'_\tau$는 사이트의 사상 $f_\tau : U_\tau \to U'_\tau$를
정의한다. 이때 얻는 붙이기 명제는 다음과 같다. $\mathcal{C}$ 위의 층
$\mathcal{F}$는 다음 자료로 주어진다.
\begin{enumerate}
\item 각 $U \in \Ob(\mathcal{C})$에 대하여 $U_\tau$ 위의 층
$\mathcal{F}_U$,
\item $\mathcal{C}$의 각 $f : U' \to U$에 대하여 사상
$c_f : f_\tau^{-1}\mathcal{F}_U \to \mathcal{F}_{U'}$.
\end{enumerate}
이 자료에는 다음 조건을 부과한다.
\begin{enumerate}
\item[(a)] $\mathcal{C}$의 $f : U' \to U$와 $g : U'' \to U'$가 주어지면
합성 $c_g \circ g_\tau^{-1}c_f$는 $c_{f \circ g}$와 같다.
\item[(b)] $f : U' \to U$가 $U_\tau$에 속하면 $c_f$는 동형사상이다.
\end{enumerate}
이 결과가 필요해지면 여기에서 정확히 진술하고 증명하겠다. (모든 사상
$c_f$가 동형사상일 것을 요구하지 않으므로 이 결과는 위의 명제들과 조금
다름에 유의하라!)
\end{remark}

\section{국소화에 관하여 더 나아가}
\label{sites-section-more-localize}

\noindent
% P01-ERRATA: P01-E0097 -- the frozen source duplicates and misorders the prepositions in this coordination.
이 절에서는 사이트 또는 국소화하는 대상에 몇 가지 추가 가정을 부과하여
국소화에 관한 몇 가지 보조정리를 증명한다.

\begin{lemma}
\label{sites-lemma-describe-j-shriek-good-site}
$\mathcal{C}$를 사이트라 하고 $U \in \Ob(\mathcal{C})$라 하자.
$\mathcal{C}$ 위의 위상이 준표준이고(Definition
\ref{sites-definition-weaker-than-canonical} 참조), $\mathcal{G}$가
$\mathcal{C}/U$ 위의 층이면
$$
j_{U!}(\mathcal{G})(V)
=
\coprod\nolimits_{\varphi \in \Mor_\mathcal{C}(V, U)}
\mathcal{G}(V \xrightarrow{\varphi} U),
$$
이다. 달리 말해 Lemma \ref{sites-lemma-describe-j-shriek}에서는 층화가 필요하지
않다.
\end{lemma}

\begin{proof}
$\mathcal{V} = \{V_i \to V\}_{i \in I}$를 사이트 $\mathcal{C}$에서 $V$의
덮개라 하자. Lemma \ref{sites-lemma-describe-j-shriek}의 준층 $\mathcal{H}$에 대한
층 조건을 직접 확인하겠다.
% P01-ERRATA: P01-E0098 -- the frozen source joins two independent sentences with a comma.
$(s_i, \varphi_i)_{i \in I} \in \prod_i \mathcal{H}(V_i)$라 하자. 이는
$\varphi_i : V_i \to U$가 $\mathcal{C}$의 사상이고
$s_i \in \mathcal{G}(V_i \xrightarrow{\varphi_i} U)$임을 뜻한다.
% P01-ERRATA: P01-E1334 -- the frozen source reverses both typed compositions pr_k and varphi.
쌍 $(s_i, \varphi_i)$의 $V_i \times_V V_j$에 대한 제한은 쌍
$(s_i|_{V_i \times_V V_j/U}, \text{pr}_1 \circ \varphi_i)$이고, 마찬가지로
쌍 $(s_j, \varphi_j)$의 $V_i \times_V V_j$에 대한 제한은 쌍
$(s_j|_{V_i \times_V V_j/U}, \text{pr}_2 \circ \varphi_j)$이다. 따라서 족
$(s_i, \varphi_i)$가 $\check{H}^0(\mathcal{V}, \mathcal{H})$에 속하면
$\text{pr}_1 \circ \varphi_i = \text{pr}_2 \circ \varphi_j$임을 알 수 있다.
$\mathcal{C}$ 위의 위상이 표준 위상보다 약하다는 조건에 의해, 각
$\varphi_i$가 $V_i \to V$와 $\varphi$의 합성이 되게 하는 유일한 사상
$\varphi : V \to U$가 존재한다. 이제 $\mathcal{G}$에 대한 층 조건에 의해
절단 $s_i$들은 유일한 절단 $s \in \mathcal{G}(V \xrightarrow{\varphi} U)$로
붙는다. 따라서 요구한 대로 $(s, \varphi) \in \mathcal{H}(V)$이다.
\end{proof}

\begin{lemma}
\label{sites-lemma-localize-given-products}
$\mathcal{C}$를 사이트라 하고 $U \in \Ob(\mathcal{C})$라 하자.
$\mathcal{C}$가 두 대상의 곱을 갖는다고 가정하자. 그러면
\begin{enumerate}
\item 함자 $j_U$는 연속인 오른쪽 수반, 즉 함자
$v(X) = X \times U / U$를 갖는다.
\item 함자 $v$는 사이트의 사상 $\mathcal{C}/U \to \mathcal{C}$를 정의하며,
이에 결부된 토포스의 사상은
$j_U : \Sh(\mathcal{C}/U) \to \Sh(\mathcal{C})$와 같다.
\item $j_{U*}\mathcal{F}(X) = \mathcal{F}(X \times U/U)$이다.
\end{enumerate}
\end{lemma}

\begin{proof}
함자 $v$가 $j_U$의 오른쪽 수반이라는 것은 $Y/U$와 $X$가 주어졌을 때
$$
\Mor_\mathcal{C}(Y, X)
=
\Mor_{\mathcal{C}/U}(Y/U, X \times U/U)
$$
라는 뜻이며, 이는 명백하다. $v$가 연속임을 확인하기 위해
$\{X_i \to X\}$를 $\mathcal{C}$의 덮개라 하자. 사이트의 세 번째 공리
(Definition \ref{sites-definition-site})에 의해
$$
\{X_i \times_X (X \times U) \to X \times_X (X \times U)\}
=
\{X_i \times U \to X \times U\}
$$
도 $\mathcal{C}$의 덮개이다. 따라서 $v$는 연속이다. 보조정리의 다른
명제들은 Lemmas \ref{sites-lemma-have-functor-other-way}와
\ref{sites-lemma-have-functor-other-way-morphism}에서 따른다.
\end{proof}

\begin{lemma}
\label{sites-lemma-relocalize-given-fibre-products}
$\mathcal{C}$를 사이트라 하고 $U \to V$를 $\mathcal{C}$의 사상이라 하자.
$\mathcal{C}$가 섬유곱을 갖는다고 가정하자. $j$를 Lemma
\ref{sites-lemma-relocalize}에서처럼 두자. 그러면
\begin{enumerate}
\item 함자 $j : \mathcal{C}/U \to \mathcal{C}/V$는 연속인 오른쪽 수반,
즉 함자 $v : (X/V) \mapsto (X \times_V U/U)$를 갖는다.
\item 함자 $v$는 사이트의 사상 $\mathcal{C}/U \to \mathcal{C}/V$를
정의하며, 이에 결부된 토포스의 사상은 $j$와 같다.
\item $j_*\mathcal{F}(X/V) = \mathcal{F}(X \times_V U/U)$이다.
\end{enumerate}
\end{lemma}

\begin{proof}
Lemma \ref{sites-lemma-localize-given-products}를 적용하면 결과를 얻는다. 실제로
Lemma \ref{sites-lemma-relocalize}에 의해 $j$를 국소화 함자로 볼 수 있다.
\end{proof}

\noindent
% P01-ERRATA: P01-E0099 -- the frozen compound subject takes a singular verb.
열린 몰입의 기본 성질 하나는 전진상의 제한과 공집합에 의한 확장의 제한이
원래의 층을 되돌려 준다는 것이다. 위의 $j_U$에 결부된 함자들에 대해서는
항상 참이지 않다. $U$가 ``종대상의 부분대상''일 때에는 참이다.

\begin{lemma}
\label{sites-lemma-restrict-back}
$\mathcal{C}$를 사이트라 하고 $U \in \Ob(\mathcal{C})$라 하자.
$\mathcal{C}$의 모든 $X$에서 $U$로 가는 사상이 많아야 하나라고 가정하자.
$\mathcal{F}$를 $\mathcal{C}/U$ 위의 층이라 하자. 표준 사상
$\mathcal{F} \to j_U^{-1}j_{U!}\mathcal{F}$ 및
$j_U^{-1}j_{U*}\mathcal{F} \to \mathcal{F}$는 동형사상이다.
\end{lemma}

\begin{proof}
% P01-ERRATA: P01-E0100 -- the frozen source says "fully faithfulness" instead of "full faithfulness".
$U$에 대한 가정은 국소화 함자 $\mathcal{C}/U \to \mathcal{C}$의
충만충실성과 동치이므로, 이는 Lemma \ref{sites-lemma-back-and-forth}의 특수한
경우이다.
\end{proof}

\begin{lemma}
\label{sites-lemma-localize-cartesian-square}
$\mathcal{C}$를 사이트라 하고
$$
\xymatrix{
U' \ar[d] \ar[r] & U \ar[d] \\
V' \ar[r] & V
}
$$
를 $\mathcal{C}$의 가환도표라 하자. Lemma \ref{sites-lemma-relocalize}의 사상들은
연속이고 쌍대연속인 함자들의 가환도표 및 토포스의 가환도표
$$
\vcenter{
\xymatrix{
\mathcal{C}/U' \ar[d]_{j_{U'/V'}} \ar[r]_{j_{U'/U}} &
\mathcal{C}/U \ar[d]^{j_{U/V}} \\
\mathcal{C}/V' \ar[r]^{j_{V'/V}} & \mathcal{C}/V
}
}
\quad\text{and}\quad
\vcenter{
\xymatrix{
\Sh(\mathcal{C}/U') \ar[d]_{j_{U'/V'}} \ar[r]_{j_{U'/U}} &
\Sh(\mathcal{C}/U) \ar[d]^{j_{U/V}} \\
\Sh(\mathcal{C}/V') \ar[r]^{j_{V'/V}} &
\Sh(\mathcal{C}/V)
}
}
$$
를 만든다. 더 나아가 $\mathcal{C}$의 처음 도표가 데카르트 도표이면
$j_{V'/V}^{-1} \circ j_{U/V, *} = j_{U'/V', *} \circ j_{U'/U}^{-1}$이다.
\end{lemma}

\begin{proof}
보조정리의 첫 번째 명제에서 왼쪽 정사각형의 가환성은 정의에서 즉시
따른다. Lemma \ref{sites-lemma-composition-cocontinuous}에 의해 이는 토포스의
도표가 가환함을 함의한다. 도표가 데카르트라고 가정하자. 수반함자의
유일성에 의해
$j_{V'/V}^{-1} \circ j_{U/V, *} = j_{U'/V', *} \circ j_{U'/U}^{-1}$임을
보이는 것은
$j_{U/V}^{-1} \circ j_{V'/V!} = j_{U'/U!} \circ j_{U'/V'}^{-1}$임을
보이는 것과 동치이다. Lemma \ref{sites-lemma-essential-image-j-shriek}의 식별을
통하여 토포스의 도표를
$$
\xymatrix{
\Sh(\mathcal{C})/h_{U'}^\# \ar[d] \ar[r] &
\Sh(\mathcal{C})/h_U^\# \ar[d] \\
\Sh(\mathcal{C})/h_{V'}^\# \ar[r] &
\Sh(\mathcal{C})/h_V^\#
}
$$
로 생각할 수 있으며, Lemma \ref{sites-lemma-relocalize-explicit}에 의해 함자
$j^{-1}$와 $j_!$를 해석하는 방법을 알고 있다. 따라서
$\mathcal{F} \to h_{V'}^\#$가 주어졌을 때, $h_U^\#$로 가는 사상을 갖는
층으로서
$$
\mathcal{F} \times_{h_{V'}^\#} h_{U'}^\# =
\mathcal{F} \times_{h_V^\#} h_U^\#
$$
임을 보여야 한다. 이는 $h_{U'} = h_{V'} \times_{h_V} h_U$이고, 층화가
완전하므로
$$
h_{U'}^\# = h_{V'}^\# \times_{h_V^\#} h_U^\#
$$
도 성립하기 때문에 참이다.
\end{proof}

\section{국소화와 사상}
\label{sites-section-localize-morphisms}

\noindent
% P01-ERRATA: P01-E0101 -- the frozen count noun "relation" lacks its determiner.
다음 보조정리는 국소화와 사이트 및 토포스의 사상 사이의 관계를 이해하는 데
중요하다.

\begin{lemma}
\label{sites-lemma-localize-morphism}
$f : \mathcal{C} \to \mathcal{D}$를 연속함자
$u : \mathcal{D} \to \mathcal{C}$에 대응하는 사이트의 사상이라 하자.
$V \in \Ob(\mathcal{D})$라 하고 $U = u(V)$로 놓자. 그러면 함자
$u' : \mathcal{D}/V \to \mathcal{C}/U$,
$V'/V \mapsto u(V')/U$는 사이트의 사상
$f' : \mathcal{C}/U \to \mathcal{D}/V$를 결정한다. 사상 $f'$는 토포스의
가환도표
$$
\xymatrix{
\Sh(\mathcal{C}/U) \ar[r]_{j_U} \ar[d]_{f'} &
\Sh(\mathcal{C}) \ar[d]^f \\
\Sh(\mathcal{D}/V) \ar[r]^{j_V} &
\Sh(\mathcal{D}).
}
$$
에 들어맞는다. Lemma \ref{sites-lemma-essential-image-j-shriek}의 식별
$\Sh(\mathcal{C}/U) = \Sh(\mathcal{C})/h_U^\#$ 및
$\Sh(\mathcal{D}/V) = \Sh(\mathcal{D})/h_V^\#$를 사용하면 함자
$(f')^{-1}$는 규칙
$$
(f')^{-1}(\mathcal{H} \xrightarrow{\varphi} h_V^\#)
=
(f^{-1}\mathcal{H} \xrightarrow{f^{-1}\varphi} h_U^\#).
$$
로 기술된다. 마지막으로 $f'_*j_U^{-1} = j_V^{-1}f_*$이다.
\end{lemma}

\begin{proof}
$u'$가 연속임은 명백하다. 따라서 함자
$f'_* = (u')^s = (u')^p$ (Sections \ref{sites-section-functoriality-PSh}와
\ref{sites-section-continuous-functors} 참조) 및 수반함자
$(f')^{-1} = (u')_s = ((u')_p\ )^\#$를 얻는다. 명제
$f'_*j_U^{-1} = j_V^{-1}f_*$는
$$
(j_V^{-1}f_*\mathcal{F})(V'/V)
= f_*\mathcal{F}(V') = \mathcal{F}(u(V'))
= (j_U^{-1}\mathcal{F})(u(V')/U)
= (f'_*j_U^{-1}\mathcal{F})(V'/V)
$$
에서 따르며, 이는 준층에 대해서도 성립한다. 선험적으로 명백하지 않은 것은
$(f')^{-1}$가 완전하다는 것, 도표가 가환한다는 것, 그리고 위의
$(f')^{-1}$에 대한 기술이 성립한다는 것이다.

\medskip\noindent
$\mathcal{H}$를 $\mathcal{D}/V$ 위의 층이라 하자.
$j_{U!}(f')^{-1}\mathcal{H}$를 계산하자. 다음을 얻는다.
\begin{eqnarray*}
j_{U!}(f')^{-1}\mathcal{H}
& =
((j_U)_p(u'_p\mathcal{H})^\#)^\# \\
& =
((j_U)_pu'_p\mathcal{H})^\# \\
& =
(u_p(j_V)_p\mathcal{H})^\# \\
& =
f^{-1}j_{V!}\mathcal{H}
\end{eqnarray*}
% P01-ERRATA: P01-E0102 -- all four frozen equality explanations are sentence fragments.
첫 번째 등식은 정의를 풀어 쓰면 얻어진다. 두 번째 등식은 Lemma
\ref{sites-lemma-technical-up}에서 따른다. 세 번째 등식은
$u \circ j_V = j_U \circ u'$이기 때문에 성립한다. 네 번째 등식은 다시
Lemma \ref{sites-lemma-technical-up}에서 따른다. 위의 모든 등식은 함자들의
동형사상이므로 다음 범주 및 함자들의 도표가 가환한다는 뜻으로 해석할 수 있다.
$$
\xymatrix{
\Sh(\mathcal{C}/U) \ar[r]_{j_{U!}} &
\Sh(\mathcal{C})/h_U^\# \ar[r] &
\Sh(\mathcal{C}) \\
\Sh(\mathcal{D}/V) \ar[r]^{j_{V!}} \ar[u]^{(f')^{-1}} &
\Sh(\mathcal{D})/h_V^\# \ar[r] \ar[u]^{f^{-1}} &
\Sh(\mathcal{D}) \ar[u]^{f^{-1}}
}
$$
가운데 화살표는 $f^{-1}h_V^\# = (h_{u(V)})^\# = h_U^\#$이므로 의미가 있다.
Lemma \ref{sites-lemma-pullback-representable-sheaf}를 보라. 특히 이로써 보조정리의
명제에 제시된 $(f')^{-1}$에 대한 기술이 증명된다. Lemma
\ref{sites-lemma-essential-image-j-shriek}에 의해 왼쪽 수평 화살표들이 동치이고,
가정에 의해 $f^{-1}$가 완전하므로 $(f')^{-1} = u'_s$가 완전하다고 결론내린다.
실제로 이는 왼쪽 수반이므로 이미 오른쪽 완전하다(Categories, Lemma
\ref{categories-lemma-adjoint-exact}). 따라서 종대상을 종대상으로 보내고
섬유곱과 가환한다는 것만 보이면 된다(Categories, Lemma
\ref{categories-lemma-characterize-left-exact}). 유도된 함자
$f^{-1} : \Sh(\mathcal{D})/h_V^\# \to \Sh(\mathcal{C})/h_U^\#$에 대해서는
둘 다 명백하다. 이로써 $f'$가 사이트의 사상임이 증명된다.

\medskip\noindent
아직 $(f')^{-1}j_V^{-1} = j_U^{-1}f^{-1}$임을 확인해야 한다. 이를 위해 위의
공식과 Lemma \ref{sites-lemma-compute-j-shriek-restrict}의 기술을 사용하자. 이들을
결합하면 $\mathcal{D}$ 위의 임의의 층 $\mathcal{G}$에 대하여
$$
j_{U!}(f')^{-1}j_V^{-1}\mathcal{G}
=
f^{-1}j_{V!}j_V^{-1}\mathcal{G}
=
f^{-1}(\mathcal{G} \times h_V^\#)
=
f^{-1}\mathcal{G} \times h_U^\#
=
j_{U!}j_U^{-1}f^{-1}\mathcal{G}.
$$
함자 $j_{U!}$는 동치
$\Sh(\mathcal{C}/U) \to \Sh(\mathcal{C})/h_U^\#$를 유도하므로 결론을
얻는다.
\end{proof}

\noindent
다음 보조정리는 위의 더 일반적인 Lemma \ref{sites-lemma-localize-morphism}의
특수한 경우이다.

\begin{lemma}
\label{sites-lemma-localize-morphism-strong}
$\mathcal{C}$, $\mathcal{D}$를 사이트라 하자.
$u : \mathcal{D} \to \mathcal{C}$를 함자라 하자.
$V \in \Ob(\mathcal{D})$라 하고 $U = u(V)$로 놓자. 다음을 가정하자.
\begin{enumerate}
\item $\mathcal{C}$와 $\mathcal{D}$는 모든 유한 극한을 갖는다.
\item $u$는 연속이다.
\item $u$는 유한 극한과 가환한다.
\end{enumerate}
다음 사이트 사상들의 가환도표가 존재한다.
$$
\xymatrix{
\mathcal{C}/U \ar[r]_{j_U} \ar[d]_{f'} & \mathcal{C} \ar[d]^f \\
\mathcal{D}/V \ar[r]^{j_V} & \mathcal{D}
}
$$
여기서 오른쪽 수직 화살표는 $u$에 대응하고, 왼쪽 수직 화살표는 함자
$u' : \mathcal{D}/V \to \mathcal{C}/U$,
$V'/V \mapsto u(V')/u(V)$에 대응하며, 수평 화살표들은 Lemma
% P01-ERRATA: P01-E0103 -- the frozen displayed functor values omit their slice structure maps.
\ref{sites-lemma-localize-given-products}에서처럼 함자
$\mathcal{C} \to \mathcal{C}/U$, $X \mapsto X \times U$ 및
$\mathcal{D} \to \mathcal{D}/V$, $Y \mapsto Y \times V$에 대응한다. 더 나아가
이에 결부된 토포스 사상들의 도표는 Lemma
\ref{sites-lemma-localize-morphism}의 도표와 같다. 특히
$f'_*j_U^{-1} = j_V^{-1}f_*$이다.
\end{lemma}

\begin{proof}
$u$는 Proposition \ref{sites-proposition-get-morphism}의 가정들을 만족하므로 그
명제에 의해 사이트의 사상 $f : \mathcal{C} \to \mathcal{D}$를 유도한다.
$u$가 제시된 함자 $u'$를 유도함은 명백하다. 이 함자도 Proposition
\ref{sites-proposition-get-morphism}의 가정들을 만족함은 명백하다. 따라서 사이트의
사상 $f' : \mathcal{C}/U \to \mathcal{D}/V$를 얻는다. 사이트 사상의 합성에
대한 우리의 정의(Definition \ref{sites-definition-composition-morphisms-sites} 참조)와
$$
u(Y \times V) = u(Y) \times u(V) = u(Y) \times U
$$
에 의해 도표는 가환한다. 실제로 이 등식은 보조정리의 도표에 반대되는 범주와
함자들의 도표가 가환함을 보인다.
\end{proof}

\noindent
이제 사이트를 국소화할 수 있고, 재국소화하는 방법도 알며, 아래쪽 사이트의
한 대상에서 사이트의 사상을 국소화할 수도 있다. 이들을 결합하면 다음과 같은
도표를 얻는다.

\begin{lemma}
\label{sites-lemma-relocalize-morphism}
$f : \mathcal{C} \to \mathcal{D}$를 연속함자
$u : \mathcal{D} \to \mathcal{C}$에 대응하는 사이트의 사상이라 하자.
$V \in \Ob(\mathcal{D})$, $U \in \Ob(\mathcal{C})$라 하고
$c : U \to u(V)$를 $\mathcal{C}$의 사상이라 하자. 토포스의 가환도표
$$
\xymatrix{
\Sh(\mathcal{C}/U) \ar[r]_{j_U} \ar[d]_{f_c} &
\Sh(\mathcal{C}) \ar[d]^f \\
\Sh(\mathcal{D}/V) \ar[r]^{j_V} &
\Sh(\mathcal{D}).
}
$$
가 존재한다. 여기서 $f' : \Sh(\mathcal{C}/u(V)) \to \Sh(\mathcal{D}/V)$는
Lemma \ref{sites-lemma-localize-morphism}에서와 같고,
$j_{U/u(V)} : \Sh(\mathcal{C}/U) \to \Sh(\mathcal{C}/u(V))$는 Lemma
\ref{sites-lemma-relocalize}에서와 같으며, $f_c = f' \circ j_{U/u(V)}$이다. Lemma
\ref{sites-lemma-essential-image-j-shriek}의 식별
$\Sh(\mathcal{C}/U) = \Sh(\mathcal{C})/h_U^\#$ 및
$\Sh(\mathcal{D}/V) = \Sh(\mathcal{D})/h_V^\#$를 사용하면 함자
$(f_c)^{-1}$는 규칙
$$
(f_c)^{-1}(\mathcal{H} \xrightarrow{\varphi} h_V^\#)
=
(f^{-1}\mathcal{H} \times_{f^{-1}\varphi, h_{u(V)}^\#, c} h_U^\#
\rightarrow h_U^\#).
$$
로 기술된다. 마지막으로 사상 $b : V' \to V$, $a : U' \to U$ 및
$c' : U' \to u(V')$가 주어지고
$$
\xymatrix{
U' \ar[r]_-{c'} \ar[d]_a & u(V') \ar[d]^{u(b)} \\
U \ar[r]^-c & u(V)
}
$$
가 가환하면, 도표
$$
\xymatrix{
\Sh(\mathcal{C}/U') \ar[r]_{j_{U'/U}} \ar[d]_{f_{c'}} &
\Sh(\mathcal{C}/U) \ar[d]^{f_c} \\
\Sh(\mathcal{D}/V') \ar[r]^{j_{V'/V}} &
\Sh(\mathcal{D}/V).
}
$$
가 가환한다.
\end{lemma}

\begin{proof}
이 보조정리는 사실상 자명하며, 이 이론의 현 단계에서 우리가 아는 것들을
모아 놓은 것에 가깝다. 예를 들어 첫 번째 정사각형의 가환성은 Diagram
(\ref{sites-equation-relocalize})의 가환성과 Lemma
\ref{sites-lemma-localize-morphism}의 도표의 가환성에서 따른다. $f_c^{-1}$에 대한
기술은 Lemma \ref{sites-lemma-relocalize-explicit}와 Lemma
\ref{sites-lemma-localize-morphism}를 결합하면 얻어진다. 마지막 정사각형의
가환성은 다음 등식에서 따른다.
% P01-ERRATA: P01-E0104 -- the frozen fibre-product subscript misplaces c' inside h_{u(V')}.
$$
f^{-1}\mathcal{H} \times_{h_{u(V)}^\#, c} h_U^\# \times_{h_U^\#} h_{U'}^\#
=
f^{-1}(\mathcal{H} \times_{h_V^\#} h_{V'}^\#)
\times_{h_{u(V'), c'}^\#} h_{U'}^\#
$$
이는 $f^{-1}h_V^\# = h_{u(V)}^\#$ 및
$f^{-1}h_{V'}^\# = h_{u(V')}^\#$를 사용하면 형식적으로 성립한다. Lemma
\ref{sites-lemma-pullback-representable-sheaf}를 보라.
\end{proof}

\noindent
다음 보조정리에서는 토포스의 사상이 쌍대연속 함자에서 나오는 경우에
국소화의 또 다른 함자성을 찾는다. 이는 Lemma
\ref{sites-lemma-localize-morphism}의 도표와 다른 종류의 도표이다. 특히 일반적으로
Lemma \ref{sites-lemma-localize-morphism}에서 본 등식
$f'_*j_U^{-1} = j_V^{-1}f_*$는 다음 보조정리의 상황에서는 성립하지 않는다.

\begin{lemma}
\label{sites-lemma-localize-cocontinuous}
$\mathcal{C}$, $\mathcal{D}$를 사이트라 하자.
$u : \mathcal{C} \to \mathcal{D}$를 쌍대연속 함자라 하자. $U$를
$\mathcal{C}$의 대상이라 하고 $V = u(U)$로 놓자. 가환도표
$$
\xymatrix{
\mathcal{C}/U \ar[r]_{j_U} \ar[d]_{u'} & \mathcal{C} \ar[d]^u \\
\mathcal{D}/V \ar[r]^-{j_V} & \mathcal{D}
}
$$
가 있으며, 여기서 왼쪽 수직 화살표는
% P01-ERRATA: P01-E0105 -- the frozen object assignment uses undefined V'.
$u' : \mathcal{C}/U \to \mathcal{D}/V$, $U'/U \mapsto V'/V$이다. 그러면
$u'$도 쌍대연속이고 토포스의 가환도표
$$
\xymatrix{
\Sh(\mathcal{C}/U) \ar[r]_{j_U} \ar[d]_{f'} &
\Sh(\mathcal{C}) \ar[d]^f \\
\Sh(\mathcal{D}/V) \ar[r]^-{j_V} &
\Sh(\mathcal{D})
}
$$
를 얻는다. 여기서 $f$ (각각 $f'$)는 $u$ (각각 $u'$)에 대응한다.
\end{lemma}

\begin{proof}
첫 번째 도표의 가환성은 명백하다. $u'$가 쌍대연속임을 보이면 이로부터 두
번째 도표의 가환성이 따른다.

\medskip\noindent
$U'/U$를 $\mathcal{C}/U$의 대상이라 하자.
$\{V_j/V \to u(U')/V\}_{j \in J}$를 $\mathcal{D}/V$에서 $u(U')/V$의
덮개라 하자. $u$가 쌍대연속이므로, 족 $\{u(U_i') \to u(U')\}$가
$\mathcal{D}$의 덮개 $\{V_j \to u(U')\}$의 세분이 되게 하는 덮개
$\{U_i' \to U'\}_{i \in I}$가 존재한다. 달리 말해 지표집합의 사상
$\alpha : I \to J$ 및 사상
% P01-ERRATA: P01-E0106 -- the frozen phrase says these D-morphisms are over U' rather than u(U').
$\phi_i : u(U_i') \to V_{\alpha(i)}$ over $U'$가 존재한다. 합성
$U'_i \to U' \to U$를 통하여 $U_i'$를 $U$ 위의 대상으로 생각하자. 그러면
$\{U'_i/U \to U'/U\}$는 $\mathcal{C}/U$의 덮개이고,
$\{u(U_i')/V \to u(U')/V\}$는 $\{V_j/V \to u(U')/V\}$의 세분이다(같은
$\alpha$와 같은 사상 $\phi_i$를 사용한다). 따라서
$u' : \mathcal{C}/U \to \mathcal{D}/V$는 쌍대연속이다.
\end{proof}

\begin{lemma}
\label{sites-lemma-localize-cocontinuous-downstairs}
$\mathcal{C}$, $\mathcal{D}$를 사이트라 하자.
$u : \mathcal{C} \to \mathcal{D}$를 쌍대연속 함자라 하자. $V$를
$\mathcal{D}$의 대상이라 하자. ${}^u_V\mathcal{I}$를 Section
\ref{sites-section-more-functoriality-PSh}에서 도입한 범주라 하자. 가환도표
$$
\vcenter{
\xymatrix{
\,_V^u\mathcal{I} \ar[r]_j \ar[d]_{u'} &
\mathcal{C} \ar[d]^u \\
\mathcal{D}/V \ar[r]^-{j_V} &
\mathcal{D}
}
}
\quad\text{where}\quad
\begin{matrix}
j : (U, \psi) \mapsto U \\
u' : (U, \psi) \mapsto (\psi : u(U) \to V)
\end{matrix}
$$
가 있다. ${}^u_V\mathcal{I}$의 사상족
$\{(U_i, \psi_i) \to (U, \psi)\}$가 덮개일 필요충분조건을
$\{U_i \to U\}$가 $\mathcal{C}$에서 덮개인 것으로 선언하자. 그러면
\begin{enumerate}
\item ${}^u_V\mathcal{I}$는 사이트이다.
\item $j$는 연속이고 쌍대연속이다.
\item $u'$는 쌍대연속이다.
\item 토포스의 가환도표
$$
\xymatrix{
\Sh({}^u_V\mathcal{I}) \ar[r]_j \ar[d]_{f'} &
\Sh(\mathcal{C}) \ar[d]^f \\
\Sh(\mathcal{D}/V) \ar[r]^-{j_V} &
\Sh(\mathcal{D})
}
$$
를 얻는다. 여기서 $f$ (각각 $f'$)는 $u$ (각각 $u'$)에 대응한다.
\item $f'_*j^{-1} = j_V^{-1}f_*$이다.
\end{enumerate}
\end{lemma}

\begin{proof}
(1), (2), (3), (4)는 정의와 함자 $j$가 섬유곱과 가환한다는 사실에서 바로
따른다. 세부사항은 생략한다. (5)를 보기 위해 $f_*$가
${}_su = {}_pu$로 주어짐을 상기하자. 따라서 $V'/V$에서
$j_V^{-1}f_*\mathcal{F}$의 값은 $V'$에서 ${}_pu\mathcal{F}$의 값이며, 이는
범주 ${}^u_{V'}\mathcal{I}$ 위의 $\mathcal{F}$의 값들의 극한이다. 명백히
범주의 동치
$$
{}^u_{V'}\mathcal{I} \to {}^{u'}_{V'/V}\mathcal{I}
$$
가 있다. $V'/V$에서 $f'_*j^{-1}\mathcal{F}$의 값은 범주
${}^{u'}_{V'/V}\mathcal{I}$ 위의 $j^{-1}\mathcal{F}$의 값들의 극한으로
주어지고, ${}^u_V\mathcal{I}$의 대상에서 $j^{-1}\mathcal{F}$의 값들은
$\mathcal{F}$의 값들일 뿐이므로($j$가 연속이고 쌍대연속인 데 대하여 Lemma
\ref{sites-lemma-when-shriek}를 적용한다), (5)가 참임을 알 수 있다.
\end{proof}

\noindent
다음 두 결과는 약간 다른 성격을 갖는다.

\begin{lemma}
\label{sites-lemma-special-square-cocontinuous}
사이트 $\mathcal{C}', \mathcal{C}, \mathcal{D}', \mathcal{D}$와 함자
$$
\xymatrix{
\mathcal{C}' \ar[r]_{v'} \ar[d]_{u'} &
\mathcal{C} \ar[d]^u \\
\mathcal{D}' \ar[r]^v &
\mathcal{D}
}
$$
가 주어졌다고 하자. 다음을 가정하자.
\begin{enumerate}
\item $u$, $u'$, $v$, $v'$는 쌍대연속이며 Lemma
\ref{sites-lemma-cocontinuous-morphism-topoi}에 의해 각각 토포스의 사상
$f$, $f'$, $g$, $g'$를 준다.
\item $v \circ u' = u \circ v'$이다.
\item $v$와 $v'$는 쌍대연속일 뿐 아니라 연속이다.
\item $\mathcal{D}'$의 임의의 대상 $V'$에 대하여 $v$가 주는 함자
${}^{u'}_{V'}\mathcal{I} \to {}^{\ \ \ u}_{v(V')}\mathcal{I}$는 공종이다.
\end{enumerate}
그러면 $f'_* \circ (g')^{-1} = g^{-1} \circ f_*$이고
$g'_! \circ (f')^{-1} = f^{-1} \circ g_!$이다.
\end{lemma}

\begin{proof}
범주 ${}^{u'}_{V'}\mathcal{I}$와 ${}^{\ \ \ u}_{v(V')}\mathcal{I}$는 Section
\ref{sites-section-more-functoriality-PSh}에서 정의되었다. 조건 (4)의 함자는
${}^{u'}_{V'}\mathcal{I}$의 대상 $\psi : u'(U') \to V'$를
${}^{\ \ \ u}_{v(V')}\mathcal{I}$의 대상
$v(\psi) : uv'(U') = vu'(U') \to v(V')$로 보낸다. $g^{-1}$는
$v^p$로 주어지고(Lemma \ref{sites-lemma-when-shriek}), $f_*$는
${}_su = {}_pu$로 주어짐을 상기하자. 따라서 $V'$에서
$g^{-1}f_*\mathcal{F}$의 값은 $v(V')$에서 ${}_pu\mathcal{F}$의 값, 즉 극한
$$
\lim_{u(U) \to v(V') \in \Ob({}^{\ \ \ u}_{v(V')}\mathcal{I}^{opp})}
\mathcal{F}(U)
$$
이다. 같은 논증으로 $V'$에서 $f'_*(g')^{-1}\mathcal{F}$의 값은 극한
$$
\lim_{u'(U') \to V' \in \Ob({}^{u'}_{V'}\mathcal{I}^{opp})} \mathcal{F}(v'(U'))
$$
으로 주어진다. 따라서 가정 (4)와 Categories, Lemma
\ref{categories-lemma-initial}에 의해 이들은 일치하고 보조정리의 첫 번째
등식이 증명된다. 두 번째 등식은 수반함자의 유일성에 의해 첫 번째 등식에서
따른다.
\end{proof}

\begin{lemma}
\label{sites-lemma-special-square-continuous}
사이트 $\mathcal{C}', \mathcal{C}, \mathcal{D}', \mathcal{D}$와 함자
$$
\xymatrix{
\mathcal{C}' \ar[r]_{v'} &
\mathcal{C} \\
\mathcal{D}' \ar[r]^v \ar[u]^{u'} &
\mathcal{D} \ar[u]_u
}
$$
가 주어졌다고 하자. Sections \ref{sites-section-morphism-sites}와
\ref{sites-section-cocontinuous-morphism-topoi}에서와 같이 표기하고 다음을 가정하자.
\begin{enumerate}
\item $u$와 $u'$는 연속이며 사이트의 사상 $f$와 $f'$를 준다.
\item $v$와 $v'$는 쌍대연속이며 토포스의 사상 $g$와 $g'$를 준다.
\item $u \circ v = v' \circ u'$이다.
\item $v$와 $v'$는 쌍대연속일 뿐 아니라 연속이다.
\end{enumerate}
그러면\footnote{이 일반성에서는 토포스의 사상으로서
$f \circ g'$가 $g \circ f'$와 같은지 알 수 없다(첫 번째 것에서 두 번째
것으로 가는 표준적인 $2$-화살표가 있지만, 이는 동형사상이 아닐 수 있다).}
$f'_* \circ (g')^{-1} = g^{-1} \circ f_*$이고
$g'_! \circ (f')^{-1} = f^{-1} \circ g_!$이다.
\end{lemma}

\begin{proof}
다음이 성립한다.
$$
f'_*(g')^{-1}\mathcal{F} = (u')^p((v')^p\mathcal{F})^\# =
(u')^p(v')^p\mathcal{F}
$$
% P01-ERRATA: P01-E0107 -- the frozen equality explanations after both displays are sentence fragments.
첫 번째 등식은 정의에서 따르고 두 번째 등식은 Lemma
\ref{sites-lemma-when-shriek}에서 따른다. 또한
$$
g^{-1}f_*\mathcal{F} = (v^pu^p\mathcal{F})^\# =
((u')^p(v')^p\mathcal{F})^\# = (u')^p(v')^p\mathcal{F}
$$
이다. 첫 번째 등식은 정의에서 따르고, 두 번째 등식은
$u \circ v = v' \circ u'$이기 때문에 성립하며, 세 번째 등식은
$(u')^p(v')^p\mathcal{F}$가 층임을 이미 보았기 때문에 성립한다. 이로써
$f'_* \circ (g')^{-1} = g^{-1} \circ f_*$가 증명되고, 등식
$g'_! \circ (f')^{-1} = f^{-1} \circ g_!$는 왼쪽 수반의 유일성에서 따른다.
\end{proof}

\section{토포스의 사상}
\label{sites-section-morphisms-topoi}

\noindent
이 절에서는 임의의 토포스의 사상이 사이트의 사상에서 유도되는 토포스의
사상과 동치임을 보인다. \cite[Expos\'e III, Proposition 4.9.4]{SGA4}와
비교하라.

\begin{lemma}
\label{sites-lemma-equivalence}
$\mathcal{C}$, $\mathcal{D}$를 사이트라 하자.
$u : \mathcal{C} \to \mathcal{D}$를 함자라 하자.
다음을 가정하자.
\begin{enumerate}
\item $u$는 쌍대연속이다.
\item $u$는 연속이다.
\item $\mathcal{C}$의 $a, b : U' \to U$에 대하여
$u(a) = u(b)$이면, $a \circ f_i = b \circ f_i$를 만족하는
$\mathcal{C}$의 덮개 $\{f_i : U'_i \to U'\}$가 존재한다.
\item $U', U \in \Ob(\mathcal{C})$ 및
$\mathcal{D}$의 사상 $c : u(U') \to u(U)$가 주어지면,
$u(c_i) = c \circ u(f_i)$를 만족하는 $\mathcal{C}$의 덮개
$\{f_i : U_i' \to U'\}$와 사상 $c_i : U_i' \to U$가 존재한다.
\item $V \in \Ob(\mathcal{D})$가 주어지면, $\mathcal{D}$에서
$\{u(U_i) \to V\}_{i \in I}$ 꼴인 $V$의 덮개가 존재한다.
\end{enumerate}
그러면 쌍대연속 함자 $u$에 Lemma
\ref{sites-lemma-cocontinuous-morphism-topoi}를 적용하여 얻는 토포스의 사상
$$
g : \Sh(\mathcal{C}) \longrightarrow \Sh(\mathcal{D})
$$
은 동치이다.
\end{lemma}

\begin{proof}
$u$가 성질 (1) -- (5)를 만족한다고 하자. 수반사상
$$
\mathcal{G} \longrightarrow g_*g^{-1}\mathcal{G}
\quad\text{and}\quad
g^{-1}g_*\mathcal{F} \longrightarrow \mathcal{F}
$$
이 동형사상임을 보이겠다.

\medskip\noindent
Lemma \ref{sites-lemma-when-shriek}를 적용할 수 있으므로
$\mathcal{D}$ 위의 임의의 층 $\mathcal{G}$에 대하여
$g^{-1}\mathcal{G}(U) = \mathcal{G}(u(U))$이다. 다음으로
$\mathcal{F}$를 $\mathcal{C}$ 위의 층이라 하고, $V$를
$\mathcal{D}$의 대상이라 하자. 정의에 의해
$g_*\mathcal{F}(V) = \lim_{u(U) \to V} \mathcal{F}(U)$이다. 따라서
$$
g^{-1}g_*\mathcal{F}(U) = \lim_{U', u(U') \to u(U)} \mathcal{F}(U')
$$
이며, 여기서 사상 $\psi : u(U') \to u(U)$는 반드시
$u(\alpha)$ 꼴일 필요가 없다.
% P01-ERRATA: P01-E0109 -- the frozen indexing-category claim names (U,id) final although preceding text says psi need not lift.
이러한 쌍 $(U', \psi)$의 범주는 종대상 $(U, \text{id})$를 가지며, 이로부터
그 극한에서 $\mathcal{F}(U)$로 가는 사상이 생긴다.
$(s_{(U', \psi)})$를 이 극한의 원소라 하자.
$s_{(U', \psi)}$가 값
$s_{(U, \text{id})} \in \mathcal{F}(U)$에 의해 유일하게 결정됨을
보이고자 한다. 성질 (4)에 의해 임의의 $(U', \psi)$에 대하여, 합성
$u(U'_i) \to u(U') \to u(U)$가 어떤
$\mathcal{C}$의 $c_i : U'_i \to U$에 대한 $u(c_i)$ 꼴이 되도록 하는
덮개 $\{U'_i \to U'\}$가 존재한다. 따라서
$$
s_{(U', \psi)}|_{U'_i} = c_i^*(s_{(U, \text{id})}).
$$
$\mathcal{F}$가 층이므로 실제로 $s_{(U', \psi)}$는
$s_{(U, \text{id})}$에 의해 결정된다. 이로써 유일성이 증명된다.
존재성을 위해 임의의
$s \in \mathcal{F}(U)$, $\psi : u(U') \to u(U)$,
$\{f_i : U_i' \to U'\}$와 위와 같이
$\psi \circ u(f_i) = u(c_i)$를 만족하는 $c_i : U_i' \to U$가
주어졌다고 하자. 다음을 만족하는 (유일한) 원소
$s_{(U', \psi)} \in \mathcal{F}(U')$가 존재한다고 주장한다.
$$
s_{(U', \psi)}|_{U'_i} = c_i^*(s).
$$
실제로 선험적으로는 원소
$c_i^*(s)|_{U_i' \times_{U'} U_j'}$와
$c_j^*(s)|_{U_i' \times_{U'} U_j'}$가 일치하는지 명확하지 않다. 도표
$$
\xymatrix{
U_i' \times_{U'} U_j' \ar[r]_-{\text{pr}_2} \ar[d]_{\text{pr}_1} &
U_j' \ar[d]^{c_j} \\
U_i' \ar[r]^{c_i} & U}
$$
가 가환할 필요가 없기 때문이다. 그러나 이 보조정리의 조건 (3)에 의해
$c_i \circ \text{pr}_1 \circ f_{ijk} = c_j \circ \text{pr}_2 \circ f_{ijk}$를
만족하는 덮개
$\{f_{ijk} : U'_{ijk} \to U_i' \times_{U'} U_j'\}_{k \in K_{ij}}$가
존재한다. 따라서
$$
f_{ijk}^* \left(c_i^*s|_{U_i' \times_{U'} U_j'}\right)
=
f_{ijk}^* \left(c_j^*s|_{U_i' \times_{U'} U_j'}\right)
$$
이다. 그러므로 $\mathcal{F}$의 층 조건에 의해
$c_i^*(s)|_{U_i' \times_{U'} U_j'} = c_j^*(s)|_{U_i' \times_{U'} U_j'}$이고,
다시 $\mathcal{F}$의 층 조건에 의해 $s_{(U', \psi)}$가 존재한다.
유일성에 의해 이렇게 얻은 족 $(s_{(U', \psi)})$는 극한의 원소이며
% P01-ERRATA: P01-E0108 -- the frozen distinguished component uses (U,psi) rather than the consistently chosen (U,id).
$s_{(U, \psi)} = s$를 만족한다. 이로써
$g^{-1}g_*\mathcal{F} \to \mathcal{F}$가 동형사상임이 증명된다.

\medskip\noindent
$\mathcal{G}$를 $\mathcal{D}$ 위의 층이라 하고, $V$를
$\mathcal{D}$의 대상이라 하자. 그러면
$$
g_*g^{-1}\mathcal{G}(V) =
\lim_{U, \psi : u(U) \to V} \mathcal{G}(u(U))
$$
이다. 앞 문단에 의해 $V = u(U)$ 꼴인 대상 $V$에서
층 $g_*g^{-1}\mathcal{G}$의 값은 $\mathcal{G}(u(U))$와 같다.
(형식적으로는 $g^{-1}g_*g^{-1} \cong g^{-1}$ 및 증명 첫머리에서 준
$g^{-1}$의 기술로부터 따르고, 비형식적으로는 여기와 위의 극한들을
비교하면 된다.) 따라서 수반사상
$\mathcal{G} \to g_*g^{-1}\mathcal{G}$는 $u(U)$ 꼴인 모든 대상 위의
절단에서 전단사이다. 공리 (5)에 의해 $u(U)$ 꼴인 대상들로 이루어진
$V$의 덮개가 존재하므로, 이 수반사상이 동형사상임을 쉽게 알 수 있다.
\end{proof}

\noindent
Lemma \ref{sites-lemma-equivalence}와 같은 쌍대연속 함자에 이름을 붙여 두는 것이
편리하다.

\begin{definition}
\label{sites-definition-special-cocontinuous-functor}
$\mathcal{C}$, $\mathcal{D}$를 사이트라 하자.
{\it $\mathcal{C}$에서 $\mathcal{D}$로 가는 특수 쌍대연속 함자 $u$}란
Lemma \ref{sites-lemma-equivalence}의 가정과 결론을 만족하는 쌍대연속 함자
$u : \mathcal{C} \to \mathcal{D}$를 말한다.
\end{definition}

\begin{lemma}
\label{sites-lemma-localize-special-cocontinuous}
$\mathcal{C}$, $\mathcal{D}$를 사이트라 하자.
$u : \mathcal{C} \to \mathcal{D}$를 특수 쌍대연속 함자라 하자.
$\mathcal{C}$의 모든 대상 $U$에 대하여 Lemma
\ref{sites-lemma-localize-cocontinuous}에서와 같은 가환도표
$$
\xymatrix{
\mathcal{C}/U \ar[r]_{j_U} \ar[d] & \mathcal{C} \ar[d]^u \\
\mathcal{D}/u(U) \ar[r]^-{j_{u(U)}} & \mathcal{D}
}
$$
가 있다. 왼쪽 수직 화살표는 특수 쌍대연속 함자이다.
따라서 토포스의 가환도표
$$
\xymatrix{
\Sh(\mathcal{C}/U) \ar[r]_{j_U} \ar[d] &
\Sh(\mathcal{C}) \ar[d]^u \\
\Sh(\mathcal{D}/u(U)) \ar[r]^-{j_{u(U)}} &
\Sh(\mathcal{D})
}
$$
에서 수직 화살표들은 동치이다.
\end{lemma}

\begin{proof}
도표들의 존재성과 가환성은 Lemma
\ref{sites-lemma-localize-cocontinuous}에서 보았다. 유도된 함자
$u : \mathcal{C}/U \to \mathcal{D}/u(U)$에 대하여 Lemma
\ref{sites-lemma-equivalence}의 가정 (1) -- (5)를 확인해야 한다.
이는 완전히 기계적이다.

\medskip\noindent
성질 (1). 이는 Lemma \ref{sites-lemma-localize-cocontinuous}이다.

\medskip\noindent
성질 (2). $\{U_i'/U \to U'/U\}_{i \in I}$를
$\mathcal{C}/U$에서 $U'/U$의 덮개라 하자. $u$가 연속이므로
$\{u(U_i')/u(U) \to u(U')/u(U)\}_{i \in I}$는
$\mathcal{D}/u(U)$에서 $u(U')/u(U)$의 덮개이다. 따라서
$u : \mathcal{C}/U \to \mathcal{D}/u(U)$에 대하여 (2)가 성립한다.

\medskip\noindent
성질 (3). $a, b : U''/U \to U'/U$를 $\mathcal{C}/U$의 사상이라 하고,
$u(a) = u(b)$가 $\mathcal{D}/u(U)$에서 성립한다고 하자.
$u$가 (3)을 만족하므로 $a \circ f_i = b \circ f_i$를 만족하는
$\mathcal{C}$의 덮개 $\{f_i : U''_i \to U''\}$가 존재한다.
이는 $a \circ f_i = b \circ f_i$를 만족하는 $\mathcal{C}/U$의 덮개
$\{f_i : U''_i/U \to U''/U\}$를 준다. 따라서
$u : \mathcal{C}/U \to \mathcal{D}/u(U)$에 대하여 (3)이 성립한다.

\medskip\noindent
성질 (4). $U''/U, U'/U \in \Ob(\mathcal{C}/U)$ 및
$\mathcal{D}/u(U)$의 사상
$c : u(U'')/u(U) \to u(U')/u(U)$가 주어졌다고 하자.
$u$가 성질 (4)를 만족하므로, $u(c_i) = c \circ u(f_i)$를 만족하는
$\mathcal{C}$의 덮개 $\{f_i : U_i'' \to U''\}$와 사상
$c_i : U_i'' \to U'$가 존재한다. 합성 $U_i'' \to U'' \to U$를
통해 $U_i''$를 $U$ 위의 대상으로 생각한다.
$c_i$가 $U$ 위의 사상일 필요는 없다! 그러나 $u(c_i)$는
$u(U)$ 위의 사상이므로 $u$에 성질 (3)을 적용하여,
$c_{ik} = c_i \circ f_{ik} : U''_{ik} \to U'$가 $U$ 위의 사상이 되게 하는
덮개 $\{f_{ik} : U''_{ik} \to U''_i\}$를 찾을 수 있다.
그러므로 $\{f_i \circ f_{ik} : U''_{ik}/U \to U''/U\}$는
$\mathcal{C}/U$의 덮개이고,
% P01-ERRATA: P01-E0110 -- the frozen refined equality omits f_i from the covering composite.
$u(c_{ik}) = c \circ u(f_{ik})$를 만족한다. 따라서
$u : \mathcal{C}/U \to \mathcal{D}/u(U)$에 대하여 (4)가 성립한다.

\medskip\noindent
성질 (5). $h : V \to u(U)$를 $\mathcal{D}/u(U)$의 대상이라 하자.
$u$가 성질 (5)를 만족하므로 $\mathcal{D}$의 덮개
$\{c_i : u(U_i) \to V\}$가 존재한다. 성질 (4)에 의해
$u(c_{ij}) = h \circ c_i \circ u(f_{ij})$를 만족하는 덮개
$\{f_{ij} : U_{ij} \to U_i\}$와 사상 $c_{ij} : U_{ij} \to U$를
찾을 수 있다. 따라서 $\{u(U_{ij})/u(U) \to V/u(U)\}$는
$\mathcal{D}/u(U)$에서 원하는 꼴의 덮개이고, 이로부터
$u : \mathcal{C}/U \to \mathcal{D}/u(U)$에 대하여 (5)가 성립한다.
\end{proof}

\begin{lemma}
\label{sites-lemma-special-equivalence}
$\mathcal{C}$를 사이트라 하자.
$\mathcal{C}' \subset \Sh(\mathcal{C})$를 다음을 만족하는
충만한 부분범주(대상들이 집합을 이룬다)라 하자.
\begin{enumerate}
\item 모든 $U \in \Ob(\mathcal{C})$에 대하여
$h_U^\# \in \Ob(\mathcal{C}')$이고,
\item $\mathcal{C}'$는 $\Sh(\mathcal{C})$에서의 섬유곱 아래 보존된다.
\end{enumerate}
$\mathcal{C}'$의 덮개를 층의 사상
$\coprod_{i \in I} \mathcal{F}_i \to \mathcal{F}$가 전사인 임의의 사상족
$\{\mathcal{F}_i \to \mathcal{F}\}_{i \in I}$로 선언하자. 그러면
\begin{enumerate}
\item $\mathcal{C}'$는 사이트이다(덮개들의 집합을 하나 선택한 뒤이며,
Sets, Lemma \ref{sets-lemma-coverings-site}를 보라).
\item $\mathcal{C}'$ 위의 표현가능 준층은 층이다(즉,
$\mathcal{C}'$ 위의 위상은 준표준이며, Definition
\ref{sites-definition-weaker-than-canonical}을 보라).
\item 함자 $v : \mathcal{C} \to \mathcal{C}'$,
$U \mapsto h_U^\#$는 특수 쌍대연속 함자이므로, 동치
$g : \Sh(\mathcal{C}) \to \Sh(\mathcal{C}')$를 유도한다.
\item 임의의 $\mathcal{F} \in \Ob(\mathcal{C}')$에 대하여
$g^{-1}h_\mathcal{F} = \mathcal{F}$이다.
\item 임의의 $U \in \Ob(\mathcal{C})$에 대하여
$g_*h_U^\# = h_{v(U)} = h_{h_U^\#}$이다.
\end{enumerate}
\end{lemma}

\begin{proof}
% P01-ERRATA: P01-E0111 -- the frozen warning contains the malformed phrase “may look be”.
경고: 위의 명제들 중 일부는 처음에는 조금 혼란스러워 보일 수 있다. 이는
$\mathcal{C}'$의 대상들을 $\mathcal{C}$ 위의 층으로도 볼 수 있기
때문이다! 보조정리에 기술한 $\mathcal{C}'$의 덮개들이 Definition
\ref{sites-definition-site}의 조건을 만족한다는 증명은 생략한다.

\medskip\noindent
$\{\mathcal{F}_i \to \mathcal{F}\}$가 층의 전사 사상족이라고 하자.
$\mathcal{G}$를 또 다른 층이라 하자. 보조정리의 (2)는
$$
\xymatrix{
\Mor_{\Sh(\mathcal{C})}(
\coprod_{i \in I} \mathcal{F}_i, \mathcal{G})
\ar@<1ex>[r] \ar@<-1ex>[r]
&
\Mor_{\Sh(\mathcal{C})}(
\coprod_{(i_0, i_1) \in I \times I}
\mathcal{F}_{i_0} \times_\mathcal{F} \mathcal{F}_{i_1}, \mathcal{G})
}
$$
의 등화자가 $\Mor_{\Sh(\mathcal{C})}(\mathcal{F}, \mathcal{G}).$라는
말이다. 이는 명백하다(예를 들어 Lemma
\ref{sites-lemma-coequalizer-surjection}을 사용하라).

\medskip\noindent
(3)을 증명하려면 Lemma \ref{sites-lemma-equivalence}의 조건 (1) -- (5)를
확인해야 한다. $v$가 쌍대연속이라는 것은 Lemma
\ref{sites-lemma-mono-epi-sheaves}의 층의 전사 사상에 대한 기술과 동치이다.
$U \mapsto h_U^\#$는 섬유곱과 가환하고 덮개를 덮개로 보내므로
(Lemma \ref{sites-lemma-sheafification-exact} 및 Lemma
\ref{sites-lemma-covering-surjective-after-sheafification}를 보라)
함자 $v$는 연속이다. Lemma \ref{sites-lemma-equivalence}의 성질 (3), (4)는
사상 $f : h_{U'}^\# \to h_U^\#$에 관한 명제이다. 이러한 사상은
$h_U^\#(U')$의 원소와 같은 것이다. 따라서 (3), (4)는 층화의 구성에서
바로 따른다. Lemma \ref{sites-lemma-equivalence}의 성질 (5)는 Lemma
\ref{sites-lemma-sheaf-coequalizer-representable}이다.
% P01-ERRATA: P01-E0112 -- the frozen naming sentence omits “by” before g.
$v$에 대응하는 토포스의 동치를 Lemma \ref{sites-lemma-equivalence}에 따라
$g : \Sh(\mathcal{C}) \to \Sh(\mathcal{C}')$로 나타내자.

\medskip\noindent
$\mathcal{F}$를 보조정리의 (4)에서와 같이 잡자. 임의의
$U \in \Ob(\mathcal{C})$에 대하여
$$
g^{-1}h_\mathcal{F}(U) = h_\mathcal{F}(v(U))
= \Mor_{\Sh(\mathcal{C})}(h_U^\#, \mathcal{F})
= \mathcal{F}(U)
$$
이다.
% P01-ERRATA: P01-E0113 -- the frozen first-equality explanation lacks a finite verb.
첫 번째 등식은 Lemma \ref{sites-lemma-when-shriek}에서 따른다.
따라서 (4)가 성립한다.

\medskip\noindent
$\mathcal{F} \in \Ob(\mathcal{C}')$이고
$U \in \Ob(\mathcal{C})$라 하자. 그러면
\begin{align*}
g_*h_U^\#(\mathcal{F})
& =
\Mor_{\Sh(\mathcal{C}')}(h_\mathcal{F}, g_*h_U^\#) \\
& =
\Mor_{\Sh(\mathcal{C})}(g^{-1}h_\mathcal{F}, h_U^\#) \\
& =
\Mor_{\Sh(\mathcal{C})}(\mathcal{F}, h_U^\#) \\
& =
\Mor_{\mathcal{C}'}(\mathcal{F}, h_U^\#)
\end{align*}
이며, 이는 원하는 바이다(세 번째 등식은 위에서 보였다).
\end{proof}

\noindent
이를 사용하면 임의의 토포스를 모든 유한 극한을 갖는 사이트 위에 놓이도록
정비할 수 있다.

\begin{lemma}
\label{sites-lemma-topos-good-site}
$\Sh(\mathcal{C})$를 토포스라 하자. $\{\mathcal{F}_i\}_{i \in I}$를
$\mathcal{C}$ 위의 층들의 집합이라 하자. 다음을 만족하는 사이트의 특수
쌍대연속 함자 $u : \mathcal{C} \to \mathcal{C}'$가 유도하는 토포스의
동치 $g : \Sh(\mathcal{C}) \to \Sh(\mathcal{C}')$가 존재한다.
\begin{enumerate}
\item $\mathcal{C}'$는 준표준 위상을 갖는다.
\item $\mathcal{C}'$의 사상족 $\{V_j \to V\}$가
$\mathcal{C}'$의 덮개와 (조합적으로) 동치일 필요충분조건은
$\coprod h_{V_j} \to h_V$가 전사인 것이다.
\item $\mathcal{C}'$는 섬유곱과 종대상을 갖는다
(즉, $\mathcal{C}'$는 모든 유한 극한을 갖는다).
\item $\mathcal{C}'$ 위의 표현가능 층의 모든 부분층은 표현가능하다.
\item 각 $g_*\mathcal{F}_i$는 표현가능 층이다.
\end{enumerate}
\end{lemma}

\begin{proof}
모든 $U \in \Ob(\mathcal{C})$에 대한 $h_U^\#$, 주어진 층
$\mathcal{F}_i$, 그리고 종층 $*$ (Example
\ref{sites-example-singleton-sheaf}를 보라)로 이루어진
충만한 부분범주 $\mathcal{C}_1 \subset \Sh(\mathcal{C})$를 생각하자.
충만한 부분범주들을 귀납적으로 정의하겠다.
% P01-ERRATA: P01-E0114 -- the frozen increasing chain repeats C_2 where C_3 is intended.
$$
\mathcal{C}_1 \subset \mathcal{C}_2 \subset \mathcal{C}_2 \subset \ldots
\subset \Sh(\mathcal{C})
$$
즉, $\mathcal{C}_n$이 주어지면 $\mathcal{C}_{n + 1}$을
$\mathcal{C}_n$의 대상들의 모든 섬유곱과 부분층으로 이루어진 충만한
부분범주로 잡는다. ($\mathcal{C}_{n + 1}$의 대상들이 집합을 이룸에
유의하라.) $\mathcal{C}' = \bigcup_{n \geq 1} \mathcal{C}_n$으로 놓자.
$\mathcal{C}'$의 덮개는 $\mathcal{C}'$의 대상들의 사상족
$\{\mathcal{G}_j \to \mathcal{G}\}_{j \in J}$ 가운데
$\coprod \mathcal{G}_j \to \mathcal{G}$가 $\mathcal{C}$ 위의 층의
사상으로서 전사인 것으로 잡는다. 함자
$v : \mathcal{C} \to \mathcal{C'}$는 $U \mapsto h_U^\#$로 주어진다.
Lemma \ref{sites-lemma-special-equivalence}를 적용하라.
\end{proof}

\noindent
다음이 이 절의 목표이다.

\begin{lemma}
\label{sites-lemma-morphism-topoi-comes-from-morphism-sites}
\begin{reference}
이 명제는 \cite[Proposition 4.9.4. Expos\'e IV]{SGA4}와 밀접하게
관련되어 있다. 전체 결과를 얻으려면
\cite[Remarque 4.7.4, Expos\'e IV]{SGA4}도 사용해야 한다.
\end{reference}
$\mathcal{C}$, $\mathcal{D}$를 사이트라 하자.
$f : \Sh(\mathcal{C}) \to \Sh(\mathcal{D})$를 토포스의 사상이라 하자.
그러면 사이트 $\mathcal{C}'$와 함자들의 도표
$$
\xymatrix{
\mathcal{C} \ar[r]_v & \mathcal{C}' & \mathcal{D} \ar[l]^u
}
$$
가 존재하여 다음을 만족한다.
\begin{enumerate}
\item 함자 $v$는 특수 쌍대연속 함자이다.
\item 함자 $u$는 섬유곱과 가환하고 연속이며, 사이트의 사상
$\mathcal{C}' \to \mathcal{D}$를 정의한다.
\item 토포스의 사상 $f$는 토포스의 사상들의 합성
$$
\Sh(\mathcal{C}) \longrightarrow
\Sh(\mathcal{C}') \longrightarrow
\Sh(\mathcal{D})
$$
과 일치한다. 여기서 첫째 화살표는 Lemma
\ref{sites-lemma-equivalence}를 통해 $v$에서 유도되고, 둘째 화살표는
Lemma \ref{sites-lemma-morphism-sites-topoi}를 통해 $u$에서 유도된다.
\end{enumerate}
\end{lemma}

\begin{proof}
모든 $U \in \Ob(\mathcal{C})$에 대한 $h_U^\#$와 모든
$V \in \Ob(\mathcal{D})$에 대한 $f^{-1}h_V^\#$로 이루어진 충만한
부분범주 $\mathcal{C}_1 \subset \Sh(\mathcal{C})$를 생각하자.
$\mathcal{C}_{n + 1}$을 $\mathcal{C}_n$의 대상들의 모든 섬유곱으로
이루어진 충만한 부분범주로 잡자.
$\mathcal{C}' = \bigcup_{n \geq 1} \mathcal{C}_n$으로 놓는다.
$\mathcal{C}'$의 덮개는 사상족
$\{\mathcal{F}_i \to \mathcal{F}\}_{i \in I}$ 가운데
$\coprod_{i \in I} \mathcal{F}_i \to \mathcal{F}$가
$\mathcal{C}$ 위의 층의 사상으로서 전사인 것으로 잡는다.
함자 $v : \mathcal{C} \to \mathcal{C'}$는
$U \mapsto h_U^\#$로 주어진다. 함자
$u : \mathcal{D} \to \mathcal{C'}$는
$V \mapsto f^{-1}h_V^\#$로 주어진다.

\medskip\noindent
(1)은 Lemma \ref{sites-lemma-special-equivalence}에서 따른다.

\medskip\noindent
보조정리의 (2), (3)을 증명하자. $V \mapsto h_V^\#$와 $f^{-1}$가 모두
섬유곱과 가환하므로 함자 $u$도 섬유곱과 가환한다. 또한 $f^{-1}$는
완전하고 임의의 쌍대극한과 가환하므로, 덮개를 층의 전사 사상족으로
보낸다. 따라서 $u$는 연속이다. 이것이 사이트의 사상을 정의함을 보이려면
$u_s$가 완전함을 더 보여야 한다. 이를 위해
$g^{-1} \circ u_s = f^{-1}$임을 보이겠다. 실제로 그러면 $g^{-1}$가
동치이고 $f^{-1}$가 완전하므로 결론을 얻는다.
$g^{-1}$는 $g_*$의 수반이고, $u_s$는 $u^s$의 수반이며,
$f^{-1}$는 $f_*$의 수반이므로 $u^s \circ g_* = f_*$임을 보이는 것으로도
충분하다. $\mathcal{F}$를 $\mathcal{C}$ 위의 층이라 하고 $V$를
$\mathcal{D}$의 대상이라 하자. 그러면
\begin{align*}
(u^s g_{\ast}\mathcal{F})(V)
& =
(g_{\ast}\mathcal{F})(f^{-1}h_V^\#) \\
& =
\Mor_{Sh(\mathcal{C}')}(h_{f^{-1}h_V^\#},g_{\ast}\mathcal{F}) \\
& =
\Mor_{Sh(\mathcal{C})}(g^{-1}h_{f^{-1}h_V^\#},\mathcal{F}) \\
& =
\Mor_{Sh(\mathcal{C})}(f^{-1}h_V^\#,\mathcal{F}) \\
& =
\Mor_{Sh(\mathcal{D})}(h_V^\#,f_{\ast}\mathcal{F}) \\
& =
f_{\ast}\mathcal{F}(V)
\end{align*}
이다.
% P01-ERRATA: P01-E0115 -- three frozen equality explanations are sentence fragments.
첫 번째 등식은 $u^s = u^p$이므로 성립한다.
두 번째 등식은 요네다 보조정리이다.
세 번째 등식은 $g^{-1}$와 $g_*$의 수반성에서 따른다.
네 번째 등식은 Lemma \ref{sites-lemma-special-equivalence} (4)에 의한다.
다섯 번째 등식은 $f^{-1}$와 $f_*$의 수반성에서 따른다.
여섯 번째 등식은 요네다 보조정리이다.
따라서 $u^s g_*\mathcal{F} = f_*\mathcal{F}$이고, 이로써 보조정리의
증명이 끝난다.
\end{proof}

\begin{remark}
\label{sites-remark-morphism-topoi-comes-from-morphism-sites}
표기와 가정은 Lemma
\ref{sites-lemma-morphism-topoi-comes-from-morphism-sites}에서와 같다.
사이트 $\mathcal{D}$가 종대상과 섬유곱을 가지면, 함자
$u : \mathcal{D} \to \mathcal{C}'$는 Proposition
\ref{sites-proposition-get-morphism}의 모든 가정을 만족한다. 실제로 보조정리에서
언급한 성질들에 더하여, $u$는 $\mathcal{D}$의 종대상을
$\mathcal{C}'$의 종대상으로 보낸다. 이는 $u$의 구성에서 명백하다.
따라서 먼저
% P01-ERRATA: P01-E0116 -- the frozen text pluralizes “Lemmas” before a single reference.
Lemma \ref{sites-lemma-topos-good-site}를 $\mathcal{D}$에 적용하고, 이어서
그 결과 얻은 토포스의 사상
$\Sh(\mathcal{C}) \to \Sh(\mathcal{D}')$에 Lemma
\ref{sites-lemma-morphism-topoi-comes-from-morphism-sites}를 적용하면 다음 명제를
얻는다. 임의의 토포스의 사상
$f : \Sh(\mathcal{C}) \to \Sh(\mathcal{D})$는 가환도표
$$
\xymatrix{
\Sh(\mathcal{C}) \ar[d]_g \ar[r]_f &
\Sh(\mathcal{D}) \ar[d]^e \\
\Sh(\mathcal{C}') \ar[r]^{f'} &
\Sh(\mathcal{D}')
}
$$
에 들어가며, 다음 성질들이 성립한다.
\begin{enumerate}
\item 사상 $e$와 $g$는 특수 쌍대연속 함자
$\mathcal{C} \to \mathcal{C}'$ 및
$\mathcal{D} \to \mathcal{D}'$에 의해 주어지는 동치이다.
\item 사이트 $\mathcal{C}'$와 $\mathcal{D}'$는 섬유곱과 종대상을 가지며
준표준 위상을 갖는다.
\item 사상 $f' : \mathcal{C}' \to \mathcal{D}'$는 Proposition
\ref{sites-proposition-get-morphism}을 적용할 수 있는 함자
$u : \mathcal{D}' \to \mathcal{C}'$에 대응하는 사이트의 사상에서
유도된다.
\item $\mathcal{C}$ 위의 임의의 층들의 집합 $\mathcal{F}_i$
(각각 $\mathcal{D}$ 위의 $\mathcal{G}_j$)이 주어지면, 이들 각각을
$\mathcal{C}'$ (각각 $\mathcal{D}'$) 위의 표현가능 층이라고 가정할 수
있다.
\end{enumerate}
$\mathcal{C}$와 $\mathcal{D}$를 $\mathcal{C}'$와 $\mathcal{D}'$로
대체하는 것이 흔히 유용하다.
\end{remark}

\begin{remark}
\label{sites-remark-equivalence-topoi-comes-from-morphism-sites}
표기와 가정은 Lemma
\ref{sites-lemma-morphism-topoi-comes-from-morphism-sites}에서와 같다.
원래의 토포스의 사상
$\Sh(\mathcal{C}) \to \Sh(\mathcal{D})$가 추가로 동치라고 하자.
그러면 Lemma \ref{sites-lemma-morphism-topoi-comes-from-morphism-sites}의
증명에서 구성한 두 함자
$$
\mathcal{C} \rightarrow \mathcal{C}' \leftarrow \mathcal{D}
$$
는 모두 특수 쌍대연속 함자이다. 따라서 이 경우 토포스의 사상을
Remark \ref{sites-remark-morphism-topoi-comes-from-morphism-sites}에서와 같이 합성
$$
\Sh(\mathcal{C}) \rightarrow
\Sh(\mathcal{C}') =
\Sh(\mathcal{D}') \leftarrow
\Sh(\mathcal{D})
$$
으로 분해할 수 있으며, 가운데 사상은 항등사상이다.
\end{remark}

\section{토포스의 국소화}
\label{sites-section-localize-topoi}

\noindent
국소화에 관한 내용 일부를 겉보기에는 더 일반적인 토포스의 경우에 대하여
되풀이한다. 실제로는 바탕 사이트를 언제나 확대하여 사이트의 대상에서
국소화한다고 가정할 수 있으므로 더 일반적인 것은 아니다.

\begin{lemma}
\label{sites-lemma-localize-topos}
$\mathcal{C}$를 사이트라 하자.
$\mathcal{F}$를 $\mathcal{C}$ 위의 층이라 하자.
그러면 범주 $\Sh(\mathcal{C})/\mathcal{F}$는 토포스이다.
Section \ref{sites-section-localize}에서와 같은 국소화인 토포스의 표준 사상
$$
j_\mathcal{F} :
\Sh(\mathcal{C})/\mathcal{F}
\longrightarrow
\Sh(\mathcal{C})
$$
이 존재하며 다음을 만족한다.
\begin{enumerate}
\item 함자 $j_\mathcal{F}^{-1}$는
$\mathcal{H} \mapsto \mathcal{H} \times \mathcal{F}/\mathcal{F}$인
함자이다.
\item 함자 $j_{\mathcal{F}!}$는 망각 함자
$\mathcal{G}/\mathcal{F} \mapsto \mathcal{G}$이다.
\end{enumerate}
\end{lemma}

\begin{proof}
Lemma \ref{sites-lemma-topos-good-site}를 적용하자. 이는 $\mathcal{C}$가
준표준 위상을 갖는 사이트이고 어떤 $U \in \Ob(\mathcal{C})$에 대하여
$\mathcal{F} = h_U = h_U^\#$라고 가정해도 됨을 뜻한다. 따라서 Section
\ref{sites-section-localize}의 내용을 적용할 수 있다. 특히 동치
$\Sh(\mathcal{C}/U) = \Sh(\mathcal{C})/h_U^\#$가 존재하여 합성
$$
\Sh(\mathcal{C}/U)
\to
\Sh(\mathcal{C})/h_U^\#
\to \Sh(\mathcal{C})
$$
은 $j_{U!}$와 같다. Lemma
\ref{sites-lemma-essential-image-j-shriek}를 보라.
% P01-ERRATA: P01-E0117 -- the frozen naming construction omits “by”.
역함자를
$a : \Sh(\mathcal{C})/h_U^\# \to \Sh(\mathcal{C}/U)$로 나타내자. 그러면
$j_{\mathcal{F}!} = j_{U!} \circ a$,
$j_\mathcal{F}^{-1} = a^{-1} \circ j_U^{-1}$, 그리고
$j_{\mathcal{F}, *} = j_{U, *} \circ a$이다.
$j_{\mathcal{F}!}$의 기술은 위에서 따르고, $j_\mathcal{F}^{-1}$의 기술은
Lemma \ref{sites-lemma-compute-j-shriek-restrict}에서 따른다.
\end{proof}

\begin{lemma}
\label{sites-lemma-localize-pushforward}
Lemma \ref{sites-lemma-localize-topos}의 상황에서 함자
% P01-ERRATA: P01-E0118 -- the frozen relative construction omits “that”.
$j_{\mathcal{F}, *}$는 $\varphi : \mathcal{G} \to \mathcal{F}$에 층
$$
U
\longmapsto
\{\alpha : \mathcal{F}|_U \to \mathcal{G}|_U
\text{ such that } \alpha \text{ is a right inverse to }\varphi|_U \}.
$$
을 대응시키는 함자이다.
\end{lemma}

\begin{proof}
임의의 $\varphi : \mathcal{G} \to \mathcal{F}$에 대하여
$\Sh(\mathcal{C})/\mathcal{F}$의 대응하는 대상을
$\mathcal{G}/\mathcal{F}$로 나타내자. 다음이 성립한다.
$$
(j_{\mathcal{F}, *}(\mathcal{G}/\mathcal{F}))(U) =
\Mor_{\Sh(\mathcal{C})}(h_U^\#, j_{\mathcal{F}, *}(\mathcal{G}/\mathcal{F})) =
\Mor_{\Sh(\mathcal{C})/\mathcal{F}}(j_{\mathcal{F}}^{-1}h_U^{\#},
(\mathcal{G}/\mathcal{F})).
$$
Lemma \ref{sites-lemma-localize-topos}에 의해 이 집합은 집합들의 사상
$$
\Mor_{\Sh(\mathcal{C})}(h_U^\# \times \mathcal{F}, \mathcal{G})
\xrightarrow{\varphi \circ}
\Mor_{\Sh(\mathcal{C})}(h_U^\# \times \mathcal{F}, \mathcal{F}).
$$
에서 원소 $h_U^\# \times \mathcal{F} \to \mathcal{F}$ 위의 섬유이다.
Lemma \ref{sites-lemma-internal-hom-sheaf}의 수반성에 의해
\begin{align*}
\Mor_{\Sh(\mathcal{C})}(h_U^{\#}\times\mathcal{F}, \mathcal{G})
& =
\Mor_{\Sh(\mathcal{C})}(h_U^{\#},\SheafHom(\mathcal{F}, \mathcal{G})) \\
& =
\Mor_{\Sh(\mathcal{C}/U)}(\mathcal{F}|_{\mathcal{C}/U},
\mathcal{G}|_{\mathcal{C}/U}), \\
\Mor_{\Sh(\mathcal{C})}(h_U^{\#} \times \mathcal{F}, \mathcal{F})
& =
\Mor_{\Sh(\mathcal{C})}(h_U^{\#},\SheafHom(\mathcal{F},\mathcal{F})) \\
& =
\Mor_{\Sh(\mathcal{C}/U)}(\mathcal{F}|_{\mathcal{C}/U},
\mathcal{F}|_{\mathcal{C}/U}),
\end{align*}
이고, 이 수반성 아래에서 사상 $\varphi\circ$는 사상
$$
\Mor_{\Sh(\mathcal{C}/U)}(\mathcal{F}|_{\mathcal{C}/U},
\mathcal{G}|_{\mathcal{C}/U})
\longrightarrow
\Mor_{\Sh(\mathcal{C}/U)}(\mathcal{F}|_{\mathcal{C}/U},
\mathcal{F}|_{\mathcal{C}/U}),\quad
\psi \longmapsto \varphi|_{\mathcal{C}/U} \circ \psi,
$$
이 되고, 원소 $h_U^\# \times \mathcal{F} \to \mathcal{F}$는
$\text{id}_{\mathcal{F}|_{\mathcal{C}/U}}$가 된다. 따라서
$(j_{\mathcal{F}, *}\mathcal{G}/\mathcal{F})(U)$는 사상
$$
\Mor_{\Sh(\mathcal{C}/U)}(\mathcal{F}|_{\mathcal{C}/U},
\mathcal{G}|_{\mathcal{C}/U})
\xrightarrow{\varphi|_{\mathcal{C}/U}\circ}
\Mor_{\Sh(\mathcal{C}/U)}(\mathcal{F}|_{\mathcal{C}/U},
\mathcal{F}|_{\mathcal{C}/U}),
$$
에서 $\text{id}_{\mathcal{F}|_{\mathcal{C}/U}}$ 위의 섬유와 동형이며,
이는 원하는 대로
\[
\{\alpha : \mathcal{F}|_U \to \mathcal{G}|_U
\text{ such that } \alpha \text{ is a right inverse to }\varphi|_U \}
\]
이다.
\end{proof}

\begin{lemma}
\label{sites-lemma-localize-topos-site}
$\mathcal{C}$를 사이트라 하고, $\mathcal{F}$를 $\mathcal{C}$ 위의
층이라 하자. $\mathcal{C}/\mathcal{F}$를 $U \in \Ob(\mathcal{C})$이고
$s \in \mathcal{F}(U)$인 쌍 $(U, s)$들의 범주라 하자.
$\mathcal{C}/\mathcal{F}$의 덮개를 $\{U_i \to U\}$가
$\mathcal{C}$의 덮개인 사상족
$\{(U_i, s_i) \to (U, s)\}$로 잡는다. 그러면
$j : \mathcal{C}/\mathcal{F} \to \mathcal{C}$는 사이트의 연속이고
쌍대연속인 함자이며 토포스의 사상
$j : \Sh(\mathcal{C}/\mathcal{F}) \to \Sh(\mathcal{C})$를 유도한다.
실제로 $j$를 $j_\mathcal{F}$로 바꾸는 동치
$\Sh(\mathcal{C}/\mathcal{F}) =
\Sh(\mathcal{C})/\mathcal{F}$가 존재한다.
\end{lemma}

\begin{proof}
$\mathcal{C}/\mathcal{F}$가 사이트라는 것과 $j$가 연속이고
쌍대연속이라는 것의 확인은 생략한다. Lemma
\ref{sites-lemma-when-shriek}에 의해 표시한 바와 같은 토포스의 사상 $j$가
존재하며, $j^{-1}\mathcal{G}(U, s) = \mathcal{G}(U)$이고
$j^{-1}$에는 왼쪽 수반 $j_!$가 존재한다.
$\mathcal{C}/\mathcal{F}$ 위의 사상
$\varphi : * \to j^{-1}\mathcal{G}$는 각 쌍 $(U, s)$에 제한사상들과
양립하는 절단 $\varphi(s) \in \mathcal{G}(U)$를 대응시키는 규칙과 같다.
따라서 이는 $\mathcal{C}$ 위의 사상
$\varphi : \mathcal{F} \to \mathcal{G}$와 같다. 이로부터
$j_!* = \mathcal{F}$를 얻는다. 특히 모든
$\mathcal{H} \in \Sh(\mathcal{C}/\mathcal{F})$에 대하여 표준 사상
$$
j_!\mathcal{H} \to j_!* = \mathcal{F}
$$
이 있다. 즉 함자
$j'_! : \Sh(\mathcal{C}/\mathcal{F}) \to \Sh(\mathcal{C})/\mathcal{F}$를
얻는다. 이 함자의 역함자는
$\Sh(\mathcal{C})/\mathcal{F}$의 대상
$\varphi : \mathcal{G} \to \mathcal{F}$에 층
$$
a(\mathcal{G}/\mathcal{F}) :
(U, s) \longmapsto \{t \in \mathcal{G}(U) \mid \varphi(t) = s\}
$$
을 대응시키는 규칙이다. $a(\mathcal{G}/\mathcal{F})$가 층이라는 것과
$a$가 $j'_!$의 역함자라는 것의 확인은 생략한다.
\end{proof}

\begin{definition}
\label{sites-definition-localize-topos}
$\mathcal{C}$를 사이트라 하자.
$\mathcal{F}$를 $\mathcal{C}$ 위의 층이라 하자.
\begin{enumerate}
\item 토포스 $\Sh(\mathcal{C})/\mathcal{F}$를
{\it $\mathcal{F}$에서 토포스 $\Sh(\mathcal{C})$의 국소화}라고 한다.
\item Lemma \ref{sites-lemma-localize-topos}의 토포스의 사상
$j_\mathcal{F} :
\Sh(\mathcal{C})/\mathcal{F}
\to
\Sh(\mathcal{C})$를 {\it 국소화 사상}이라고 한다.
\end{enumerate}
\end{definition}

\noindent
층 $\mathcal{F}$가 사이트의 어떤 대상 $U$에 대한 $h_U^\#$와 같을 때마다
토포스의 국소화가 $U$에서 사이트를 국소화한 것 위의 층들의 범주와 같음을
보이겠다. 또한 모든 함자성이 이 식별과 양립함을 확인하겠다.

\begin{lemma}
\label{sites-lemma-localize-compare}
$\mathcal{C}$를 사이트라 하자. $\mathcal{F} = h_U^\#$이고 $U$가
$\mathcal{C}$의 대상이라고 하자. Lemma \ref{sites-lemma-localize-topos}에서
구성한 $j_\mathcal{F} : \Sh(\mathcal{C})/\mathcal{F}
\to \Sh(\mathcal{C})$는 Section \ref{sites-section-localize}에서 구성한
토포스의 사상 $j_U : \Sh(\mathcal{C}/U) \to \Sh(\mathcal{C})$와
Lemma \ref{sites-lemma-essential-image-j-shriek}의 식별
$\Sh(\mathcal{C}/U) = \Sh(\mathcal{C})/h_U^\#$을 통하여 일치한다.
\end{lemma}

\begin{proof}
Lemma \ref{sites-lemma-essential-image-j-shriek}에서 합성
$\Sh(\mathcal{C}/U) \to \Sh(\mathcal{C})/h_U^\#
\to \Sh(\mathcal{C})$가 $j_{U!}$임을 보았다. 함자
$\Sh(\mathcal{C})/h_U^\# \to \Sh(\mathcal{C})$는 Lemma
\ref{sites-lemma-localize-topos}에 의해 $j_{\mathcal{F}!}$이다. 따라서 이 식별
아래 $j_{\mathcal{F}!} = j_{U!}$이다. 그러므로 (수반성에 의해)
$j_\mathcal{F}^{-1} = j_U^{-1}$이고, (다시 수반성에 의해)
$j_{\mathcal{F}, *} = j_{U, *}$이다.
\end{proof}

\begin{lemma}
\label{sites-lemma-relocalize-topos}
$\mathcal{C}$를 사이트라 하자.
$s : \mathcal{G} \to \mathcal{F}$가 $\mathcal{C}$ 위의 층들의
사상이면 토포스의 사상들의 자연스러운 가환도표
$$
\xymatrix{
\Sh(\mathcal{C})/\mathcal{G} \ar[rd]_{j_\mathcal{G}} \ar[rr]_j & &
\Sh(\mathcal{C})/\mathcal{F} \ar[ld]^{j_\mathcal{F}} \\
& \Sh(\mathcal{C}) &
}
$$
가 존재한다. 여기서 $j = j_{\mathcal{G}/\mathcal{F}}$는 토포스
$\Sh(\mathcal{C})/\mathcal{F}$를 대상
$\mathcal{G}/\mathcal{F}$에서 국소화한 것이다. 특히
$$
j^{-1}(\mathcal{H} \to \mathcal{F}) =
(\mathcal{H} \times_\mathcal{F} \mathcal{G} \to \mathcal{G})
$$
이고
$$
j_!(\mathcal{E} \xrightarrow{e} \mathcal{G}) =
(\mathcal{E} \xrightarrow{s \circ e} \mathcal{F}).
$$
이다.
\end{lemma}

\begin{proof}
$j^{-1}$와 $j_!$의 기술은 Lemma
\ref{sites-lemma-localize-topos}에 있는 이 함자들의 기술에서 따른다.
함자들의 등식 $j_{\mathcal{G}!} = j_{\mathcal{F}!} \circ j_!$은
이 함자들을 망각 함자로 기술한 것에서 명백하다. 수반성에 의해 등식
$j_\mathcal{G}^{-1} = j^{-1} \circ j_\mathcal{F}^{-1}$ 및
$j_{\mathcal{G}, *} = j_{\mathcal{F}, *} \circ j_*$도 얻는다.
\end{proof}

\begin{lemma}
\label{sites-lemma-relocalize-compare}
$\mathcal{C}$와 $s : \mathcal{G} \to \mathcal{F}$에 대한 가정은
Lemma \ref{sites-lemma-relocalize-topos}에서와 같다고 하자.
$\mathcal{G} = h_V^\#$, $\mathcal{F} = h_U^\#$이고
$s : \mathcal{G} \to \mathcal{F}$가 $\mathcal{C}$의 사상
$V \to U$에서 유도되면, Lemma \ref{sites-lemma-relocalize-topos}의 도표는
diagram (\ref{sites-equation-relocalize})와 Lemma
\ref{sites-lemma-essential-image-j-shriek}의 식별
$\Sh(\mathcal{C}/V) = \Sh(\mathcal{C})/h_V^\#$ 및
$\Sh(\mathcal{C}/U) = \Sh(\mathcal{C})/h_U^\#$을 통하여
식별된다.
\end{lemma}

\begin{proof}
$j^{-1}$의 기술들이 일치하므로 참이다. Lemma
\ref{sites-lemma-relocalize-explicit} 및 Lemma
\ref{sites-lemma-relocalize-topos}를 보라.
\end{proof}

\section{국소화와 토포스의 사상}
\label{sites-section-localize-morphisms-topoi}

\noindent
이 절은 토포스의 사상에 대한 Section
\ref{sites-section-localize-morphisms}의 유사물이다.

\begin{lemma}
\label{sites-lemma-localize-morphism-topoi}
$f : \Sh(\mathcal{C}) \to \Sh(\mathcal{D})$를 토포스의 사상이라 하자.
$\mathcal{G}$를 $\mathcal{D}$ 위의 층이라 하고
$\mathcal{F} = f^{-1}\mathcal{G}$로 놓자. 그러면 토포스의 가환도표
$$
\xymatrix{
\Sh(\mathcal{C})/\mathcal{F} \ar[r]_{j_\mathcal{F}} \ar[d]_{f'} &
\Sh(\mathcal{C}) \ar[d]^f \\
\Sh(\mathcal{D})/\mathcal{G} \ar[r]^{j_\mathcal{G}} &
\Sh(\mathcal{D}).
}
$$
가 존재한다. 사상 $f'$는 다음 성질로 특징지어진다.
$$
(f')^{-1}(\mathcal{H} \xrightarrow{\varphi} \mathcal{G})
=
(f^{-1}\mathcal{H} \xrightarrow{f^{-1}\varphi} \mathcal{F})
$$
또한 $f'_*j_\mathcal{F}^{-1} = j_\mathcal{G}^{-1}f_*$이다.
\end{lemma}

\begin{proof}
이 명제는 토포스에 관한 것이고 바탕 사이트들을 언급하지 않으므로 사이트를
마음대로 바꾸어도 된다. 따라서 Remark
\ref{sites-remark-morphism-topoi-comes-from-morphism-sites}의 논의에 의해,
모든 유한 극한과 준표준 위상을 갖는 사이트들 사이에서
Proposition \ref{sites-proposition-get-morphism}의 가정을 만족하는 연속 함자
$u : \mathcal{D} \to \mathcal{C}$가 $f$를 주고, 어떤
$\mathcal{D}$의 대상 $V$에 대하여 $\mathcal{G} = h_V$라고 가정해도
된다. 그러면 Lemma \ref{sites-lemma-pullback-representable-sheaf}에 의해
$\mathcal{F} = f^{-1}h_V = h_{u(V)}$이다.
Lemma \ref{sites-lemma-localize-morphism}에 의해 토포스의 사상들의 가환도표
% P01-ERRATA: P01-E0119 -- the frozen proof localizes at undefined U rather than explicitly setting U=u(V).
$$
\xymatrix{
\Sh(\mathcal{C}/U) \ar[r]_{j_U} \ar[d]_{f'} &
\Sh(\mathcal{C}) \ar[d]^f \\
\Sh(\mathcal{D}/V) \ar[r]^{j_V} &
\Sh(\mathcal{D}),
}
$$
를 얻고, $f'_*j_U^{-1} = j_V^{-1}f_*$이다. Lemma
\ref{sites-lemma-localize-compare}에 의해 $j_\mathcal{F}$와 $j_U$를,
$j_\mathcal{G}$와 $j_V$를 각각 식별할 수 있다. $(f')^{-1}$의 기술은
Lemma \ref{sites-lemma-localize-morphism}에 주어져 있다.
\end{proof}

\begin{lemma}
\label{sites-lemma-localize-morphism-compare}
$f : \mathcal{C} \to \mathcal{D}$를 연속 함자
$u : \mathcal{D} \to \mathcal{C}$가 주는 사이트의 사상이라 하자.
$V$를 $\mathcal{D}$의 대상이라 하고 $U = u(V)$로 놓자.
$\mathcal{G} = h_V^\#$ 및
$\mathcal{F} = h_U^\# = f^{-1}h_V^\#$로 놓는다(Lemma
\ref{sites-lemma-pullback-representable-sheaf}를 보라).
그러면 Lemma \ref{sites-lemma-localize-morphism-topoi}의 토포스의 사상들의
도표는 Lemma \ref{sites-lemma-localize-morphism}의 토포스의 사상들의 도표와
Lemma \ref{sites-lemma-localize-compare}의 식별
$j_\mathcal{F}= j_U$ 및 $j_\mathcal{G} = j_V$를 통하여 일치한다.
\end{lemma}

\begin{proof}
Lemma \ref{sites-lemma-localize-morphism-topoi}의 증명에서 $f$를 유도하도록
선택한 사이트의 사상은 이 보조정리에서 주어진 사이트의 사상과 다를 수
있으므로 이는 완전히 자명하지는 않다. 그러나 두 경우 모두 함자
$(f')^{-1}$는 같은 규칙으로 기술된다. 따라서 이들은 일치하고 대응하는
토포스의 사상도 같다. 몇 가지 세부사항은 생략한다.
\end{proof}

\begin{lemma}
\label{sites-lemma-relocalize-morphism-topoi}
$f : \Sh(\mathcal{C}) \to \Sh(\mathcal{D})$를 토포스의 사상이라 하자.
$\mathcal{G} \in \Sh(\mathcal{D})$,
$\mathcal{F} \in \Sh(\mathcal{C})$이고
$s : \mathcal{F} \to f^{-1}\mathcal{G}$를 층의 사상이라 하자.
그러면 토포스의 가환도표
$$
\xymatrix{
\Sh(\mathcal{C})/\mathcal{F} \ar[r]_{j_\mathcal{F}} \ar[d]_{f_s} &
\Sh(\mathcal{C}) \ar[d]^f \\
\Sh(\mathcal{D})/\mathcal{G} \ar[r]^{j_\mathcal{G}} &
\Sh(\mathcal{D}).
}
$$
가 존재한다. 다음이 성립한다.
$f_s = f' \circ j_{\mathcal{F}/f^{-1}\mathcal{G}}$. 여기서
% P01-ERRATA: P01-E0120 -- the frozen codomain slices Sh(D) over the C-sheaf F instead of G.
$f' :
\Sh(\mathcal{C})/f^{-1}\mathcal{G}
\to
\Sh(\mathcal{D})/\mathcal{F}$
는 Lemma \ref{sites-lemma-localize-morphism-topoi}에서와 같고,
$j_{\mathcal{F}/f^{-1}\mathcal{G}} :
\Sh(\mathcal{C})/\mathcal{F}
\to
\Sh(\mathcal{C})/f^{-1}\mathcal{G}$
는 Lemma \ref{sites-lemma-relocalize-topos}에서와 같다.
함자 $(f_s)^{-1}$는 다음 규칙으로 기술된다.
$$
(f_s)^{-1}(\mathcal{H} \xrightarrow{\varphi} \mathcal{G})
=
(f^{-1}\mathcal{H} \times_{f^{-1}\varphi, f^{-1}\mathcal{G}, s} \mathcal{F}
\rightarrow \mathcal{F}).
$$
마지막으로 사상 $b : \mathcal{G}' \to \mathcal{G}$,
$a : \mathcal{F}' \to \mathcal{F}$ 및
$s' : \mathcal{F}' \to f^{-1}\mathcal{G}'$가 주어져서 도표
$$
\xymatrix{
\mathcal{F}' \ar[r]_-{s'} \ar[d]_a & f^{-1}\mathcal{G}' \ar[d]^{f^{-1}b} \\
\mathcal{F} \ar[r]^-s & f^{-1}\mathcal{G}
}
$$
가 가환하면 도표
$$
\xymatrix{
\Sh(\mathcal{C})/\mathcal{F}'
\ar[r]_{j_{\mathcal{F}'/\mathcal{F}}} \ar[d]_{f_{s'}} &
\Sh(\mathcal{C})/\mathcal{F} \ar[d]^{f_s} \\
\Sh(\mathcal{D})/\mathcal{G}' \ar[r]^{j_{\mathcal{G}'/\mathcal{G}}} &
\Sh(\mathcal{D})/\mathcal{G}.
}
$$
도 가환한다.
\end{lemma}

\begin{proof}
첫 번째 정사각형의 가환성은 Lemma \ref{sites-lemma-relocalize-topos}의
도표의 가환성과 Lemma \ref{sites-lemma-localize-morphism-topoi}의 도표의
가환성에서 따른다. $f_s^{-1}$의 기술은 Lemma
\ref{sites-lemma-localize-morphism-topoi}의 $(f')^{-1}$에 대한 기술과 Lemma
\ref{sites-lemma-relocalize-topos}의
$(j_{\mathcal{F}/f^{-1}\mathcal{G}})^{-1}$에 대한 기술을 결합하면
얻어진다. 마지막 정사각형의 가환성은 형식적인 등식
$$
f^{-1}\mathcal{H} \times_{f^{-1}\mathcal{G}, s} \mathcal{F}
\times_\mathcal{F} \mathcal{F}'
=
f^{-1}(\mathcal{H} \times_\mathcal{G} \mathcal{G}')
\times_{f^{-1}\mathcal{G}', s'} \mathcal{F}'
$$
에서 따른다.
\end{proof}

\begin{lemma}
\label{sites-lemma-relocalize-morphism-compare}
$f : \mathcal{C} \to \mathcal{D}$를 연속 함자
$u : \mathcal{D} \to \mathcal{C}$가 주는 사이트의 사상이라 하자.
$V$를 $\mathcal{D}$의 대상이라 하고 $c : U \to u(V)$를 사상이라 하자.
% P01-ERRATA: P01-E0121 -- the frozen equality identifies h_U# with h_{u(V)}# for a general c:U->u(V).
$\mathcal{G} = h_V^\#$ 및
$\mathcal{F} = h_U^\# = f^{-1}h_V^\#$로 놓자.
$s : \mathcal{F} \to f^{-1}\mathcal{G}$를 $c$가 유도하는 사상이라 하자.
그러면 Lemma \ref{sites-lemma-relocalize-morphism}의 토포스의 사상들의 도표는
Lemma \ref{sites-lemma-relocalize-morphism-topoi}의 토포스의 사상들의 도표와
Lemma \ref{sites-lemma-localize-compare}의 식별
$j_\mathcal{F} = j_U$ 및 $j_\mathcal{G} = j_V$를 통하여 일치한다.
\end{lemma}

\begin{proof}
Lemma \ref{sites-lemma-relocalize-compare}와 Lemma
\ref{sites-lemma-localize-morphism-compare}를 결합하면 따른다.
\end{proof}

\section{점}
\label{sites-section-points}

\begin{definition}
\label{sites-definition-point-topos}
$\mathcal{C}$를 사이트라 하자.
{\it 토포스 $\Sh(\mathcal{C})$의 점}이란
$\Sh(pt)$에서 $\Sh(\mathcal{C})$로 가는 토포스의 사상 $p$이다.
\end{definition}

\noindent
사이트의 점을 함자 $u : \mathcal{C} \to \textit{Sets}$로 정의하겠다.
나중에 $u$가 사이트의 사상을 정의하고, 이것이 $\mathcal{C}$에 대응하는
토포스의 점을 유도함을 알게 될 것이다. Lemma
\ref{sites-lemma-site-point-morphism}을 보라.

\medskip\noindent
$\mathcal{C}$를 사이트라 하자. $p = u$를 함자
$u : \mathcal{C} \to \textit{Sets}$라 하자.
함자의 근방을 말하는 것은 우스워 보이므로 이 색다른 용어를 도입한다.
즉 “점” $p$를 범주론적 자료인 함자 $u$로 주어지는 기하학적 대상으로
생각한다. $p$가 실제로 $u$와 같다는 사실은 중요하지 않다.
$p$의 {\it 근방}이란 $U \in \Ob(\mathcal{C})$이고
$x \in u(U)$인 쌍 $(U, x)$이다.
{\it 근방의 사상} $(V, y) \to (U, x)$은
$u(\alpha)(y) = x$를 만족하는 $\mathcal{C}$의 사상
$\alpha :V \to U$로 주어진다. 근방들의 범주는 “큰” 범주가 아님에
유의하라.

\medskip\noindent
준층 $\mathcal{F}$의 $p$에서의 {\it 줄기}를 다음과 같이 정의한다.
\begin{equation}
\label{sites-equation-stalk}
\mathcal{F}_p = \colim_{\{(U, x)\}^{opp}} \mathcal{F}(U).
\end{equation}
쌍대극한은 $p$의 근방들의 범주의 반대범주 위에서 취한다. 달리 말해
$\mathcal{F}_p$의 원소는 삼중항 $(U, x, s)$로 주어지며, 여기서
$(U, x)$는 $p$의 근방이고 $s \in \mathcal{F}(U)$이다. 삼중항들의
동등성은 위와 같은 $\alpha$에 대한 관계
$(U, x, s) \sim (V, y, \alpha^*s)$가 생성하는 동치관계이다.

\medskip\noindent
$\varphi : \mathcal{F} \to \mathcal{G}$가 집합의 준층들의 사상이면 줄기들의
표준 사상 $\varphi_p : \mathcal{F}_p \to \mathcal{G}_p$를 얻는다.
따라서 {\it 줄기 함자}
$$
\textit{PSh}(\mathcal{C}) \longrightarrow \textit{Sets}, \quad
\mathcal{F} \longmapsto \mathcal{F}_p.
$$
를 얻는다. 줄기 함자는 임의의 함자
$p = u : \mathcal{C} \to \textit{Sets}$를 사용하여 정의했다.
이 정의가 성립하는 데에는 아무 조건도 필요하지
않다\footnote{모든 $U$에 대하여 $u(U) = \emptyset$인 경우는 피해야 한다.}.
한편 일반적으로 줄기 함자가 좋은 성질을 갖지 않으므로, $p$가 실제로
점인 경우(아래 정의를 보라)가 아니면 이 개념을 사용하지 않는 편이
아마도 낫다.

\begin{definition}
\label{sites-definition-point}
$\mathcal{C}$를 사이트라 하자. {\it 사이트 $\mathcal{C}$의 점 $p$}이란
다음을 만족하는 함자 $u : \mathcal{C} \to \textit{Sets}$로 주어지는
것이다.
\begin{enumerate}
\item $\mathcal{C}$의 모든 덮개 $\{U_i \to U\}$에 대하여 사상
$\coprod u(U_i) \to u(U)$는 전사이다.
\item $\mathcal{C}$의 모든 덮개 $\{U_i \to U\}$와 모든 사상
$V \to U$에 대하여 사상
$u(U_i \times_U V) \to u(U_i) \times_{u(U)} u(V)$는 전단사이다.
\item 줄기 함자 $\Sh(\mathcal{C}) \to \textit{Sets}$,
$\mathcal{F} \mapsto \mathcal{F}_p$는 좌완전이다.
\end{enumerate}
\end{definition}

\noindent
이 조건들은 Definitions \ref{sites-definition-continuous} 및
\ref{sites-definition-morphism-sites}의 조건들을 본떠 만든 것이므로 익숙할
것이다. (3)에 의해 $*_p = \{*\}$임에 유의하라. Example
\ref{sites-example-singleton-sheaf}를 보라. 따라서 적어도 어떤 $U$에 대해서는
$u(U) \not = \emptyset$이다(빈 쌍대극한은 공집합을 만들기 때문이다).
아래에서 이것이 실제로 토포스 $\Sh(\mathcal{C})$의 점을 유도함을 보일
것이다(Lemma \ref{sites-lemma-point-site-topos}). 그 전에 일반 함자 $u$에
대한 몇 가지 보조정리를 증명한다.

\begin{lemma}
\label{sites-lemma-points-recover}
$\mathcal{C}$를 사이트라 하자.
$p = u : \mathcal{C} \to \textit{Sets}$를 함자라 하자.
$U \in \Ob(\mathcal{C})$에 대하여 함자적 동형사상
$(h_U)_p = u(U)$가 존재한다.
\end{lemma}

\begin{proof}
$(h_U)_p$의 원소는 삼중항 $(V, y, f)$로 주어지며, 여기서
$V \in \Ob(\mathcal{C})$, $y\in u(V)$이고
$f \in h_U(V) = \Mor_\mathcal{C}(V, U)$이다. 이러한 두 삼중항
$(V, y, f)$, $(V', y', f')$에 대하여
$u(\phi)(y) = y'$ 및 $f' \circ \phi = f$를 만족하는 사상
$\phi : V \to V'$가 존재하면 같은 원소를 정하고, 일반적으로는 올바른
동치관계를 얻기 위해 이와 같은 항등식들의 사슬을 취해야 한다.
% P01-ERRATA: P01-E0122 -- the frozen proof says “same object” and “same triple” where it proves equality of stalk elements.
어쨌든 모든 $(V, y, f)$는 원소
$(U, u(f)(y), \text{id}_U)$와 동치이다. 위와 같은 $\phi$가 존재하면
삼중항 $(V, y, f)$, $(V', y', f')$는 같은 원소
$(U, u(f)(y), \text{id}_U) = (U, u(f')(y'), \text{id}_U)$를 정한다.
따라서 사상
$u(U) \to (h_U)_p$, $x \mapsto \text{class of }(U, x, \text{id}_U)$는
전단사이다.
\end{proof}

\noindent
$\mathcal{C}$를 사이트라 하고
$p = u : \mathcal{C} \to \textit{Sets}$를 함자라 하자.
Section \ref{sites-section-functoriality-PSh}의 구성과 유사하게, 집합 $E$가
주어지면 다음 규칙으로 준층 $u^pE$를 정의한다.
\begin{equation}
\label{sites-equation-skyscraper}
U
\longmapsto
u^pE(U) = \Mor_{\textit{Sets}}(u(U), E) = \text{Map}(u(U), E).
\end{equation}
이는 함자
$u^p : \textit{Sets} \to \textit{PSh}(\mathcal{C})$, $E \mapsto u^pE$를
정의한다.

\begin{lemma}
\label{sites-lemma-adjoint-point-push-stalk}
% P01-ERRATA: P01-E0123 -- the frozen statement separates its quantifier phrase from the predicate with a period.
임의의 함자 $u : \mathcal{C} \to \textit{Sets}$에 대하여 함자 $u^p$는
준층 위의 줄기 함자의 오른쪽 수반이다.
\end{lemma}

\begin{proof}
$\mathcal{F}$를 $\mathcal{C}$ 위의 준층이라 하고 $E$를 집합이라 하자.
사상 $\mathcal{F} \to u^pE$는 사상들의 양립하는 계
$\mathcal{F}(U) \to \text{Map}(u(U), E)$, 즉 사상들의 양립하는 계
$\mathcal{F}(U) \times u(U) \to E$로 주어진다. 그리고
$\mathcal{F}_p$의 정의에 의해 사상 $\mathcal{F}_p \to E$는 각 삼중항
$(U, x, \sigma)$에 동치인 삼중항들이 같은 원소로 가도록 $E$의 원소를
대응시키는 규칙으로 주어진다. Equation
(\ref{sites-equation-stalk}) 주위의 논의를 보라. 이것도 사상들의 양립하는 계
$\mathcal{F}(U) \times u(U) \to E$를 뜻한다.
\end{proof}

\noindent
Section \ref{sites-section-continuous-functors}와 유사하게 다음 보조정리를 얻는다.

\begin{lemma}
\label{sites-lemma-point-pushforward-sheaf}
$\mathcal{C}$를 사이트라 하고
$p = u : \mathcal{C} \to \textit{Sets}$를 함자라 하자.
$\mathcal{C}$의 모든 덮개 $\{U_i \to U\}$에 대하여 다음을 가정하자.
\begin{enumerate}
\item 사상 $\coprod u(U_i) \to u(U)$는 전사이다.
\item 사상
$u(U_i \times_U U_j) \to u(U_i) \times_{u(U)} u(U_j)$는 전사이다.
\end{enumerate}
그러면 다음이 성립한다.
\begin{enumerate}
\item 모든 집합 $E$에 대하여 준층 $u^pE$는 층이며, 이를 $u^sE$로
나타낸다.
\item 줄기 함자 $\Sh(\mathcal{C}) \to \textit{Sets}$와 함자
$u^s: \textit{Sets} \to \Sh(\mathcal{C})$는 수반이다.
\item 모든 집합의 준층 $\mathcal{F}$에 대하여
$\mathcal{F}_p = \mathcal{F}^\#_p$이다.
\end{enumerate}
\end{lemma}

\begin{proof}
첫 번째 주장은 층의 정의, 가정 (1)과 (2), 그리고 $u^pE$의 정의에서
곧바로 따른다. 두 번째 주장은 $u^p$와 줄기 함자의 수반성
(Lemma \ref{sites-lemma-adjoint-point-push-stalk})을 층들에 제한한 것을
다시 말한 것이다. 세 번째 주장은 임의의 집합 $E$에 대하여
$$
\text{Map}(\mathcal{F}_p, E) =
\Mor_{\textit{PSh}(\mathcal{C})}(\mathcal{F}, u^pE) =
\Mor_{\Sh(\mathcal{C})}(\mathcal{F}^\#, u^sE) =
\text{Map}(\mathcal{F}^\#_p, E)
$$
가 층화의 수반 성질에 의해 성립하는 것에서 따른다.
\end{proof}

\noindent
특히 $p = u$가 점이면 Lemma \ref{sites-lemma-point-pushforward-sheaf}가
성립한다. 이 경우 층 $u^sE$를 $p$에서 값 $E$를 갖는 “마천루” 층으로
생각한다.

\begin{definition}
\label{sites-definition-pushforward-point}
% P01-ERRATA: P01-E0124 -- the frozen definition runs together the definiendum and its appositive without “to be” or punctuation.
$p$를 함자 $u$로 주어지는 사이트 $\mathcal{C}$의 점이라 하자.
집합 $E$에 대하여 $p_*E = u^sE$를 위의 Lemma
\ref{sites-lemma-point-pushforward-sheaf}에서 기술한 층으로 정의한다.
이를 때때로 {\it 마천루 층}이라 부른다.
\end{definition}

\noindent
특히 마천루 층과 줄기 사이에는 다음 수반 성질이 있다.
$$
\Mor_{\Sh(\mathcal{C})}(\mathcal{F}, p_*E)
=
\text{Map}(\mathcal{F}_p, E)
$$
이것은 때때로 사용할 표기
$p^{-1}\mathcal{F} = \mathcal{F}_p$의 동기가 된다.

\begin{lemma}
\label{sites-lemma-point-site-topos}
$\mathcal{C}$를 사이트라 하자.
\begin{enumerate}
\item $p$를 사이트 $\mathcal{C}$의 점이라 하자.
그러면 위에서 도입한 함자들의 쌍 $(p_*, p^{-1})$은 토포스의 사상
$\Sh(pt) \to \Sh(\mathcal{C})$를 정의한다.
\item $p = (p_*, p^{-1})$를 토포스 $\Sh(\mathcal{C})$의 점이라 하자.
그러면 함자 $u : U \mapsto p^{-1}(h_U^\#)$는 사이트
$\mathcal{C}$의 점 $p'$를 유도하고, 이에 대응하는 토포스의 사상
$(p'_*, (p')^{-1})$은 $p$와 같다.
\end{enumerate}
\end{lemma}

\begin{proof}
(1)의 증명. 위에서 보았듯이 함자 $p_*$와 $p^{-1}$은 수반이다.
함자 $p^{-1}$은 Definition \ref{sites-definition-point}에 의해 완전이어야
한다. 따라서 Definition \ref{sites-definition-topos}에서 부과한 조건들이
모두 만족되므로 (1)이 성립한다.

\medskip\noindent
(2)의 증명. $\{U_i \to U\}$를 $\mathcal{C}$의 덮개라 하자.
그러면 Lemma \ref{sites-lemma-covering-surjective-after-sheafification}에
의해 $\coprod (h_{U_i})^\# \to h_U^\#$는 전사이다.
$p^{-1}$은 토포스의 사상의 정의에 의해 완전이므로
$\coprod u(U_i) \to u(U)$가 전사임을 얻는다.
이로써 Definition \ref{sites-definition-point}의 (1)을 증명했다.
층화는 완전이다. Lemma \ref{sites-lemma-sheafification-exact}를 보라.
따라서 $\mathcal{C}$에서 $U \times_V W$가 존재하면
$$
h_{U \times_V W}^\# =  h_U^\# \times_{h_V^\#} h_W^\#
$$
이고, $p^{-1}$이 완전이므로
$u(U \times_V W) = u(U) \times_{u(V)} u(W)$임을 알 수 있다.
이로써 Definition \ref{sites-definition-point}의 (2)를 증명했다.
$p' = u$라 하고, $u$를 사용하여 Equation
(\ref{sites-equation-stalk})으로 정의한 줄기 함자를
$\mathcal{F}_{p'}$라 하자. 표준 비교 사상
$c : \mathcal{F}_{p'} \to \mathcal{F}_p = p^{-1}\mathcal{F}$가
존재한다. 실제로 $\mathcal{F}_{p'}$의 원소 $\xi$를 나타내는 삼중항
$(U, x, \sigma)$가 주어졌다고 하자. $\sigma$를 사상
$\sigma : h_U^\# \to \mathcal{F}$로 생각하고,
$x \in u(U) = p^{-1}(h_U^\#)$이므로
$c(\xi) = p^{-1}(\sigma)(x)$로 놓을 수 있다. Lemma
\ref{sites-lemma-points-recover}에 의해 $(h_U)_{p'} = u(U)$이다.
Definition \ref{sites-definition-point}의 조건 (1)과 (2)가 $p'$에 대하여
성립하므로 Lemma \ref{sites-lemma-point-pushforward-sheaf}에 의해
$(h_U^\#)_{p'} = (h_U)_{p'}$도 성립한다. 따라서
$$
(h_U^\#)_{p'} = (h_U)_{p'} = u(U) = p^{-1}(h_U^\#)
$$
이다. 이 전단사가 비교 사상
$c : (h_U^\#)_{p'} \to p^{-1}(h_U^\#)$와 같다고 주장한다
(확인은 생략한다). $\mathcal{C}$ 위의 임의의 층은
Lemma \ref{sites-lemma-sheaf-coequalizer-representable}에 의해
$h_U^\#$ 꼴의 층들의 쌍대합들 사이의 사상들의 쌍대등화자이다.
줄기 함자 $\mathcal{F} \mapsto \mathcal{F}_{p'}$와 함자 $p^{-1}$은
둘 다 왼쪽 수반이므로 임의의 쌍대극한과 가환한다.
따라서 모든 층 $\mathcal{F}$에 대하여 $c$는 동형사상이다.
그러므로 줄기 함자 $\mathcal{F} \mapsto \mathcal{F}_{p'}$는
$p^{-1}$과 동형이고, 특히 완전임을 알 수 있다.
이로써 Definition \ref{sites-definition-point}의 조건 (3)이 성립하고
$p'$가 점임을 증명했다. 위에서 $p^{-1} = (p')^{-1}$임을 보았으므로
마지막 주장도 이미 증명되었다.
\end{proof}

\noindent
실제로 점은 언제나 사이트의 사상에 대응한다. 이를 다음 보조정리에서
보인다.

\begin{lemma}
\label{sites-lemma-site-point-morphism}
$\mathcal{C}$를 사이트라 하고, $p$를
$u : \mathcal{C} \to \textit{Sets}$로 주어지는 $\mathcal{C}$의
점이라 하자. 모든 $U \in \Ob(\mathcal{C})$에 대하여
$u(U) \subset S_0$인 무한집합 $S_0$를 택하자. $\mathcal{S}$를
Remark \ref{sites-remark-pt-topos}에서 멱집합 $S = \mathcal{P}(S_0)$로부터
구성한 사이트라 하자. 그러면
\begin{enumerate}
\item 동치 $i : \Sh(pt) \to \Sh(\mathcal{S})$가 존재하고,
\item 함자 $u : \mathcal{C} \to \mathcal{S}$는 사이트의 사상
$f : \mathcal{S} \to \mathcal{C}$를 유도하며,
\item 합성
$$
\Sh(pt) \to
\Sh(\mathcal{S}) \to
\Sh(\mathcal{C})
$$
은 Lemma \ref{sites-lemma-point-site-topos}의 토포스의 사상
$(p_*, p^{-1})$이다.
\end{enumerate}
\end{lemma}

\begin{proof}
(1)은 Remark \ref{sites-remark-pt-topos}에서 보았다.
또한 이 동치는 집합 $E$에 $\mathcal{S}$ 위의 층 $i_*E$를 대응시키며,
이 층은 규칙 $V \mapsto \Mor_{\textit{Sets}}(V, E)$로 정의됨을
상기하라. (2)는 $\mathcal{C}$의 점의 정의
(Definition \ref{sites-definition-point})와 사이트의 사상의 정의
(Definition \ref{sites-definition-morphism-sites})에서 명백하다.
마지막으로 $f_*i_*E$를 생각하자. 구성에 의해
$$
f_*i_*E(U) = i_*E(u(U)) = \Mor_{\textit{Sets}}(u(U), E)
$$
이고, 이는 Equation (\ref{sites-equation-skyscraper})에 의해
$p_*E(U)$와 같다. 이로써 (3)을 증명했다.
\end{proof}

\noindent
위상적인 경우와 달리 값 $E$를 갖는 마천루 층의 줄기가 언제나 $E$인
것은 아니다. 일반적으로 참인 것은 다음과 같다.

\begin{lemma}
\label{sites-lemma-stalk-skyscraper}
$\mathcal{C}$를 사이트라 하자.
$p : \Sh(pt) \to \Sh(\mathcal{C})$를 $\mathcal{C}$에 대응하는
토포스의 점이라 하자. 임의의 집합 $E$에 대하여 합성이
$\text{id}_E$인 표준 사상들
$$
E \longrightarrow (p_*E)_p \longrightarrow E
$$
가 존재한다.
\end{lemma}

\begin{proof}
언제나 수반 사상 $(p_*E)_p = p^{-1}p_*E \to E$가 존재한다.
$p_*$와 $p^{-1}$은 둘 다 좌완전이어서 종대상을 종대상으로 보내므로,
$E = \{*\}$일 때 이 사상은 동형사상이다. 따라서 $e \in E$가 주어지면
사상 $i_e : \{*\} \to E$를 생각하여
$$
\xymatrix{
p^{-1}p_*\{*\} \ar[rr]_{p^{-1}p_*i_e} \ar[d]_{\cong} & & p^{-1}p_*E \ar[d] \\
\{*\} \ar[rr]^{i_e} & & E
}
$$
를 얻고, 이로부터 원하는 사상
$E \to (p_*E)_p = p^{-1}p_*E$를 얻는다.
\end{proof}

\begin{lemma}
\label{sites-lemma-skyscraper-functor-exact}
$\mathcal{C}$를 사이트라 하고,
$p : \Sh(pt) \to \Sh(\mathcal{C})$를 $\mathcal{C}$에 대응하는
토포스의 점이라 하자. 함자 $p_* : \textit{Sets} \to \Sh(\mathcal{C})$는
다음 성질들을 갖는다. 임의의 극한과 가환하고, 좌완전이며, 충실하고,
전사를 전사로 보내며, 쌍대등화자와 가환하고, 단사성, 전사성 및
동형사상을 반영한다.
\end{lemma}

\begin{proof}
$p_*$는 오른쪽 수반이므로 임의의 극한과 가환하고 좌완전이다.
$p^{-1}p_*E \to E$가 표준적으로 분할되는 전사라는 사실로부터
$p_*$가 충실하고 단사성, 전사성 및 동형사상을 반영함이 따른다.
Lemma \ref{sites-lemma-point-site-topos}에 의해 $p$가 바탕 사이트
$\mathcal{C}$의 점 $u : \mathcal{C} \to \textit{Sets}$에서 온다고
가정해도 된다. 이 경우 층 $p_*E$는
$$
p_*E(U) = \Mor_{\textit{Sets}}(u(U), E)
$$
로 주어진다. Equation (\ref{sites-equation-skyscraper}) 및
Definition \ref{sites-definition-pushforward-point}를 보라.
이 공식에서 $p_*$가 전사를 전사로 보내고 쌍대등화자를 쌍대등화자로
보냄이 곧바로 따른다.
\end{proof}

\section{점의 구성}
\label{sites-section-construct-points}

\noindent
이 절에서는 사이트에서 집합의 범주로 가는 함자가 언제 그 사이트의 점을
정의하는지에 대한 판정법을 제시한다.

\begin{lemma}
\label{sites-lemma-neighbourhoods-cofiltered}
$\mathcal{C}$를 사이트라 하고
$p = u : \mathcal{C} \to \textit{Sets}$를 함자라 하자.
$p$의 근방들의 범주가 쌍대여과이면 줄기 함자
(\ref{sites-equation-stalk})는 좌완전이다.
\end{lemma}

\begin{proof}
$\mathcal{I} \to \Sh(\mathcal{C})$,
$i \mapsto \mathcal{F}_i$를 층들의 유한 도표라 하자.
이 계의 극한의 줄기가 줄기들의 극한과 일치함을 보여야 한다.
$\mathcal{F}$를 이 계의 {\it 준층}으로서의 극한이라 하자.
Lemma \ref{sites-lemma-limit-sheaf}에 따르면 이는 층이며 층의 범주에서의
극한이다. 따라서
$\mathcal{F}_p = \lim_\mathcal{I} \mathcal{F}_{i, p}$임을 보이면 된다.
또한 $\mathcal{F}$에는 간단한 기술이 있음을 상기하라. Section
\ref{sites-section-limits-colimits-PSh}를 보라. 그러므로 다음을 보이면 된다.
$$
\lim_i \colim_{\{(U, x)\}^{opp}} \mathcal{F}_i(U)
=
\colim_{\{(U, x)\}^{opp}} \lim_i \mathcal{F}_i(U).
$$
가정에 의해 근방들의 범주의 반대범주가 여과이므로, 이는 Categories,
Lemma \ref{categories-lemma-directed-commutes}에 의해 성립한다.
\end{proof}

\begin{lemma}
\label{sites-lemma-neighbourhoods-directed}
$\mathcal{C}$를 사이트라 하자. $\mathcal{C}$가 종대상 $X$와
섬유곱들을 갖는다고 가정하자.
$p = u : \mathcal{C} \to \textit{Sets}$가 다음을 만족하는 함자라고 하자.
\begin{enumerate}
\item $u(X)$는 한원소집합이다.
\item 공통의 공역을 갖는 모든 두 사상 $U \to W$와 $V \to W$에 대하여
사상 $u(U \times_W V) \to u(U) \times_{u(W)} u(V)$는 전단사이다.
\end{enumerate}
% P01-ERRATA: P01-E0125 -- the frozen source duplicates the determiner in “Then the the category”.
그러면 $p$의 근방들의 범주는 쌍대여과이고, 따라서 줄기 함자
$\Sh(\mathcal{C}) \to \textit{Sets}$,
$\mathcal{F} \to \mathcal{F}_p$는 유한 극한과 가환한다.
\end{lemma}

\begin{proof}
Lemma \ref{sites-lemma-directed}와의 유사성에 유의하라.
$\mathcal{C}$에 대한 가정은 $\mathcal{C}$가 유한 극한을 가짐을 함의한다.
Categories, Lemma \ref{categories-lemma-finite-limits-exist}를 보라.
가정 (1)에 의해 근방들의 범주는 비어 있지 않다.
$(U, x)$와 $(V, y)$가 근방이라고 하자. (2)에 의해
$u(U \times V) = u(U \times_X V) =
u(U) \times_{u(X)} u(V) = u(U) \times u(V)$이다.
따라서 $(U, x)$와 $(V, y)$ 둘 다로 가는 근방
$(U \times_X V, z)$가 존재한다.
$a, b : (V, y) \to (U, x)$를 근방들의 범주에서의 두 사상이라 하자.
$W$를 $a, b : V \to U$의 등화자라 하자. Categories, Lemma
\ref{categories-lemma-finite-limits-exist}의 증명에서와 같이 $W$를
섬유곱들로 다음처럼 쓸 수 있다.
$$
W = (V \times_{a, U, b} V) \times_{(pr_1, pr_2), V \times V, \Delta} V
$$
(2)의 전단사성은 $((y, y), y)$로 가는 원소 $z \in u(W)$가 존재함을
보장한다. 그러면 $(W, z) \to (V, y)$는 원하는 대로 $a, b$를
등화한다. “따라서”라는 절은 Lemma
\ref{sites-lemma-neighbourhoods-cofiltered}이다.
\end{proof}

\begin{proposition}
\label{sites-proposition-point-limits}
$\mathcal{C}$를 사이트라 하자. $\mathcal{C}$에 유한 극한이 존재한다고
가정하자(즉, $\mathcal{C}$가 섬유곱들과 종대상을 갖는다고 하자).
이러한 사이트 $\mathcal{C}$의 점 $p$는 다음을 만족하는 함자
$u : \mathcal{C} \to \textit{Sets}$로 주어진다.
\begin{enumerate}
\item $u$는 유한 극한과 가환한다.
\item $\{U_i \to U\}$가 덮개이면
$\coprod_i u(U_i) \to u(U)$는 전사이다.
\end{enumerate}
\end{proposition}

\begin{proof}
먼저 $p$가 함자 $u$로 주어지는 점이라고 하자
(Definition \ref{sites-definition-point}). 조건 (2)는 점의 정의에서
직접 성립한다. Lemma \ref{sites-lemma-points-recover}에 의해
$(h_U)_p = u(U)$이고, Lemma \ref{sites-lemma-point-pushforward-sheaf}에
의해 $(h_U^\#)_p = (h_U)_p$이다. 따라서 $u$는 함자들의 합성
$$
\mathcal{C} \xrightarrow{h}
\textit{PSh}(\mathcal{C}) \xrightarrow{{}^\#}
\Sh(\mathcal{C}) \xrightarrow{()_p}
\textit{Sets}
$$
과 같다. 이 함자들은 각각 좌완전이므로 $u$가 (1)을 만족함을 알 수 있다.

\medskip\noindent
역으로 $u$가 (1)과 (2)를 만족한다고 하자. 이 경우 $u$가 Definition
\ref{sites-definition-point}의 처음 두 조건을 만족함은 곧바로 알 수 있다.
그리고 그 줄기 함자는 Lemma \ref{sites-lemma-point-pushforward-sheaf}에
의해 왼쪽 수반이고 Lemma \ref{sites-lemma-neighbourhoods-directed}에 의해
유한 극한과 가환하므로 완전이다.
\end{proof}

\begin{remark}
\label{sites-remark-improve-proposition-points-limits}
실제로 $\mathcal{C}$를 사이트라 하자. $\mathcal{C}$가 종대상
$X$와 섬유곱들을 갖는다고 가정하자.
$p = u: \mathcal{C} \to \textit{Sets}$가 다음을 만족하는 함자라고 하자.
\begin{enumerate}
% P01-ERRATA: P01-E0126 -- the frozen item omits the copula between the equality and “a singleton”.
\item $u(X) = \{*\}$는 한원소집합이다.
\item 공통의 공역을 갖는 모든 두 사상 $U \to W$와 $V \to W$에 대하여
사상 $u(U \times_W V) \to u(U) \times_{u(W)} u(V)$는 전사이다.
\item 모든 덮개 $\{U_i \to U\}$에 대하여 사상
$\coprod u(U_i) \to u(U)$는 전사이다.
\end{enumerate}
그러나 일반적으로 $p$는 $\mathcal{C}$의 점이 {\bf 아니다}.
두 대상 $\{U, X\}$와 항등사상이 아닌 사상 하나 $U \to X$만을 갖는
범주 $\mathcal{C}$가 한 예이다. $\mathcal{C}$에는 자명한 위상, 즉
덮개가 $\{U \to U\}$와 $\{X \to X\}$뿐인 위상을 준다.
층 $\mathcal{F}$는 준층과 같은 것이며 삼중항 $(A, B, A \to B)$로
이루어진다. 즉 $A = \mathcal{F}(X)$, $B = \mathcal{F}(U)$이고
$A \to B$는 $U \to X$에 대응하는 제한 사상이다.
$U \times_X U = U$이므로 섬유곱들이 존재함에 유의하라.
$u(X) = \{*\}$이고 $u(U) = \{*_1, *_2\}$인 함자 $u = p$를 생각하자.
이는 (1), (2), (3)을 만족하지만 대응하는 줄기 함자
(\ref{sites-equation-stalk})는 함자
$$
(A, B, A \to B) \longmapsto B \amalg_A B
$$
이며 완전이 아니다. 실제로 층들의 단사 사상
$(\emptyset, \{1\}, \emptyset \to \{1\}) \to (\{1\}, \{1\}, \{1\} \to \{1\})$
을 생각하자. 줄기 함자는 이를 집합들의 단사가 아닌 사상
$$
\{1\} \amalg \{1\} \longrightarrow \{1\} \amalg_{\{1\}} \{1\}
$$
으로 보낸다.
\end{remark}

\begin{example}
\label{sites-example-point-topological}
$X$를 위상공간이라 하고, $X_{Zar}$를 Example
\ref{sites-example-site-topological}의 사이트라 하자.
$x \in X$를 점이라 하자. 함자
$$
u : X_{Zar} \longrightarrow \textit{Sets}, \quad
U \mapsto
\left\{
\begin{matrix}
\emptyset & \text{if} & x \not \in U \\
\{*\} & \text{if} & x \in U
\end{matrix}
\right.
$$
를 생각하자. 이 함자는 곱과 섬유곱과 가환하고 덮개들을 전사인
사상족으로 보낸다. 따라서 사이트 $X_{Zar}$의 점 $p$를 얻는다.
줄기 함자가 Sheaves, Section \ref{sheaves-section-stalks}에서 정의한
$x$에서의 줄기와 일치함은 곧바로 확인된다.
\end{example}

\begin{example}
\label{sites-example-point-topology}
$X$를 위상공간이라 하자. 토포스 $\Sh(X)$의 점들은 무엇인가?
이를 알아보기 위해 $X_{Zar}$를 Example
\ref{sites-example-site-topological}의 사이트라 하자.
Lemma \ref{sites-lemma-point-site-topos}에 의해 $\Sh(X)$의 점은 이 사이트의
점에 대응한다. $p$를 함자
$u : X_{Zar} \to \textit{Sets}$로 주어지는 사이트 $X_{Zar}$의
점이라 하자. Proposition \ref{sites-proposition-point-limits}에 주어진
이러한 $u$의 특징짓기를 사용하겠다. 이로부터 곧바로
$u(\emptyset) = \emptyset$ 및
$u(U \cap V) = u(U) \times u(V)$를 얻는다. 특히 대각사상을 통하여
$u(U) = u(U) \times u(U)$이고, 이는 $u(U)$가 한원소집합이거나
공집합임을 함의한다. 또한 $U = \bigcup U_i$가 열린 덮개이면
$$
u(U) = \emptyset \Rightarrow \forall i, \ u(U_i) = \emptyset
\quad\text{and}\quad
u(U) \not = \emptyset \Rightarrow \exists i, \ u(U_i) \not = \emptyset.
$$
따라서 $u(W) = \emptyset$인 유일한 최대 열린집합 $W \subset X$가
존재하며, 이는 $u(U) = \emptyset$인 모든 열린집합 $U$의 합집합이다.
$Z = X \setminus W$로 놓자. $Z = Z_1 \cup Z_2$이고
$Z_i \subset Z$가 닫힌집합이면
$W = (X \setminus Z_1) \cap (X \setminus Z_2)$이므로
$\emptyset = u(W) = u(X \setminus Z_1) \times u(X \setminus Z_2)$이고,
$u(X \setminus Z_1) = \emptyset$이거나
$u(X \setminus Z_2) = \emptyset$이라고 결론내린다.
이는 $X \setminus Z_1 = W$이거나 $X \setminus Z_2 = W$임을 뜻한다.
달리 말해 $Z$는 기약이다. 이제 $u$가 다음 규칙으로 기술됨을 알 수 있다.
$$
u : X_{Zar} \longrightarrow \textit{Sets}, \quad
U \mapsto
\left\{
\begin{matrix}
\emptyset & \text{if} & Z \cap U = \emptyset \\
\{*\} & \text{if} & Z \cap U \not = \emptyset
\end{matrix}
\right.
$$
임의의 기약 닫힌집합 $Z \subset X$에 대하여 이 함자는 Proposition
\ref{sites-proposition-point-limits}의 가정 (1), (2)를 만족하므로 점을
정의함에 유의하라. 달리 말해 사이트 $X_{Zar}$의 점들은 $X$의 기약
닫힌 부분집합들과 일대일 대응한다. 특히 $X$가 소버 위상공간이면
$X_{Zar}$의 점들과 $X$의 점들이 일대일 대응한다. Example
\ref{sites-example-point-topological}을 보라.
\end{example}

\begin{example}
\label{sites-example-point-G-sets}
Example \ref{sites-example-site-on-group} 및 Section
\ref{sites-section-example-sheaf-G-sets}에서 기술한 사이트
$\mathcal{T}_G$를 생각하자. 망각 함자
$u : \mathcal{T}_G \to \textit{Sets}$는 곱 및 섬유곱과 가환하고
덮개들을 전사인 사상족으로 보낸다. 따라서 이는 $\mathcal{T}_G$의
점을 정의한다. $\Sh(\mathcal{T}_G)$와 $G\textit{-Sets}$를 식별한다.
줄기 함자
$$
p^{-1} :
\Sh(\mathcal{T}_G) = G\textit{-Sets}
\longrightarrow
\textit{Sets}
$$
는 망각 함자이다. 직접상 $p_*$는 집합 $S$를
$G$-집합 $\text{Map}(G, S)$로 보내는 함자
$$
\textit{Sets}
\longrightarrow
\Sh(\mathcal{T}_G) = G\textit{-Sets}
$$
이다. 여기서 작용은 $g \cdot \psi = \psi \circ R_g$이고
$R_g$는 오른쪽 곱셈이다. 특히 집합으로서
$p^{-1}p_*S = \text{Map}(G, S)$이고, Lemma
\ref{sites-lemma-stalk-skyscraper}의 사상들
$S \to \text{Map}(G, S) \to S$는 명백한 사상들이다.
\end{example}

\begin{example}
\label{sites-example-indiscrete-points}
$\mathcal{C}$를 카오스 위상을 준 범주라 하자
(Example \ref{sites-example-indiscrete}). $\mathcal{C}$의 모든 대상 $U_0$에
대하여 함자 $u : U \mapsto \Mor_\mathcal{C}(U_0, U)$는
$\mathcal{C}$의 점 $p$를 정의한다. 실제로 유일한 덮개들은 항등사상들로
주어지므로 Definition \ref{sites-definition-point}의 조건 (1)과 (2)는
곧바로 성립한다.
% P01-ERRATA: P01-E0127 -- the frozen source repeats condition (2), although exactness establishes condition (3); the canonical Korean rendering preserves the frozen number.
조건 (2)는 $\mathcal{F}_p = \mathcal{F}(U_0)$이고
% P01-ERRATA: P01-E0128 -- the frozen source says “discrete” after specifying the chaotic topology; the canonical Korean rendering preserves the frozen word.
위상이 이산 위상이므로 $U_0$ 위에서 단면을 취하는 것이 완전함자라는
사실에서 성립한다.
\end{example}

\section{점과 토포스의 사상}
\label{sites-section-functorial-points}

\noindent
이 절에서는 점과 토포스의 사상에 관한 몇 가지 설명을 한다.

\begin{lemma}
\label{sites-lemma-point-functor}
$u : \mathcal{C} \to \mathcal{D}$를 함자라 하자.
$v : \mathcal{D} \to \textit{Sets}$를 함자라 하고
$w = v \circ u$로 놓자. $v$와 $w$에 각각 대응하는 줄기 함자
(\ref{sites-equation-stalk})를 각각 $q$와 $p$로 나타내자.
그러면 $\mathcal{C}$ 위의 준층 $\mathcal{F}$에 함자적으로
$(u_p\mathcal{F})_q = \mathcal{F}_p$이다.
\end{lemma}

\begin{proof}
이는 간단한 범주론적 사실이다. 다음이 성립한다.
\begin{align*}
(u_p\mathcal{F})_q
& =
\colim_{(V, y)} \colim_{U, \phi : V \to u(U)} \mathcal{F}(U) \\
& = \colim_{(V, y, U, \phi : V \to u(U))} \mathcal{F}(U) \\
& = \colim_{(U, x)} \mathcal{F}(U) \\
& = \mathcal{F}_p
\end{align*}
첫 번째 등식은 $u_p$의 정의와 줄기 함자의 정의에서 성립한다.
$y \in v(V)$임을 관찰하라. 두 번째 등식에서는 단순히 쌍대극한들을
결합했다. 세 번째 등식을 보기 위해 규칙
$$
F((V, y, U, \phi : V \to u(U))) = (U, v(\phi)(y)).
$$
으로 정의한 도표 범주들의 함자 $F$에 Categories, Lemma
\ref{categories-lemma-colimit-constant-connected-fibers}를 적용한다.
$w(U) = v(u(U))$이므로 이 정의는 뜻이 있다. Categories, Lemma
\ref{categories-lemma-colimit-constant-connected-fibers}의 가정을
확인하자. $F$는 오른쪽 역
$(U, x) \mapsto (u(U), x, U, \text{id} : u(U) \to u(U))$를 가짐에
유의하라. $x \in w(U) = v(u(U))$이므로 이 또한 뜻이 있다.
한편 $(U, x)$ 위의 섬유범주에는 언제나 사상
$$
(V, y, U, \phi : V \to u(U))
\longrightarrow
(u(U), v(\phi)(y), U, \text{id} : u(U) \to u(U))
$$
가 존재하므로 섬유범주들은 연결이다. 네 번째이자 마지막 등식은
명백하다.
\end{proof}

\begin{lemma}
\label{sites-lemma-point-morphism-sites}
\begin{slogan}
사이트의 사상은 점들 사이의 사상을 정의하며, 당김은 줄기를 보존한다.
\end{slogan}
$f : \mathcal{D} \to \mathcal{C}$를 연속 함자
$u : \mathcal{C} \to \mathcal{D}$로 주어지는 사이트의 사상이라 하자.
$q$를 함자 $v : \mathcal{D} \to \textit{Sets}$로 주어지는
$\mathcal{D}$의 점이라 하자. Definition \ref{sites-definition-point}를 보라.
그러면 함자 $v \circ u : \mathcal{C} \to \textit{Sets}$는
$\mathcal{C}$의 점 $p$를 정의하고, 더 나아가 $\mathcal{C}$ 위의
임의의 층 $\mathcal{F}$에 대하여 표준 식별
$$
(f^{-1}\mathcal{F})_q = \mathcal{F}_p
$$
이 존재한다.
\end{lemma}

\begin{proof}[Lemma \ref{sites-lemma-point-morphism-sites}의 첫 번째 증명]
$u$가 연속이고 $v$가 점을 정의하므로 $v \circ u$가 Definition
\ref{sites-definition-point}의 조건 (1)과 (2)를 만족함은 즉시 알 수 있다.
표시된 등식을 증명하자. $\mathcal{F}$를 $\mathcal{C}$ 위의 층이라 하자.
그러면
$$
(f^{-1}\mathcal{F})_q = (u_s\mathcal{F})_q =
(u_p \mathcal{F})_q = \mathcal{F}_p
$$
이다. 첫 번째 등식은 $f^{-1} = u_s$이므로 성립하고, 두 번째 등식은
Lemma \ref{sites-lemma-point-pushforward-sheaf}에 의해, 세 번째 등식은
Lemma \ref{sites-lemma-point-functor}에 의해 성립한다. 이제 $p$가
Definition \ref{sites-definition-point}의 조건 (3)도 만족함을 알 수 있는데,
이는 완전함자들의 합성이기 때문이다. 이로써 증명이 끝난다.
\end{proof}

\begin{proof}[Lemma \ref{sites-lemma-point-morphism-sites}의 두 번째 증명]
Lemma \ref{sites-lemma-site-point-morphism}에 의해 $(q_*, q^{-1})$을
$$
\Sh(pt) \xrightarrow{i} \Sh(\mathcal{S})
\xrightarrow{h} \Sh(\mathcal{D})
$$
로 분해할 수 있다. 여기서 두 번째 토포스의 사상은 함자
$v : \mathcal{D} \to \mathcal{S}$가 유도하는 사이트의 사상
$h : \mathcal{S} \to \mathcal{D}$에서 온다
($\mathcal{S} \subset \textit{Sets}$가 $v$의 상에 있는 모든 대상을
포함하는 충만 부분범주이므로 이는 뜻이 있다). Lemma
\ref{sites-lemma-composition-morphisms-sites}에 의해 합성
$v \circ u : \mathcal{C} \to \mathcal{S}$는 사이트의 사상
$g : \mathcal{S} \to \mathcal{C}$를 정의한다. 특히 함자
$v \circ u : \mathcal{C} \to \mathcal{S}$는 연속이고,
Remark \ref{sites-remark-pt-topos}에 있는 $\mathcal{S}$의 덮개들의 정의에
따르면 이는 $v \circ u$가 Definition \ref{sites-definition-point}의
조건 (1)과 (2)를 만족함을 뜻한다. 한편 Remark
\ref{sites-remark-pt-topos}에서의 $i$의 구성에 의해
% P01-ERRATA: P01-E0129 -- the frozen display is missing one closing parenthesis in i_*E(v(u(U))).
$$
g_*i_*E(U) = i_*E(v(u(U)) = \Mor_{\textit{Sets}}(v(u(U)), E)
$$
이다.
% P01-ERRATA: P01-E0130 -- the frozen prose has an objectless “formula for” followed by a malformed relative clause.
이것은 그 공식과 같고, 이는 Equation
(\ref{sites-equation-skyscraper})에 의해 $(v \circ u)^pE$와 같다.
Lemma \ref{sites-lemma-point-pushforward-sheaf}에 의해 함자
$g_*i_* = (v \circ u)^p = (v \circ u)^s$는 줄기 함자
% P01-ERRATA: P01-E0131 -- the frozen inline formula uses the q-stalk although the adjunction here constructs the p-stalk; it is preserved.
$\mathcal{F} \mapsto \mathcal{F}_q$의 오른쪽 수반이다.
따라서 줄기 함자 $p^{-1}$이 $i^{-1} \circ g^{-1}$과 표준적으로
동형임을 알 수 있다. 그러므로 이는 완전이고 $p$가 점이라고 결론내린다.
마지막으로 구성에 의해 $g = f \circ h$이므로
$p^{-1} = i^{-1} \circ h^{-1} \circ f^{-1} = q^{-1} \circ f^{-1}$임을
알 수 있다. 즉 보조정리의 표시된 공식을 얻는다.
\end{proof}

\begin{lemma}
\label{sites-lemma-point-morphism-topoi}
$f : \Sh(\mathcal{D}) \to \Sh(\mathcal{C})$를 토포스의 사상이라 하고
$q : \Sh(pt) \to \Sh(\mathcal{D})$를 점이라 하자.
그러면 $p = f \circ q$는 토포스 $\Sh(\mathcal{C})$의 점이고,
$\mathcal{C}$ 위의 임의의 층 $\mathcal{F}$에 대하여 표준 식별
$$
(f^{-1}\mathcal{F})_q = \mathcal{F}_p
$$
이 존재한다.
\end{lemma}

\begin{proof}
이는 정의와 토포스의 사상들을 합성할 수 있다는 사실에서 곧바로 따른다.
\end{proof}

\section{국소화와 점}
\label{sites-section-localize-points}

\noindent
이 절에서는 국소화 $\mathcal{C}/U$의 점들이 $\mathcal{C}$의 점들로부터
간단히 구성됨을 보인다.

\begin{lemma}
\label{sites-lemma-point-localize}
$\mathcal{C}$를 사이트라 하고, $p$를 함자
$u : \mathcal{C} \to \textit{Sets}$로 주어지는 $\mathcal{C}$의 점이라
하자. $U$를 $\mathcal{C}$의 대상이라 하고 $x \in u(U)$라 하자. 함자
$$
v : \mathcal{C}/U \longrightarrow \textit{Sets}, \quad
(\varphi : V \to U) \longmapsto \{y \in u(V) \mid u(\varphi)(y) = x\}
$$
는 다음 도표가 가환하도록 사이트 $\mathcal{C}/U$의 점 $q$를 정의한다.
$$
\xymatrix{
& \Sh(pt) \ar[d]^p \ar[ld]_q \\
\Sh(\mathcal{C}/U) \ar[r]^{j_U} &
\Sh(\mathcal{C})
}
$$
달리 말해 $\mathcal{C}$ 위의 임의의 층에 대하여
% P01-ERRATA: P01-E0132 -- the frozen source omits the quantified sheaf symbol after “any sheaf”.
$\mathcal{F}_p = (j_U^{-1}\mathcal{F})_q$이다.
\end{lemma}

\begin{proof}
Lemma \ref{sites-lemma-site-point-morphism}에서와 같이 $S$와 $\mathcal{S}$를
선택하자. 그 보조정리에서처럼 $\Sh(pt) = \Sh(\mathcal{S})$로 식별하고,
$p = f : \Sh(\mathcal{S}) \to \Sh(\mathcal{C})$를 $u$가 유도하는
토포스의 사상으로 쓸 수 있다. Lemma
\ref{sites-lemma-localize-morphism}에 의해 토포스의 가환도표
$$
\xymatrix{
\Sh(\mathcal{S}/u(U)) \ar[r]_-{j_{u(U)}} \ar[d]_{p'} &
\Sh(\mathcal{S}) \ar[d]^p \\
\Sh(\mathcal{C}/U) \ar[r]^{j_U} &
\Sh(\mathcal{C}),
}
$$
를 얻는다. 여기서 $p'$는 함자
$u' : \mathcal{C}/U \to \mathcal{S}/u(U)$,
$V/U \mapsto u(V)/u(U)$로 주어진다. 집합 $E$에 집합 $E$를 값 $x$인
상수사상 $E \to u(U)$로 갖추어 대응시키는 함자
$j_x : \mathcal{S} \cong \mathcal{S}/x$를 생각하자.
그러면 $j_x$는 연속인 오른쪽 수반
$v_x : (\psi : E \to u(U)) \mapsto \psi^{-1}(\{x\})$를 갖는
충실충만 쌍대연속 함자이다. 또한
$j_{u(U)} \circ j_x = \text{id}_\mathcal{S}$이고
$v_x \circ u' = v$임에 유의하라. 이 관찰들로부터 다음 토포스의
가환도표를 얻는다.
$$
\xymatrix{
\Sh(\mathcal{S}) \ar[rd]^a \ar[rdd]_q
\ar `r[rrr] `d[dd]^p [rrdd] & & & \\
& \Sh(\mathcal{S}/u(U)) \ar[r]_-{j_{u(U)}} \ar[d]^{p'} &
\Sh(\mathcal{S}) \ar[d]^p & \\
& \Sh(\mathcal{C}/U) \ar[r]^{j_U} &
\Sh(\mathcal{C}) &
}
$$
구체적으로 다음과 같다.
\begin{enumerate}
\item 사상
$a : \Sh(\mathcal{S}) \to \Sh(\mathcal{S}/u(U))$는
쌍대연속 함자 $j_x$에 대응하는 토포스의 사상이고, 이는 연속 함자
$v_x$에 대응하는 사상과 같다. Lemma
\ref{sites-lemma-cocontinuous-morphism-topoi} 및 Section
\ref{sites-section-cocontinuous-adjoint}를 보라.
\item $j_{u(U)} \circ j_x = \text{id}_\mathcal{S}$이므로 합성
$p \circ j_{u(U)} \circ a = p$이다.
\item 합성 $p' \circ a$는 토포스의 사상을 준다. 더 나아가 이는 연속 함자
$v_x \circ u' = v$에 대응하는 토포스의 사상이다. 따라서 $v$는 실제로
$\mathcal{C}/U$의 점 $q$를 정의하며, 구성에 의해 위 도표에 들어맞는다.
\end{enumerate}
이로써 보조정리의 증명이 끝난다.
\end{proof}

\begin{lemma}
\label{sites-lemma-points-above-point}
$\mathcal{C}$, $p$, $u$, $U$는 Lemma
\ref{sites-lemma-point-localize}에서와 같다고 하자. Lemma
\ref{sites-lemma-point-localize}의 구성은 $p$ 위에 놓이는
$\mathcal{C}/U$의 점 $q$들과 $u(U)$의 원소 $x$들 사이의 일대일
대응을 준다.
\end{lemma}

\begin{proof}
$q$를 함자 $v : \mathcal{C}/U \to \textit{Sets}$로 주어지는
$\mathcal{C}/U$의 점이라 하고, 토포스의 사상으로서
$j_U \circ q = p$라고 하자. $\mathcal{C}$의 임의의 대상 $V$에 대하여
$u(V) = p^{-1}(h_V^\#)$임을 상기하라. Lemma
\ref{sites-lemma-point-site-topos}를 보라. 마찬가지로 $\mathcal{C}/U$의
임의의 대상 $V/U$에 대하여
$v(V/U) = q^{-1}(h_{V/U}^\#)$이다. 다음 두 도표를 생각하자.
$$
\vcenter{
\xymatrix{
\Mor_{\mathcal{C}/U}(W/U, V/U) \ar[r] \ar[d] &
\Mor_\mathcal{C}(W, V) \ar[d] \\
\Mor_{\mathcal{C}/U}(W/U, U/U) \ar[r] &
\Mor_\mathcal{C}(W, U)
}
}
\quad
\vcenter{
\xymatrix{
h_{V/U}^\# \ar[r] \ar[d] & j_U^{-1}(h_V^\#) \ar[d] \\
h_{U/U}^\# \ar[r] & j_U^{-1}(h_U^\#)
}
}
$$
오른쪽 도표는 $\mathcal{C}/U$ 위의 준층들의 도표를 층화한 것이며,
이 준층 도표는 $W/U$를 왼쪽 집합 도표로 보낸다
(여기에는 작은 기술적 사항이 있다. 즉
$(j_U^{-1}h_V)^\# = j_U^{-1}(h_V^\#)$이고 $h_U$에 대해서도 마찬가지이다.
Lemma \ref{sites-lemma-technical-pu}를 보라). 왼쪽 집합 도표는
데카르트임에 유의하라. 층화는 완전이므로
(Lemma \ref{sites-lemma-sheafification-exact}) 오른쪽 도표도
데카르트라고 결론내린다.

\medskip\noindent
완전함자 $q^{-1}$을 오른쪽 도표에 적용하면 집합들의 데카르트 도표
$$
\xymatrix{
v(V/U) \ar[r] \ar[d] & u(V) \ar[d] \\
v(U/U) \ar[r] & u(U)
}
$$
를 얻는다. 여기서는
% P01-ERRATA: P01-E0133 -- the frozen formula uses undefined j rather than the localization morphism j_U; it is preserved.
$q^{-1} \circ j^{-1} = p^{-1}$을 사용했다.
$U/U$는 $\mathcal{C}/U$의 종대상이므로 $v(U/U)$는 한원소집합이다.
따라서 $v(U/U)$의 $u(U)$에서의 상은 원소 $x$이고, 위쪽 수평사상은
원하는 전단사
$v(V/U) \to  \{y \in u(V) \mid y \mapsto x\text{ in }u(U)\}$를 준다.
\end{proof}

\begin{lemma}
\label{sites-lemma-stalk-j-shriek}
$\mathcal{C}$를 사이트라 하고, $p$를 함자
$u : \mathcal{C} \to \textit{Sets}$로 주어지는 $\mathcal{C}$의 점이라
하자. $U$를 $\mathcal{C}$의 대상이라 하자.
$\mathcal{C}/U$ 위의 임의의 층 $\mathcal{G}$에 대하여
$$
(j_{U!}\mathcal{G})_p =
\coprod\nolimits_q \mathcal{G}_q
$$
이다. 여기서 쌍대합은 Lemma \ref{sites-lemma-point-localize}에서처럼
원소 $x \in u(U)$에 대응하는 $\mathcal{C}/U$의 점 $q$들
위에서 취한다.
\end{lemma}

\begin{proof}
Lemma \ref{sites-lemma-describe-j-shriek}에 있는, 준층
$V \mapsto
\coprod\nolimits_{\varphi \in \Mor_\mathcal{C}(V, U)}
\mathcal{G}(V/_\varphi U)$에 대응하는 층으로서의
$j_{U!}\mathcal{G}$의 기술을 사용한다. 또한 Lemma
\ref{sites-lemma-point-pushforward-sheaf}에 의해
$p$에서의 $j_{U!}\mathcal{G}$의 줄기는 이 준층의 줄기와 같다.
따라서
$$
(j_{U!}\mathcal{G})_p =
\colim_{(V, y)} \coprod\nolimits_{\varphi : V \to U} \mathcal{G}(V/_\varphi U)
$$
임을 알 수 있다. 이 쌍대극한의 각 원소 $(V, y, \varphi, s)$에
$x = u(\varphi)(y) \in u(U)$를 대응시킬 수 있다. 따라서
$$
(j_{U!}\mathcal{G})_p =
\coprod\nolimits_{x \in u(U)}
\colim_{(\varphi : V \to U, y), \ u(\varphi)(y) = x} \mathcal{G}(V/_\varphi U).
$$
이는 점 $q$들의 구성에 의해 보조정리의 식과 같다.
\end{proof}

\begin{remark}
\label{sites-remark-not-pushforward}
주의: Lemma \ref{sites-lemma-stalk-j-shriek}의 결과에는 $j_{U, *}$에 대한
유사물이 없다.
\end{remark}

\section{토포스의 2-사상}
\label{sites-section-2-category}

\noindent
이 절에서는 토포스의 $2$-사상이라는 개념을 간략히 다룬다.

\begin{definition}
\label{sites-definition-2-morphism-topoi}
$f, g : \Sh(\mathcal{C}) \to \Sh(\mathcal{D})$를 토포스의 두 사상이라
하자. {\it $f$에서 $g$로 가는 2-사상}은 함자들의 변환
$t : f_* \to g_*$로 주어진다.
\end{definition}

\noindent
그림으로는 $t$를 때때로 다음과 같이 나타낸다.
$$
\xymatrix{
\Sh(\mathcal{C})
\rrtwocell^f_g{t}
&
&
\Sh(\mathcal{D})
}
$$
$f^{-1}$가 $f_*$의 수반이고 $g^{-1}$가 $g_*$의 수반이므로,
$t$는 함자들의 변환
$t : g^{-1} \to f^{-1}$도 유도한다(보통 같은 기호로 나타낸다).
이는 다음 도표가 가환한다는 조건으로 유일하게 특징지어진다.
$$
\xymatrix{
\Mor_{\Sh(\mathcal{D})}(\mathcal{G}, f_*\mathcal{F})
\ar[d]_{t \circ -} \ar@{=}[r] &
\Mor_{\Sh(\mathcal{C})}(f^{-1}\mathcal{G}, \mathcal{F})
\ar[d]^{- \circ t}
\\
\Mor_{\Sh(\mathcal{D})}(\mathcal{G}, g_*\mathcal{F})
\ar@{=}[r] &
\Mor_{\Sh(\mathcal{C})}(g^{-1}\mathcal{G}, \mathcal{F})
}
$$
집합론적 어려움(Remark \ref{sites-remark-morphism-topoi-big} 참조) 때문에
토포스의 2-범주를 얻지는 못한다. 그러나 수평 합성과 수직 합성을 정의하고
Categories, Section \ref{categories-section-2-categories}에 열거된
엄밀한 2-범주의 공리들이 성립함을 보일 수는 있다.
실제로 2-사상의 수직 합성은 분명하고(함자들의 변환을 합성하면 된다),
1-사상의 합성은 Definition \ref{sites-definition-topos}에서 정의했으며, 다음
$$
\xymatrix{
\Sh(\mathcal{C})
\rtwocell^f_g{t}
&
\Sh(\mathcal{D})
\rtwocell^{f'}_{g'}{s}
&
\Sh(\mathcal{E})
}
$$
의 수평 합성은 Categories, Definition
\ref{categories-definition-horizontal-composition}에서 도입한
함자들의 변환 $s \star t$로 정의한다.
구체적으로 $s \star t$는 다음 두 합성 가운데 어느 것으로도 주어진다.
$$
\xymatrix{
f'_*f_*\mathcal{F}
\ar[r]^{f'_*t} &
f'_*g_*\mathcal{F}
\ar[r]^s &
g'_*g_*\mathcal{F}
}
\quad\text{or}\quad
\xymatrix{
f'_*f_*\mathcal{F}
\ar[r]^s &
g'_*f_*\mathcal{F}
\ar[r]^{g'_*t} &
g'_*g_*\mathcal{F}
}
$$
(두 사상은 같다). 이 정의들이 Categories, Section
\ref{categories-section-formal-cat-cat}의 정의들과 일치하므로,
Categories, Lemma \ref{categories-lemma-properties-2-cat-cats}에 의해
이 정의들에 대하여 엄밀한 2-범주의 공리들이 성립한다.




\section{점 사이의 사상}
\label{sites-section-morphisms-points}

\begin{lemma}
\label{sites-lemma-maps-u-points}
$\mathcal{C}$를 사이트라 하자.
$u, u' : \mathcal{C} \to \textit{Sets}$를 두 함자라 하고
$t : u' \to u$를 함자들의 변환이라 하자.
그러면 줄기 함자들의 표준적 변환
$t_{stalk} : \mathcal{F}_{p'} \to \mathcal{F}_p$를 얻으며,
이는 Lemma \ref{sites-lemma-points-recover}의 식별을 통해 $t$와 일치한다.
\end{lemma}

\begin{proof}
생략한다.
\end{proof}

\begin{definition}
\label{sites-definition-morphism-points}
$\mathcal{C}$를 사이트라 하자. 함자
$u, u' : \mathcal{C} \to \textit{Sets}$로 주어지는
$\mathcal{C}$의 점 $p, p'$를 생각하자.
{\it 사상 $f : p \to p'$}는 함자들의 변환
$$
f_u : u' \to u.
$$
로 주어진다.
\end{definition}

\noindent
함자들의 변환이 반대 방향으로 간다는 점에 유의하자.
뒤에서 보겠지만, 사상 $f$를 그림으로 다음과 같은 일종의 $2$-화살표로
생각하면 이는 자연스럽다.
$$
\xymatrix{
\textit{Sets}
=
\Sh(pt)
\rrtwocell^p_{p'}{f}
&
&
\Sh(\mathcal{C})
}
$$
실제로 뒤에서 $f_u$가 마천루 층 구성들 사이의 표준적 함자 변환
$p_* \to p'_*$를 유도함을 볼 것이다.

\medskip\noindent
이는 상당히 중요한 개념이므로 여기서 더 완전하게 다룰 가치가 있다.
필요한 사항은 다음과 같다.
\begin{enumerate}
\item Example \ref{sites-example-point-G-sets}에서 설명한
$\mathcal{T}_G$의 점의 자기동형사상들을 기술한다.
\item $\Mor(p, p')$를 $\Mor(p_*, p'_*)$로 기술한다.
\item 위상공간에서 점의 특수화를 다룬다.
위상공간 $X$에서 $x' \in \overline{\{x\}}$이면 사상
$p \to p'$가 있음을 보인다. 여기서 $p$ (각각 $p'$)는
$x$ (각각 $x'$)에 대응하는 $X_{Zar}$의 점이다.
\end{enumerate}




\section{충분한 점을 갖는 사이트}
\label{sites-section-sites-enough-points}

\begin{definition}
\label{sites-definition-enough-points}
$\mathcal{C}$를 사이트라 하자.
\begin{enumerate}
\item 점들의 족 $\{p_i\}_{i\in I}$이 {\it 보수적}이라는 것은,
층의 임의의 사상 $\phi : \mathcal{F} \to \mathcal{G}$가
% P01-ERRATA: P01-E0134
모든 섬유 $\mathcal{F}_{p_i} \to \mathcal{G}_{p_i}$에서
동형사상이면 그 자체가 동형사상이라는 뜻이다.
\item 보수적인 점들의 족이 존재하면 $\mathcal{C}$가
{\it 충분한 점을 갖는다}고 한다.
\end{enumerate}
\end{definition}

\noindent
이 경우 ``완전성''을 줄기에서 확인할 수 있음이 밝혀진다.

\begin{lemma}
\label{sites-lemma-exactness-stalks}
$\mathcal{C}$를 사이트라 하고 $\{p_i\}_{i\in I}$를
보수적인 점들의 족이라 하자. 그러면 다음이 성립한다.
\begin{enumerate}
\item 층의 임의의 사상 $\varphi : \mathcal{F} \to \mathcal{G}$에
대하여, $\forall i, \varphi_{p_i}$가 단사이면 $\varphi$도 단사이다.
\item 층의 임의의 사상 $\varphi : \mathcal{F} \to \mathcal{G}$에
대하여, $\forall i, \varphi_{p_i}$가 전사이면 $\varphi$도 전사이다.
\item 층의 임의의 두 사상
$\varphi_1, \varphi_2 : \mathcal{F} \to \mathcal{G}$에 대하여,
$\forall i, \varphi_{1, p_i} = \varphi_{2, p_i}$이면
$\varphi_1 = \varphi_2$이다.
\item 유한 도표 $\mathcal{G} : \mathcal{J}
\to \Sh(\mathcal{C})$, 층 $\mathcal{F}$, 그리고 사상들
$q_j : \mathcal{F} \to \mathcal{G}_j$이 주어졌다고 하자.
그러면 $(\mathcal{F}, q_j)$가 이 도표의 극한일 필요충분조건은,
각 $i$에 대하여 줄기
$(\mathcal{F}_{p_i}, (q_j)_{p_i})$가 극한인 것이다.
\item 유한 도표 $\mathcal{F} : \mathcal{J}
\to \Sh(\mathcal{C})$, 층 $\mathcal{G}$, 그리고 사상들
$e_j : \mathcal{F}_j \to \mathcal{G}$가 주어졌다고 하자.
그러면 $(\mathcal{G}, e_j)$가 이 도표의 쌍대극한일 필요충분조건은,
각 $i$에 대하여 줄기
$(\mathcal{G}_{p_i}, (e_j)_{p_i})$가 쌍대극한인 것이다.
\end{enumerate}
\end{lemma}

\begin{proof}
모든 줄기 함자가 임의의 유한 극한 및 쌍대극한, 따라서 곱과 섬유곱
등과 가환한다는 사실을 반복해서 사용할 것이다. 또한 단사 사상은
단사사상이고 전사 사상은 전사사상이라는 사실도 사용할 것이다.
층의 사상 $\varphi : \mathcal{F} \to \mathcal{G}$가 단사일
필요충분조건은
$\mathcal{F} \to \mathcal{F} \times_\mathcal{G} \mathcal{F}$가
동형사상인 것이다. 따라서 (1)이 성립한다.
마찬가지로 $\varphi : \mathcal{F} \to \mathcal{G}$가 전사일
필요충분조건은
$\mathcal{G} \amalg_\mathcal{F} \mathcal{G} \to \mathcal{G}$가
동형사상인 것이다. 따라서 (2)가 성립한다.
사상 $a, b : \mathcal{F} \to \mathcal{G}$가 같을 필요충분조건은
$\text{Equalizer}(a, b : \mathcal{F} \to \mathcal{G}) \to \mathcal{F}$가
동형사상인 것이다\footnote{등화자는 곱과 섬유곱으로 나타낼 수 있다.
Categories, Lemma \ref{categories-lemma-finite-limits-exist}의 증명을
보라.}. 따라서 (3)이 성립한다.
(4)와 (5)는 정의와 이 증명 첫머리의 설명에서 곧바로 따른다.
\end{proof}

\begin{lemma}
\label{sites-lemma-enough}
$\mathcal{C}$를 사이트라 하고 $\{(p_i, u_i)\}_{i\in I}$를
점들의 족이라 하자. 이 족이 보수적일 필요충분조건은 다음과 같다.
임의의 층 $\mathcal{F}$, 임의의 $U\in \Ob(\mathcal{C})$, 그리고
서로 다른 임의의 두 절단 $s, s' \in \mathcal{F}(U)$,
$s \not = s'$에 대하여, 삼중항 $(U, x, s)$와 $(U, x, s')$가
$\mathcal{F}_{p_i}$의 서로 다른 원소를 정의하도록 하는
$i$와 $x\in u_i(U)$가 존재한다.
\end{lemma}

\begin{proof}
먼저 이 족이 보수적이고 $\mathcal{F}$, $U$, $s, s'$가
보조정리에서와 같다고 하자. 절단 $s$, $s'$는 서로 다른 사상
$a, a' : (h_U)^\# \to \mathcal{F}$를 정의한다. 따라서 Lemma
\ref{sites-lemma-exactness-stalks}에 의해 $a_{p_i} \not = a'_{p_i}$인
$i$가 존재한다. Lemmas \ref{sites-lemma-points-recover}와
\ref{sites-lemma-point-pushforward-sheaf}에 의해
$(h_U)^\#_{p_i} = u_i(U)$임을 상기하자.
따라서 $\mathcal{F}_{p_i}$에서 $a_{p_i}(x)
\not = a'_{p_i}(x)$가 되게 하는 $x \in u_i(U)$가 존재한다.
정의를 풀어 쓰면 $(U, x, s)$와 $(U, x, s')$가 보조정리의
주장과 같음을 알 수 있다.

\medskip\noindent
역을 보이기 위해 보조정리의 점 존재 조건을 가정하자.
모든 줄기에서 동형사상인 층의 사상
$\phi : \mathcal{F} \to \mathcal{G}$를 생각하자.
Lemma \ref{sites-lemma-mono-epi-sheaves}에 따라 $\phi$가 단사이면서
전사임을 보여야 한다. 단사성은 가정에서 즉시 따른다.
전사성을 보이려면
$\mathcal{G} \amalg_\mathcal{F} \mathcal{G} \to \mathcal{G}$가
동형사상임을 보여야 한다
(Categories, Lemma \ref{categories-lemma-characterize-mono-epi}).
이 사상은 명백히 전사이므로 단사성만 확인하면 충분한데, 이는 가정에 의해
$\mathcal{G} \amalg_\mathcal{F} \mathcal{G} \to \mathcal{G}$가
모든 줄기에서 단사라는 사실에서 따른다.
\end{proof}

\noindent
다음 보조정리의 점 $q_{i, x}$들은 Lemma
\ref{sites-lemma-points-above-point}의 의미에서 점 $p_i$ 위에 놓이는
$\mathcal{C}/U$의 모든 점과 정확히 같다.

\begin{lemma}
\label{sites-lemma-localize-enough}
$\mathcal{C}$를 사이트라 하자. $U$를 $\mathcal{C}$의 대상이라 하자.
% P01-ERRATA: P01-E0135
$\{(p_i, u_i)\}_{i\in I}$를 $\mathcal{C}$의 점들의 족이라 하자.
$x \in u_i(U)$에 대하여 $q_{i, x}$를
Lemma \ref{sites-lemma-point-localize}에서 구성한 $\mathcal{C}/U$의 점이라
하자. $\{p_i\}$가 보수적인 점들의 족이면
$\{q_{i, x}\}_{i \in I, x \in u_i(U)}$는
$\mathcal{C}/U$의 보수적인 점들의 족이다.
특히 $\mathcal{C}$가 충분한 점을 가지면 모든 국소화
$\mathcal{C}/U$도 충분한 점을 갖는다.
\end{lemma}

\begin{proof}
Lemma \ref{sites-lemma-essential-image-j-shriek}에 의해
$j_{U!}$가 동치
$j_{U!} : \Sh(\mathcal{C}/U) \to \Sh(\mathcal{C})/h_U^\#$를
유도함을 안다. 또한 Lemma \ref{sites-lemma-stalk-j-shriek}에 의해
$(j_{U!}\mathcal{G})_{p_i} = \coprod_x \mathcal{G}_{q_{i, x}}$임을
안다. 따라서 결과는 형식적으로 따른다.
\end{proof}

\noindent
다음 보조정리는 점의 존재를 사이트 위에서 국소적으로 확인할 수 있음을
말해 준다.

\begin{lemma}
\label{sites-lemma-enough-points-local}
$\mathcal{C}$를 사이트라 하자. $\{U_i\}_{i \in I}$를
$\mathcal{C}$의 대상들의 족이라 하자. 다음을 가정한다.
\begin{enumerate}
\item $\coprod h_{U_i}^\# \to *$는 층의 전사 사상이고,
\item 각 국소화 $\mathcal{C}/U_i$는 충분한 점을 갖는다.
\end{enumerate}
그러면 $\mathcal{C}$는 충분한 점을 갖는다.
\end{lemma}

\begin{proof}
각 $i \in I$에 대하여 $\{p_j\}_{j \in J_i}$를
$\mathcal{C}/U_i$의 보수적인 점들의 족이라 하자.
% P01-ERRATA: P01-E0136
$j \in J_i$에 대하여 $q_j : \Sh(pt) \to \Sh(\mathcal{C})$를
$p_j$와 국소화 사상
$\Sh(\mathcal{C}/U_i) \to \Sh(\mathcal{C})$의 합성이라 하자.
그러면 Lemma \ref{sites-lemma-point-morphism-topoi}에 의해 $q_j$는 점이다.
점들의 족 $\{q_j\}_{j \in \coprod J_i}$가 보수적임을 주장한다.
실제로 $\mathcal{C}$ 위의 층의 사상
$\mathcal{F} \to \mathcal{G}$가 주어져서 모든
$j \in \coprod J_i$에 대하여
$\mathcal{F}_{q_j} \to \mathcal{G}_{q_j}$가 동형사상이라고 하자.
$W$를 $\mathcal{C}$의 대상이라 하자.
가정 (1)에 의해 덮개 $\{W_a \to W\}$와 사상들
$W_a \to U_{i(a)}$가 존재한다.
Lemma \ref{sites-lemma-point-morphism-topoi}에 의해
$(\mathcal{F}|_{\mathcal{C}/U_{i(a)}})_{p_j} = \mathcal{F}_{q_j}$이고
$(\mathcal{G}|_{\mathcal{C}/U_{i(a)}})_{p_j} = \mathcal{G}_{q_j}$이므로,
점들의 족 $\{p_j\}_{j \in J_{i(a)}}$의 보수성에서
$\mathcal{F}|_{U_{i(a)}} \to \mathcal{G}|_{U_{i(a)}}$가 동형사상임을
알 수 있다. 따라서 각 $a$에 대하여
$\mathcal{F}(W_a) \to \mathcal{G}(W_a)$는 전단사이다.
마찬가지로 각 $a, b$에 대하여
$\mathcal{F}(W_a \times_W W_b) \to \mathcal{G}(W_a \times_W W_b)$도
전단사이다. 층 조건에 의해
$\mathcal{F}(W) \to \mathcal{G}(W)$는 전단사, 즉
$\mathcal{F} \to \mathcal{G}$는 동형사상이다.
\end{proof}

\begin{lemma}
\label{sites-lemma-check-morphism-sites}
$u : \mathcal{C} \to \mathcal{D}$를 사이트의 연속 함자라 하자.
$\{(q_i, v_i)\}_{i\in I}$를 $\mathcal{D}$의 보수적인 점들의 족이라
하자. 각 함자 $u_i = v_i \circ u$가 $\mathcal{C}$의 점을 정의한다면,
$u$는 사이트의 사상 $f : \mathcal{D} \to \mathcal{C}$를 정의한다.
\end{lemma}

\begin{proof}
함자 $u_i$에 대응하는 $\textit{PSh}(\mathcal{C})$ 위의 줄기 함자
(\ref{sites-equation-stalk})를 $p_i$로 나타내자.
% P01-ERRATA: P01-E0137
다음이 성립한다.
$$
(f^{-1}\mathcal{F})_{q_i} =
(u_s\mathcal{F})_{q_i} =
(u_p\mathcal{F})_{q_i} =
\mathcal{F}_{p_i}
$$
첫째 등식은 $f^{-1} = u_s$이기 때문이고, 둘째 등식은
Lemma \ref{sites-lemma-point-pushforward-sheaf}, 셋째 등식은
Lemma \ref{sites-lemma-point-functor}에 의한 것이다.
따라서 $p_i$가 점이면 $f$로 당긴 다음 $q_i$에서 줄기를 취하는 함자는
완전하다. 점들의 족 $\{q_i\}$가 보수적이므로 이는 $f^{-1}$가
완전한 함자임을 뜻하고, 따라서 Definition
\ref{sites-definition-morphism-sites}에 의해 $f$가 사이트의 사상임을 알 수 있다.
\end{proof}





\section{점의 존재 판정법}
\label{sites-section-criterion-points}

\noindent
이 절은 \cite[Expos\'e VI]{SGA4}에 실린 Deligne의 부록에 해당한다.
실제로 거의 글자 그대로 같다.

\medskip\noindent
$\mathcal{C}$를 사이트라 하자.
$(I, \geq)$가 유향 집합이고
$(U_i, f_{ii'})$가 $I$ 위의 역계라고 하자. Categories, Definition
\ref{categories-definition-system-over-poset}을 보라.
데이터 $(I, \geq, U_i, f_{ii'})$가 주어지면 다음을 정의한다.
$$
u : \mathcal{C} \longrightarrow \textit{Sets}, \quad
u(V) = \colim_i \Mor_\mathcal{C}(U_i , V)
$$
Section \ref{sites-section-points}에서와 같이 $u$에 대응하는 줄기 함자를
$\mathcal{F} \mapsto \mathcal{F}_p$라 하자.
이 특별한 경우에는 정의에서 곧바로 실제로
$$
\mathcal{F}_p = \colim_i \mathcal{F}(U_i)
$$
임을 알 수 있다.
각 함자 $V \mapsto \Mor_\mathcal{C}(U_i , V)$가 모든 유한 극한
(즉 $\mathcal{C}$에서 표현 가능한 것들)과 가환하므로 $u$도 그러하다.
% P01-ERRATA: P01-E0138
Categories, Lemma \ref{categories-lemma-directed-commutes}를 보라.

\medskip\noindent
계 $(I, \geq, U_i, f_{ii'})$가
$(J, \geq, V_j, g_{jj'})$의 {\it 세분}이라는 것은
% P01-ERRATA: P01-E0139
$J \subset I$이고, $J$ 위의 순서가 $I$의 순서에서 유도되며,
$V_j = U_j$, $g_{jj'} = f_{jj'}$라는 뜻이다(달리 말하면 $J$ 위의
역계가 $I$ 위의 역계에서 유도된다). $(I, \geq, U_i, f_{ii'})$에
대응하는 함자를 $u$라 하고 $(J, \geq, V_j, g_{jj'})$에 대응하는
함자를 $u'$라 하자. 그러면 함자들의 변환
$$
u' \longrightarrow u
$$
를 얻는다. 이는 단순히 $u'$의 쌍대극한들이 $u$의 쌍대극한들에 쓰이는
계의 부분계 위에서 취해지기 때문이다.
특히 대응하는 줄기 함자들의 변환
$\mathcal{F}_{p'} \to \mathcal{F}_p$를 얻는다.
Lemma \ref{sites-lemma-maps-u-points}를 보라.

\begin{lemma}
\label{sites-lemma-refine}
$\mathcal{C}$를 사이트라 하자.
$(J, \geq, V_j, g_{jj'})$를 위와 같은 계라 하고, 이에 대응하는
함자들의 쌍을 $(u', p')$라 하자.
$\mathcal{F}$를 $\mathcal{C}$ 위의 층이라 하자.
$s, s' \in \mathcal{F}_{p'}$를 서로 다른 원소라 하자.
$\{W_k \to W\}$를 $\mathcal{C}$의 유한 덮개라 하자.
$f \in u'(W)$라 하자.
$(J, \geq, V_j, g_{jj'})$의 세분
$(I, \geq, U_i, f_{ii'})$가 존재하여, $s, s'$는
$\mathcal{F}_p$의 서로 다른 원소로 가고 $u(W)$ 안에서 $f$의 상은
$u(W_k)$들 가운데 하나의 상에 속한다.
\end{lemma}

\begin{proof}
$f$가 $f' : V_{j_0} \to W$로 주어지게 하는 $j_0 \in J$가 존재한다.
% P01-ERRATA: P01-E0140
$j \geq j_0$에 대하여
$V_{j, k} = V_j \times_{f'\circ f_{j j_0}, W} W_k$로 놓는다.
그러면 $\{V_{j, k} \to V_j\}$는 사이트 $\mathcal{C}$의 유한 덮개이다.
따라서
$\mathcal{F}(V_j) \subset \prod_k \mathcal{F}(V_{j, k})$이다.
Categories, Lemma \ref{categories-lemma-directed-commutes}를 다시 적용하면
$$
\mathcal{F}_{p'} =
\colim_j \mathcal{F}(V_j)
\longrightarrow
\prod\nolimits_k \colim_j \mathcal{F}(V_{j, k})
$$
가 단사임을 알 수 있다. 따라서 $s$와 $s'$가
$\colim_j \mathcal{F}(V_{j, k})$에서 서로 다른 상을 갖게 하는
$k$가 존재한다. $J_0 = \{j \in J, j \geq j_0\}$ 및
$I = J \amalg J_0$로 놓는다. 두 번째 합성분의 어떤 원소도
첫 번째 합성분의 어떤 원소보다 작지 않게 하고, 그 밖에는 $J$의 순서를
사용하여 $I$에 순서를 준다. $j \in I$가 첫 번째 합성분에 속하면
$V_j$를 사용하고, $j \in I$가 두 번째 합성분에 속하면
$V_{j, k}$를 사용한다.
역계의 전이 사상들의 정의는 생략한다. 위에서 본 바에 의해
$s, s'$는 $\mathcal{F}_p$에서 서로 다른 상을 갖는다.
더욱이 구성에 의해 $f'$의 $V_{j, k}$로의 제한은 $W_k$를 거쳐
인수분해된다.
\end{proof}

\begin{lemma}
\label{sites-lemma-refine-all-at-once}
$\mathcal{C}$를 사이트라 하자.
$(J, \geq, V_j, g_{jj'})$를 위와 같은 계라 하고, 이에 대응하는
함자들의 쌍을 $(u', p')$라 하자.
$\mathcal{F}$를 $\mathcal{C}$ 위의 층이라 하자.
$s, s' \in \mathcal{F}_{p'}$를 서로 다른 원소라 하자.
$(J, \geq, V_j, g_{jj'})$의 세분
$(I, \geq, U_i, f_{ii'})$가 존재하여, $s, s'$는
$\mathcal{F}_p$의 서로 다른 원소로 가고, 사이트 $\mathcal{C}$의
모든 유한 덮개 $\{W_k \to W\}$ 및 임의의 $f \in u'(W)$에 대하여
$u(W)$ 안에서 $f$의 상은 $u(W_k)$들 가운데 하나의 상에 속한다.
\end{lemma}

\begin{proof}
$E$를 쌍 $(\{W_k \to W\}, f\in u'(W))$들의 집합이라 하자.
다음을 만족하는 쌍
$(E' \subset E, (I, \geq, U_i, f_{ii'}))$들을 생각하자.
\begin{enumerate}
% P01-ERRATA: P01-E0141
\item $(I, \geq, U_i, g_{ii'})$는 $(J, \geq, V_j, g_{jj'})$의 세분이다.
\item $s, s'$는 $\mathcal{F}_p$의 서로 다른 원소로 간다.
\item 모든 쌍 $(\{W_k \to W\}, f\in u'(W)) \in E'$에 대하여,
$u(W)$ 안에서 $f$의 상은 $u(W_k)$들 가운데 하나의 상에 속한다.
\end{enumerate}
이러한 쌍들에 첫 번째 성분에서는 포함관계로, 두 번째 성분에서는
세분관계로 순서를 준다.
% P01-ERRATA: P01-E0142
위와 같은 모든 쌍 $(E' \subset E, (I, \geq, U_i, f_{ii'}))$들의
클래스를 $\mathcal{S}$라 하자.
$\mathcal{S}$가 Zorn 보조정리의 가정을 만족한다고 주장한다.
실제로 $(E'_a, (I_a, \geq, U_i, f_{ii'}))_{a \in A}$를
$\mathcal{S}$의 전순서 부분클래스라 하자. 그러면
$E' = \bigcup_{a \in A} E'_a$와 $I = \bigcup_{a \in A} I_a$를
정의할 수 있다. 이에 대응하는 쌍
$(E' , (I, \geq, U_i, f_{ii'}))$가 $\mathcal{S}$의 원소라고
주장한다. 조건 (1)과 (3)은 분명하다. 조건 (2)에 대해서는
$$
u = \colim_{a \in A} u_a
\text{ and correspondingly }
\mathcal{F}_p = \colim_{a \in A} \mathcal{F}_{p_a}
$$
임에 유의하자. 이 줄기에서 $s, s'$의 상이 서로 다르다는 것은
집합의 유향 쌍대극한에 대한 기술에서 따른다. Categories, Section
\ref{categories-section-directed-colimits}을 보라. 이 상황에서는 간단히
$(E', (I, \ldots)) = \bigcup_{a \in A}(E'_a, (I_a, \ldots))$라고
쓰겠다.

\medskip\noindent
이제 $\mathcal{S}$가 집합이었다면 Zorn 보조정리를 적용할 수 있고,
위의 Lemma \ref{sites-lemma-refine}와 결합하여 보조정리를 쉽게 증명할 수
있었을 것이다. 그러나 $\mathcal{S}$는 클래스이므로 그렇지 않다.
이를 피하기 위해 $E$에 정렬순서를 하나 택한다.
$e \in E$에 대하여 $E'_e = \{e' \in E \mid e' \leq e\}$로 놓는다.
초한귀납을 사용하여 다음을 만족하는 쌍
$(E'_e, (I_e, \ldots)) \in \mathcal{S}$들을 구성한다.
$e_1 \leq e_2 \Rightarrow (E'_{e_1}, (I_{e_1}, \ldots))
\leq (E'_{e_2}, (I_{e_2}, \ldots))$.
$e \in E$, 이를테면 $e = (\{W_k \to W\}, f\in u'(W))$라 하자.
$e$가 선행자 $e - 1$을 가지면, 계
$e = (\{W_k \to W\}, f\in u'(W))$에 관하여
Lemma \ref{sites-lemma-refine}에서와 같이 $(I_e, \ldots)$를
$(I_{e - 1}, \ldots)$의 세분으로 잡는다.
$e$가 선행자를 갖지 않으면, 계
$e = (\{W_k \to W\}, f\in u'(W))$에 관하여
$(I_e, \ldots)$를 $\bigcup_{e' < e} (I_{e'}, \ldots)$의 세분으로
잡는다. 마지막으로 합집합 $\bigcup_{e \in E} I_e$가
보조정리에서 제기한 문제의 해가 된다.
\end{proof}

\begin{proposition}
\label{sites-proposition-criterion-points}
\begin{reference}
\cite[Expos\'e VI, Appendix by Deligne, Proposition 9.0]{SGA4}
\end{reference}
$\mathcal{C}$를 사이트라 하자. 다음을 가정한다.
\begin{enumerate}
\item $\mathcal{C}$에서 유한 극한들이 존재하고,
\item 모든 덮개 $\{U_i \to U\}_{i \in I}$는
$\mathcal{C}$의 유한 덮개에 의한 세분을 갖는다.
\end{enumerate}
그러면 $\mathcal{C}$는 충분한 점을 갖는다.
\end{proposition}

\begin{proof}
$\mathcal{C}$ 위의 임의의 층 $\mathcal{F}$, 임의의
$U \in \Ob(\mathcal{C})$, 그리고 서로 다른 임의의 절단
$s, s' \in \mathcal{F}(U)$가 주어졌을 때, $s, s'$가
$\mathcal{F}_p$에서 서로 다른 상을 갖게 하는 점 $p$가 존재함을
보여야 한다. Lemma \ref{sites-lemma-enough}를 보라.
$J = \{1\}$, $V_1 = U$, $g_{11} = \text{id}_U$인 계
$(J, \geq, V_j, g_{jj'})$를 생각하자.
Lemma \ref{sites-lemma-refine-all-at-once}를 적용한다.
그 보조정리에 의해 주어진 계를 세분하는 계
$(I, \geq, U_i, f_{ii'})$를 얻는데, 이는
$s_p \not = s'_p$이고 사이트 $\mathcal{C}$의 모든 유한 덮개
$\{W_k \to W\}$에 대하여 사상
$\coprod_k u(W_k) \to u(W)$가 전사이도록 한다.
$\mathcal{C}$의 모든 덮개는 유한 덮개로 세분될 수 있으므로,
사이트 $\mathcal{C}$의 {\it 임의의} 덮개
$\{W_k \to W\}$에 대하여
$\coprod_k u(W_k) \to u(W)$가 전사임을 결론 내린다.
이는 $u = p$가 점임을 뜻한다. Proposition
\ref{sites-proposition-point-limits} 및 이 절의 시작에서 $u$가 유한
극한들과 가환함을 보장한 논의를 보라.
\end{proof}

\begin{lemma}
\label{sites-lemma-criterion-points}
$\mathcal{C}$를 사이트라 하자. $I$를 집합이라 하고, 각
$i \in I$에 대하여 $U_i$를 다음을 만족하는 $\mathcal{C}$의
대상이라 하자.
\begin{enumerate}
% P01-ERRATA: P01-E0143
\item $\coprod h_{U_i}$는 $\Sh(\mathcal{C})$의 종대상 위로 전사하고,
\item $\mathcal{C}/U_i$는 Proposition
\ref{sites-proposition-criterion-points}의 가정들을 만족한다.
\end{enumerate}
그러면 $\mathcal{C}$는 충분한 점을 갖는다.
\end{lemma}

\begin{proof}
가정 (2)와 명제에 의해 $\mathcal{C}/U_i$는 충분한 점을 갖는다.
$\mathcal{C}/U_i$의 점들은 Lemma
\ref{sites-lemma-point-morphism-sites}의 절차를 통해 $\mathcal{C}$의
점들을 준다. 따라서 다음을 보이면 충분하다.
$\phi : \mathcal{F} \to \mathcal{G}$가 $\mathcal{C}$ 위의 층의
사상이고 모든 $i$에 대하여
$\phi|_{\mathcal{C}/U_i}$가 동형사상이면, $\phi$도 동형사상이다.
가정 (1)에 의해 $\mathcal{C}$의 모든 대상 $W$에 대하여
덮개 $\{W_j \to W\}_{j \in J}$가 존재하여, 각 $j \in J$에
대해 어떤 $i \in I$와 사상 $f_j : W_j \to U_i$가 존재한다.
그러면 사상들
$\mathcal{F}(W_j) \to \mathcal{G}(W_j)$는 전단사이고, 마찬가지로
$\mathcal{F}(W_j \times_W W_{j'}) \to \mathcal{G}(W_j \times_W W_{j'})$도
전단사이다. 층 조건에 의해 원하는 대로
$\mathcal{F}(W) \to \mathcal{G}(W)$가 전단사이다.
\end{proof}




\section{약하게 축약 가능한 대상}
\label{sites-section-w-contractible}

\noindent
사이트의 대상이 {\it 약하게 축약 가능하다}는 것은 다음 보조정리의
서로 동치인 조건들을 만족한다는 뜻이다.

\begin{lemma}
\label{sites-lemma-w-contractible}
$\mathcal{C}$를 사이트라 하고 $U$를 $\mathcal{C}$의 대상이라 하자.
다음 조건들은 서로 동치이다.
\begin{enumerate}
\item 모든 덮개 $\{U_i \to U\}$에 대하여, 사상
$\coprod h_{U_i} \to h_U$를 층화한 사상의 오른쪽 역인 층의 사상
$h_U^\# \to \coprod h_{U_i}^\#$가 존재한다.
\item 집합의 층의 모든 전사 사상
$\mathcal{F} \to \mathcal{G}$에 대하여 사상
$\mathcal{F}(U) \to \mathcal{G}(U)$가 전사이다.
\end{enumerate}
\end{lemma}

\begin{proof}
(1)을 가정하고 집합의 층의 전사 사상
$\mathcal{F} \to \mathcal{G}$를 생각하자.
$s \in \mathcal{G}(U)$에 대하여, $s|_{U_i}$로 가는 원소
$t_i \in \mathcal{F}(U_i)$와 덮개 $\{U_i \to U\}$가 존재한다.
Definition \ref{sites-definition-sheaves-injective-surjective}를 보라.
(\ref{sites-equation-map-representable-into-sheaf})를 통해 $t_i$를 사상
$t_i : h_{U_i}^\# \to \mathcal{F}$로 생각하자.
% P01-ERRATA: P01-E0144
그러면 $\coprod t_i : \coprod h_{U_i}^\# \to \mathcal{F}$에
(1)에서 얻은 사상 $h_U^\# \to \coprod h_{U_i}^\#$를 앞에서 합성하여,
$s$로 가는 절단 $t \in \mathcal{F}(U)$를 얻는다.
따라서 (2)가 성립한다.

\medskip\noindent
(2)를 가정하자. $\{U_i \to U\}$를 덮개라 하자.
그러면 Lemma \ref{sites-lemma-covering-surjective-after-sheafification}에 의해
$\coprod h_{U_i}^\# \to h_U^\#$는 전사이다.
따라서 (2)에 의해, $\coprod h_{U_i}^\#$의 절단 $s$ 가운데
$h_U^\#$의 절단 $\text{id}_U$로 가는 것이 존재한다.
이 절단은 $\coprod h_{U_i} \to h_U$를 층화한 사상의 오른쪽 역인
사상 $h_U^\# \to \coprod h_{U_i}^\#$에 대응한다.
이로써 (1)이 증명된다.
\end{proof}

\begin{definition}
\label{sites-definition-w-contractible}
$\mathcal{C}$를 사이트라 하자.
\begin{enumerate}
\item $\mathcal{C}$의 대상 $U$가 Lemma
\ref{sites-lemma-w-contractible}의 서로 동치인 조건들을 만족하면
{\it 약하게 축약 가능하다}고 한다.
\item $\mathcal{C}$의 모든 대상 $U$가, 모든 $i$에 대하여
$U_i$가 약하게 축약 가능한 덮개 $\{U_i \to U\}$를 가지면
이 사이트에는 {\it 약하게 축약 가능한 대상이 충분히 많다}고 한다.
\item 더 일반적으로 $P$가 $\mathcal{C}$의 대상들의 성질이라고 하자.
$\mathcal{C}$의 모든 대상 $U$가, 모든 $i$에 대하여 $U_i$가
$P$를 갖는 덮개 $\{U_i \to U\}$를 가지면
$\mathcal{C}$에는 {\it $P$-대상이 충분히 많다}고 한다.
\end{enumerate}
\end{definition}

\noindent
$\mathbf{A}^1_\mathbf{C}$의 작은 \'etale 사이트에는 약하게 축약 가능한
대상이 하나도 없다. 반면 임의의 스킴의 작은 pro-\'etale 사이트에는
축약 가능한 대상이 충분히 많다.






\section{직접상의 완전성 성질}
\label{sites-section-pushforward}

\noindent
$f$를 토포스의 사상이라 하자. 일반적으로 함자 $f_*$는 좌완전할 뿐이다.
이 함자에 여러 추가 조건을 부과하여 토포스의 특정한 사상 부류들을
골라낼 수 있다. 여기서는 그 조건들을 모으고 논리적 의존관계 몇 가지를
기록한다. 다음 보조정리의 일부는 순전히 범주론적이다. 즉 토포스의
사상이 필요하지 않고, 수반함자 한 쌍만 있으면 충분하다.

\begin{lemma}
\label{sites-lemma-exactness-properties}
$f : \Sh(\mathcal{C}) \to \Sh(\mathcal{D})$를 토포스의 사상이라 하자.
(집합의 층에 관한) 다음 성질들을 생각하자.
\begin{enumerate}
\item $f_*$는 충실하다.
\item $f_*$는 충실충만하다.
\item $\Sh(\mathcal{C})$의 모든 $\mathcal{F}$에 대하여
$f^{-1}f_*\mathcal{F} \to \mathcal{F}$는 전사이다.
\item $f_*$는 전사를 전사로 보낸다.
\item $f_*$는 쌍대등화자와 가환한다.
\item $f_*$는 밂과 가환한다.
\item $\Sh(\mathcal{C})$의 모든 $\mathcal{F}$에 대하여
$f^{-1}f_*\mathcal{F} \to \mathcal{F}$는 동형사상이다.
\item $f_*$는 단사성을 반영한다.
\item $f_*$는 전사성을 반영한다.
\item $f_*$는 전단사성을 반영한다.
\item 임의의 전사 $\mathcal{F} \to f^{-1}\mathcal{G}$에 대하여,
$f^{-1}\mathcal{G}' \to f^{-1}\mathcal{G}$가
$\mathcal{F} \to f^{-1}\mathcal{G}$를 거쳐 인수분해되게 하는 전사
$\mathcal{G}' \to \mathcal{G}$가 존재한다.
\end{enumerate}
그러면 다음 함의들이 성립한다.
\begin{enumerate}
\item[(a)] (2) $\Rightarrow$ (1),
\item[(b)] (3) $\Rightarrow$ (1),
\item[(c)] (7) $\Rightarrow$ (1), (2), (3), (8), (9), (10).
\item[(d)] (3) $\Leftrightarrow$ (9),
\item[(e)] (6) $\Rightarrow$ (4) and (5) $\Rightarrow$ (4),
\item[(f)] (4) $\Leftrightarrow$ (11),
\item[(g)] (9) $\Rightarrow$ (8), (10), and
\item[(h)] (2) $\Leftrightarrow$ (7).
\end{enumerate}
그림으로 나타내면 다음과 같다.
$$
\xymatrix{
(6) \ar@{=>}[rd] & & & & & (9) \ar@{=>}[r] \ar@{=>}[rd] & (8) \\
& (4) \ar@{<=>}[r] & (11) &
(2) \ar@{<=>}[r] &
(7) \ar@{=>}[ru] \ar@{=>}[rd] & & (10) \\
(5) \ar@{=>}[ur] & & & & & (3) \ar@{=>}[r] & (1)
}
$$
\end{lemma}

\begin{proof}
(a)의 증명: 정의에서 즉시 따른다.

\medskip\noindent
(b)의 증명. $a, b : \mathcal{F} \to \mathcal{F}'$를
$\mathcal{C}$ 위의 층의 사상들이라 하자.
$f_*a = f_*b$이면 $f^{-1}f_*a = f^{-1}f_*b$이다.
다음 가환도표를 생각하자.
$$
\xymatrix{
\mathcal{F} \ar@<-1ex>[r] \ar@<1ex>[r] & \mathcal{F}' \\
f^{-1}f_*\mathcal{F} \ar@<-1ex>[r] \ar@<1ex>[r] \ar[u] &
f^{-1}f_*\mathcal{F}' \ar[u]
}
$$
아래쪽 두 화살표가 같고 수직 화살표들이 전사이면 위쪽 두 화살표도 같다.
따라서 (b)가 따른다.

\medskip\noindent
(c)의 증명. $a : \mathcal{F} \to \mathcal{F}'$를
$\mathcal{C}$ 위의 층의 사상이라 하자. 다음 가환도표를 생각하자.
$$
\xymatrix{
\mathcal{F} \ar[r] & \mathcal{F}' \\
f^{-1}f_*\mathcal{F} \ar[r] \ar[u] &
f^{-1}f_*\mathcal{F}' \ar[u]
}
$$
(7)이 성립하면 수직 화살표들은 동형사상이다.
따라서 $f_*a$가 단사(각각 전사, 각각 전단사)이면 아래쪽 화살표가
단사(각각 전사, 각각 전단사)이고, 그러므로 위쪽 화살표도
단사(각각 전사, 각각 전단사)이다. 따라서 (7)이 (8), (9), (10)을
함의함을 알 수 있다. (7)이 (3)을 함의하는 것은 분명하다.
(7) $\Rightarrow$ (2), (1)은 (a)와 아래에서 보일 (h)에서 따른다.

\medskip\noindent
(d)의 증명. (3)을 가정하자. $a : \mathcal{F} \to \mathcal{F}'$를
$f_*a$가 전사인 $\mathcal{C}$ 위의 층의 사상이라 하자.
$f^{-1}$가 완전하므로
$f^{-1}f_*a : f^{-1}f_*\mathcal{F} \to f^{-1}f_*\mathcal{F}'$는
전사이다. 이를 (3)과 결합하면 $a$가 전사임을 알 수 있다.
이는 (9)가 성립한다는 뜻이다.
(9)를 가정하자. $\mathcal{F}$를 $\mathcal{C}$ 위의 층이라 하자.
사상 $f^{-1}f_*\mathcal{F} \to \mathcal{F}$가 전사임을 보여야 한다.
$f_*f^{-1}f_*\mathcal{F} \to f_*\mathcal{F}$가 전사임을 보이면
충분하다. 이는 한쪽 역인 표준적 사상
$f_*\mathcal{F} \to f_*f^{-1}f_*\mathcal{F}$가 있기 때문에 참이다.

\medskip\noindent
(e)의 증명. 이후 따로 언급하지 않고 Categories, Lemma
\ref{categories-lemma-characterize-mono-epi}를 사용한다.
$\mathcal{F} \to \mathcal{F}'$가 전사이면
$\mathcal{F}' \amalg_\mathcal{F} \mathcal{F}' \to \mathcal{F}'$는
동형사상이다. 따라서 (6)은
$$
f_*\mathcal{F}' \amalg_{f_*\mathcal{F}} f_*\mathcal{F}' =
f_*(\mathcal{F}' \amalg_\mathcal{F} \mathcal{F}')
\longrightarrow
f_*\mathcal{F}'
$$
도 동형사상임을 함의한다. 이는 다시
$f_*\mathcal{F} \to f_*\mathcal{F}'$가 전사임을 함의한다.
따라서 (6)이 (4)를 함의함을 알 수 있다.
$\mathcal{F} \to \mathcal{F}'$가 전사이면 Lemma
\ref{sites-lemma-coequalizer-surjection}에 의해 $\mathcal{F}'$는
두 사영
$\mathcal{F} \times_{\mathcal{F}'} \mathcal{F} \to \mathcal{F}$의
쌍대등화자이다. 따라서 (5)가 성립하면 $f_*\mathcal{F}'$는
두 사영
$$
f_*(\mathcal{F} \times_{\mathcal{F}'} \mathcal{F}) =
f_*\mathcal{F} \times_{f_*\mathcal{F}'} f_*\mathcal{F}
\longrightarrow
f_*\mathcal{F}
$$
의 쌍대등화자이다. 이는 분명히
$f_*\mathcal{F} \to f_*\mathcal{F}'$가 전사라는 뜻이다.
따라서 (5)도 (4)를 함의한다.

\medskip\noindent
(f)의 증명. (4)를 가정하자. $\mathcal{F} \to f^{-1}\mathcal{G}$를
$\mathcal{C}$ 위의 층의 전사 사상이라 하자. (4)에 의해
$f_*\mathcal{F} \to f_*f^{-1}\mathcal{G}$가 전사이다.
$\mathcal{G}'$를 다음 섬유곱으로 정의하자.
$$
\xymatrix{
f_*\mathcal{F} \ar[r] & f_*f^{-1}\mathcal{G} \\
\mathcal{G}' \ar[u] \ar[r] & \mathcal{G} \ar[u]
}
$$
그러면 $\mathcal{G}' \to \mathcal{G}$도 전사이다.
다음 가환도표를 생각하면 필요한 결과를 얻는다.
$$
\xymatrix{
\mathcal{F} \ar[r] & f^{-1}\mathcal{G} \\
f^{-1}f_*\mathcal{F} \ar[r] \ar[u] & f^{-1}f_*f^{-1}\mathcal{G} \ar[u] \\
f^{-1}\mathcal{G}' \ar[u] \ar[r] & f^{-1}\mathcal{G} \ar[u]
}
$$
역으로 (11)을 가정하자.
% P01-ERRATA: P01-E0145
$a : \mathcal{F} \to \mathcal{F}'$를 $\mathcal{C}$ 위의 층의
전사 사상이라 하자. 다음 섬유곱 도표를 생각하자.
$$
\xymatrix{
\mathcal{F} \ar[r] & \mathcal{F}' \\
\mathcal{F}'' \ar[u] \ar[r] & f^{-1}f_*\mathcal{F}' \ar[u]
}
$$
아래쪽 수평 화살표가 전사이므로, (11)에 의해 전사
$\gamma : \mathcal{G}' \to f_*\mathcal{F}'$를 찾아서
$f^{-1}\gamma$가
$\mathcal{F}'' \to f^{-1}f_*\mathcal{F}'$를 거쳐 인수분해되게 할
수 있다.
$$
\xymatrix{
& \mathcal{F} \ar[r] & \mathcal{F}' \\
f^{-1}\mathcal{G}' \ar[r] & \mathcal{F}'' \ar[u] \ar[r] &
f^{-1}f_*\mathcal{F}' \ar[u]
}
$$
여기에 $f_*$를 적용하면 다음 가환도표를 얻는다.
$$
\xymatrix{
& f_*\mathcal{F} \ar[r] & f_*\mathcal{F}' \\
f_*f^{-1}\mathcal{G}' \ar[r] & f_*\mathcal{F}'' \ar[u] \ar[r] &
f_*f^{-1}f_*\mathcal{F}' \ar[u] \\
\mathcal{G}' \ar[u] \ar[rr] & & f_*\mathcal{F}' \ar[u]
}
$$
이는 (4)가 성립함을 증명한다.

\medskip\noindent
(g)의 증명. (9)를 가정하자. 이후 따로 언급하지 않고 Categories,
Lemma \ref{categories-lemma-characterize-mono-epi}를 사용한다.
$a : \mathcal{F} \to \mathcal{F}'$를 $f_*a$가 단사인
$\mathcal{C}$ 위의 층의 사상이라 하자. 이는
$f_*\mathcal{F} \to f_*\mathcal{F} \times_{f_*\mathcal{F}'} f_*\mathcal{F} =
f_*(\mathcal{F} \times_{\mathcal{F}'} \mathcal{F})$가
동형사상이라는 뜻이다. 따라서 (9)에 의해
$\mathcal{F} \to \mathcal{F} \times_{\mathcal{F}'} \mathcal{F}$가
전사, 즉 동형사상임을 알 수 있다. 그러므로 $a$는 단사이고,
즉 (8)이 성립한다. (10)은 자명하게 (8) $+$ (9)와 동치이므로
(g)의 증명이 끝난다.

\medskip\noindent
(h)의 증명. 이는 Categories, Lemma
\ref{categories-lemma-adjoint-fully-faithful}이다.
\end{proof}

\noindent
다음은 사이트의 사상에 관한 조건으로, 함자 $f_*$가 전사 사상을
전사 사상으로 보냄을 보장한다.

\begin{lemma}
\label{sites-lemma-weaker}
$f : \mathcal{D} \to \mathcal{C}$를 연속 함자
$u : \mathcal{C} \to \mathcal{D}$에 대응하는 사이트의 사상이라 하자.
$\mathcal{C}$의 임의의 대상 $U$와 $\mathcal{D}$의 임의의 덮개
$\{V_j \to u(U)\}$에 대하여, 층의 사상
$$
\coprod h_{u(U_i)}^\# \to h_{u(U)}^\#
$$
가 층의 사상
$$
\coprod h_{V_j}^\# \to h_{u(U)}^\#.
$$
를 거쳐 인수분해되도록 하는 $\mathcal{C}$의 덮개
$\{U_i \to U\}$가 존재한다고 가정하자.
그러면 $f_*$는 층의 전사 사상을 층의 전사 사상으로 보낸다.
\end{lemma}

\begin{proof}
$a : \mathcal{F} \to \mathcal{G}$를 $\mathcal{D}$ 위의 층의
전사 사상이라 하자. $U$를 $\mathcal{C}$의 대상이라 하고
$s \in f_*\mathcal{G}(U) = \mathcal{G}(u(U))$라 하자.
가정에 의해 덮개 $\{V_j \to u(U)\}$와 절단
$s_j \in \mathcal{F}(V_j)$들이 존재하여
$a(s_j) = s|_{V_j}$이다. 이제 절단 $s$, $s_j$ 및 $a$가
다음 층의 사상들의 가환도표를 준다고 생각할 수 있다.
$$
\xymatrix{
\coprod h_{V_j}^\# \ar[r]_-{\coprod s_j} \ar[d] & \mathcal{F} \ar[d]^a \\
h_{u(U)}^\# \ar[r]^s & \mathcal{G}
}
$$
가정에 의해 덮개 $\{U_i \to U\}$가 존재하여 위 가환도표를
다음과 같이 확장할 수 있다.
$$
\xymatrix{
& \coprod h_{V_j}^\# \ar[r]_-{\coprod s_j} \ar[d] & \mathcal{F} \ar[d]^a \\
\coprod h_{u(U_i)}^\# \ar[r] \ar[ur] &
h_{u(U)}^\# \ar[r]^s & \mathcal{G}
}
$$
$\mathcal{F}$가 층이므로 왼쪽 아래 꼭짓점에서 오른쪽 위 꼭짓점으로 가는
사상은 절단족 $s_i \in \mathcal{F}(u(U_i))$, 즉 절단들
$s_i \in f_*\mathcal{F}(U_i)$에 대응한다.
도표의 가환성은 $a(s_i)$가 $s$의 $U_i$로의 제한과 같음을 뜻한다.
달리 말하면 $f_*a$가 층의 전사 사상임을 보였다.
\end{proof}

\begin{example}
\label{sites-example-weaker-not-coequalizer}
$f : \mathcal{D} \to \mathcal{C}$가 Lemma \ref{sites-lemma-weaker}의
가정들을 만족한다고 하자. 일반적으로 $f_*$가 쌍대등화자 또는 밂과
가환하는 것은 아니다. 실제로 $f$가 위상공간의 사상
$X = \{1, 2\} \to Y = \{*\}$에 대응하는 사이트의 사상이라고 하자
(Example \ref{sites-example-continuous-map} 참조). 여기서 $Y$는 한 점 공간이고
$X = \{1, 2\}$는 두 점을 갖는 이산공간이다.
$X$ 위의 층 $\mathcal{F}$는 집합의 쌍 $(A_1, A_2)$로 주어진다.
그러면 $f_*\mathcal{F}$는 집합 $A_1 \times A_2$에 대응한다.
따라서
$a = (a_1, a_2), b = (b_1, b_2) : (A_1, A_2) \to (B_1, B_2)$가
$X$ 위의 층의 사상들이라면 $a, b$의 쌍대등화자는
$(C_1, C_2)$이다. 여기서 $C_i$는 $a_i, b_i$의 쌍대등화자이다.
한편 $f_*a, f_*b$의 쌍대등화자는
$$
a_1 \times a_2, b_1 \times b_2 :
A_1 \times A_2 \longrightarrow B_1 \times B_2
$$
의 쌍대등화자이며, 이는 일반적으로 $C_1 \times C_2$와 다르다.
실제로 $A_2 = \emptyset$이면 $A_1 \times A_2 = \emptyset$이므로
표시된 화살표들의 쌍대등화자는 $B_1 \times B_2$이지만,
일반적으로 $C_1 \not = B_1$이다. 밂에 대해서도 비슷한 예가 성립한다.
\end{example}

\noindent
다음 보조정리는 사이트의 사상에 대응하는 함자 $f_*$가 단사성과
전사성을 반영하기 위한 판정법을 준다. 이는 또한 $f_*$가 충실하고
사상 $f^{-1}f_*\mathcal{F} \to \mathcal{F}$가 항상 전사임을 뜻한다.

\begin{lemma}
\label{sites-lemma-cover-from-below}
$f : \mathcal{D} \to \mathcal{C}$를 함자
$u : \mathcal{C} \to \mathcal{D}$로 주어지는 사이트의 사상이라 하자.
$\mathcal{D}$의 모든 대상 $V$에 대하여 $\mathcal{C}$의 대상들
$U_i$와 사상들 $u(U_i) \to V$가 존재하여
$\{u(U_i) \to V\}$가 $\mathcal{D}$의 덮개라고 가정하자.
이 경우 함자 $f_* : \Sh(\mathcal{D}) \to \Sh(\mathcal{C})$는
단사성과 전사성을 반영한다.
\end{lemma}

\begin{proof}
% P01-ERRATA: P01-E0146
$\alpha : \mathcal{F} \to \mathcal{G}$를 $\mathcal{D}$ 위의 층의
사상이라 하자. 가정과 층 조건에 의해 $\mathcal{D}$의 모든 대상
$V$에 대하여, 어떤 $U_i \in \Ob(\mathcal{C})$가 존재해서
$\mathcal{F}(V) \subset \prod \mathcal{F}(u(U_i)) =
\prod f_*\mathcal{F}(U_i)$이고, $\mathcal{G}$에 대해서도 마찬가지이다.
따라서 $f_*\alpha$가 단사이면 $\alpha$가 단사임이 분명하다.
달리 말하면 $f_*$는 단사성을 반영한다.

\medskip\noindent
$f_*\alpha$가 전사라고 하자. 위와 같은 $V, U_i, u(U_i) \to V$와
절단 $s \in \mathcal{G}(V)$에 대하여, $s|_{u(U_{ij})}$가
$\mathcal{F}(u(U_{ij}))$의 상에 속하도록 하는 덮개들
$\{U_{ij} \to U_i\}$가 존재한다. $u$가 연속이고 사이트의 공리들이
성립하므로 $\{u(U_{ij}) \to V\}$는 덮개이다. 따라서 $s$는
국소적으로 상에 속한다. 그러므로 $\alpha$는 전사이다.
달리 말하면 $f_*$는 전사성을 반영한다.
\end{proof}

\begin{example}
\label{sites-example-cover-from-below}
Lemma \ref{sites-lemma-cover-from-below}의 가정들을 만족하는 사상
$f : \mathcal{D} \to \mathcal{C}$를 구성하자.
$\varphi : G \to H$를 유한군들의 사상이라 하자.
가산 $G$-집합과 $H$-집합을 각각 대상으로 하고, 공동으로 전사인
사상들의 가산족을 덮개로 하는 사이트
$\mathcal{D} = \mathcal{T}_G$와 $\mathcal{C} = \mathcal{T}_H$를
생각하자(Example \ref{sites-example-site-on-group}).
$u : \mathcal{T}_H \to \mathcal{T}_G$를 Section
\ref{sites-section-G-sets-morphisms}에서 설명한 함자라 하고
$f : \mathcal{T}_G \to \mathcal{T}_H$를 이에 대응하는 사이트의
사상이라 하자. $\varphi$가 단사이면 모든 가산 $G$-집합은
$G$-집합으로서 가산 $H$-집합의 몫이다
($\varphi$가 단사가 아니면 이는 성립하지 않는다).
따라서 $f$는 Lemma \ref{sites-lemma-cover-from-below}의 가정을 만족한다.
$\mathcal{T}_G$ 위의 층 $\mathcal{F}$가 $G$-집합 $S$에 대응하면
표준적 사상
$$
f^{-1}f_*\mathcal{F} \longrightarrow \mathcal{F}
$$
은 사상
$$
\text{Map}_G(H, S) \longrightarrow S,\quad
a \longmapsto a(1_H)
$$
에 대응한다. $\varphi$가 단사이지만 전사가 아니면 이 사상은
전사이지만(Lemma \ref{sites-lemma-cover-from-below}에 따라 그래야 한다)
일반적으로 단사는 아니다(예를 들어
$G = \{1\}$, $H = \{1, \sigma\}$, $S = \{1, 2\}$로 놓는다).
더욱이 함자 $f_*$는 쌍대등화자 또는 밂과 가환하지 않는다
($G = \{1\}$이고 $H = \{1, \sigma\}$인 경우).
\end{example}





\section{거의 쌍대연속인 함자}
\label{sites-section-almost-cocontinuous}

\noindent
$\mathcal{C}$를 사이트라 하자. 범주 $\textit{PSh}(\mathcal{C})$는
시초대상을 갖는데, 이는 $\mathcal{C}$의 각 대상에 공집합을 대응시키는
준층이다. 이 준층을 $\emptyset$로 나타내자. 층화의 성질에 의해
$\emptyset$의 층화 $\emptyset^\#$는 $\mathcal{C}$ 위의 층들의 범주
$\Sh(\mathcal{C})$의 시초대상이다.

\begin{definition}
\label{sites-definition-empty}
$\mathcal{C}$를 사이트라 하자. $\mathcal{C}$의 대상 $U$에 대하여
$\emptyset^\# \to h_U^\#$가 층의 동형사상이면 그 대상이
{\it 층론적으로 공이다}라고 한다.
\end{definition}

\noindent
다음 보조정리는 이 개념을 더 명시적으로 나타낸다.

\begin{lemma}
\label{sites-lemma-characterize-empty}
$\mathcal{C}$를 사이트라 하고 $U$를 $\mathcal{C}$의 대상이라 하자.
다음 조건들은 서로 동치이다.
\begin{enumerate}
\item $U$는 층론적으로 공이다.
\item 모든 층 $\mathcal{F}$에 대하여 $\mathcal{F}(U)$는 한원소집합이다.
\item $\emptyset^\#(U)$는 한원소집합이다.
\item $\emptyset^\#(U)$는 공집합이 아니다.
\item 공족은 $\mathcal{C}$에서 $U$의 덮개이다.
\end{enumerate}
더욱이 $U$가 층론적으로 공이면, $\mathcal{C}$의 임의의 사상
$U' \to U$에 대하여 대상 $U'$도 층론적으로 공이다.
\end{lemma}

\begin{proof}
임의의 층 $\mathcal{F}$에 대하여
$\mathcal{F}(U) = \Mor_{\Sh(\mathcal{C})}(h_U^\#, \mathcal{F})$이다.
따라서 (1)과 (2)가 동치임을 알 수 있다.
(2)는 (3)을 함의하고 (3)은 (4)를 함의함이 분명하다.
$U$의 모든 덮개가 공집합이 아닌 족으로 주어진다면, 플러스 구성의
정의에 의해 $\emptyset^+(U)$는 공집합이다.
$\emptyset$는 분리된 준층이므로
$\emptyset^+ = \emptyset^\#$임에 유의하자. Theorem
\ref{sites-theorem-plus}를 보라.
따라서 (4)가 (5)를 함의함을 알 수 있다. (5)가 성립하면
$\mathcal{F}$에 대한 층 조건에 의해 모든 층 $\mathcal{F}$에 대하여
$\mathcal{F}(U)$는 한원소집합이다. Remark
\ref{sites-remark-sheaf-condition-empty-covering}을 보라.
따라서 (5)는 (2)를 함의하고 (1) -- (5)는 서로 동치이다.
% P01-ERRATA: P01-E0147
보조정리의 마지막 주장은 Definition \ref{sites-definition-site}의
공리 (3)을 $U$의 공 덮개에 적용하면 따른다.
\end{proof}

\begin{definition}
\label{sites-definition-almost-cocontinuous}
$\mathcal{C}$와 $\mathcal{D}$를 사이트라 하고
$u : \mathcal{C} \to \mathcal{D}$를 함자라 하자.
$\mathcal{C}$의 모든 대상 $U$와 모든 덮개
$\{V_j \to u(U)\}_{j \in J}$에 대하여 $\mathcal{C}$의 덮개
$\{U_i \to U\}_{i \in I}$가 존재해서 각 $I$ 안의 $i$에 대하여
다음 두 조건 가운데 적어도 하나가 성립하면 $u$가
{\it 거의 쌍대연속이다}라고 한다.
\begin{enumerate}
\item $u(U_i)$는 층론적으로 공이다.
\item 어떤 $j \in J$에 대하여 사상 $u(U_i) \to u(U)$는
$V_j$를 거쳐 인수분해된다.
\end{enumerate}
\end{definition}

\noindent
이 정의의 동기는 위상공간의 닫힌 몰입 $i : Z \to X$에서 나온다.
Example \ref{sites-example-closed-map-cocontinuous-false}에서 논의했듯이
연속 함자 $X_{Zar} \to Z_{Zar}$, $U \mapsto Z \cap U$는
쌍대연속이 아니다. 그러나 위에서 정의한 의미에서는 거의 쌍대연속이다.
% P01-ERRATA: P01-E0148
$i_*$는 집합의 층에서는 완전하지 않지만 아벨군의 층에서는 완전함을
안다. Sheaves, Remark \ref{sheaves-remark-i-star-not-exact}를 보라.
이는 일반적인 연속이고 거의 쌍대연속인 함자에 대해서도 성립한다.

\begin{lemma}
\label{sites-lemma-almost-cocontinuous-sheafification}
$\mathcal{C}$와 $\mathcal{D}$를 사이트라 하고
$u : \mathcal{C} \to \mathcal{D}$를 함자라 하자.
$u$가 연속이고 거의 쌍대연속이라고 가정하자.
$\mathcal{G}$를 $\mathcal{D}$ 위의 준층이라 하고, $V$가
층론적으로 공일 때마다 $\mathcal{G}(V)$가 한원소집합이라고 하자.
그러면 $(u^p\mathcal{G})^\# = u^p(\mathcal{G}^\#)$이다.
\end{lemma}

\begin{proof}
$U \in \Ob(\mathcal{C})$라 하자.
$(u^p\mathcal{G})^\#(U) = u^p(\mathcal{G}^\#)(U)$임을 보여야 한다.
$\mathcal{G}^+$도 보조정리의 가정을 만족하는 준층이므로
$(u^p\mathcal{G})^+(U) = u^p(\mathcal{G}^+)(U)$임을 보이면 충분하다.
다음이 성립한다.
$$
u^p(\mathcal{G}^+)(U) =
\mathcal{G}^+(u(U)) =
\colim_\mathcal{V} \check H^0(\mathcal{V}, \mathcal{G})
$$
여기서 쌍대극한은 $\mathcal{D}$ 안의 $u(U)$의 덮개
$\mathcal{V}$들 위에서 취한다. 한편 다음을 알 수 있다.
$$
u^p(\mathcal{G})^+(U) =
\colim_\mathcal{U} \check H^0(u(\mathcal{U}), \mathcal{G})
$$
여기서 쌍대극한은 $\mathcal{C}$ 안의 $U$의 덮개
$\mathcal{U} = \{U_i \to U\}_{i \in I}$들의 범주 위에서 취하고,
$u(\mathcal{U}) = \{u(U_i) \to u(U)\}_{i \in I}$이다.
$u$가 연속이라는 조건은 각 $u(\mathcal{U})$가 덮개라는 뜻이다.
$I = I_1 \amalg I_2$로 쓰자. 여기서
$$
I_2 = \{i \in I \mid u(U_i)\text{ is sheaf theoretically empty}\}
$$
이다.
% P01-ERRATA: P01-E0149
그러면 $u(\mathcal{U})' = \{u(U_i) \to u(U)\}_{i \in I_1}$은
여전히 그 대상의 덮개이다. 다른 각 조각은 공족으로 덮을 수 있으므로
Definition \ref{sites-definition-site}의 공리 (2)에 의해 버릴 수 있기 때문이다.
더욱이 $\mathcal{G}$에 대한 가정에 의해
$\check H^0(u(\mathcal{U}), \mathcal{G}) =
\check H^0(u(\mathcal{U})', \mathcal{G})$이다.
마지막으로 $u$가 거의 쌍대연속이라는 조건에 의해, $u(U)$의 모든 덮개
$\mathcal{V}$에 대하여 $U$의 덮개 $\mathcal{U}$가 존재해서
$u(\mathcal{U})'$가 $\mathcal{V}$를 세분한다.
따라서 위에 표시한 두 쌍대극한은 원하는 대로 같은 값을 갖는다.
\end{proof}

\begin{lemma}
\label{sites-lemma-continuous-almost-cocontinuous}
$\mathcal{C}$와 $\mathcal{D}$를 사이트라 하고
$u : \mathcal{C} \to \mathcal{D}$를 함자라 하자.
$u$가 연속이고 거의 쌍대연속이라고 가정하자.
그러면 $u^s = u^p : \Sh(\mathcal{D}) \to \Sh(\mathcal{C})$는
밂 및 쌍대등화자와 가환한다(더 일반적으로 유한 연결 쌍대극한과 가환한다).
\end{lemma}

\begin{proof}
$\mathcal{I}$를 유한 연결 지표범주라 하자.
% P01-ERRATA: P01-E0150
$\mathcal{I} \to \Sh(\mathcal{D})$,
$i \mapsto \mathcal{G}_i$를 도표라 하자.
이 도표의 쌍대극한은 준층의 범주에서 취한 쌍대극한을 층화한 것임을 안다.
Lemma \ref{sites-lemma-colimit-sheaves}를 보라.
준층의 범주에서의 쌍대극한을 $\colim^{Psh}$로 나타내자.
$\mathcal{I}$가 유한이고 연결이므로
$\colim^{Psh}_i \mathcal{G}_i$는 Lemma
\ref{sites-lemma-almost-cocontinuous-sheafification}의 가정들을 만족하는
준층이다(한원소집합들의 유한 연결 쌍대극한은 한원소집합이기 때문이다).
% P01-ERRATA: P01-E0151
따라서 그 보조정리에 의해 다음을 얻는다.
\begin{align*}
u^s(\colim_i \mathcal{G}_i) & =
u^s((\colim^{Psh}_i \mathcal{G}_i)^\#) \\
& = (u^p(\colim^{Psh}_i \mathcal{G}_i))^\# \\
& = (\colim^{PSh}_i u^p(\mathcal{G}_i))^\# \\
& = \colim_i u^s(\mathcal{G}_i)
\end{align*}
이는 원하는 바이다.
\end{proof}

\begin{lemma}
\label{sites-lemma-morphism-of-sites-almost-cocontinuous}
$f : \mathcal{D} \to \mathcal{C}$를 연속 함자
$u : \mathcal{C} \to \mathcal{D}$에 대응하는 사이트의 사상이라 하자.
$u$가 거의 쌍대연속이면 $f_*$는 밂 및 쌍대등화자와 가환한다
(더 일반적으로 유한 연결 쌍대극한과 가환한다).
\end{lemma}

\begin{proof}
이는 Lemma \ref{sites-lemma-continuous-almost-cocontinuous}의 특별한 경우이다.
\end{proof}








\section{부분토포스}
\label{sites-section-subtopoi}

\noindent
정의는 다음과 같다.

\begin{definition}
\label{sites-definition-embedding}
$\mathcal{C}$와 $\mathcal{D}$를 사이트라 하자.
토포스의 사상 $f : \Sh(\mathcal{D}) \to \Sh(\mathcal{C})$에 대하여
$f_*$가 충실충만이면 이를 {\it 매장}이라 한다.
\end{definition}

\noindent
Lemma \ref{sites-lemma-exactness-properties}에 의하면 이는
$\mathcal{D}$ 위의 모든 층 $\mathcal{F}$에 대하여 수반사상
$f^{-1}f_*\mathcal{F} \to \mathcal{F}$가 동형사상이라고 요구하는
것과 동치이다.

\begin{definition}
\label{sites-definition-subtopos}
$\mathcal{C}$를 사이트라 하자. 엄밀히 충만한 부분범주
$E \subset \Sh(\mathcal{C})$에 대하여, 어떤 토포스의 매장
$f : \Sh(\mathcal{D}) \to \Sh(\mathcal{C})$가 존재해서
$E$가 함자 $f_*$의 본질적 상과 같으면 그 부분범주를
{\it 부분토포스}라 한다.
\end{definition}

\noindent
다음 보조정리에서 구성하는 부분토포스는 뒤의 정의에서 “열린”
부분토포스라고 부를 것이다.

\begin{lemma}
\label{sites-lemma-open-subtopos}
$\mathcal{C}$를 사이트라 하고 $\mathcal{F}$를 $\mathcal{C}$ 위의
층이라 하자. 다음은 서로 동치이다.
\begin{enumerate}
\item $\mathcal{F}$는 $\Sh(\mathcal{C})$의 종대상의 부분대상이다.
\item 토포스 $\Sh(\mathcal{C})/\mathcal{F}$는
$\Sh(\mathcal{C})$의 부분토포스이다.
\end{enumerate}
\end{lemma}

\begin{proof}
Lemma \ref{sites-lemma-localize-topos}에서
$\Sh(\mathcal{C})/\mathcal{F}$가 토포스임을 보았다.
그 증명을 상기하자. 먼저 Lemma \ref{sites-lemma-topos-good-site}를
적용하면, $\mathcal{C}$가 준표준 위상, 섬유곱, 종대상 $X$,
그리고 $\mathcal{F} = h_U$인 대상 $U$를 갖는 사이트라고 가정할 수 있다.
Lemma \ref{sites-lemma-localize-topos}의 증명에 따르면 토포스의 사상
$j_\mathcal{F} : \Sh(\mathcal{C})/\mathcal{F} \to \Sh(\mathcal{C})$는
(어떤 식별들을 거치면) 국소화 사상
$j_U : \Sh(\mathcal{C}/U) \to \Sh(\mathcal{C})$와 같다.

\medskip\noindent
(2)를 가정하자. 이는 $\mathcal{C}/U$ 위의 모든 층 $\mathcal{G}$에
대하여 $j_U^{-1}j_{U, *}\mathcal{G} \to \mathcal{G}$가
동형사상이라는 뜻이다. $\mathcal{C}/U$의 임의의 대상 $Z/U$에 대하여,
Lemma \ref{sites-lemma-localize-given-products}에 의해
$$
(j_{U, *}h_{Z/U})(U) = \Mor_{\mathcal{C}/U}(U \times_X U/U, Z/U)
$$
이다. 위 등식에서 $\mathcal{G} = h_{Z/U}$로 놓으면 모든 $Z/U$에
대하여
$$
\Mor_{\mathcal{C}/U}(U \times_X U/U, Z/U) = \Mor_{\mathcal{C}/U}(U, Z/U)
$$
를 얻는다. 요네다 보조정리(Categories, Lemma
\ref{categories-lemma-yoneda})에 의해 이는 $U \times_X U = U$를
함의한다. Categories, Lemma
\ref{categories-lemma-characterize-mono-epi}에 의해
$U \to X$는 단사사상이다. 달리 말하면 (1)이 성립한다.

\medskip\noindent
(1)을 가정하자. 그러면 Lemma \ref{sites-lemma-restrict-back}에 의해
$j_U^{-1} j_{U, *} = \text{id}$이다.
\end{proof}

\begin{definition}
\label{sites-definition-open-subtopos}
$\mathcal{C}$를 사이트라 하자. 엄밀히 충만한 부분범주
$E \subset \Sh(\mathcal{C})$에 대하여, $\Sh(\mathcal{C})$의
종대상의 부분층 $\mathcal{F}$가 존재해서 $E$가 Lemma
\ref{sites-lemma-open-subtopos}에서 설명한 부분토포스
$\Sh(\mathcal{C})/\mathcal{F}$와 같으면 그 부분범주를
{\it 열린 부분토포스}라 한다.
\end{definition}

\noindent
이는 $\Sh(\mathcal{C})$의 열린 부분토포스들의 모임과
$\Sh(\mathcal{C})$의 종대상의 부분대상들의 집합 사이에 전단사가
있다는 뜻이다. 열린 부분토포스가 주어지면 “닫힌” 여집합이 있다.

\begin{lemma}
\label{sites-lemma-closed-subtopos}
$\mathcal{C}$를 사이트라 하자. $\mathcal{F}$를
$\Sh(\mathcal{C})$의 종대상 $*$의 부분층이라 하자.
$\mathcal{F} \times \mathcal{G} \to \mathcal{F}$가 동형사상인
층 $\mathcal{G}$들로 이루어진 충만한 부분범주는
$\Sh(\mathcal{C})$의 부분토포스이다.
\end{lemma}

\begin{proof}
Lemma \ref{sites-lemma-topos-good-site}를 적용하면 $\mathcal{C}$가 그
보조정리에 열거된 성질들을 갖는 사이트라고 가정할 수 있다.
특히 $\mathcal{C}$는 종대상 $X$를 가지며(따라서 $* = h_X$),
$\mathcal{F} = h_U$인 대상 $U$를 갖는다.

\medskip\noindent
범주로서 $\mathcal{D} = \mathcal{C}$로 놓되, 사상족
% P01-ERRATA: P01-E0152
$\{V_j \to V\}$가 덮개라는 것을
$\{V_i \to V\} \cup \{U \times_X V \to V\}$가 덮개라는 뜻으로
정의하자. $\mathcal{C}$의 선택에 의해 이는 정확히
$$
h_{U \times_X V} \amalg \coprod h_{V_i} \longrightarrow h_V
$$
가 전사라는 뜻이다. $\mathcal{D}$가 사이트, 즉 덮개들이 Definition
\ref{sites-definition-site}의 조건 (1), (2), (3)을 만족한다고 주장한다.
조건 (1)은 성립한다. 조건 (2)에 대해서는
$\{V_i \to V\}$와 $\{V_{ij} \to V_i\}$가 $\mathcal{D}$의
덮개라고 하자. 그러면 합성
$$
\coprod \left(
h_{U \times_X V_i} \amalg \coprod h_{V_{ij}}
\right) \longrightarrow
h_{U \times_X V} \amalg \coprod h_{V_i} \longrightarrow h_V
$$
은 전사이다. 각 사상 $U \times_X V_i \to V$가
$U \times_X V$를 거쳐 인수분해되므로
$$
h_{U \times_X V} \amalg \coprod h_{V_{ij}} \longrightarrow h_V
$$
는 전사이다. 즉 $\{V_{ij} \to V\}$는 $\mathcal{D}$에서 $V$의
덮개이다. 층의 전사 사상을 밑변환하면 다시 전사이므로
조건 (3)도 같은 방식으로 따른다.

\medskip\noindent
(항등) 함자 $u : \mathcal{C} \to \mathcal{D}$는 연속이며 섬유곱 및
종대상과 가환함에 유의하자. 따라서 사이트의 사상
$f : \mathcal{D} \to \mathcal{C}$를 얻는다
(Proposition \ref{sites-proposition-get-morphism}).
$f_*$는 바탕 준층 위의 항등 함자이므로 충실충만하다.
증명을 끝내려면 $f_*$의 본질적 상이 보조정리에서 골라낸 충만한 부분범주
$E \subset \Sh(\mathcal{C})$임을 보여야 한다.
이를 위해 $\mathcal{G} \in \Ob(\Sh(\mathcal{C}))$가 $E$에 속할
필요충분조건이 $\mathcal{C}$의 모든 대상 $V$에 대하여
$\mathcal{G}(U \times_X V)$가 한원소집합인 것임에 유의하자.
따라서 이러한 층은 $\mathcal{D}$의 모든 덮개에 대한 층 조건을
만족한다(증명 생략). 역으로 $\mathcal{G}$가 $\mathcal{D}$의
모든 덮개에 대한 층 조건을 만족하면
$\mathcal{G}(U \times_X V)$는 한원소집합이다.
$\mathcal{D}$에서 대상 $U \times_X V$는 공 덮개로 덮이기 때문이다.
\end{proof}

\begin{definition}
\label{sites-definition-closed-subtopos}
$\mathcal{C}$를 사이트라 하자. 엄밀히 충만한 부분범주
$E \subset \Sh(\mathcal{C})$에 대하여, $\Sh(\mathcal{C})$의
종대상의 부분층 $\mathcal{F}$가 존재해서 $E$가 Lemma
\ref{sites-lemma-closed-subtopos}에서 설명한 부분토포스이면
그 부분범주를 {\it 닫힌 부분토포스}라 한다.
\end{definition}

\noindent
이제 토포스의 닫힌 몰입과 열린 몰입의 의미를 정의할 수 있다.

\begin{definition}
\label{sites-definition-immersion-topoi}
$f : \Sh(\mathcal{D}) \to \Sh(\mathcal{C})$를 토포스의 사상이라 하자.
\begin{enumerate}
\item $f$가 매장이고 $f_*$의 본질적 상이 열린 부분토포스이면
$f$를 {\it 열린 몰입}이라 한다.
\item $f$가 매장이고 $f_*$의 본질적 상이 닫힌 부분토포스이면
$f$를 {\it 닫힌 몰입}이라 한다.
\end{enumerate}
\end{definition}

\begin{lemma}
\label{sites-lemma-closed-immersion}
$i : \Sh(\mathcal{D}) \to \Sh(\mathcal{C})$를 토포스의 닫힌 몰입이라
하자. 그러면 $i_*$는 충실충만하고, 전사를 전사로 보내며,
쌍대등화자 및 밂과 가환하고, 단사성, 전사성 및 전단사성을 반영한다.
\end{lemma}

\begin{proof}
$\mathcal{F}$를 $\Sh(\mathcal{C})$의 종대상 $*$의 부분층이라 하고
$E \subset \Sh(\mathcal{C})$를
$\mathcal{F} \times \mathcal{G} \to \mathcal{F}$가 동형사상인
$\mathcal{G}$들로 이루어진 충만한 부분범주라 하자.
Lemma \ref{sites-lemma-closed-subtopos}에 의해 함자 $i_*$는 포함함자
$\iota : E \to \Sh(\mathcal{C})$와 동형이다.

\medskip\noindent
$j_{\mathcal{F}} : \Sh(\mathcal{C})/\mathcal{F} \to \Sh(\mathcal{C})$를
국소화 함자라 하자(Lemma \ref{sites-lemma-localize-topos}).
$E$는 $j_\mathcal{F}^{-1}\mathcal{G} = *$인 층
$\mathcal{G}$들의 모임으로도 기술할 수 있음에 유의하자.

\medskip\noindent
% P01-ERRATA: P01-E0153
$a, b : \mathcal{G}_1 \to \mathcal{G}_2$를 $E$의 두 사상이라 하자.
$\iota$가 쌍대등화자와 가환함을 증명하려면
$\Sh(\mathcal{C})$에서 $a$, $b$의 쌍대등화자가 $E$에 속함을 보이면
충분하다. 이는 두 사상 $* \to *$의 쌍대등화자가 $*$이고
$j_\mathcal{F}^{-1}$가 완전하기 때문에 분명하다. 밂도 마찬가지이다.

\medskip\noindent
따라서 $i_*$는 Lemma \ref{sites-lemma-exactness-properties}의 성질
(5), (6), (7)을 만족하고, 그러므로 사상 $i$는 그 보조정리에서
언급한 모든 성질, 특히 이 보조정리에서 언급한 성질들을 만족한다.
\end{proof}







\section{대수적 구조의 층}
\label{sites-section-sheaves-algebraic-structures}

\noindent
Sheaves, Section \ref{sheaves-section-algebraic-structures}에서
대수적 구조의 유형을 쌍 $(\mathcal{A}, s)$로 도입했다.
여기서 $\mathcal{A}$는 범주이고
$s : \mathcal{A} \to \textit{Sets}$는 다음을 만족하는 함자이다.
\begin{enumerate}
\item $s$는 충실하다.
\item $\mathcal{A}$는 극한들을 가지며 $s$는 극한들과 가환한다.
\item $\mathcal{A}$는 여과 쌍대극한들을 가지며 $s$는 이들과 가환한다.
\item $s$는 동형사상을 반영한다.
\end{enumerate}
이러한 대수적 구조의 유형에 대하여, 공간 $X$ 위에서
$\mathcal{A}$에 값을 갖는 준층 $\mathcal{F}$가 층일 필요충분조건은
그에 대응하는 집합의 준층이 층인 것이다. 또한 줄기의 개념을 구체화했고,
연속 사상 $f : X \to Y$가 주어졌을 때 대수적 구조의 층에 대한
수반함자 직접상과 역상을 정의했다.
% P01-ERRATA: P01-E0154
이들은 바탕 집합의 층에 대한 직접상 및 역상과 일치한다.
그 밖에도 대수적 구조의 층을 기저에서 공간의 모든 열린집합으로 확장하는
과정은 예상대로 작동한다.

\medskip\noindent
이 내용의 일부는 사이트와 층의 설정에서도 그대로 작동한다.
$(\mathcal{A}, s)$를 대수적 구조의 유형이라 하자.
$\mathcal{C}$를 사이트라 하자.
$\textit{PSh}(\mathcal{C}, \mathcal{A})$ 및
$\Sh(\mathcal{C}, \mathcal{A})$로 각각 $\mathcal{C}$ 위에서
$\mathcal{A}$에 값을 갖는 준층 및 층의 범주를 나타내자.
\begin{itemize}
\item[] ($\alpha$) $\mathcal{A}$에 값을 갖는 준층이 층일
필요충분조건은 그 바탕 집합의 준층이 층인 것이다.
Sheaves, Lemma \ref{sheaves-lemma-sheaves-structure}의 증명을 보라.
\item[] ($\beta$) $\mathcal{A}$에 값을 갖는 준층 $\mathcal{F}$가
주어지면 준층 ${\mathcal{F}}^\# = (\mathcal{F}^+)^+$는 층이다.
이는 층화 과정의 쌍대극한들이 여과되어 있고, 더 나아가 유향 집합 위의
쌍대극한이기 때문이다(Section \ref{sites-section-sheafification}, 특히
Lemma \ref{sites-lemma-sheafification-exact}의 증명 참조).
또한 $s$는 여과 쌍대극한과 가환한다.
\item[] ($\gamma$) 다음 가환도표를 얻는다.
$$
\xymatrix{
\Sh(\mathcal{C}, \mathcal{A}) \ar@<1ex>[r] \ar[d]_s &
\textit{PSh}(\mathcal{C}, \mathcal{A}) \ar@<1ex>[l]^{{}^\#} \ar[d]^s\\
\Sh(\mathcal{C}) \ar@<1ex>[r] &
\textit{PSh}(\mathcal{C}) \ar@<1ex>[l]
}
$$
\item[] ($\delta$) $\mathcal{F} = \mathcal{F}^\#$일 필요충분조건은
$\mathcal{F}$가 대수적 구조의 층인 것이다.
\item[] ($\epsilon$) 함자 ${}^\#$은 포함함자의 수반이다.
$$
\Mor_{\textit{PSh}(\mathcal{C}, \mathcal{A})}(\mathcal{G}, \mathcal{F})
=
\Mor_{\Sh(\mathcal{C}, \mathcal{A})}(\mathcal{G}^\#, \mathcal{F})
$$
증명은 Proposition \ref{sites-proposition-sheafification-adjoint}의 증명과 같다.
\item[] ($\zeta$) 함자 $\mathcal{F} \mapsto \mathcal{F}^\#$은
좌완전하다. 증명은 Lemma \ref{sites-lemma-sheafification-exact}의 증명과 같다.
\end{itemize}

\begin{definition}
\label{sites-definition-pushforward-algebraic-structures}
$f : \mathcal{D} \to \mathcal{C}$를 함자
$u : \mathcal{C} \to \mathcal{D}$로 주어지는 사이트의 사상이라 하자.
대수적 구조의 준층에 대한 {\it 직접상} 함자를
$u^p\mathcal{F}(U) = \mathcal{F}(uU)$로 정의하고,
대수적 구조의 층에 대해서도 같은 규칙, 즉
$f_*\mathcal{F}(U) = \mathcal{F}(uU)$로 정의한다.
\end{definition}

\noindent
% P01-ERRATA: P01-E0155
문제는 역상을 정의하려고 할 때 생긴다.
그 이유는 Section \ref{sites-section-functoriality-PSh}에서 함자 $u_p$를
정의하는 쌍대극한들이 여과되어 있지 않을 수 있기 때문이다.
따라서 일반적으로는 위의 공리들만으로 대수적 구조의 (준)층의 역상을
정의하기에 충분하지 않다. 그럼에도 거의 모든 경우에는 다음 보조정리로
대수적 구조의 (준)층의 직접상과 역상을 정의하기에 충분하다.

\begin{lemma}
\label{sites-lemma-push-pull-good-case}
함자 $u : \mathcal{C} \to \mathcal{D}$가 Proposition
\ref{sites-proposition-get-morphism}의 가정들을 만족하여 사이트의 사상
$f : \mathcal{D} \to \mathcal{C}$를 준다고 하자.
이 경우 역상 함자 $f^{-1}$ (각각 $u_p$)와 직접상 함자
$f_*$ (각각 $u^p$)는 대수적 구조의 층(각각 준층)의 범주 위의
수반함자쌍으로 확장된다. 또한 이 함자들은 바탕 집합의 층
(각각 준층)을 취하는 것과 가환한다.
\end{lemma}

\begin{proof}
위에서 $f_* = u^p$를 정의했다.
Proposition \ref{sites-proposition-get-morphism}의 증명 과정에서, 그 명제의
가정 아래에는 $u_p$를 정의하는 모든 쌍대극한이 여과되어 있음을 보았다.
따라서 대수적 구조의 유형의 정의에 의해 대수적 구조의 준층 위의 함자
$u_p$를 집합의 준층에서와 정확히 같은 쌍대극한으로 정의할 수 있다.
$u_p$와 $u^p$의 수반성은 Lemma \ref{sites-lemma-adjoints-u}의 증명과
정확히 같은 방식으로 증명한다. 위의 대수적 구조의 준층의 층화에 관한
논의에 따라 $f^{-1}(\mathcal{F}) = (u_p\mathcal{F})^\#$로
정의할 수 있다.
\end{proof}

\noindent
일반적인 사이트의 사상, 더 일반적으로 토포스의 임의의 사상에 대하여
역상과 직접상을 다루는 방법을 간략히 논의하자.

\medskip\noindent
$\mathcal{C}$를 사이트라 하자.
$\mathcal{A} = \textit{Ab}$인 경우에는 $\mathcal{C}$ 위의
아벨 (준)층을 사중항 $(\mathcal{F}, +, 0, i)$로 생각할 수 있다.
여기서 데이터는 다음과 같다.
\begin{enumerate}
\item[(D1)] $\mathcal{F}$는 집합의 층이다.
\item[(D2)] $+ : \mathcal{F} \times \mathcal{F} \to \mathcal{F}$는
집합의 층의 사상이다.
\item[(D3)] $0 : * \to \mathcal{F}$는 한원소 층(Example
\ref{sites-example-singleton-sheaf} 참조)에서 $\mathcal{F}$로 가는 사상이다.
\item[(D4)] $i : \mathcal{F} \to \mathcal{F}$는 집합의 층의 사상이다.
\end{enumerate}
이 데이터는 다음 공리들을 만족해야 한다.
\begin{enumerate}
\item[(A1)] $+$는 결합적이고 가환적이다.
\item[(A2)] $0$은 $+$의 단위원이다.
\item[(A3)] $+ \circ (1, i) = 0 \circ (\mathcal{F} \to *)$이다.
\end{enumerate}
Sheaves, Lemma \ref{sheaves-lemma-abelian-presheaves}와 비교하라.
$f : \mathcal{D} \to \mathcal{C}$를 사이트의 사상이라 하자.
$f^{-1}$가 완전하므로
$f^{-1}* = *$이고
$f^{-1}(\mathcal{F} \times \mathcal{F}) =
f^{-1}\mathcal{F} \times f^{-1}\mathcal{F}$임에 유의하자.
따라서 $f^{-1}\mathcal{F}$를 단순히 사중항
$(f^{-1}\mathcal{F}, f^{-1}+, f^{-1}0, f^{-1}i)$로 정의할 수 있다.
$f^{-1}$는 유한 극한들과 가환하는 함자이므로 공리들이 보존된다.
마지막으로 $f_*$와 $f^{-1}$가 평소와 같이 수반임을 확인하는 것은
어렵지 않다.

\medskip\noindent
\cite{SGA4}에서는 이 방법을 사용한다. 거기서는
% P01-ERRATA: P01-E0156
``{\it esp\`ece the structure alg\'ebrique $\ll$d\'efinie
par limites projectives finie$\gg$}''라고 부르는 것을 도입한다.
이러한 esp\`ece에 대해서는 위에 설명한 방법을 사용하여 위와 같은
수반함자쌍 $f^{-1}$와 $f_*$를 정의할 수 있다.
이는 실제로 마주치는 대부분의 대수적 구조에 분명히 적용된다.
이 구성을 형식화하는 대신, 이 방법이 적용되는 대수적 구조를
(각 경우에 확인한다는 전제 아래) 열거하겠다.
실제로 이 방법은 토포스의 임의의 사상에 적용된다.

\begin{proposition}
\label{sites-proposition-functoriality-algebraic-structures-topoi}
\begin{slogan}
토포스의 사상은 대수적 구조를 보존한다.
\end{slogan}
$\mathcal{C}$와 $\mathcal{D}$를 사이트라 하자.
$f = (f^{-1}, f_*)$를 $\Sh(\mathcal{D}) \to \Sh(\mathcal{C})$인
토포스의 사상이라 하자. 위에서 도입한 방법은 다음 대수적 구조의
유형들에 대하여, 바탕 집합의 층을 취하는 것과 호환되는
수반함자쌍 $(f^{-1}, f_*)$를 대수적 구조의 층 위에 준다.
\begin{enumerate}
\item 점붙인 집합,
\item 아벨군,
\item 군,
\item 모노이드,
\item 환,
\item 고정된 환 위의 가군,
\item 고정된 체 위의 Lie 대수.
\end{enumerate}
더욱이 각 경우에 위에서 ($\alpha$), ($\beta$), ($\gamma$),
($\delta$), ($\epsilon$), ($\zeta$)로 표시한 결과들이 성립한다.
\end{proposition}

\begin{proof}
마지막 주장은 열거한 각 범주가 자명한 망각 함자와 함께, 이 절의
시작에서 설명한 의미의 대수적 구조의 유형이기 때문에 성립한다.
Sheaves, Lemma \ref{sheaves-lemma-list-algebraic-structures}를 보라.

\medskip\noindent
(2)의 증명. 명제 앞의 논의에서 설명했듯이 아벨군의 층을 사중항
$(\mathcal{F}, +, 0, i)$로 생각한다.
$(\mathcal{F}, +, 0, i)$가 $\mathcal{C}$ 위에 있으면 그 역상은
$(f^{-1}\mathcal{F}, f^{-1}+, f^{-1}0, f^{-1}i)$로 정의한다.
$(\mathcal{G}, +, 0, i)$가 $\mathcal{D}$ 위에 있으면 그 직접상은
$(f_*\mathcal{G}, f_*+, f_*0, f_*i)$로 정의한다.
$f_*(\mathcal{G} \times \mathcal{G}) =
f_*\mathcal{G} \times f_*\mathcal{G}$이므로 이는 잘 정의된다.
수반성은 집합에 바탕을 둔 함자들의 수반성에서 따른다. 실제로
$$
\Mor_{\textit{Ab}(\mathcal{C})}
((\mathcal{F}_1, +, 0, i), (\mathcal{F}_2, +, 0, i))
=
\left\{
\begin{matrix}
\varphi \in \Mor_{\Sh(\mathcal{C})}
(\mathcal{F}_1, \mathcal{F}_2) \\
\varphi \text{ is compatible with }+, 0, i
\end{matrix}
\right\}
$$
이다. 세부사항은 독자에게 맡긴다.

\medskip\noindent
이 방법은 단위원을 갖는 환의 층을, 임의의 기초 대수학 책에서 찾을 수
있는 공리들의 목록을 만족하는 육중항
$(\mathcal{O}, + , 0, i, \cdot, 1)$로 생각하면 환의 층에도 적용된다.

\medskip\noindent
점붙인 집합의 층은 쌍 $(\mathcal{F}, p)$이다.
여기서 $\mathcal{F}$는 집합의 층이고
$p : * \to \mathcal{F}$는 집합의 층의 사상이다.

\medskip\noindent
군의 층은 적절한 공리들을 갖춘 사중항
$(\mathcal{F}, \cdot, 1, i)$로 주어진다.

\medskip\noindent
% P01-ERRATA: P01-E0157
모노이드의 층은 적절한 공리들을 갖춘 쌍
$(\mathcal{F}, \cdot)$로 주어진다.

\medskip\noindent
% P01-ERRATA: P01-E0158
$R$을 환이라 하자. $R$-가군의 층은 오중항
$(\mathcal{F}, +, 0, i, \{\lambda_r\}_{r \in R})$로 주어진다.
여기서 사중항 $(\mathcal{F}, +, 0, i)$는 위와 같은 아벨군의 층이고,
$\lambda_r : \mathcal{F} \to \mathcal{F}$는 다음을 만족하는
집합의 층의 사상족이다.
$\lambda_r \circ 0 = 0$,
$\lambda_r \circ + = + \circ (\lambda_r, \lambda_r)$,
% P01-ERRATA: P01-E0159
$\lambda_{r + r'} =
+ \circ \lambda_r \times \lambda_{r'} \circ (\text{id}, \text{id})$,
$\lambda_{rr'} = \lambda_r \circ \lambda_{r'}$,
$\lambda_1 = \text{id}$, $\lambda_0 = 0 \circ (\mathcal{F} \to *)$.
\end{proof}

\noindent
환의 층 위의 가군의 층의 범주는 Modules on Sites, Section
\ref{sites-modules-section-sheaves-modules}에서 논의할 것이다.

\begin{remark}
\label{sites-remark-no-pullback-presheaves}
$\mathcal{C}$와 $\mathcal{D}$를 사이트라 하자.
$u : \mathcal{D} \to \mathcal{C}$를 사이트의 사상
$\mathcal{C} \to \mathcal{D}$를 주는 연속 함자라 하자.
아벨군의 경우에도 아벨군의 준층에 대한 역상 함자를 정의하지 않았음에
유의하자. 아벨군의 범주에서는 모든 쌍대극한이 표현 가능하므로,
집합의 준층 위의 $u_p$를 정의할 때 사용한 것과 같은 쌍대극한으로
아벨 준층 위의 함자 $u_p^{ab}$를 정의할 수 있음은 확실하다.
또한 $u_p^{ab}$는 아벨 준층의 범주에서 $u^p$의 수반이다.
그러나 $u_p^{ab}$가 바탕 집합의 준층 위에서 $u_p$와 항상
일치하는 것은 아니다.
\end{remark}












\section{당김 사상}
\label{sites-section-pullback}

\noindent
사이트 $\mathcal{C}$가 종대상을 갖지 않는 경우가 종종 있다.
이 경우 대역 절단 함자를 다음과 같이 정의한다.

\begin{definition}
\label{sites-definition-global-sections}
% P01-ERRATA: P01-E0160
사이트 $\mathcal{C}$ 위의 집합의 준층 $\mathcal{F}$의
{\it 대역 절단들의 집합}은
$$
\Gamma(\mathcal{C}, \mathcal{F}) =
\Mor_{\textit{PSh}(\mathcal{C})}(*, \mathcal{F})
$$
이다. 여기서 $*$는 $\mathcal{C}$ 위의 준층의 범주의 종대상,
즉 모든 대상에 한원소집합을 대응시키는 준층이다.
\end{definition}

\noindent
물론 같은 정의를 층에도 적용할 수 있다.
대역 절단을 계산하는 한 가지 방법은 다음과 같다.

\begin{lemma}
\label{sites-lemma-compute-global-sections}
$\mathcal{C}$를 사이트라 하자.
% P01-ERRATA: P01-E0161
$a, b : V \to U$를 $\mathcal{C}$의 사상들이라 하고,
$$
\xymatrix{
h_V^\# \ar@<1ex>[r] \ar@<-1ex>[r] & h_U^\# \ar[r] & {*}
}
$$
가 $\Sh(\mathcal{C})$에서 쌍대등화자라고 하자.
그러면 $\Gamma(\mathcal{C}, \mathcal{F})$는
$a^*, b^* : \mathcal{F}(U) \to \mathcal{F}(V)$의 등화자이다.
\end{lemma}

\begin{proof}
$\Mor_{\Sh(\mathcal{C})}(h_U^\#, \mathcal{F}) = \mathcal{F}(U)$이므로
정의에서 분명히 따른다.
\end{proof}

\noindent
이제 $f : \Sh(\mathcal{D}) \to \Sh(\mathcal{C})$를 토포스의 사상이라
하자. 그러면 $\mathcal{C}$ 위의 임의의 층 $\mathcal{F}$에 대하여
당김 사상
$$
f^{-1} :
\Gamma(\mathcal{C}, \mathcal{F})
\longrightarrow
\Gamma(\mathcal{D}, f^{-1}\mathcal{F})
$$
이 있다. 실제로 $f^{-1}$가 완전하므로 $*$를 $*$로 보낸다.
따라서 $\mathcal{C}$ 위의 $\mathcal{F}$의 대역 절단 $s$, 즉 층의 사상
$s : * \to \mathcal{F}$를
$f^{-1}s : * = f^{-1}* \to f^{-1}\mathcal{F}$로 당길 수 있다.

\medskip\noindent
이를 조금 일반화하자. $\mathcal{C}$, $\mathcal{D}$ 위의 층들의 쌍
$\mathcal{F}$, $\mathcal{G}$와 사상
$f^{-1}\mathcal{F} \to \mathcal{G}$가 주어졌다고 하자.
그러면 위 구성에 자명한 사상
$\Gamma(\mathcal{D}, f^{-1}\mathcal{F}) \to \Gamma(\mathcal{D}, \mathcal{G})$를
합성하여 사상
$$
\Gamma(\mathcal{C}, \mathcal{F})
\longrightarrow
\Gamma(\mathcal{D}, \mathcal{G})
$$
을 얻는다. 이 사상도 때때로 당김 사상이라 부른다.

\medskip\noindent
자연스럽게 자주 나타나는 조금 더 일반적인 구성은 다음과 같다.
토포스의 사상들의 가환도표
$$
\xymatrix{
\Sh(\mathcal{D}) \ar[rd]_h \ar[rr]_f & &
\Sh(\mathcal{C}) \ar[ld]^g \\
& \Sh(\mathcal{B})
}
$$
가 있다고 하자. 다음으로 $\mathcal{C}$ 위의 층 $\mathcal{F}$가
주어졌다고 하자. 그러면 {\it 당김 사상}
$$
f^{-1} : g_*\mathcal{F} \longrightarrow h_*f^{-1}\mathcal{F}
$$
이 있다. 이는 식별
$g_*f_*f^{-1}\mathcal{F} = h_*f^{-1}\mathcal{F}$와
표준적 사상 $\mathcal{F} \to f_*f^{-1}\mathcal{F}$에 $g_*$를
적용한 사상에서 나오는 사상이다.
$g$가 항등이면 대역 절단 위의 이 사상은 위의 당김 사상과 일치한다.

\medskip\noindent
앞 문단의 상황에서 $\mathcal{C}$, $\mathcal{D}$ 위의 층들의 쌍
$\mathcal{F}$, $\mathcal{G}$와 사상
$f^{-1}\mathcal{F} \to \mathcal{G}$가 주어졌다고 하자.
그러면 위 당김 사상에
$f^{-1}\mathcal{F} \to \mathcal{G}$에 $h_*$를 적용한 사상을
합성하여 사상
$$
g_*\mathcal{F} \longrightarrow h_*\mathcal{G}
$$
을 얻는다. $\mathcal{B}$의 한 대상 위의 절단으로 제한하면
(사이트를 적절히 선택했을 때) 위에서 논의한 대역 절단 위의
“당김 사상”을 다시 얻는다.

\medskip\noindent
더 일반적으로 토포스의 가환도표
$$
\xymatrix{
\Sh(\mathcal{D}) \ar[d]_h \ar[r]_f &
\Sh(\mathcal{C}) \ar[d]^g \\
\Sh(\mathcal{B}) \ar[r]^e & \Sh(\mathcal{A})
}
$$
와 $\mathcal{C}$ 위의 층 $\mathcal{G}$가 주어진 경우가 있다.
그러면 {\it 밑변환 사상}
$$
e^{-1}g_*\mathcal{G} \longrightarrow h_*f^{-1}\mathcal{G}.
$$
이 있다. 실제로 이 사상은 사상
$g_*\mathcal{G} \to e_*h_*f^{-1}\mathcal{G} = (e \circ h)_*f^{-1}\mathcal{G}$의
수반인데, 후자는 방금 설명한 당김 사상이다.

\begin{remark}
\label{sites-remark-compose-base-change}
토포스의 가환도표
$$
\xymatrix{
\Sh(\mathcal{B}') \ar[r]_k \ar[d]_{f'} &
\Sh(\mathcal{B}) \ar[d]^f \\
\Sh(\mathcal{C}') \ar[r]^l \ar[d]_{g'} &
\Sh(\mathcal{C}) \ar[d]^g \\
\Sh(\mathcal{D}') \ar[r]^m &
\Sh(\mathcal{D})
}
$$
를 생각하자. 두 정사각형의 밑변환 사상들을 합성하면 바깥쪽 직사각형의
밑변환 사상을 얻는다. 더 정확히는 합성
\begin{align*}
m^{-1} \circ (g \circ f)_*
& =
m^{-1} \circ g_* \circ f_* \\
& \to g'_* \circ l^{-1} \circ f_* \\
& \to g'_* \circ f'_* \circ k^{-1} \\
& = (g' \circ f')_* \circ k^{-1}
\end{align*}
이 그 직사각형의 밑변환 사상이다. 확인은 생략한다.
\end{remark}

\begin{remark}
\label{sites-remark-compose-base-change-horizontal}
% P01-ERRATA: P01-E0162
환 달린 토포스의 가환도표
$$
\xymatrix{
\Sh(\mathcal{C}'') \ar[r]_{g'} \ar[d]_{f''} &
\Sh(\mathcal{C}') \ar[r]_g \ar[d]_{f'} &
\Sh(\mathcal{C}) \ar[d]^f \\
\Sh(\mathcal{D}'') \ar[r]^{h'} &
\Sh(\mathcal{D}') \ar[r]^h &
\Sh(\mathcal{D})
}
$$
를 생각하자. 두 정사각형의 밑변환 사상들을 합성하면 바깥쪽 직사각형의
밑변환 사상을 얻는다. 더 정확히는 합성
\begin{align*}
(h \circ h')^{-1} \circ f_*
& =
(h')^{-1} \circ h^{-1} \circ f_* \\
& \to (h')^{-1} \circ f'_* \circ g^{-1} \\
& \to f''_* \circ (g')^{-1} \circ g^{-1} \\
& = f''_* \circ (g \circ g')^{-1}
\end{align*}
이 그 직사각형의 밑변환 사상이다. 확인은 생략한다.
\end{remark}








\section{SGA4와의 비교}
\label{sites-section-notation-comparison}

\noindent
절 \ref{sites-section-functoriality-PSh}의 함자 $u^p$와 $u_p$ 및
절 \ref{sites-section-continuous-functors}의 함자 $u^s$와 $u_s$에 대한
우리의 표기법은 \cite[pages 14 and 42]{ArtinTopologies}에서 가져왔다.
이렇게 선택하고 나면 절 \ref{sites-section-more-functoriality-PSh}의 함자
${}_pu$와 절 \ref{sites-section-cocontinuous-functors}의 함자 ${}_su$에 대한
표기법도 자연스럽다.
이 절에서는 우리의 표기법을 SGA4의 표기법과 비교한다.

\medskip\noindent
{\bf 준층:} $u : \mathcal{C} \to \mathcal{D}$가 범주 사이의 함자라고 하자.
함자 $u^p$는 \cite[Exposee I, Section 5]{SGA4}에서 $u^*$로 표기된다.
함자 $u_p$는 \cite[Exposee I, Proposition 5.1]{SGA4}에서 $u_!$로 표기된다.
함자 ${}_pu$는 \cite[Exposee I, Proposition 5.1]{SGA4}에서 $u_*$로 표기된다.
다시 말해,
$$
u_p, u^p, {}_pu\quad(SP)
\quad\text{versus}\quad
u_!, u^*, u_*\quad(SGA4)
$$
이다. SGA4의 서로 다른 부분에서는 이 함자들에 서로 다른 표기법을
사용한다는 점에 주의해야 한다.

\medskip\noindent
{\bf 층과 연속 함자:}
$\mathcal{C}$와 $\mathcal{D}$가 사이트이고
$u : \mathcal{C} \to \mathcal{D}$가 연속 함자라고 하자
(정의 \ref{sites-definition-continuous}).
함자 $u^s$는 \cite[Exposee III, 1.11]{SGA4}에서 $u_s$로 표기된다.
함자 $u_s$는 \cite[Exposee III, Proposition 1.2]{SGA4}에서 $u^s$로
표기된다. 다시 말해,
$$
u_s, u^s\quad(SP)
\quad\text{versus}\quad
u^s, u_s\quad(SGA4)
$$
이다. $u$가 사이트의 사상
$f : \mathcal{D} \to \mathcal{C}$를 정의할 때
(정의 \ref{sites-definition-morphism-sites}), 이에 연관된 토포스의 사상
(보조정리 \ref{sites-lemma-morphism-sites-topoi})은
\cite[Exposee IV, (4.9.1.1)]{SGA4}의 것과 같다.

\medskip\noindent
{\bf 층과 쌍대연속 함자:}
$\mathcal{C}$와 $\mathcal{D}$가 사이트이고
$u : \mathcal{C} \to \mathcal{D}$가 쌍대연속 함자라고 하자
(정의 \ref{sites-definition-cocontinuous}).
함자 ${}_su$ (보조정리 \ref{sites-lemma-pu-sheaf})는
\cite[Exposee III, Proposition 2.3]{SGA4}에서 $u_*$로 표기된다.
함자 $(u^p\ )^\#$는
\cite[Exposee III, Proposition 2.3]{SGA4}에서 $u^*$로 표기된다.
다시 말해,
$$
(u^p\ )^\#, {}_su\quad(SP)
\quad\text{versus}\quad
u^*, u_*\quad(SGA4)
$$
이다. 따라서 보조정리 \ref{sites-lemma-cocontinuous-morphism-topoi}에서
$u$에 연관된 토포스의 사상은 \cite[Exposee IV, 4.7]{SGA4}의 것과 같다.

\medskip\noindent
{\bf 토포스의 사상:} $f$가 함자 $(f^{-1}, f_*)$로 주어진 토포스의
사상이면, 함자 $f^{-1}$는
\cite[Exposee IV, Definition 3.1]{SGA4}에서 $f^*$로 표기된다.
집합의 층이나, 더 일반적으로 대수적 구조의 층의 역상을 나타낼 때에는
$f^{-1}$를 사용한다
(절 \ref{sites-section-sheaves-algebraic-structures}).
환 달린 토포스의 사상에 대한 가군층의 당김을 나타낼 때에는
$f^*$를 사용한다
(Modules on Sites, Definition \ref{sites-modules-definition-pushforward}).



\section{위상}
\label{sites-section-topologies}

\noindent
이 절에서는 \cite{SGA4}에서 정의한 범주 위의 위상을 정의한다.
이 언어를 사용하여 층과 토포스에 관한 모든 이론을 전개할 수 있다.
이 내용에 대한 현대적인 설명은 \cite{KS}에서 찾을 수 있다.
그러나 대수기하학에서 가장 중요한 경우는 사이트가 그 밑범주 위에
정의하는 위상이다. 따라서 처음 읽는 독자에게는
{\bf 이 절과 이 장의 나머지 모든 절을 건너뛸 것을 강력히 권한다}!

\begin{definition}
\label{sites-definition-sieve}
$\mathcal{C}$를 범주라 하고 $U \in \Ob(\mathcal{C})$라 하자.
{\it $U$ 위의 시브 $S$}란 부분준층 $S \subset h_U$를 말한다.
\end{definition}

\noindent
다시 말해, $U$ 위의 시브는 각 대상 $T \in \Ob(\mathcal{C})$에 대하여
모든 사상 $T \to U$의 집합에서 부분집합 $S(T)$를 골라낸다.
실제로 부분집합족
$S(T) \subset h_U(T) = \Mor_\mathcal{C}(T, U)$에 부과되는 유일한
조건은 다음 규칙이다.
\begin{equation}
\label{sites-equation-property-sieve}
\left.
\begin{matrix}
(\alpha : T \to U) \in S(T) \\
g : T' \to T
\end{matrix}
\right\} \Rightarrow
(\alpha \circ g : T' \to U) \in S(T')
\end{equation}
유용한 직관적 그림은 사상 $S \to h_U$를
“$S$에서 $U$로 가는 사상”으로 생각하는 것이다.

\begin{lemma}
\label{sites-lemma-sieves-set}
$\mathcal{C}$를 범주라 하고 $U \in \Ob(\mathcal{C})$라 하자.
\begin{enumerate}
\item $U$ 위의 시브들은 집합을 이룬다.
\item 포함관계는 이 집합 위에 부분순서를 정의한다.
\item 시브들의 합집합과 교집합은 시브이다.
\item
\label{sites-item-sieve-generated}
공역이 $U$인 $\mathcal{C}$의 사상족
$\{U_i \to U\}_{i\in I}$가 주어지면, 각 $U_i \to U$가
$S(U_i)$에 속하는 $U$ 위의 유일한 최소 시브 $S$가 존재한다.
\item 시브 $S = h_U$는 최대 시브이다.
\item 공 부분준층은 최소 시브이다.
\end{enumerate}
\end{lemma}

\begin{proof}
부분준층의 정의에 따라 준층 $\mathcal{F}$의 모든 부분준층의 모임은
$\prod_{U \in \Ob(\mathcal{C})} \mathcal{P}(\mathcal{F}(U))$의
부분집합이다. 그리고 이는 집합이다. (여기서 $\mathcal{P}(A)$는
$A$의 멱집합을 뜻한다.) 따라서 $U$ 위의 시브들의 모임은 집합이다.

\medskip\noindent
부분순서는 다음과 같이 정의한다. 모든 $T \to U$에 대하여
$S(T) \subset S'(T)$일 때 그리고 그때에만 $S \leq S'$로 둔다.
표기는 $S \subset S'$이다.

\medskip\noindent
$U$ 위의 시브족 $S_i$, $i \in I$가 주어지면, 모든
$T \in \Ob(\mathcal{C})$에 대하여
$(\bigcup S_i)(T) = \bigcup S_i(T)$인 시브로 $\bigcup S_i$를
정의할 수 있다. 교집합 $\bigcap S_i$도 같은 방식으로 정의한다.

\medskip\noindent
명제의 $\{U_i \to U\}_{i\in I}$가 주어졌다고 하자.
준층의 사상 $h_{U_i} \to h_U$를 생각한다. 이 준층 사상들의
상들(정의 \ref{sites-definition-image})의 합집합으로 $S$를 정의하면 된다.

\medskip\noindent
보조정리의 마지막 두 명제는 자명하다.
\end{proof}

\begin{definition}
\label{sites-definition-sieve-generated}
$\mathcal{C}$를 범주라 하자. 공역이 $U$인 $\mathcal{C}$의 사상족
$\{f_i : U_i \to U\}_{i\in I}$가 주어졌다고 하자.
보조정리 \ref{sites-lemma-sieves-set}의 (\ref{sites-item-sieve-generated})에서
설명한 $U$ 위의 시브 $S$를
{\it 사상들 $f_i$에 의해 생성된 $U$ 위의 시브}라고 한다.
\end{definition}

\begin{definition}
\label{sites-definition-pullback-sieve}
$\mathcal{C}$를 범주라 하고 $f : V \to U$를 $\mathcal{C}$의
사상이라 하자. $S \subset h_U$를 시브라 하자.
$f$에 의한 $S$의 {\it 당김}을 다음 규칙으로 정의되는 $V$의 시브
$S \times_U V$로 정의한다.
$$
(\alpha : T \to V) \in (S \times_U V)(T)
\Leftrightarrow
(f \circ \alpha : T \to U) \in S(T)
$$
\end{definition}

\noindent
이것이 실제로 시브임을 확인하는 것은 독자에게 맡긴다
(힌트: 식 \ref{sites-equation-property-sieve}을 사용하라).
$S \times_U V$를 $f : V \to U$에 의한 $S$의
{\it 밑변환}이라고 부르기도 한다.

\begin{lemma}
\label{sites-lemma-pullback-sieve-section}
$\mathcal{C}$를 범주라 하고 $U \in \Ob(\mathcal{C})$라 하자.
$S$를 $U$ 위의 시브라 하자. $f : V \to U$가 $S$에 속하면
$S \times_U V = h_V$는 최대 시브이다.
\end{lemma}

\begin{proof}
정의에서 자명하다.
\end{proof}

\begin{definition}
\label{sites-definition-topology}
$\mathcal{C}$를 범주라 하자. $\mathcal{C}$ 위의 {\it 위상}이란
각 $U \in \Ob(\mathcal{C})$에 $U$ 위의 모든 시브들의 집합의
부분집합 $J(U)$를 대응시키면서 다음 조건들을 만족하는 규칙이다.
\begin{enumerate}
\item $\mathcal{C}$의 모든 사상 $f : V \to U$와 모든 원소
$S \in J(U)$에 대하여 당김 $S \times_U V$는 $J(V)$의 원소이다.
\item $S$와 $S'$가 $U \in \Ob(\mathcal{C})$ 위의 시브이고,
$S \in J(U)$이며, 모든 $f \in S(V)$에 대하여 당김
$S' \times_U V$가 $J(V)$에 속하면, $S'$는 $J(U)$에 속한다.
\item 모든 $U \in \Ob(\mathcal{C})$에 대하여 최대 시브
$S = h_U$는 $J(U)$에 속한다.
\end{enumerate}
\end{definition}

\noindent
이 경우 $J(U)$에 속하는 시브를 {\it 덮개 시브}라고 한다.

\begin{lemma}
\label{sites-lemma-topology-basic}
$\mathcal{C}$를 범주라 하고 $J$를 $\mathcal{C}$ 위의 위상이라 하자.
$U \in \Ob(\mathcal{C})$라 하자.
\begin{enumerate}
\item $J(U)$의 원소들의 유한 교집합은 $J(U)$에 속한다.
\item $S \in J(U)$이고 $S' \supset S$이면 $S' \in J(U)$이다.
\end{enumerate}
\end{lemma}

\begin{proof}
(1)의 증명. 공족의 교집합에 대한 명제는
정의 \ref{sites-definition-topology}의 조건 (3)에서 나온다.
$S, S' \in J(U)$라 하고 $S'' = S \cap S'$를 생각하자.
$S(V)$에 속하는 모든 $V \to U$에 대하여
$$
S' \times_U V = S'' \times_U V
$$
이다. 이는 $V \to U$가 이미 $S$에 속하기 때문이다.
따라서 정의의 두 번째 공리에 의해 $S'' \in J(U)$이다.

\medskip\noindent
(2)의 증명. $S \in J(U)$이고 $S' \supset S$라 하자.
$S(V)$에 속하는 모든 $V \to U$에 대하여 보조정리
\ref{sites-lemma-pullback-sieve-section}에 의해
$S' \times_U V = h_V$이다. 따라서 세 번째 공리에 의해
$S' \times_U V \in J(V)$이고, 두 번째 공리에 의해 $S' \in J(U)$이다.
\end{proof}

\begin{definition}
\label{sites-definition-finer}
$\mathcal{C}$를 범주라 하고 $J$, $J'$를 $\mathcal{C}$ 위의 두 위상이라 하자.
$\mathcal{C}$의 모든 대상 $U$에 대하여 $J'(U) \subset J(U)$일 때
그리고 그때에만 $J$가 $J'$보다 {\it 더 세밀하다} 또는
{\it 더 강하다}고 한다. 이 경우 $J'$가 $J$보다
{\it 더 거칠다} 또는 {\it 더 약하다}고도 한다.
\end{definition}

\noindent
다시 말해, $J'$의 모든 덮개 시브는 $J$의 덮개 시브이다.
$\mathcal{C}$ 위에는 가장 세밀한 위상이 존재한다. 즉, 모든 시브가
덮개 시브인 위상이며, 이를 $\mathcal{C}$의 {\it 이산 위상}이라고 한다.
가장 거친 위상도 존재한다. 즉, 모든 대상 $U$에 대하여
$J(U) = \{h_U\}$인 위상이며, 이를 {\it 카오스 위상} 또는
{\it 비이산 위상}이라고 한다.

\begin{lemma}
\label{sites-lemma-play-with-topologies}
$\mathcal{C}$를 범주라 하고 $\{J_i\}_{i\in I}$를 위상들의 집합이라 하자.
\begin{enumerate}
\item 규칙 $J(U) = \bigcap J_i(U)$는 $\mathcal{C}$ 위의 위상을 정의한다.
\item 모든 위상 $J_i$보다 더 세밀한 위상들 가운데 가장 거친 것이 존재한다.
\end{enumerate}
\end{lemma}

\begin{proof}
첫 번째 명제는 정의에서 바로 나온다.
두 번째 명제는 모든 $J_i$보다 더 세밀한 모든 위상의 교집합을 취하면 나온다.
\end{proof}

\noindent
이제 아무런 동기 설명 없이 층을 정의할 수 있다.

\begin{definition}
\label{sites-definition-sheaf-sets-topology}
$\mathcal{C}$가 위상 $J$를 갖춘 범주이고 $\mathcal{F}$가
$\mathcal{C}$ 위의 집합의 준층이라고 하자.
모든 $U \in \Ob(\mathcal{C})$와 $U$의 모든 덮개 시브 $S$에 대하여
표준 사상
$$
\Mor_{\textit{PSh}(\mathcal{C})}(h_U, \mathcal{F})
\longrightarrow
\Mor_{\textit{PSh}(\mathcal{C})}(S, \mathcal{F})
$$
이 전단사이면 $\mathcal{F}$를 $\mathcal{C}$ 위의 {\it 층}이라고 한다.
\end{definition}

\noindent
표시된 식의 왼쪽이 $\mathcal{F}(U)$와 같음을 상기하자.
다시 말해, $\mathcal{F}$가 층이라는 것은, $U$ 위의 $\mathcal{F}$의 절단 하나를
주는 일이 $(\alpha : T \to U) \in S(T)$로 매개화된 양립 가능한
절단들 $s_{T, \alpha} \in \mathcal{F}(T)$의 모임을 주는 일과
같다는 뜻이며, 이는 $U$ 위의 모든 덮개 시브 $S$에 대하여 성립해야 한다.

\begin{lemma}
\label{sites-lemma-topology-presheaves-sheaves}
$\mathcal{C}$를 범주라 하고 $\{ \mathcal{F}_i \}_{i\in I}$를
$\mathcal{C}$ 위의 집합의 준층들의 모임이라 하자. 각
$U \in \Ob(\mathcal{C})$에 대하여 다음 성질을 갖는 $U$ 위의 시브
$S$들의 집합을 $J(U)$라 하자. 모든 사상 $V \to U$에 대하여 사상들
% P01-ERRATA: P01-E0163
$$
\Mor_{\textit{PSh}(\mathcal{C})}(h_V, \mathcal{F}_i)
\longrightarrow
\Mor_{\textit{PSh}(\mathcal{C})}(S \times_U V, \mathcal{F}_i)
$$
이 모든 $i \in I$에 대하여 전단사라고 하자. 그러면 $J$는
$\mathcal{C}$ 위의 위상을 정의한다. 이 위상은 모든 $\mathcal{F}_i$가
층이 되는 가장 세밀한 위상이다.
\end{lemma}

\begin{proof}
먼저 $J$가 위상이면 보조정리의 마지막 명제가 곧바로 나온다는 점에 유의하자.

\medskip\noindent
각 $i \in I$와 $U \in \Ob(\mathcal{C})$에 대하여 다음 성질을 갖는
$U$ 위의 시브 $S$들의 집합을 $J_i(U)$라 하자. 모든 사상
$V \to U$에 대하여 사상
$$
\Mor_{\textit{PSh}(\mathcal{C})}(h_V, \mathcal{F}_i)
\longrightarrow
\Mor_{\textit{PSh}(\mathcal{C})}(S \times_U V, \mathcal{F}_i)
$$
이 전단사라고 하자. $J(U) = \bigcap J_i(U)$임에 유의하자.
그러므로 $J_i$가 위상을 정의함을 보이면 보조정리
\ref{sites-lemma-play-with-topologies}에 의해 결론을 얻는다.
따라서 다음 문단에서 논의하는 경우로 귀착된다.

\medskip\noindent
우리의 모임이 하나의 준층 $\mathcal{F}$로 주어질 때 $J$가
위상임을 증명하자. 위상의 첫 번째 공리와 세 번째 공리는 바로 확인된다.
따라서 대상 $U$와 $U$ 위의 시브 $S, S'$가 있어서 $S \in J(U)$이고,
$S(V)$에 속하는 모든 $V \to U$에 대하여
$S' \times_U V \in J(V)$라고 가정하자. $S' \in J(U)$임을 보여야 한다.
다시 말해, 임의의 $f : W \to U$에 대하여 사상
$$
\mathcal{F}(W) =
\Mor_{\textit{PSh}(\mathcal{C})}(h_W, \mathcal{F})
\longrightarrow
\Mor_{\textit{PSh}(\mathcal{C})}(S' \times_U W, \mathcal{F})
$$
이 전단사임을 보여야 한다. 원소
$\varphi \in
\Mor_{\textit{PSh}(\mathcal{C})}(S' \times_U W, \mathcal{F})$를 고르자.
$\varphi$로 가는 절단 $s \in \mathcal{F}(W)$를 구성할 것이다.

\medskip\noindent
$\alpha : V \to W$가 $S \times_U W$의 원소라고 하자.
당김의 정의에 따라 합성 $f \circ \alpha : V \to W  \to U$는
$S$에 속한다. 따라서 시브의 쌍 $S, S'$에 대한 가정에 의해
$S' \times_U V$는 $J(V)$에 속한다. 이제 준층의 가환도표
$$
\xymatrix{
S' \times_U V \ar[r] \ar[d] & h_V \ar[d] \\
S' \times_U W \ar[r] & h_W
}
$$
가 있다. $\varphi$를 $S' \times_U V$로 제한한 것은 원소
$s_{V, \alpha} \in \mathcal{F}(V)$에 대응한다.
이는 $J$의 정의와 $S' \times_U V$가 $J(V)$에 속한다는 사실에서 나온다.
규칙 $(V, \alpha) \mapsto s_{V, \alpha}$가 원소
$\psi \in
\Mor_{\textit{PSh}(\mathcal{C})}(S \times_U W, \mathcal{F})$를 정의함을
확인하는 것은 독자에게 맡긴다. $S \in J(U)$이므로 $J$의 정의에서
$\psi$가 $\mathcal{F}(W)$의 한 원소 $s$에 대응함을 곧바로 알 수 있다.

\medskip\noindent
구성 $\varphi \mapsto s$가 위에서 표시한 자연사상의 역임을
확인하는 것은 독자에게 맡긴다.
\end{proof}

\begin{definition}
\label{sites-definition-canonical-topology}
$\mathcal{C}$를 범주라 하자. 모든 표현 가능한 준층이 층이 되게 하는
$\mathcal{C}$ 위의 가장 세밀한 위상
(보조정리 \ref{sites-lemma-topology-presheaves-sheaves} 참조)을
$\mathcal{C}$의 {\it 표준 위상}이라고 한다.
\end{definition}













\section{사이트가 정의하는 위상}
\label{sites-section-topology-site}

\noindent
$\mathcal{C}$가 범주이고, $\text{Cov}_1(\mathcal{C})$와
$\text{Cov}_2(\mathcal{C})$가 $\mathcal{C}$ 위에 사이트의 구조를
정의하는 덮개들의 집합이라고 하자. 이 상황에서는
$\text{Cov}_1(\mathcal{C})$와 $\text{Cov}_2(\mathcal{C})$에 대한
(집합의) 층들의 범주가 같을 수 있다. 예를 들어 보조정리
\ref{sites-lemma-refine-same-topology}을 보라.

\medskip\noindent
어떤 위상 $J$에 관한 $\mathcal{C}$ 위의 층들의 범주는
위상 $J$를 결정하며 그 위상에 의해 결정된다는 사실이 성립한다.
이는 자명하지 않은 명제이며, 뒤의 정리 \ref{sites-theorem-topology-and-topos}에서
다룰 것이다.

\medskip\noindent
당장은 이 사실을 받아들이면, 사이트가 결정하는 위상을 연구하는 것이 자연스럽다.

\begin{lemma}
\label{sites-lemma-site-gives-topology}
$\mathcal{C}$가 덮개들 $\text{Cov}(\mathcal{C})$를 갖는 사이트라고 하자.
$\mathcal{C}$의 모든 대상 $U$에 대하여, 다음 성질을 갖는 $U$ 위의
시브 $S$들의 집합을 $J(U)$라 하자. 덮개
$\{f_i : U_i \to U\}_{i\in I} \in \text{Cov}(\mathcal{C})$가 존재하여,
$f_i$들에 의해 생성된 시브 $S'$
(정의 \ref{sites-definition-sieve-generated} 참조)가 $S$에 포함된다.
\begin{enumerate}
\item 이 $J$는 $\mathcal{C}$ 위의 위상이다.
\item 준층 $\mathcal{F}$가 이 위상의 층
(정의 \ref{sites-definition-sheaf-sets-topology} 참조)일 필요충분조건은
그것이 이 사이트 위의 층인 것이다
(정의 \ref{sites-definition-sheaf-sets} 참조).
\end{enumerate}
\end{lemma}

\begin{proof}
첫 번째 명제를 증명하려면 사이트의 정의
(정의 \ref{sites-definition-site})의 공리 (1), (2), (3)이 위상의 정의
(정의 \ref{sites-definition-topology})의 공리 (3), (2), (1)을 각각
바로 함의함을 관찰하면 된다. 예로 $J$가 성질 (2)를 가짐을 증명하자.
$U$가 $\mathcal{C}$의 대상이고, $S, S'$가 $U$ 위의 시브들이며,
$S \in J(U)$이고, $S(V)$에 속하는 모든 $V \to U$에 대하여
$S' \times_U V \in J(V)$라고 하자. $J(U)$의 정의에 의해 사이트의
덮개 $\{f_i : U_i \to U\}$를 찾아서, $S$, 곧
$h_{U_i} \to h_U$의 상이 $S$에 포함되게 할 수 있다.
% P01-ERRATA: P01-E0164
각 $S'\times_U U_i$는 $J(U_i)$에 속하므로, 사이트의 덮개들
$\{U_{ij} \to U_i\}$가 존재하여 $h_{U_{ij}} \to h_{U_i}$가
$S' \times_U U_i$에 포함된다. 밑변환의 정의에 의해 이는
$h_{U_{ij}} \to h_U$가 부분준층 $S' \subset h_U$에 포함된다는
뜻이다. 사이트에 대한 공리 (2)에 의해 $\{U_{ij} \to U\}$는
$U$의 덮개이고, $J$의 정의에 의해 $S' \in J(U)$라고 결론 내린다.

\medskip\noindent
$\mathcal{F}$를 준층이라 하자. $\mathcal{F}$가 위상 $J$의 층이라고
가정하자. $\mathcal{F}$가 사이트 위의 층이기도 함을 보이겠다.
$\{f_i : U_i \to U\}_{i\in I}$를 사이트의 덮개라 하자.
$s_i \in \mathcal{F}(U_i)$가 모든 $i, j$에 대하여
$s_i|_{U_i \times_U U_j} = s_j|_{U_i \times_U U_j}$를 만족하는
절단들의 족이라고 하자. 각 $U_i$로 제한하면 $s_i$가 되는 유일한
절단 $s \in \mathcal{F}(U)$가 존재함을 보여야 한다.
$S \subset h_U$를 $f_i$들에 의해 생성된 시브라 하자.
정의에 의해 $S \in J(U)$임에 유의하자. $s$를 직접 구성하는 대신,
위상에서의 층 조건에 의해 원소
% P01-ERRATA: P01-E0165
$$
\varphi \in \Mor_{\textit{PSh}(\mathcal{C})}(S, \mathcal{F}).
$$
를 구성하면 충분하다. 어떤 대상 $T \in \mathcal{U}$에 대하여
$\alpha \in S(T)$를 택하자.
% P01-ERRATA: P01-E0166
이는 정확히 $\alpha : T \to U$가 어떤 $i\in I$에 대하여 $f_i$를
통해 분해되는 사상이라는 뜻이다(아마 $1$개보다 더 많은 것을 통해
분해될 수도 있다). 그런 지표 $i$와 분해
$\alpha = f_i \circ \alpha_i$를 하나 택하자.
$\varphi(\alpha) = \alpha_i^* s_i$로 정의한다. $i'$와
$\alpha = f_i \circ \alpha_{i'}'$가 두 번째 선택이면,
% P01-ERRATA: P01-E0167
$\alpha_i^* s_i = (\alpha_{i'}')^* s_{i'}$이다. 이는 모든
$i, j$에 대한 우리의 조건
$s_i|_{U_i \times_U U_j} = s_j|_{U_i \times_U U_j}$에서 정확히
따라온다. 따라서 $\varphi(\alpha)$는 잘 정의된다.
차례로 $s$를 결정하는 $\varphi$가 올바르다는 것, 즉 $s$를
제한하면 $s_i$가 된다는 것을 확인하는 일은 독자에게 맡긴다.

\medskip\noindent
$\mathcal{F}$를 준층이라 하자. $\mathcal{F}$가 사이트
$(\mathcal{C}, \text{Cov}(\mathcal{C}))$ 위의 층이라고 가정하자.
$\mathcal{F}$가 위상 $J$의 층이기도 함을 보이겠다.
$U$를 $\mathcal{C}$의 대상이라 하자. $S$를 위상 $J$에 관한
$U$ 위의 덮개 시브라고 하자. 다음 원소를 잡자.
$$
\varphi \in \Mor_{\textit{PSh}(\mathcal{C})}(S, \mathcal{F}).
$$
$\varphi$로 제한되는 유일한 원소가
$\mathcal{F}(U) = \Mor_{\textit{PSh}(\mathcal{C})}(h_U, \mathcal{F})$에
존재함을 보여야 한다. 정의에 의해 덮개
$\{f_i : U_i \to U\}_{i\in I} \in \text{Cov}(\mathcal{C})$가 존재하여
$f_i : U_i \to U$가 $S(U_i)$에 속한다. 따라서
$s_i = \varphi(f_i) \in \mathcal{F}(U_i)$로 둘 수 있다.
모든 $i, j$에 대하여
$s_i|_{U_i \times_U U_j} = s_j|_{U_i \times_U U_j}$임을 확인하는 것은
즐거운 연습이다. 따라서 사이트
$(\mathcal{C}, \text{Cov}(\mathcal{C}))$ 위의 $\mathcal{F}$에 대한
층 조건으로 원하는 절단 $s$를 얻는다.
세부사항은 독자에게 맡긴다.
\end{proof}

\begin{definition}
\label{sites-definition-topology-associated-site}
$\mathcal{C}$가 덮개들 $\text{Cov}(\mathcal{C})$를 갖는 사이트라고 하자.
위의 보조정리 \ref{sites-lemma-site-gives-topology}에서 구성한 위상 $J$를
{\it $\mathcal{C}$에 연관된 위상}이라고 한다.
\end{definition}

\noindent
$\mathcal{C}$를 범주라 하자. $\text{Cov}_1(\mathcal{C})$와
$\text{Cov}_2(\mathcal{C})$가 $\mathcal{C}$ 위에 사이트의 구조를
정의하는 두 덮개라고 하자. 이들이 정의하는 위상은 얼마든지 같을 수 있다.
그런 경우 $\text{Cov}_1(\mathcal{C})$와
$\text{Cov}_2(\mathcal{C})$가 $\mathcal{C}$ 위에
{\it 같은 위상을 정의한다}고 한다. 이 경우 보조정리
\ref{sites-lemma-site-gives-topology}에 의해 층들의 범주가 같다.
% P01-ERRATA: P01-E0168

\medskip\noindent
보통 우리는 덮개들의 모임이 정의하는 위상에만 관심이 있으며,
서로 다른 덮개들의 집합을 고를 수 있다는 점을 그 위상을 연구하는
도구로 여긴다.

\begin{remark}
\label{sites-remark-enlarge-coverings}
덮개류를 확대하기. 분명히 $\text{Cov}(\mathcal{C})$가
$\mathcal{C}$ 위에 사이트의 구조를 정의하면,
$\text{Cov}(\mathcal{C})$의 원소들과 자명하게 동치인
(정의 \ref{sites-definition-combinatorial-tautological} 참조)
공역이 고정된 사상족들의 임의의 집합을 $\mathcal{C}$에 추가해도
위상은 바뀌지 않는다.
% P01-ERRATA: P01-E0169
\end{remark}

\begin{remark}
\label{sites-remark-shrink-coverings}
덮개류를 축소하기. $\mathcal{C}$를 사이트라 하자.
집합
$$
\mathcal{S} = P(\text{Arrows}(\mathcal{C})) \times \Ob(\mathcal{C})
$$
를 생각하자. 여기서 $P(\text{Arrows}(\mathcal{C}))$는 사상들의
집합의 멱집합, 즉 사상들로 이루어진 모든 집합들의 집합이다.
$\mathcal{S}_\tau \subset \mathcal{S}$를 다음 조건을 만족하는
$(T, U) \in \mathcal{S}$들로 이루어진 부분집합이라 하자.
(a) 모든 $\varphi \in T$의 공역은 $U$이고,
(b) 모임 $\{\varphi\}_{\varphi \in T}$는
$\text{Cov}(\mathcal{C})$의 어떤 덮개와 자명하게 동치이다
(정의 \ref{sites-definition-combinatorial-tautological} 참조).
$\mathcal{S}_\tau$의 원소들을 덮개로 보면, 정의
\ref{sites-definition-site}에서 정의한 사이트 개념을 정확히 얻지는 못함이
분명하다. 얻은 구조 $(\mathcal{C}, \mathcal{S}_\tau)$는 조금 수정된
조건들을 만족한다. 수정된 조건들은 다음과 같다.
\begin{enumerate}
\item[(0')] $\text{Cov}(\mathcal{C}) \subset
P(\text{Arrows}(\mathcal{C})) \times \Ob(\mathcal{C})$,
\item[(1')] $V \to U$가 동형이면
$(\{V \to U\}, U) \in \text{Cov}(\mathcal{C})$이다.
\item[(2')] $(T, U) \in \text{Cov}(\mathcal{C})$이고 $T$에 속하는
$f : U' \to U$에 대하여
$(T_f, U') \in \text{Cov}(\mathcal{C})$가 주어졌다고 하자.
$T' = \{f \circ f' \mid f \in T,\ f' \in T_f\}$로 두면
$(T', U) \in \text{Cov}(\mathcal{C})$이다.
\item[(3')] $(T, U) \in \text{Cov}(\mathcal{C})$이고
$g : V \to U$가 $\mathcal{C}$의 사상이면,
\begin{enumerate}
\item $T$에 속하는 $f : U' \to U$에 대하여
$U' \times_{f, U, g} V$가 존재하고,
\item 섬유곱을 어떻게 선택했을 때
$T' = \{\text{pr}_2 : U' \times_{f, U, g} V \to V \mid f : U' \to U \in T\}$로
두면 $(T', V) \in \text{Cov}(\mathcal{C})$이다.
\end{enumerate}
\end{enumerate}
위의 (0') -- (3')을 만족하는 구조가 주어졌을 때
$\text{Cov}(\mathcal{C})$를 적절히 확대하면
(Sets, Section \ref{sets-section-coverings-site} 참조) 사이트를
얻는다는 것은 쉽게 확인된다. 적어도 위상의 관점에서는 이 개념과
실제 사이트의 개념 사이에 차이가 거의 없음이 분명하다.
두 가지 이점이 있다. 위의 조건 (0') 때문에 덮개들이 자동으로
집합을 이루고, (0') 때문에 이 유형의 모든 구조의 전체도 집합을
이룬다. 그 대가는 어디에서나 “자명하게 동치”라고 계속 써야 한다는 것이다.
\end{remark}

\section{위상에서의 층화}
\label{sites-section-sheafification-topology}

\noindent
이 절에서는 위상에서 층화 구성에 대응하는 구성을 설명한다.

\medskip\noindent
$\mathcal{C}$를 범주라 하자.
$J$를 $\mathcal{C}$ 위의 위상이라 하자.
$\mathcal{F}$를 집합의 준층이라 하자.
모든 $U \in \Ob(\mathcal{C})$에 대하여
$$
L\mathcal{F}(U)
=
\colim_{S \in J(U)^{opp}}
\Mor_{\textit{PSh}(\mathcal{C})}(S, \mathcal{F})
$$
를 여극한으로 정의한다. 여기서 $J(U)$는 포함관계로 순서가 주어진
부분순서집합으로 생각한다. 보조정리 \ref{sites-lemma-sieves-set}을 보라.
이 계의 전이사상은 다음과 같이 정의한다.
$S \subset S'$가 $J(U)$에 속하면 $S \to S'$는 준층의 사상이다.
따라서 자연스러운 제한 사상
$$
\Mor_{\textit{PSh}(\mathcal{C})}(S', \mathcal{F})
\longrightarrow
\Mor_{\textit{PSh}(\mathcal{C})}(S, \mathcal{F}).
$$
이 있다. 그러므로 $J(U)$ 위의 순서를 뒤집으면
(위첨자 ${}^{opp}$는 바로 이것을 나타낸다),
$S \mapsto \Mor_{\textit{PSh}(\mathcal{C})}(S, \mathcal{F})$는
Categories, Definition \ref{categories-definition-system-over-poset}의
유향계이다. 특히 $h_U \in J(U)$이므로 식별
$\mathcal{F}(U) = \Mor_{\textit{PSh}(\mathcal{C})}(h_U, \mathcal{F})$에서
나오는 표준 사상
$$
\ell : \mathcal{F}(U) \longrightarrow L\mathcal{F}(U)
$$
이 있다. 또한 $L\mathcal{F}(U)$를 정의하는 여극한은 유향이다.
실제로 $U$ 위의 임의의 덮개 시브의 쌍 $S, S'$에 대하여
시브 $S \cap S'$도 덮개 시브이다. 보조정리
\ref{sites-lemma-sieves-set}을 보라.
% P01-ERRATA: P01-E0170

\medskip\noindent
$f : V \to U$를 $\mathcal{C}$의 사상이라 하고 $S \in J(U)$라 하자.
가환도표
$$
\xymatrix{
S \times_U V \ar[r] \ar[d] & h_V \ar[d] \\
S \ar[r] & h_U
}
$$
가 있다. 왼쪽 수직사상을 사용하여 표준적인 제한 사상
$$
\Mor_{\textit{PSh}(\mathcal{C})}(S, \mathcal{F})
\to \Mor_{\textit{PSh}(\mathcal{C})}(S \times_U V, \mathcal{F}).
$$
을 얻을 수 있다. 밑변환 $S \mapsto S \times_U V$는
순서보존 사상 $J(U) \to J(V)$를 유도한다. 제한 사상들은
Categories, Lemma {categories-lemma-functorial-colimit}의 의미에서
함자들의 변환을 정의한다.
% P01-ERRATA: P01-E0171
따라서 자연스러운 제한 사상
$$
L\mathcal{F}(U) \longrightarrow L\mathcal{F}(V).
$$
을 얻는다.

\begin{lemma}
\label{sites-lemma-L-presheaf}
위 상황에서는 다음이 성립한다.
% P01-ERRATA: P01-E0172
\begin{enumerate}
\item 대응 $U \mapsto L\mathcal{F}(U)$와 위에서 정의한 제한 사상들을
합치면 준층을 이룬다.
\item 사상들 $\ell$은 붙어서 준층의 사상
$\ell : \mathcal{F} \to L\mathcal{F}$을 준다.
\item 규칙 $\mathcal{F} \mapsto (\mathcal{F} \xrightarrow{\ell}
L\mathcal{F})$는 함자이다.
\item $\mathcal{F}$가 $\mathcal{G}$의 부분준층이면
$L\mathcal{F}$는 $L\mathcal{G}$의 부분준층이다.
\item 사상 $\ell : \mathcal{F} \to L\mathcal{F}$은 다음 성질을 갖는다.
모든 절단 $s \in L\mathcal{F}(U)$에 대하여 $U$ 위의 덮개 시브
$S$와 원소
$\varphi \in \Mor_{\textit{PSh}(\mathcal{C})}(S, \mathcal{F})$가
존재하여 $\ell(\varphi)$는 $s$를 $S$에 제한한 것과 같다.
\end{enumerate}
\end{lemma}

\begin{proof}
생략한다.
\end{proof}

\begin{definition}
\label{sites-definition-presheaf-separated-topology}
$\mathcal{C}$를 범주라 하고 $J$를 $\mathcal{C}$ 위의 위상이라 하자.
집합의 준층 $\mathcal{F}$가 {\it 분리되었다}는 것은, 모든 대상
$U$와 $U$ 위의 모든 덮개 시브 $S$에 대하여 표준 사상
$\mathcal{F}(U) \to \Mor_{\textit{PSh}(\mathcal{C})}(S, \mathcal{F})$가
단사라는 뜻이다.
\end{definition}

\begin{theorem}
\label{sites-theorem-L-topology}
$\mathcal{C}$를 범주라 하고 $J$를 $\mathcal{C}$ 위의 위상이라 하자.
$\mathcal{F}$를 집합의 준층이라 하자.
\begin{enumerate}
\item 준층 $L\mathcal{F}$는 분리되어 있다.
\item $\mathcal{F}$가 분리되어 있으면 $L\mathcal{F}$는 층이고,
준층의 사상 $\mathcal{F} \to L\mathcal{F}$은 단사이다.
\item $\mathcal{F}$가 층이면 $\mathcal{F} \to L\mathcal{F}$은 동형이다.
\item 준층 $LL\mathcal{F}$는 항상 층이다.
\end{enumerate}
\end{theorem}

\begin{proof}
(3)은 $L$의 정의와 층의 정의
(정의 \ref{sites-definition-sheaf-sets-topology})에서 자명하다.
(4)는 나머지 명제들에서 형식적으로 따라온다.

\medskip\noindent
(1)의 증명을 개략적으로 설명하자. $S$가 대상 $U$의 덮개 시브라고 하자.
$\varphi_i \in L\mathcal{F}(U)$, $i = 1, 2$가
$\Mor_{\textit{PSh}(\mathcal{C})}(S, L\mathcal{F})$의 같은 원소로
간다고 하자. 두 $\varphi_i$가 모두 원소
$\varphi_i \in \Mor_{\textit{PSh}(\mathcal{C})}(S', \mathcal{F})$로
표현되는 $U$ 위의 하나의 덮개 시브 $S'$를 찾을 수 있다.
$S$와 $S'$를 역시 덮개 시브인 $S' \cap S$로 함께 바꾸면
$S' = S$라고 가정할 수 있다. 보조정리 \ref{sites-lemma-sieves-set}을 보라.
$V\in \Ob(\mathcal{C})$이고 $\alpha : V \to U$가 $S(V)$에
속한다고 하자. 그러면 $S \times_U V = h_V$이다.
보조정리 \ref{sites-lemma-pullback-sieve-section}을 보라.
따라서 $\varphi_i$를 $V \to U$를 통해 제한한 것은
$V$ 위의 $\mathcal{F}$의 절단 $s_{i, V, \alpha}$에 대응한다.
가정에 따르면 $V$의 덮개 시브 $S_{V, \alpha}$가 존재하여
$s_{i, V, \alpha}$들이
$\Mor_{\textit{PSh}(\mathcal{C})}(S_{V, \alpha}, \mathcal{F})$의
같은 원소로 제한된다. 다음 규칙으로 정의되는 $U$ 위의 시브
$S''$를 생각하자.
\begin{eqnarray}
\label{sites-equation-S-prime-prime}
(f : T \to U) \in S''(T)
& \Leftrightarrow &
\exists\ V , \ \alpha : V \to U, \ \alpha \in S(V), \nonumber \\
& &
\exists\ g : T \to V, \ g \in S_{V, \alpha}(T), \\
& &
f = \alpha \circ g \nonumber
\end{eqnarray}
위상의 공리 (2)에 의해 $S''$는 $U$ 위의 덮개 시브이다.
구성에 의해 $\varphi_1$과 $\varphi_2$가
$\Mor_{\textit{PSh}(\mathcal{C})}(S'', L\mathcal{F})$의
같은 원소로 제한됨을 알 수 있다. 이는 원하던 것이다.

\medskip\noindent
(2)의 증명을 개략적으로 설명하자. $\mathcal{F}$가 위상 $J$에 관하여
$\mathcal{C}$ 위에서 분리된 집합의 준층이라고 가정하자.
$S$를 $\mathcal{C}$의 대상 $U$의 덮개 시브라 하자.
$\varphi \in \Mor_\mathcal{C}(S, L\mathcal{F})$라고 가정하자.
% P01-ERRATA: P01-E0173
$\varphi$로 제한되는 원소 $s \in L\mathcal{F}(U)$를 찾아야 한다.
$V\in \Ob(\mathcal{C})$이고 $\alpha : V \to U$가 $S(V)$에
속한다고 하자. 값 $\varphi(\alpha) \in L\mathcal{F}(V)$는
$V$ 위의 덮개 시브 $S_{V, \alpha}$와 준층의 사상
$\varphi_{V, \alpha} : S_{V, \alpha} \to \mathcal{F}$로 주어진다.
위의 증명에서처럼 식 (\ref{sites-equation-S-prime-prime})으로 $U$ 위의
덮개 시브 $S''$를 정의한다. 다음의 간단한 규칙으로
$$
\varphi'' : S'' \longrightarrow \mathcal{F}
$$
를 정의한다. 모든 $f : T \to U$에 대하여,
$f \in S''(T)$이면 식 (\ref{sites-equation-S-prime-prime})에서처럼
$V, \alpha, g$를 고른다. 그리고
$$
\varphi''(f) = \varphi_{V, \alpha}(g).
$$
로 둔다. 이것이 $V, \alpha, g$의 선택과 무관하다고 주장한다.
그런 두 번째 선택 $V', \alpha', g'$를 생각하자.
% P01-ERRATA: P01-E0174
$\varphi_{V, \alpha}$와 $\varphi_{V', \alpha'}$를 $T$ 위의
다음 두 덮개 시브의 교집합으로 제한하면 서로 같다.
$$
(S_{V, \alpha} \times_{V, g} T) \cap (S_{V', \alpha'} \times_{V', g'} T)
$$
실제로 이 두 제한은 모두 $f$를 통한 $\varphi$의 $T$로의 제한에
대응하고, 원하는 등식은 $\mathcal{F}$가 분리되어 있기 때문에 나온다.
그 공통 제한을 $\psi$라고 하자.
% P01-ERRATA: P01-E0175
선택 무관성은
$\varphi_{V, \alpha}(g) = \psi(\text{id}_T) =
\varphi_{V', \alpha'}(g')$에서 나온다. 이제 $\varphi''$는 원소
$s \in L\mathcal{F}(U)$를 준다. $s$가 $\varphi$로 제한됨을 확인하는
것은 독자에게 맡긴다.
\end{proof}

\begin{definition}
\label{sites-definition-associated-sheaf-topology}
$\mathcal{C}$가 위상 $J$를 갖춘 범주이고 $\mathcal{F}$가
$\mathcal{C}$ 위의 집합의 준층이라고 하자.
층 $\mathcal{F}^\# := LL\mathcal{F}$과 표준 사상
$\mathcal{F} \to \mathcal{F}^\#$을 함께
{\it $\mathcal{F}$에 연관된 층}이라고 한다.
\end{definition}

\begin{proposition}
\label{sites-proposition-sheafification-adjoint-topology}
$\mathcal{C}$가 위상을 갖춘 범주이고 $\mathcal{F}$가
$\mathcal{C}$ 위의 집합의 준층이라고 하자.
표준 사상 $\mathcal{F} \to \mathcal{F}^\#$은 다음의 보편성질을 갖는다.
$\mathcal{G}$가 집합의 층일 때, 임의의 사상
$\mathcal{F} \to \mathcal{G}$에 대하여 유일한 사상
$\mathcal{F}^\# \to \mathcal{G}$가 존재하여
$\mathcal{F} \to \mathcal{F}^\# \to \mathcal{G}$가 주어진 사상과 같다.
\end{proposition}

\begin{proof}
명제 \ref{sites-proposition-sheafification-adjoint}의 증명과 같다.
\end{proof}

\section{위상과 층}
\label{sites-section-topology-and-sheaves}

\begin{lemma}
\label{sites-lemma-sieve-sheafification}
$\mathcal{C}$가 위상 $J$를 갖춘 범주라고 하자.
$U$를 $\mathcal{C}$의 대상이라 하자.
$S$를 $U$ 위의 시브라 하자. 다음 조건들은 동치이다.
\begin{enumerate}
\item 시브 $S$는 덮개 시브이다.
\item 사상 $S \to h_U$의 층화 $S^\# \to h_U^\#$는 동형이다.
\end{enumerate}
\end{lemma}

\begin{proof}
먼저 두 가지 일반적인 관찰을 하자.
$S^\# = LLS$와 $h_U^\# = LLh_U$를 사용할 것이다.
특히 보조정리 \ref{sites-lemma-L-presheaf}에 의해
$S^\# \to h_U^\#$가 단사임을 알 수 있다.
$\text{id}_U \in h_U(U)$임에 유의하자. 따라서 이는 $Lh_U$와
$h_U^\# = LLh_U$의 $U$ 위 절단들을 주며, 이 절단들도
$\text{id}_U$로 표기하겠다.

\medskip\noindent
$S$가 덮개 시브라고 가정하자. 합성
$h_U \to h_U^\#$가 표준 사상이 되게 하는 사상
$h_U \to S^\#$를 찾으면 충분함이 분명하다.
그런 사상을 찾으려면 $\text{id}_U$로 제한되는 절단
$s \in S^\#(U)$를 찾으면 충분하다. 그런데 $S$가 덮개 시브이므로
원소 $\text{id}_S \in \Mor_{\textit{PSh}(\mathcal{C})}(S, S)$는
$U$ 위의 $LS$의 절단을 주며, 이는 $Lh_U$에서 $\text{id}_U$로
제한된다. 따라서 원하는 결론을 얻는다.

\medskip\noindent
$S^\# \to h_U^\#$가 동형이라고 가정하자.
$1 \in S^\#(U)$을 $h_U^\#(U)$의 $\text{id}_U$에 대응하는
원소라 하자. $S^\# = LLS$이므로 $U$ 위의 덮개 시브 $S'$가
존재하여 $1$은 다음 원소에서 나온다.
$$
\varphi \in \Mor_{\textit{PSh}(\mathcal{C})}(S', LS).
$$
이는 다시 모든 $\alpha : V \to U$와 $\alpha\in S'(V)$에 대하여
$V$ 위의 덮개 시브 $S_{V, \alpha}$가 존재하여
$\varphi(\alpha)$가 준층의 사상 $S_{V, \alpha} \to S$에
대응한다는 뜻이다. 다시 말해 $S_{V, \alpha}$는
$S \times_U V$에 포함된다. 위상의 두 번째 공리에 의해
$S$가 덮개 시브임을 알 수 있다.
\end{proof}

\begin{theorem}
\label{sites-theorem-topology-and-topos}
$\mathcal{C}$를 범주라 하고 $J$, $J'$를 $\mathcal{C}$ 위의
위상들이라 하자. 다음 조건들은 동치이다.
\begin{enumerate}
\item $J = J'$,
\item 위상 $J$에 대한 층들은 위상 $J'$에 대한 층들과 같다.
\end{enumerate}
\end{theorem}

\begin{proof}
$J = J'$이면 층의 개념들이 같다는 것은 자명하다.
역으로 보조정리 \ref{sites-lemma-sieve-sheafification}은 층화 함자를
사용하여 덮개 시브들을 특징짓는다. 그런데 층화 함자
$\textit{PSh}(\mathcal{C}) \to \Sh(\mathcal{C}, J)$는 포함함자
$\Sh(\mathcal{C}, J) \to \textit{PSh}(\mathcal{C})$의 왼쪽
수반이다. 따라서 부분범주 $\Sh(\mathcal{C}, J)$와
$\Sh(\mathcal{C}, J')$가 같으면 층화 함자들도 같고,
따라서 덮개 시브들의 모임들도 같다.
\end{proof}

\begin{lemma}
\label{sites-lemma-finer-topology}
가정들과 표기법은 정리 \ref{sites-theorem-topology-and-topos}에서와 같다.
그때 $J \subset J'$일 필요충분조건은 위상 $J'$의 모든 층이
위상 $J$의 층인 것이다.
% P01-ERRATA: P01-E0176
\end{lemma}

\begin{proof}
한 방향은 분명하다. 다른 방향을 위해
$\Sh(\mathcal{C}, J') \subset \Sh(\mathcal{C}, J)$라고 가정하자.
형식적 논의에 의해 다음이 따라온다. $\mathcal{F}$가 집합의
준층이고, $\mathcal{F} \to \mathcal{F}^\#$ 및 각각
$\mathcal{F} \to \mathcal{F}^{\#, \prime}$가 각각 $J$ 및 $J'$에
관한 층화이면, 표준 사상
$\mathcal{F}^\# \to \mathcal{F}^{\#, \prime}$가 존재하여
$\mathcal{F} \to \mathcal{F}^\# \to \mathcal{F}^{\#, \prime}$가
표준 사상 $\mathcal{F} \to \mathcal{F}^{\#, \prime}$와 같다.
물론 $\mathcal{F}^\# \to \mathcal{F}^{\#, \prime}$는 두 번째 층을
첫 번째 층의 위상 $J'$에 관한 층화와 식별한다.
이를 보조정리 \ref{sites-lemma-sieve-sheafification}의 사상
$S \to h_U$에 적용하자. 가환도표
$$
\xymatrix{
S \ar[r] \ar[d] &
S^\# \ar[r] \ar[d] &
S^{\#, \prime} \ar[d] \\
h_U \ar[r] &
h_U^\# \ar[r] &
h_U^{\#, \prime}
}
$$
를 얻는다. $S$가 위상 $J$의 덮개 시브이면 중간 수직사상은
(보조정리에 의해) 동형임이 분명하고, 오른쪽 수직사상도
중간 사상의 $J'$에 관한 층화이므로 동형이라고 결론 내린다.
다시 보조정리에 의해 $S$는 $J'$의 덮개 시브이기도 하다.
\end{proof}




\section{위상과 연속 함자}
\label{sites-section-topologies-continuous-functors}

\noindent
연속 함자가 층들 위에 한 쌍의 수반함자를 어떻게 주는지 설명한다.





\section{점과 위상}
\label{sites-section-points-topologies}

\noindent
절 \ref{sites-section-points}을 상기하자. 함자
$p = u : \mathcal{C} \to \textit{Sets}$가 주어지면 줄기 함자
$$
\textit{PSh}(\mathcal{C}) \longrightarrow \textit{Sets},
\mathcal{F} \longmapsto \mathcal{F}_p.
$$
를 정의할 수 있다.

\begin{definition}
\label{sites-definition-point-topology}
$\mathcal{C}$를 범주라 하고 $J$를 $\mathcal{C}$ 위의 위상이라 하자.
이 위상의 {\it 점 $p$}는 다음 조건들을 만족하는 함자
$u : \mathcal{C} \to \textit{Sets}$로 주어진다.
\begin{enumerate}
\item $U$ 위의 모든 덮개 시브 $S$에 대하여 사상
$S_p \to (h_U)_p$는 전사이다.
\item 줄기 함자 $\Sh(\mathcal{C}) \to \textit{Sets}$,
$\mathcal{F} \to \mathcal{F}_p$는 완전하다.
\end{enumerate}
% P01-ERRATA: P01-E0177
\end{definition}
